\documentclass[11pt]{article}
\pdfvariable minorversion=7\relax
\usepackage[margin=1in]{geometry}
\usepackage{amsmath,amssymb,amsthm,mathtools}
\usepackage{microtype,lmodern}
\usepackage[T1]{fontenc}
\usepackage{xcolor}
\usepackage[colorlinks=true,allcolors=blue]{hyperref}
\usepackage[nameinlink]{cleveref}
% Semantic, multi-shade mathematical notation for the Davis--Kahan paper.
% The six core families follow visualizations/dkvis/palette.py (LIGHT).
% Subshades distinguish objects that are RELATED but not IDENTICAL.
% Values were chosen for legibility against white manuscript paper; some
% amber, green, and vermilion anchors are a little darker than the slides.
%
% Color control (independent of hover reminders):
%   \DKMathColorsfalse   (after % Semantic, multi-shade mathematical notation for the Davis--Kahan paper.
% The six core families follow visualizations/dkvis/palette.py (LIGHT).
% Subshades distinguish objects that are RELATED but not IDENTICAL.
% Values were chosen for legibility against white manuscript paper; some
% amber, green, and vermilion anchors are a little darker than the slides.
%
% Color control (independent of hover reminders):
%   \DKMathColorsfalse   (after % Semantic, multi-shade mathematical notation for the Davis--Kahan paper.
% The six core families follow visualizations/dkvis/palette.py (LIGHT).
% Subshades distinguish objects that are RELATED but not IDENTICAL.
% Values were chosen for legibility against white manuscript paper; some
% amber, green, and vermilion anchors are a little darker than the slides.
%
% Color control (independent of hover reminders):
%   \DKMathColorsfalse   (after \input{math_symbol_colors.tex} in paper.tex)
% or pass \def\DKMonochromeMath{1} before \input{paper.tex}, without editing.
% EVERY \sym, \like, and \dkrole obeys this switch, including the glossary.
% No color is required to identify a mathematical object: the glyphs, indices,
% definitions and adjacent prose also distinguish them.
%
% Families (base hue): exact/target=blue, trial/source=amber, exact unwanted=
% cyan, angles and their operators=magenta, residual/displacement=green,
% gaps=violet. Norm gauges=slate; provenance retains brown/steel/vermilion.
% P_ex is the exact-target projection in sine-theta; P is the source
% projection in transport. R denotes a residual and D a direct rotation.
% Semantic keys retain their roles even when the printed notation changes.
\newif\ifDKMathColors
\DKMathColorstrue
\ifdefined\DKMonochromeMath
  \DKMathColorsfalse
\fi

% Family anchors (print-compatible versions of the Manim light palette).
\definecolor{dkfg}{HTML}{1B1F24}
\definecolor{dkmuted}{HTML}{6B7280}
\definecolor{dkwanted}{HTML}{1F6FD1}
\definecolor{dktrial}{HTML}{A66705}
\definecolor{dkunwanted}{HTML}{007C91}
\definecolor{dksine}{HTML}{D12A7B}
\definecolor{dkresid}{HTML}{087D43}
\definecolor{dkgap}{HTML}{6E4BD8}
\definecolor{dkold}{HTML}{705519}
\definecolor{dkcurrent}{HTML}{054669}
\definecolor{dkperturb}{HTML}{BB491B}
\definecolor{dkrefuted}{HTML}{C62828}

% Related yet distinguishable shades. These are NOT new unrelated categories:
% objects inherit their parent family; a different shade signals a different
% mathematical object (e.g., a frame versus its subspace or its projection).
\definecolor{dkwantedframe}{HTML}{1254A9}
\definecolor{dkwantedprojection}{HTML}{08618D}
\definecolor{dkwantedrotation}{HTML}{4248B0}
\definecolor{dkwantedspectrum}{HTML}{345F8F}
\definecolor{dkwantedvector}{HTML}{1468AA}
\definecolor{dktrialframe}{HTML}{905600}
\definecolor{dktrialprojection}{HTML}{79511B}
\definecolor{dktrialcompetitor}{HTML}{A04D10}
\definecolor{dktrialspectrum}{HTML}{79591C}
\definecolor{dktrialvector}{HTML}{9B522A}
\definecolor{dkunwantedframe}{HTML}{006477}
\definecolor{dkunwantedspectrum}{HTML}{00546B}
\definecolor{dkunwantedprojection}{HTML}{145B6D}
\definecolor{dkunwantedvector}{HTML}{256879}
\definecolor{dkangleoperator}{HTML}{923478}
\definecolor{dksineoperator}{HTML}{AA275A}
\definecolor{dksineblock}{HTML}{B43F70}
\definecolor{dksylvester}{HTML}{804074}
\definecolor{dkanglemax}{HTML}{B32185}
\definecolor{dkanglecos}{HTML}{824276}
\definecolor{dkresidvector}{HTML}{086038}
\definecolor{dkdisplacement}{HTML}{226C50}
\definecolor{dkgaporiginal}{HTML}{4E4D89}
\definecolor{dkgapstat}{HTML}{7039A7}
\definecolor{dkgapinterval}{HTML}{574A82}
\definecolor{dkambient}{HTML}{36516A}
\definecolor{dknorm}{HTML}{475875}
\definecolor{dknormaux}{HTML}{66526E}
\definecolor{dkintertwiner}{HTML}{4A6280}
\definecolor{dkturn}{HTML}{47647D}
\definecolor{dkmix}{HTML}{505B65}

% Coloring primitives; all descendants (including indices and adjoints) can
% be colored as ONE semantic expression by writing \like{key}{expression}.
\newcommand{\dkrole}[2]{\ifDKMathColors\textcolor{dk#1}{#2}\else#2\fi}
% An optional plain-text reminder supplies definitions and types on hover,
% while the final argument stays short enough for the printed notation table.
\newcommand{\dkDeclareSymbol}[5][]{%
  \expandafter\def\csname dkgloss@#2\endcsname{#5}%
  \if\relax\detokenize{#1}\relax
    \expandafter\pdfstringdef\csname dkpdfgloss@#2\endcsname{#5}%
  \else
    \expandafter\pdfstringdef\csname dkpdfgloss@#2\endcsname{\detokenize{#1}}%
  \fi
  \expandafter\def\csname dkglyph@#2\endcsname{#3}%
  \expandafter\def\csname dkrole@#2\endcsname{#4}%
  \expandafter\def\csname dksym@#2\endcsname{\dkrole{#4}{#3}}%
}
\DeclareRobustCommand{\sym}[1]{%
  \ifcsname dksym@#1\endcsname
    \dkSymbolTooltip{#1}{\csname dksym@#1\endcsname}%
  \else
    \PackageError{dk-math-colors}{Unknown symbol key #1}{See math_symbol_colors.tex}%
  \fi
}
\DeclareRobustCommand{\like}[2]{%
  \ifcsname dkrole@#1\endcsname
    \dkSymbolTooltip{#1}{\dkrole{\csname dkrole@#1\endcsname}{#2}}%
  \else
    \PackageError{dk-math-colors}{Unknown role key #1}{See math_symbol_colors.tex}%
  \fi
}

% A short description belongs to the semantic key, not the bare letter.
% The same descriptions are available to the notation appendix.
\DeclareRobustCommand{\dkSymbolDescription}[1]{\csname dkgloss@#1\endcsname}
\newif\ifDKSymbolTooltips
\DKSymbolTooltipstrue
\ifdefined\DKNoSymbolTooltips\DKSymbolTooltipsfalse\fi
\newsavebox{\DKTooltipBox}
\newcount\DKTooltipCount
% Register only shipped widgets through hyperref's existing form-field
% machinery. Poppler/Okular need the catalog's field list, not just page
% annotations. The deferred aux write also discards measuring-only copies.
\makeatletter
\newcommand{\dkRegisterTooltipField}{%
  \expandafter\HyField@@AddToFields\expandafter{\number\pdffeedback lastannot}%
}
\makeatother
\AtBeginDocument{\ifDKSymbolTooltips\Form[NeedAppearances=false]\fi}
\AtEndDocument{\ifDKSymbolTooltips\endForm\fi}
% Native LuaLaTeX PDF widgets: no navigation action or JavaScript. Avoid
% pdfcomment's soulpos dependency, which previously broke local Vim builds.
% Measure in the CURRENT math style, then emit the original expression as-is:
% boxing the visible math would change spacing, scripts and line breaking.
\newcommand{\dkSymbolTooltip}[2]{%
  \ifDKSymbolTooltips
    \begingroup
      \DKSymbolTooltipsfalse % suppress nested annotations while measuring
      \ifmmode
        \edef\DKTooltipStyle{\ifcase\mathstyle
          \displaystyle\or\displaystyle\or\textstyle\or\textstyle\or
          \scriptstyle\or\scriptstyle\or\scriptscriptstyle\or\scriptscriptstyle\fi}%
        \sbox{\DKTooltipBox}{$\DKTooltipStyle#2$}%
      \else
        \leavevmode\sbox{\DKTooltipBox}{#2}%
      \fi
      \global\advance\DKTooltipCount by 1
      \pdfextension annot width\wd\DKTooltipBox
        height\ht\DKTooltipBox depth\dp\DKTooltipBox
        {/Subtype /Widget /FT /Btn /Ff 65536 /F 512
         /T (dk-symbol-\the\DKTooltipCount) /TU (\csname dkpdfgloss@#1\endcsname)
         /H /N /BS << /W 0 >> /Border [0 0 0]}%
      \dkRegisterTooltipField
      #2%
    \endgroup
  \else
    #2%
  \fi
}
\pdfstringdefDisableCommands{%
  \def\sym#1{\csname dkglyph@#1\endcsname}%
  \def\like#1#2{#2}%
}

% Declarations: [optional detailed hover reminder], key, glyph, role, short description.
% This registry includes local proof notation. The notation appendix selects
% the symbols useful for reference; registration does not add a printed row.
% Manim's original symbol keys, now split into semantic subshades.
% Ambient and perturbed operators (provenance is independent of geometry).
\dkDeclareSymbol[Reference self-adjoint operator A on the ambient Hilbert space; the perturbed operator is Ahat = A + H. In a single-operator eigenproblem, A itself is the operator being approximated. An unbounded A acts only on domain(A).]
  {A}{A}{old}{Reference self-adjoint operator before perturbation.}
\dkDeclareSymbol{A1}{A_1}{old}{Complementary block of the reference operator.}
\dkDeclareSymbol{a1}{a_1}{old}{Reference eigenvalue or complementary spectral value.}
\dkDeclareSymbol[Reference block A_0: the restriction of A to the range of the reducing source projection P, expressed in source coordinates. In Theorem 8.1 it lies at or below alpha; in Theorem 8.2 its spectrum lies in the enlarged central interval. The perturbed blocks are denoted Lambda_0 and Lambda_1.]
  {ReferenceBlock0}{A_0}{old}{Reference operator restricted to the source subspace in Section 8.}
\dkDeclareSymbol[Reference block A_1: the restriction of A to the kernel of the reducing source projection P, expressed in complementary coordinates. In Theorem 8.1 it lies at or above alpha + delta in quadratic-form order.]
  {ReferenceBlock1}{A_1}{old}{Reference operator restricted to the source complement in Section 8.}
\dkDeclareSymbol[Self-adjoint ambient operator Ahat being studied; in the perturbation setting Ahat = A + H. Its action on a trial vector is compared with the proposed eigenvalue times that vector.]
  {Ahat}{\widehat A}{current}{Operator under study; in the perturbation setting, the reference operator plus its perturbation.}
\dkDeclareSymbol[Self-adjoint ambient operator Ahat being studied. In the perturbation setting Ahat = A + H. The trial residual is R = Ahat E_0 - E_0 A_0, and the exact complementary block satisfies Ahat F_1 = F_1 Lambda_1 on its domain.]
  {Acurrent}{\widehat A}{current}{Operator under study; in the perturbation setting, the reference operator plus its perturbation.}
\dkDeclareSymbol[Additive self-adjoint perturbation H = Ahat - A. A bound on its norm can control the residual of an eigenspace of one operator relative to the other.]
  {H}{H}{perturb}{Additive self-adjoint perturbation of the reference operator.}
\dkDeclareSymbol{H0}{H_0}{perturb}{Wanted diagonal block of the perturbation.}
\dkDeclareSymbol{H1}{H_1}{perturb}{Unwanted diagonal block of the perturbation.}
\dkDeclareSymbol{B}{B}{perturb}{Off-diagonal perturbation or coupling block.}
\dkDeclareSymbol{b}{b}{perturb}{Scalar perturbation or coupling coefficient.}
% Trial/source family: subspace, frame, compressed operator, projector, vector.
\dkDeclareSymbol{U}{U}{wanted}{Desired exact subspace in the local geometric notation.}% Manim's U remains wanted in local slide notation.
\dkDeclareSymbol{V}{V}{trial}{Trial subspace in the local geometric notation.}% Manim's V remains trial in local slide notation.
\dkDeclareSymbol[Source subspace U to be carried onto the target V by an orthogonal or unitary map W. In the transport sections P projects onto U and Q onto V. The matching theorem also names the source subspace K.]
  {Usource}{U}{trial}{Source subspace to be transported onto the target subspace.}
\dkDeclareSymbol[Trial frame E_0: F -> E, a linear isometry from trial-coordinate space F into ambient space E. Its range is the trial subspace; E_0* returns trial coordinates and E_0* E_0 = I_F. In rank one, E_0(x) = x v.]
  {E0}{E_0}{trialframe}{Isometric embedding of the trial subspace; its adjoint returns trial coordinates.}
\dkDeclareSymbol[Self-adjoint trial operator A_0 on trial-coordinate space F. The residual compares Ahat E_0 with E_0 A_0. Rayleigh--Ritz chooses A_0 = E_0* Ahat E_0 when this compression is defined; the theorem allows other trial operators.]
  {A0}{A_0}{trialspectrum}{Self-adjoint operator on trial coordinates; its eigenvalues are the trial spectral values.}
\dkDeclareSymbol[Trial vector: a proposed approximate eigenvector. In the rank-one geometry v has unit length and E_0(x) = x v. Its residual for trial scalar rho is r = Ahat v - rho v.]
  {v}{v}{trialvector}{Trial vector or vector in the trial subspace.}
\dkDeclareSymbol[Scalar trial operator in rank one, A_0(x) = rho x. Figure 1 uses the Rayleigh value rho = <v, Ahat v> for unit v; then rho v is the projection of Ahat v onto the trial line and r is perpendicular to v.]
  {rho}{\rho}{trialspectrum}{Trial eigenvalue or scalar trial operator.}
\dkDeclareSymbol[Approximate eigenvalue assigned to a trial vector v. For the local operator A, its residual is A v - mu v. For unit v, the Rayleigh--Ritz choice is mu = <A v, v>.]
  {mu}{\mu}{trialspectrum}{Trial spectral value.}
\dkDeclareSymbol{a0}{a_0}{trialspectrum}{Wanted trial eigenvalue or trial spectral value.}
\dkDeclareSymbol[Ambient orthogonal or unitary transporter W with W(U) = V, equivalently W P W* = Q. Its full displacement is I - W on the whole ambient space. The real matching theorem minimizes its trace norm over every such orthogonal W.]
  {Wcompetitor}{W}{trialcompetitor}{Orthogonal or unitary competitor carrying the source subspace onto the target.}
\dkDeclareSymbol[Orthogonal projection P onto the source subspace U in the transport problem; P = P* = P squared. A feasible orthogonal or unitary W satisfies W P W* = Q. Restricted displacement (I - W)P tests only source vectors.]
  {Psource}{P}{trialprojection}{Orthogonal projection onto the source subspace in the transport problem.}
\dkDeclareSymbol{ProjSource}{P_U}{trialprojection}{Orthogonal projection onto the source subspace in the transport problem.}
\dkDeclareSymbol[Vector x in the source subspace U. For principal vectors, unit x_j in U and y_j in V are chosen with <x_i, y_j> = 0 for i different from j and <x_j, y_j> = cos(theta_j).]
  {SourceVector}{x}{trialvector}{Vector in the source subspace.}
\dkDeclareSymbol[Orthonormal coordinate frame X for the source subspace, so X* X = I. Together with a target frame Y, the singular values of X* Y are the principal cosines. In the four-dimensional counterexample X is the vertical stack of I_2 and 0.]
  {Xframe}{X}{trialframe}{Orthonormal coordinate matrix for the source subspace in the counterexample.}
\dkDeclareSymbol[Hilbert space F of trial coordinates. E_0 embeds F isometrically into ambient space E; A_0 and the positive sine operator sin(Theta_0) act on F, while the residual R maps F into E.]
  {Ftrial}{F}{trialframe}{Hilbert space of trial coordinates.}
\dkDeclareSymbol[Ambient Hilbert space E over the real or complex scalars. It contains the trial subspace range(E_0) and the orthogonal exact subspaces range(F_0) and range(F_1). The source theorem permits infinite-dimensional separable E.]
  {Eambient}{E}{ambient}{Ambient Hilbert space containing the trial and exact subspaces.}
% Exact/wanted family: the target subspace, its frame, projection and rotation.
\dkDeclareSymbol[Target subspace V in the transport problem: feasible maps satisfy W(U) = V, or equivalently W P W* = Q. The matching theorem also names this subspace L.]
  {Vtarget}{V}{wanted}{Target subspace in the transport problem.}
\dkDeclareSymbol[Exact wanted frame F_0: F_ex -> E, a linear isometry onto the desired exact spectral subspace. F_0* returns its coordinates. P_ex = F_0 F_0* is its ambient orthogonal projection; in Figure 1, F_0(h) = (h, 0).]
  {F0}{F_0}{wantedframe}{Isometric coordinates for the desired exact spectral subspace; star denotes the adjoint.}
\dkDeclareSymbol[Self-adjoint operator block on coordinates for the desired exact spectral subspace. Its spectrum consists of the wanted exact spectral values; in the corresponding decomposition Ahat F_0 = F_0 Lambda_0 on the operator domain.]
  {Lambda0}{\Lambda_0}{wantedspectrum}{Self-adjoint block on the desired exact spectral subspace.}
\dkDeclareSymbol{u}{u}{wantedvector}{Vector in the desired exact subspace.}
\dkDeclareSymbol{f0}{f_0}{wantedvector}{Component or coordinate vector in the desired exact subspace.}
\dkDeclareSymbol{lambda0}{\lambda_0}{wantedspectrum}{Wanted exact eigenvalue or wanted spectral value.}
\dkDeclareSymbol{lambdau}{\lambda_u}{wantedspectrum}{Eigenvalue associated with a desired exact vector.}
\dkDeclareSymbol[Exact wanted orthogonal projection P_ex: E -> E, with P_ex = F_0 F_0*. It retains the desired exact spectral component. Its complementary projection is P_ex_perp = I - P_ex = F_1 F_1*; the trial leakage is S = P_ex_perp E_0.]
  {Pexact}{P_{\mathrm{ex}}}{wantedprojection}{Orthogonal projection onto the desired exact subspace in the sine-theta theorem.}
\dkDeclareSymbol[Exact complementary orthogonal projection P_ex_perp = I - P_ex = F_1 F_1*: E -> E. It retains the unwanted exact component. In Figure 1, P_ex_perp v = S(1), whereas P_ex_perp r is the unwanted component of the residual.]
  {PexactPerp}{P_{\mathrm{ex}}^\perp}{unwantedprojection}{Orthogonal projection onto the exact complementary subspace.}
\dkDeclareSymbol{ReflectionSigma}{\Sigma}{wantedprojection}{Reflection associated with the desired exact spectral projection.}
\dkDeclareSymbol[Orthogonal projection Q onto the target subspace V in the transport problem. A feasible orthogonal or unitary W satisfies W P W* = Q, where P projects onto the source U.]
  {Qtarget}{Q}{wantedprojection}{Orthogonal projection onto the target subspace in the transport problem.}
\dkDeclareSymbol{ProjTarget}{P_V}{wantedprojection}{Orthogonal projection onto the target subspace in the transport problem.}
\dkDeclareSymbol[Canonical direct rotation D from the source U to the target V for an acute pair. It is the unitary polar factor of T = QP + (I - Q)(I - P), and acts by the principal angle on each principal plane. R is reserved for the trial residual.]
  {Rrotation}{D}{wantedrotation}{Canonical direct rotation carrying the source subspace onto the target.}
\dkDeclareSymbol[Direct rotation D between the ranges of orthogonal projections P and Q, including the nonacute case. Relative to range(P) and kernel(P), its diagonal blocks are positive and its off-diagonal blocks are opposite adjoints. It exists when the two crossed defect spaces have equal dimension, and is unique when both defects vanish.]
  {SourceDirectRotation}{D}{wantedrotation}{Direct rotation, including the nonacute case in the source-result table.}
\dkDeclareSymbol[Orthonormal coordinate frame Y for the target subspace, so Y* Y = I. In the counterexample Y is the vertical stack of cos(theta) I_2 and sin(theta) J_0; hence X* Y = cos(theta) I_2 and both principal angles equal theta.]
  {Yframe}{Y}{wantedframe}{Orthonormal coordinate matrix for the target subspace in the counterexample.}
\dkDeclareSymbol[Vector y in the target subspace V. For principal vectors, unit x_j in U and y_j in V are paired so that <x_j, y_j> = cos(theta_j); distinct pairs have zero cross inner product.]
  {TargetVector}{y}{wantedvector}{Vector in the target subspace.}
\dkDeclareSymbol[Hilbert space F_ex of exact wanted coordinates. The isometry F_0: F_ex -> E identifies it with the desired exact spectral subspace. This space differs from H, the additive perturbation.]
  {Hexact}{F_{\mathrm{ex}}}{wantedframe}{Hilbert space of coordinates for the desired exact spectral subspace.}% Coordinate Hilbert space (not perturbation).
% Exact/unwanted family; the trial error operators are MAGENTA, not cyan.
\dkDeclareSymbol{Uperp}{U^{\perp}}{unwanted}{Orthogonal complement of the desired exact subspace in the local notation.}
\dkDeclareSymbol[Exact complementary frame F_1: G -> E, a linear isometry onto the orthogonal complement of range(F_0). F_1* returns unwanted coordinates; F_1 F_1* = I - P_ex. In Figure 1, F_1(g) = (0, g).]
  {F1}{F_1}{unwantedframe}{Isometric coordinates for the exact complementary subspace; star denotes the adjoint.}
\dkDeclareSymbol[Self-adjoint unwanted block Lambda_1 on G, satisfying Ahat F_1 = F_1 Lambda_1 on its domain. Its spectrum is separated from the trial spectrum of A_0 by the specified positive delta. Figure 1 uses the scalar block 0.7.]
  {Lambda1}{\Lambda_1}{unwantedspectrum}{Self-adjoint operator on exact complementary coordinates: the unwanted spectral block.}
% Spectral measure and set (distinct from the ambient space E and sine block S).
\dkDeclareSymbol[Projection-valued spectral measure E associated with a self-adjoint operator A: A = integral lambda dE(lambda). For a Borel set S, E(S) is the orthogonal projection onto the corresponding spectral subspace.]
  {SpectralMeasure}{E}{wantedprojection}{Projection-valued spectral measure of a self-adjoint operator.}
\dkDeclareSymbol[Borel subset S of the real spectral axis. The spectral projection E(S) selects the part of a self-adjoint operator with spectral values in S; its range is a reducing subspace.]
  {SpectralSet}{S}{wantedspectrum}{Borel spectral set used to select a spectral subspace.}
% Generic Sylvester coefficients and right side, by source/target role.
\dkDeclareSymbol[Left coefficient B of the Sylvester map X -> B X - X C; B acts on the output space of X. For the sine bound, B = Lambda_1 on the unwanted-coordinate space G.]
  {SylvesterLeft}{B}{unwantedspectrum}{Left coefficient of the Sylvester equation; corresponds to the unwanted spectral block.}
\dkDeclareSymbol[Right coefficient C of the Sylvester map X -> B X - X C; C acts on the input space of X. For the sine bound, C = A_0 on trial space F. The historical equation also denotes this coefficient by A.]
  {SylvesterRight}{C}{trialspectrum}{Right coefficient of the Sylvester equation; corresponds to the trial operator.}
\dkDeclareSymbol[Right-hand side Y in B X - X C = Y, with the same input and output spaces as X. For the sine bound, Y = F_1* R: F -> G, the residual expressed in unwanted exact coordinates.]
  {SylvesterRHS}{Y}{resid}{Right-hand side of the Sylvester equation, obtained by projecting the residual.}
\dkDeclareSymbol{w}{w}{unwantedvector}{Vector in an unwanted or complementary spectral subspace.}
\dkDeclareSymbol{f1}{f_1}{unwantedvector}{Component or coordinate vector in the unwanted exact subspace.}
\dkDeclareSymbol{lambda1}{\lambda_1}{unwantedspectrum}{Unwanted exact eigenvalue or complementary spectral value.}
\dkDeclareSymbol{lambdaw}{\lambda_w}{unwantedspectrum}{Eigenvalue associated with an unwanted exact vector.}
\dkDeclareSymbol[Hilbert space G of exact complementary coordinates. F_1: G -> E embeds it onto the unwanted exact subspace; Lambda_1 acts on G and X = F_1* E_0 maps trial coordinates into G.]
  {Gunwanted}{G}{unwantedframe}{Hilbert space of exact complementary coordinates.}
\dkDeclareSymbol{E1}{E_1}{unwantedframe}{Isometric coordinates for the complementary trial subspace.}
\dkDeclareSymbol{Vperp}{V^{\perp}}{unwanted}{Orthogonal complement of the trial subspace in the local notation.}
% Angle and error geometry, including distinct rectangular and positive maps.
\dkDeclareSymbol[Individual principal angle, between 0 and pi/2. For orthonormal subspace frames X and Y, the singular values of X* Y are the principal cosines. In Figure 1, theta is the angle between the unit trial vector v and the exact wanted line.]
  {theta}{\theta}{sine}{Principal angle between the subspaces under comparison.}
\dkDeclareSymbol[Full positive angle operator Theta = diag(Theta_0, Theta_1) for the source/target decompositions. Theta_0 compares the source with the target; Theta_1 compares their orthogonal complements. Its operator norm is the largest spectral angle. For positive c, the conditions Theta <= c and Theta < c mean operator norm at most c and strictly below c, respectively.]
  {Theta}{\Theta}{angleoperator}{Full angle operator, the diagonal sum of the source and complementary angle blocks.}
\dkDeclareSymbol[Directed principal-angle operator Theta_0 on source or trial coordinates. For a trial frame E_0 and exact wanted projection P_ex, cos(Theta_0) = sqrt(E_0* P_ex E_0), and sin(Theta_0) = sqrt(I_F - E_0* P_ex E_0). In the two-projection geometry it comes from the positive cosine compression P Q P restricted to range(P).]
  {Theta0}{\Theta_0}{angleoperator}{Directed principal-angle operator acting on source or trial coordinates.}
\dkDeclareSymbol[Complementary directed angle operator Theta_1, obtained from the projections I-P and I-Q by the same cosine-compression construction as Theta_0. It acts on the source complement, or its isometric coordinates. The full angle operator is diag(Theta_0, Theta_1).]
  {Theta1}{\Theta_1}{angleoperator}{Directed angle operator comparing the orthogonal complements.}
\dkDeclareSymbol[Right-angle source defect D_0 = range(P) intersect kernel(Q). These are source vectors orthogonal to the target. The full-angle conclusion in the Section 8 smallness criterion assumes dim(D_0) = dim(D_1); in Lean their identification is CrossedDefectsEquivalent.]
  {CrossedSourceDefect}{D_0}{fg}{Source directions perpendicular to the target subspace.}
\dkDeclareSymbol[Right-angle target defect D_1 = kernel(P) intersect range(Q). These are target vectors orthogonal to the source. Equality of its Hilbert dimension with that of D_0 is assumed for the full-angle conclusion of Theorem 8.2.]
  {CrossedTargetDefect}{D_1}{fg}{Target directions perpendicular to the source subspace.}
\dkDeclareSymbol[Positive complementary cosine block C_1 = cos(Theta_1). The reference complement is identified with the perturbed complement by the direct rotation D, so A_1 and Lambda_1 act on common coordinates. Theorem 8.1 compares A_1 - alpha I with C_1 (Lambda_1 - alpha I) C_1 as quadratic forms.]
  {ComplementCosineBlock}{C_1}{anglecos}{Cosine of the complementary angle block in the spectral-repulsion inequality.}
\dkDeclareSymbol[Scalar sine of a principal angle. For a unit principal trial vector v and exact wanted projection P_ex, sin(theta) = ||(I - P_ex)v||. In Figure 1 it is the length of the pink leg S(1).]
  {sintheta}{\sin\theta}{sine}{Sine of an individual principal angle: the length of the complementary component of a unit principal trial vector.}
\dkDeclareSymbol[Positive sine operator sin(Theta_0) = sqrt(I_F - E_0* P_ex E_0): F -> F. It equals |S| for S = (I - P_ex)E_0: F -> E. Its eigenvalues in finite dimensions are the principal sines; the sine theorem bounds its ideal norm by N(R)/delta.]
  {sinTheta0}{\sin\Theta_0}{sineoperator}{Positive sine operator on trial coordinates; the modulus of the rectangular complementary error.}
\dkDeclareSymbol[Rectangular complementary error S = (I - P_ex) E_0: F -> E. It records how much each trial vector leaks outside the exact wanted subspace. Its modulus |S| = sqrt(S* S) is sin(Theta_0): F -> F. In rank one S(1) = P_ex_perp v.]
  {S}{S}{sineblock}{Rectangular complementary error: the trial frame projected onto the unwanted exact subspace.}
\dkDeclareSymbol[Complementary-coordinate map X = F_1* E_0: F -> G, with S = F_1 X. It satisfies Lambda_1 X - X A_0 = F_1* R on domain(A_0). Isometric multiplication gives N(X) = N(S); the separate modulus identity abs(S) = sin(Theta_0) gives N(S) = N(sin(Theta_0)), whenever these norms are finite.]
  {X}{X}{sylvester}{Solution of the Sylvester equation; for the sine bound, the trial frame in unwanted exact coordinates.}
\dkDeclareSymbol{alphaAngle}{\alpha}{anglemax}{Largest principal angle.}
\dkDeclareSymbol[Scalar c = cos(theta) in the counterexample. In the general two-projection decomposition, the same role denotes a positive cosine operator C; the cosine and sine blocks commute and satisfy C squared + S squared = I.]
  {CosAngle}{c}{anglecos}{Cosine of a principal angle or the corresponding positive cosine block.}
\dkDeclareSymbol[Scalar s = sin(theta) in the counterexample. In the general two-projection decomposition, the same role denotes a positive sine operator S; it commutes with the positive cosine block C and satisfies C squared + S squared = I.]
  {SinAngle}{s}{sineoperator}{Sine of a principal angle or the corresponding positive sine block.}
\dkDeclareSymbol[Positive small-angle parameter in the local subspace comparison. It controls the allowed angular deviation; it is an angle parameter rather than a perturbation-operator norm.]
  {EpsilonAngle}{\varepsilon}{anglemax}{Angle parameter controlling the maximum subspace deviation.}
\dkDeclareSymbol{PhiAngle}{\varphi}{angleoperator}{Auxiliary angle or angle operator in the two-subspace comparison.}
\dkDeclareSymbol{Bblock}{B}{sineblock}{Off-diagonal block in the two-subspace decomposition.}
\dkDeclareSymbol{Cblock}{C}{sylvester}{Complementary or Sylvester block in the two-subspace decomposition.}
\dkDeclareSymbol{Ablock}{A}{anglecos}{Positive cosine-related block in the two-subspace decomposition.}
% Residual/error-displacement family.
\dkDeclareSymbol[Residual vector r = Ahat v - rho v = R(1) for the rank-one frame E_0(x) = x v. Figure 1 draws it from rho v to Ahat v. For unit v and the Rayleigh value rho = <v, Ahat v>, <v, r> = 0; a general trial scalar need not give orthogonality.]
  {r}{r}{residvector}{Residual vector of an approximate eigenpair.}
\dkDeclareSymbol[Bounded residual R: F -> E, agreeing with Ahat E_0 - E_0 A_0 on domain(A_0). The domain condition requires E_0(domain(A_0)) inside domain(Ahat). R = 0 means that E_0 intertwines A_0 with Ahat; its norm measures the defect.]
  {R}{R}{resid}{Bounded residual of the trial eigenproblem, agreeing with the operator difference on the trial domain.}
\dkDeclareSymbol[Full displacement Delta_W = I - W on the ambient space. In the failed Ky Fan comparison W = W_(pi/4); Delta_W has singular values sqrt(2), sqrt(2), 0, 0. Multiplying by a principal-plane projection restricts its input to that plane.]
  {Kdisplacement}{\Delta_W}{displacement}{Full displacement of a transporter from the identity.}
\dkDeclareSymbol[Squared full displacement Dcal(W) = (I-W)* (I-W), a positive operator. Proposition 4.3 compares its unitarily invariant norm for a direct rotation D with that for every feasible transporter W.]
  {SquaredDisplacement}{\mathcal D}{displacement}{Squared full displacement as a function of the transporter.}
% Matching proof: the gradient and its spectral level are not a space or gap.
\dkDeclareSymbol[Skew-symmetric spectral gradient G of Phi_p. It commutes with W and represents first variations by d/dt Phi_p(W exp(tB)) at t = 0 equal to -tr(G B) for skew-symmetric B. Feasible stationarity gives J_src G J_src = -G.]
  {SpectralGradient}{G}{displacement}{Skew-symmetric spectral gradient of the displacement power sum in the matching proof.}
\dkDeclareSymbol[Positive spectral level g of the skew-symmetric gradient G: on this invariant subspace, -G squared = g squared I. Consequently G/g squares to -I and supplies the complex structure used in the real matching proof.]
  {GradientLevel}{g}{displacement}{Positive gradient level; the gradient squared is minus this level squared times the identity.}
\dkDeclareSymbol[Complex structure I_g = G/g on the positive gradient-level space V_g. It is a real isometry with square -I. Multiplication by a+ib means a times identity plus b times I_g; source reflection J_src reverses I_g and acts as conjugation.]
  {GradientComplexStructure}{\mathcal I_g}{displacement}{Complex multiplication constructed from a positive gradient level.}
\dkDeclareSymbol[Real skew-symmetric matrix B, with B transpose = -B, giving the orthogonal variation W exp(tB). If B commutes with the source projection P, this curve preserves the constraint W(U) = V.]
  {SkewDirection}{B}{ambient}{Skew-symmetric direction for a first variation of an orthogonal transporter.}
% Spectral separation family: gap and original gap are NOT identical.
\dkDeclareSymbol[Positive spectral separation in the stated configuration. In the sine theorem it separates the trial operator from the exact complementary block; Theorem 8.1 uses it for the ordered reference gap, and Theorem 8.2 for the perturbed central-band separation. Figure 1 uses delta = abs(rho - 0.7).]
  {delta}{\delta}{gap}{Positive spectral separation in the stated gap configuration.}
\dkDeclareSymbol{g}{g}{gaporiginal}{Original spectral gap of the reference operator; distinct from the certified sine-theta separation.}
\dkDeclareSymbol{Delta}{\Delta}{gapstat}{Population or statistical spectral gap.}
\dkDeclareSymbol[Spectral endpoint alpha. In the interval/exterior configuration, one spectrum lies between beta and alpha, the other at or below beta - delta or at or above alpha + delta. Theorem 8.2 places the selected perturbed block in the central interval. In Theorem 8.1, alpha bounds the lower reference block and alpha + delta bounds the upper block from below.]
  {alphaEndpoint}{\alpha}{gapinterval}{Upper endpoint of the central band, or lower-block cut in the ordered gap.}
\dkDeclareSymbol[Lower endpoint beta of the finite interval from beta to alpha in the gap hypothesis. One spectrum lies in that interval; the other lies at or below beta - delta or at or above alpha + delta.]
  {betaEndpoint}{\beta}{gapinterval}{Lower endpoint of the finite interval in the interval/exterior gap hypothesis.}
% Spectral/projector and auxiliary matrix objects, context-scoped.
\dkDeclareSymbol{P}{P}{fg}{Orthogonal projection; its source or exact-target role is specified locally.}% Neutral historical Manim key; prefer Pexact/Psource.
\dkDeclareSymbol{Q}{Q}{fg}{Orthogonal projection; its role is specified locally.}% Neutral historical Manim key; prefer Qtarget.
\dkDeclareSymbol[Two-projection intertwiner T = QP + (I - Q)(I - P), satisfying TP = QT and T* T = I - (P - Q) squared. For ||P - Q|| < 1 it is invertible; its unitary polar factor T |T| inverse is the canonical direct rotation.]
  {IntertwinerT}{T}{intertwiner}{Two-projection intertwiner; its unitary polar factor is the acute direct rotation.}
\dkDeclareSymbol[Planar quarter-turn J_0 = ((0, -1), (1, 0)), with J_0 squared = -I_2 and J_0 transpose = -J_0. It defines the rotation C = cos(2 theta) I_2 + sin(2 theta) J_0 used to construct the real competitor.]
  {Jturn}{J_0}{turn}{Planar quarter-turn, with square minus the identity.}
\dkDeclareSymbol[Partial isometry V_pol in the polar decomposition S = V_pol |S|, where |S| = sqrt(S* S). It is isometric on the orthogonal complement of kernel(S) and zero on kernel(S). For an invertible operator its polar factor is unitary.]
  {PolarU}{V_{\mathrm{pol}}}{intertwiner}{Partial isometry in the polar decomposition; unitary when the operator is invertible.}
\dkDeclareSymbol[Planar rotation C = cos(2 theta) I_2 + sin(2 theta) J_0 through twice the principal angle. The competitor W_theta is orthogonally conjugate to C direct-sum I_2, so it rotates one two-plane and fixes the other.]
  {Cturn}{C}{turn}{Planar rotation through twice the principal angle in the counterexample.}
\dkDeclareSymbol[Orthogonal mixing matrix H_mix = (1/sqrt(2)) times the block matrix ((I_2, I_2), (I_2, -I_2)). It constructs W_theta = H_mix (C direct-sum I_2) H_mix transpose.]
  {Hmix}{H_{\mathrm{mix}}}{mix}{Orthogonal mixing matrix in the counterexample.}
\dkDeclareSymbol[Source reflection J_src = 2P - I: it acts as +I on the source subspace and -I on its orthogonal complement, so its square is I. In the matching proof stationarity gives J_src G J_src = -G for the skew-symmetric spectral gradient G.]
  {Jreflect}{J_{\mathrm{src}}}{trialprojection}{Reflection equal to twice the source projection minus the identity.}
\dkDeclareSymbol[Skew angle generator J_ang for a direct rotation D = exp(J_ang Theta). On each principal plane it is the oriented quarter-turn; it vanishes on the common fixed directions. It commutes with the full angle operator Theta.]
  {AngleGenerator}{J_{\mathrm{ang}}}{turn}{Skew angle generator in the direct-rotation result summaries.}
\dkDeclareSymbol{Dblock}{D}{mix}{Auxiliary diagonal or block matrix in the counterexample.}
\dkDeclareSymbol[Orthogonal projection onto the first principal plane span(e_0, e_3) in the four-dimensional counterexample. For Delta_W = I - W_(pi/4), the Ky Fan 2-norm of Delta_W Omega_1 equals 2.]
  {OmegaOne}{\Omega_1}{turn}{First principal-plane projection in the failed Ky Fan comparison.}
\dkDeclareSymbol[Orthogonal projection onto the second principal plane span(e_1, e_2) in the four-dimensional counterexample. For Delta_W = I - W_(pi/4), the Ky Fan 2-norm of Delta_W Omega_2 equals 2, although the trace norm of Delta_W is only 2 sqrt(2).]
  {OmegaTwo}{\Omega_2}{intertwiner}{Second principal-plane projection in the failed Ky Fan comparison.}
\dkDeclareSymbol{Dplanar}{D}{wantedrotation}{Planar direct-rotation block.}
\dkDeclareSymbol{Hreflection}{H}{trialcompetitor}{Planar reflection carrying the source line to the target.}
% Unitarily invariant gauge family; N, M are NOT a norm versus its input.
\dkDeclareSymbol[Unitarily invariant norm N, or a symmetric operator-ideal norm on its domain of finiteness. It satisfies N(U T V) = N(T) for unitaries and N(B T) <= ||B||op N(T). In the sine theorem it is normalized by N(u v*) = ||u|| ||v|| and has Ky Fan dominance where both norms are finite. Both R and sin(Theta_0) must belong to its ideal.]
  {NormN}{N}{norm}{Unitarily invariant norm on its domain of finiteness; in the sine theorem, normalized by $N(uv^*)=\|u\|\|v\|$.}
\dkDeclareSymbol[Sine-theta separation G_delta(B,C): finite interval/exterior separation in either spectral order, or ordered separation B <= a I and C >= (a+delta) I, or the reverse. The ordered alternatives allow both spectra to be unbounded.]
  {SineGap}{\mathcal G_\delta}{gap}{Interval/exterior separation or either ordered separation, as in the sine-theta theorem.}
\dkDeclareSymbol[Auxiliary unitarily invariant norm M witnessing the Q-norm property N(T) squared = M(T* T). For the Frobenius norm N, one can take M to be the trace norm.]
  {NormM}{M}{normaux}{Auxiliary unitarily invariant norm used to define a Q-norm through the squared modulus.}
\dkDeclareSymbol[Q-norm property: a unitarily invariant norm N is a Q-norm when N(T) squared = M(T* T) for some unitarily invariant norm M. The Frobenius norm is an example, with M the trace norm; the repaired full-displacement theorem uses this norm class.]
  {NormQ}{Q}{normaux}{Q-norm: its square is a unitarily invariant norm of the squared modulus.}
\dkDeclareSymbol[Singleton trace cost s_i = min(4 sin(theta_i/2), 2) for angle label i. It transports one real principal plane independently, choosing a rotation of cost 4 sin(theta_i/2) or a reflection of cost 2.]
  {SingletonCost}{s}{norm}{Minimum trace displacement for one real principal plane: choose a rotation or a reflection.}
\dkDeclareSymbol[Paired trace cost e_ij = 4 sin((theta_i + theta_j)/2) for distinct angle labels i,j. It is attained by coupling their principal planes in a real four-dimensional block. The edge weight s_i + s_j - e_ij is the difference between the cost of two singletons and this paired cost.]
  {PairCost}{e}{norm}{Trace displacement cost of coupling two principal planes in a four-dimensional block.}
\dkDeclareSymbol[Scalar phase cost f_p(t) = (2 sin(t/2)) to the power p, for 0 <= t <= pi. A real planar rotation through t contributes two copies, so its displacement power cost is 2 f_p(t).]
  {MatchingPhaseCost}{f_p}{norm}{Displacement power cost for one scalar phase.}
\dkDeclareSymbol[Singleton power cost a_i(p) = min(2 f_p(theta_i), 2 to the power p). These are the rotation and reflection choices for one principal plane. At p=1 it equals the trace-cost weight s_i.]
  {SingletonPowerCost}{a}{norm}{Minimum displacement power cost for one angle.}
\dkDeclareSymbol[Paired power cost b_ij(p) = min over d <= tau <= h of 2(f_p(h+tau) + f_p(h-tau)), where h=(theta_i+theta_j)/2 and d=abs(theta_i-theta_j)/2. This is the positive-determinant pair minimum; unrestricted block optimization also compares two singletons. At p=1 it equals e_ij.]
  {PairPowerCost}{b}{norm}{Positive-determinant displacement power cost for a pair of angles.}
\dkDeclareSymbol[Matching M on the finite general graph of positive principal-angle labels. Each label belongs to at most one edge; unmatched labels use singleton transporters. Minimizing paired and singleton costs is equivalent to maximum-weight matching with edge weights g_ij = s_i + s_j - e_ij. The empty matching is allowed.]
  {Matching}{M}{normaux}{Matching of positive principal-angle labels: each label is used in at most one pair.}
\dkDeclareSymbol[Displacement power sum Phi_p(W) = sum_j s_j(I - W)^p. The matching proof first minimizes it for 1 < p < 2, then lets p decrease to 1, where it becomes the trace norm ||I - W||_1.]
  {PowerCost}{\Phi}{norm}{Sum of the p-th powers of displacement singular values.}
% Auxiliary matching-proof notation: neutral ink, with the same hover registry.
\dkDeclareSymbol[Off-diagonal gradient block G_01 mapping the source complement K perpendicular into K. Stationarity makes the diagonal blocks of G zero, so G = ((0,G_01),(-G_01 transpose,0)) in the source/complement decomposition.]
  {GradientOffDiagonal}{G_{01}}{fg}{Complement-to-source block of the skew-symmetric gradient.}
\dkDeclareSymbol[Orthogonal group O(n): real n by n matrices W satisfying W transpose times W = I. Its elements preserve Euclidean inner products and lengths. The notation O(2) refers to the same group in dimension two.]
  {OrthogonalGroup}{O(n)}{fg}{Group of real orthogonal matrices.}
\dkDeclareSymbol[Feasible transporters F = {W in O(n) : W(K) = L}. Equivalently W P_K W transpose = P_L. This nonempty compact set is the optimization domain throughout the matching proof.]
  {FeasibleTransport}{\mathcal F}{fg}{Orthogonal maps carrying the source subspace onto the target.}
\dkDeclareSymbol[Singular value s_j(T), the nonnegative square root of an eigenvalue of T transpose times T, counted with multiplicity. Summing s_j(I-W) gives the trace displacement; summing its p-th powers gives Phi_p(W). These singular values differ from the singleton weights s_i.]
  {TransportSingularValues}{s_j}{fg}{A displacement singular value in the matching proof.}
\dkDeclareSymbol[Scalar singular-value magnitude s between zero and two for the displacement I-W of an orthogonal matrix. Uniform convergence of s to the power p to s on this interval permits the trace-norm limit. This variable differs from the rotation parameter s used earlier in the proof.]
  {DisplacementSize}{s}{fg}{Scalar displacement singular-value magnitude in the uniform limit.}
\dkDeclareSymbol[All-singleton cost C_0 = sum_i s_i, before any two angle labels are paired. A matching has cost C_0 minus the sum of its selected edge weights.]
  {SingletonBaseline}{C_0}{fg}{Trace displacement when every angle uses its singleton transporter.}
\dkDeclareSymbol[Matching power cost C_p(M): sum of b_ij(p) over selected pairs plus a_i(p) over unmatched labels. Its minimum equals the minimum of Phi_p over feasible transporters. At p=1 it is C(M).]
  {MatchingPowerCost}{C_p}{fg}{Total paired and singleton power cost of a matching.}
\dkDeclareSymbol[Invariant orthogonal block V_a preserved by the source projection, target projection, and minimizing transporter. The displacement cost adds over these blocks, and their principal-angle lists combine into the original list.]
  {MatchingBlock}{V_a}{fg}{Simultaneous invariant block in the matching proof.}
\dkDeclareSymbol[Gradient-level space V_g: the eigenspace of -G squared for a positive eigenvalue g squared. It is preserved by G, W, and the source reflection J_src. Multiplication by i on it is G/g.]
  {GradientLevelSpace}{V_g}{fg}{Positive gradient-level eigenspace.}
\dkDeclareSymbol[Eigenphase varphi of W on V_g viewed as a complex space with multiplication by i = G/g. It lies between zero and pi and satisfies f_p prime(varphi) = g. There are at most two such phases.]
  {GradientEigenphase}{\varphi}{fg}{One of the at most two phases at a positive gradient level.}
\dkDeclareSymbol[Orthogonal projection E onto the first phase space of W on V_g. Here W = exp(i varphi_1) E + exp(i varphi_2) (I-E). In coordinates where the source reflection is conjugation, E = X + iY.]
  {GradientPhaseProjection}{E}{fg}{Projection onto the first phase space in a gradient-level space.}
\dkDeclareSymbol[Real part X of the phase projection E = X + iY. It is a real symmetric matrix with spectrum in the closed interval from zero to one, satisfying X squared - Y squared = X and XY + YX = Y. This X differs from the source-coordinate and Sylvester matrices elsewhere in the paper.]
  {ProjectionRealPart}{X}{fg}{Real symmetric part of the phase projection.}
\dkDeclareSymbol[Imaginary part Y of the phase projection E = X + iY. It is a real skew-symmetric matrix. It maps the lambda eigenspace of X to the 1-lambda eigenspace and constructs the small invariant planes.]
  {ProjectionImaginaryPart}{Y}{fg}{Real skew-symmetric imaginary part of the phase projection.}
\dkDeclareSymbol[Eigenvalue lambda in the closed interval from zero to one of the real part X of the phase projection. It is local to the block construction and differs from the original operator's spectral parameter.]
  {ProjectionEigenvalue}{\lambda}{fg}{Eigenvalue of the real part of the phase projection.}
\dkDeclareSymbol[Plane coupling eta = sqrt(lambda (1-lambda)) for 0 < lambda < 1. If u is a unit lambda eigenvector of X, then v = Yu/eta completes its invariant real plane.]
  {ProjectionCoupling}{\eta}{fg}{Normalization used to construct an invariant real plane.}
\dkDeclareSymbol[Pair mean angle h = (theta_i + theta_j)/2. Together with d = (theta_i - theta_j)/2 after ordering theta_i >= theta_j, it parametrizes the two-angle block cost.]
  {PairMeanAngle}{h}{fg}{Mean of the two ordered principal angles.}
\dkDeclareSymbol[Pair half-difference d = (theta_i - theta_j)/2 after ordering theta_i >= theta_j. It satisfies 0 <= d <= h and h+d <= pi/2.]
  {PairHalfDifference}{d}{fg}{Half the difference of the two ordered principal angles.}
\dkDeclareSymbol[Scalar pair parameter tau in the closed interval from d to h. The paired power cost is the minimum of 2(f_p(h+tau) + f_p(h-tau)) over this interval. At p=1 a minimum is attained at tau=h.]
  {PairPhaseParameter}{\tau}{fg}{Scalar parameter in the two-angle power-cost minimum.}
\dkDeclareSymbol[Attaining rotation angle gamma = arccos(cos(tau)/cos(d)). Choosing U_0 = U_1 = R_gamma realizes phase parameters x=h and y=tau in the four-dimensional block.]
  {PairAttainingAngle}{\gamma}{fg}{Rotation angle realizing a chosen scalar pair parameter.}
\dkDeclareSymbol[Four-dimensional direct rotation D = ((C,-S),(S,C)) in source and complementary principal coordinates, with C the diagonal principal cosines and S the diagonal principal sines. Every transporter on this block is D diag(U_0,U_1).]
  {PairDirectRotation}{D}{fg}{Direct rotation on the two principal planes.}
\dkDeclareSymbol[Diagonal cosine block C = diag(cos(theta_i), cos(theta_j)) in the four-dimensional direct rotation D. This matrix differs from the scalar displacement cost C(M).]
  {PairCosineBlock}{C}{fg}{Diagonal principal-cosine block.}
\dkDeclareSymbol[Diagonal sine block S = diag(sin(theta_i), sin(theta_j)) in the four-dimensional direct rotation D. This matrix is local to the matching proof.]
  {PairSineBlock}{S}{fg}{Diagonal principal-sine block.}
\dkDeclareSymbol[Free planar orthogonal matrices U_0 and U_1 acting on the two-dimensional source and its two-dimensional complement, respectively. Multiplying D by diag(U_0,U_1) parametrizes every transporter on this four-dimensional block.]
  {PairFreeOrthogonal}{U}{fg}{Free source or complementary planar orthogonal factor.}
\dkDeclareSymbol[Phase parameter x in the closed interval from h to pi-h for a positive-determinant four-dimensional transporter. Its eigenvalues are exp(plus or minus i(x+y)) and exp(plus or minus i(x-y)).]
  {PairPhaseX}{x}{fg}{First phase parameter of a four-dimensional transporter.}
\dkDeclareSymbol[Phase parameter y in the closed interval from d to pi-d for a positive-determinant four-dimensional transporter. Together with x it gives rotational phases x+y and x-y, folded into the closed interval from zero to pi for the displacement cost.]
  {PairPhaseY}{y}{fg}{Second phase parameter of a four-dimensional transporter.}
\dkDeclareSymbol[Scalar a = cos(x) in the four-dimensional phase calculation. Together with b = cos(y), it expresses the trace as 4ab and determines the reciprocal characteristic polynomial.]
  {PairPhaseCosineX}{a}{fg}{Cosine of the first phase parameter in the trace calculation.}
\dkDeclareSymbol[Scalar b = cos(y) in the four-dimensional phase calculation, paired with a = cos(x). These scalars are local to the characteristic-polynomial calculation.]
  {PairPhaseCosineY}{b}{fg}{Cosine of the second phase parameter in the trace calculation.}
\dkDeclareSymbol[Rotation parameter r: the two free orthogonal factors use planar rotation angles r+s and r-s. It appears in cos(y) = cos(d) cos(r) for two rotations.]
  {PairRotationR}{r}{fg}{Half the sum of the two free planar rotation angles.}
\dkDeclareSymbol[Rotation parameter s: the two free orthogonal factors use planar rotation angles r+s and r-s. It appears in cos(x) = cos(h) cos(s) for two rotations. This scalar differs from the singular values and singleton costs.]
  {PairRotationS}{s}{fg}{Half the difference of the two free planar rotation angles.}
\dkDeclareSymbol[Planar reflection J_2 = diag(1,-1). Every two-dimensional orthogonal matrix of determinant -1 can be written as a planar rotation times J_2.]
  {PairPlanarReflection}{J_2}{fg}{Fixed planar reflection used to parametrize determinant -1 factors.}
\dkDeclareSymbol[Real planar rotation R_t = ((cos(t),-sin(t)),(sin(t),cos(t))). Its two displacement singular values are both 2 times abs(sin(t/2)). In the matching proof a rotation through t in the closed interval from zero to pi contributes 2 f_p(t).]
  {MatchingPlanarRotation}{R_t}{fg}{Planar rotation through the indicated angle.}
% Statistical application / operator provenance.
\dkDeclareSymbol{Sigma}{\Sigma}{old}{Population covariance matrix.}
\dkDeclareSymbol{SigmaHat}{\widehat\Sigma}{current}{Estimated or perturbed covariance matrix.}
\dkDeclareSymbol{YWSV}{V}{wanted}{Orthonormal coordinates for the population eigenspace in the statistical application.}
\dkDeclareSymbol{YWSVhat}{\widehat V}{trialframe}{Estimated eigenspace frame in the statistical application.}

% Figure-specific objects retain their own meanings when bare letters overlap.
\dkDeclareSymbol[Signed rotation angle phi of the displayed map on its invariant real two-plane. At theta = pi/4 the direct rotation turns its two principal planes by opposite signed angles of magnitude pi/4; the competitor turns its moving plane by pi/2 and fixes its other plane.]
  {PlaneRotationAngle}{\varphi}{sine}{Signed rotation angle of the orthogonal map on the displayed invariant plane.}
\dkDeclareSymbol[Unit vector q in the displayed invariant plane. The map T is the direct rotation D or the competitor W in that panel; Tq is the image and (I - T)q = q - Tq is the displacement.]
  {FigureInput}{q}{ambient}{Unit input to the orthogonal map shown in this invariant-plane panel.}
\dkDeclareSymbol[Image Tq of the panel unit input q under T = D or W. Both q and Tq lie on the unit circle; the green displacement arrow joins Tq to q, so it represents q - Tq.]
  {FigureImage}{Tq}{wantedvector}{Image of the unit input under the orthogonal map shown in this panel.}
\dkDeclareSymbol[Displacement (I - T)q = q - Tq, drawn from the image Tq to the unit input q. On a plane rotated through phi its length is 2 |sin(phi/2)|; it vanishes on a fixed plane.]
  {DisplacementVector}{(I-T)q}{displacement}{Difference between a unit input and its image under the displayed orthogonal map.}
\dkDeclareSymbol[First principal plane Pi_1 = span(e_0, e_3) for the source-target pair in the real four-dimensional counterexample. At theta = pi/4 the direct rotation acts on this plane by a quarter of pi, with displacement singular value sqrt(2 - sqrt(2)).]
  {CounterPrincipalOne}{\Pi_1}{fg}{First principal plane in the real counterexample: the span of e0 and e3.}
\dkDeclareSymbol[Second principal plane Pi_2 = span(e_1, e_2) for the source-target pair in the real four-dimensional counterexample. The direct rotation acts by the same angle magnitude as on Pi_1, with the opposite sign in the displayed plane coordinates.]
  {CounterPrincipalTwo}{\Pi_2}{fg}{Second principal plane in the real counterexample: the span of e1 and e2.}
\dkDeclareSymbol[Moving plane M = range(W - I) = span(e_0 + e_2, e_1 + e_3) for the competitor. W rotates it by 2 theta; at theta = pi/4 the two displacement singular values contributed by this plane are both sqrt(2).]
  {CounterMovingPlane}{M}{ambient}{Moving plane of the competing transporter: the range of W minus the identity.}
\dkDeclareSymbol[Fixed plane Z = kernel(W - I) = span(e_0 - e_2, e_1 - e_3) for the competitor. W acts as the identity here, so this plane contributes two zero displacement singular values.]
  {CounterFixedPlane}{Z}{ambient}{Fixed plane of the competing transporter: the kernel of W minus the identity.}
\dkDeclareSymbol[Direct-rotation displacement singular value a = 2 sin(pi/8) = sqrt(2 - sqrt(2)) at theta = pi/4. Each of its two principal planes contributes two copies, giving trace displacement 4a.]
  {CounterChordA}{a}{displacement}{Direct-rotation displacement singular value in the counterexample: sqrt(2 - sqrt(2)).}
\dkDeclareSymbol[Competitor displacement singular value b = 2 sin(pi/4) = sqrt(2) at theta = pi/4. Its moving plane contributes two copies; the fixed plane contributes two zeros, giving trace displacement 2b.]
  {CounterChordB}{b}{displacement}{Competing-transporter displacement singular value in the counterexample: square root of 2.}
\dkDeclareSymbol[Singular value s_j(I - T) of the full ambient displacement, ordered from largest to smallest. Each rotating real two-plane contributes two values 2 |sin(phi/2)|; fixed directions contribute zeros. Their sum is the trace cost.]
  {DisplacementSingular}{s_j}{displacement}{An individual singular value of the full displacement from the identity.}
\dkDeclareSymbol[Edge weight g_ij = s_i + s_j - e_ij, the difference between the cost of two singletons and the paired cost. A matching has total cost C_0 minus the sum of its edge weights, with C_0 = sum_i s_i. Maximizing this sum gives the minimum transport cost; edges of negative weight may be omitted.]
  {MatchingEdgeWeight}{g_{ij}}{current}{Edge weight: the cost of two singletons minus their paired cost.}
\dkDeclareSymbol[Total trace-norm displacement cost. For a transporter W, C(W) = ||I - W||_1, the sum of its displacement singular values. For a matching M, C(M) is the sum of its paired and unmatched singleton costs.]
  {TraceCost}{C}{norm}{Trace-norm displacement cost of the indicated orthogonal transporter or matching.}

% Spectral values have their own key, distinct from an operator block.
\dkDeclareSymbol[Spectral parameter lambda of a self-adjoint operator. In a finite-dimensional eigenvalue list lambda_j(A) it denotes an individual eigenvalue; in A = integral lambda dE(lambda), it is the real integration variable.]{SpectralValue}{\lambda}{wantedspectrum}{Real spectral parameter or individual eigenvalue of a self-adjoint operator.}
 in paper.tex)
% or pass \def\DKMonochromeMath{1} before \documentclass[11pt]{article}
% PDF version used for the manuscript and its symbol annotations.
\pdfvariable minorversion=7\relax

\usepackage[margin=1in]{geometry}
\usepackage{amsmath,amssymb,amsthm,mathtools}
\usepackage{booktabs,tabularx,array,longtable}
\usepackage[section]{placeins}
\usepackage{float}
\usepackage{microtype}
\usepackage{graphicx}
\usepackage{pgf} % Live vector figures with LaTeX symbol labels and tooltips
% Editing citations: switch to the venue's numeric style at submission.
\usepackage[authoryear,round]{natbib}
\usepackage{xurl}
\usepackage[colorlinks=true,allcolors=blue]{hyperref}
\usepackage[nameinlink]{cleveref}
\usepackage[T1]{fontenc}
\usepackage{lmodern} % Scalable LaTeX lettering, shared with the Manim slides
\usepackage{xcolor}
% Draft controls. Set these independently for review, print, or arXiv versions.
\newif\ifDKAnonymous
\DKAnonymousfalse
\ifdefined\DKAnonymousReview\DKAnonymoustrue\fi
\newif\ifDKColoredSymbols
\DKColoredSymbolstrue
\ifdefined\DKMonochromeMath\DKColoredSymbolsfalse\fi
% The colors retain the Manim glossary's light-theme mathematical roles.
\input{math_symbol_colors.tex}
\ifDKColoredSymbols\else\DKMathColorsfalse\fi

%\DeclareMathOperator{\DKRange}{range}
%\DeclareMathOperator{\DKKernel}{kernel}
%\DeclareMathOperator{\DKDom}{domain}
\DeclareMathOperator{\DKRange}{ran}
\DeclareMathOperator{\DKKernel}{ker}
\DeclareMathOperator{\DKDom}{dom}
\usepackage[most]{tcolorbox}
\usepackage{listings}
\usepackage{listingsutf8}
\usepackage{authblk}
\usepackage{emoji} % Requires LuaLaTeX compiler

\setlength{\emergencystretch}{3em}

% ArXiv-friendly Lean code: no shell escape, minted, or external highlighter.
\definecolor{leankeyword}{HTML}{243B6B}
\definecolor{leantactic}{HTML}{70477A}
\definecolor{leantype}{HTML}{1F6473}
\definecolor{leancomment}{HTML}{39734D}
\definecolor{leanstring}{HTML}{8A4B2A}
\definecolor{leannumber}{HTML}{777777}
\definecolor{leanbackground}{HTML}{F7F7F3}
\lstdefinelanguage{lean}{%
  sensitive=true,
  morecomment=[l]{--},
  morecomment=[s]{/-}{-/},
  morestring=[b]",
  morekeywords=[1]{import,open,namespace,section,end,noncomputable,variable,variables,
    universe,universes,set_option,attribute,axiom,opaque,theorem,lemma,def,abbrev,
    example,instance,structure,class,inductive,notation,local,deriving,extends,
    by,where,let,have,show,fun,match,with,if,then,else,do,return,in,forall},
  morekeywords=[2]{simp,simpa,rw,erw,dsimp,unfold,ring,ring_nf,nlinarith,linarith,
    positivity,norm_num,field_simp,omega,exact,refine,apply,assumption,intro,intros,
    ext,funext,cases,rcases,obtain,constructor,left,right,classical,by_contra,
    contrapose,convert,using,suffices,calc},
  morekeywords=[3]{Prop,Type,Nat,Int,Real,Fin,Set,Submodule,Matrix,EuclideanSpace,
    LinearMap,ContinuousLinearMap,LinearPMap,RCLike,IsTrialResidual,
    IsExactSpectralDecomposition,FormBoundedSylvesterGap,
    NormalizedSymmetricOperatorIdealFamily},
}
\lstdefinestyle{leanpaper}{%
  language=lean,
  columns=fullflexible,
  backgroundcolor=\color{leanbackground},
  basicstyle=\ttfamily\scriptsize,
  commentstyle=\color{leancomment}\itshape,
  keywordstyle=[1]\color{leankeyword}\bfseries,
  keywordstyle=[2]\color{leantactic},
  keywordstyle=[3]\color{leantype},
  stringstyle=\color{leanstring},
  numberstyle=\tiny\color{leannumber},
  breakatwhitespace=false,
  breaklines=true,
  captionpos=b,
  keepspaces=true,
  numbers=left,
  numbersep=6pt,
  showspaces=false,
  showstringspaces=false,
  showtabs=false,
  tabsize=2,
  aboveskip=0.8em,
  belowskip=0.8em,
  literate=
    {\\KField}{{$\Bbbk$}}1
    {∃}{{$\exists$}}1
    {∧}{{$\land$}}1
    {∀}{{$\forall$}}1
    {≠}{{$\neq$}}1
    {⊤}{{$\top$}}1
}
\crefname{lstlisting}{formalization}{formalizations}
\Crefname{lstlisting}{Formalization}{Formalizations}
\renewcommand{\lstlistingname}{Formalization}

\newtheorem{theorem}{Theorem}
\newtheorem{proposition}[theorem]{Proposition}
\newtheorem{lemma}[theorem]{Lemma}
\newtheorem{remark}[theorem]{Remark}
\newtheorem{corollary}[theorem]{Corollary}

\newcommand{\Lean}{\textsc{Lean}}
\newcommand{\DK}{Davis--Kahan}
\newcommand{\TauCeti}{Tau Ceti}
% Use \ie{} and \eg{}; supply surrounding punctuation at the call site.
\newcommand{\ie}{i.e.\@}
\newcommand{\eg}{e.g.\@}
\newcommand{\DKResultCount}{29}
\newcommand{\DKProvedCount}{28}
\newcommand{\DKRefutedCount}{1}

% Usage: \margincomment{Check this formula!}
%\newcommand{\jonsays}[1]{%
%  \marginpar{%
%    \raggedright\color{red}\tiny% Change to \footnotesize if text is too small
%    \emoji{person-raising-hand} \textbf{Comment:}\\#1%
%  }%
%}

\geometry{
  % Keeps your main text layout intact while expanding the side margin box
  marginparwidth=2.2cm,
  marginparsep=0.3cm
}

% Usage: \margincomment{Your comment that now wraps perfectly!}
\newcommand{\jonsays}[1]{%
  \marginpar{%
    \raggedright% Forces text to wrap beautifully without ugly hyphenation
    \color{red}\scriptsize% Using \scriptsize makes the text tighter
    \emoji{person-raising-hand}: \\#1%
  }%
}
%\newcommand{\jonsays}[1]{{\color{red}\small\textbf{[\emoji{person-raising-hand}: #1]}}}

% Usage: \margincomment{Your comment that now wraps perfectly!}
%\newcommand{\agentsays}[1]{%
%  \marginpar{%
%    \raggedright% Forces text to wrap beautifully without ugly hyphenation
%    \color{red}\scriptsize% Using \scriptsize makes the text tighter
%    \emoji{robot}: \\#1%
%  }%
%}
\newcommand{\agentsays}[1]{{\color{red}\small\textbf{[\emoji{robot}: #1]}}}

\newcommand{\jonrsp}[1]{{\color{red}\small\textbf{[\emoji{person-raising-hand}: #1]}}}

\title{Formalizing Davis--Kahan Perturbation Theory in Lean 4}
% Additional authors are expected; keep the confirmed ordering below stable.

\ifDKAnonymous
\author{Anonymous authors}
\else
\author[1]{Jonathan Crall}
\author[2]{Edward Wang}
\author[1]{Brian Hu}
\author[2]{Carey E. Priebe}

\affil[1]{%
  Kitware, Inc.\\
  \texttt{jon.crall@kitware.com},
  \texttt{brian.hu@kitware.com}
}

\affil[2]{%
  Johns Hopkins University\\
  \texttt{ewang43@jhu.edu},
  \texttt{cep@jhu.edu}
}
\fi

\date{28 October 2026}

\begin{document}
\maketitle

\begin{abstract}

We formalize the Davis--Kahan theory of spectral-subspace perturbations.
From Davis and Kahan's landmark 1970 paper we identify \DKResultCount{}
  results as the targets of the formalization: of these
  \DKProvedCount{} are proved in Lean~4 under their stated hypotheses, and the remaining
  statement --- as originally written --- is refuted.
The formal development includes the sine and tangent theorems, direct rotations,
  principal-angle operators, Sylvester separation, unitarily invariant
  operator-ideal norms, and unbounded self-adjoint operators.
We place these constructions in their mathematical lineage, from the geometry of
  principal angles through the theory of projections and the spectral methods
  underlying Davis and Kahan's estimates.


The refuted statement \citep[Proposition~4.4]{davis1970rotation} concerns
  minimization of full ambient displacement by direct rotation.
We construct real four-dimensional examples in which direct rotation fails
  to minimize the trace norm at every prescribed positive acute angle.
We locate the failed singular-value comparison and prove a corrected optimality
  statement for Q-norms, including the Frobenius norm.

The formalization builds on Mathlib and adapts the Spectra project's
  projection-valued-measure and Borel-calculus constructions.
Portions of the reusable operator theory developed here were subsequently
  upstreamed into Tau Ceti.
The resulting infrastructure also supports a theorem reducing trace-norm
  optimal real subspace transport to weighted matching on principal angles.
These outcomes show how formalizing the foundations of a classical mathematical
  theory can support validating results, identifying historical errors, and
  developing new formal research.

\end{abstract}

\noindent\textbf{Keywords:} Davis--Kahan theory; spectral perturbation;
formal verification; Lean 4; operator theory; direct rotation; principal angles.

\smallskip

\noindent\textbf{Mathematics Subject Classification (2020):}
\href{https://mathscinet.ams.org/msc/msc2020.html?t=68V20}{68V20},
\href{https://mathscinet.ams.org/msc/msc2020.html?t=47A55}{47A55},
\href{https://mathscinet.ams.org/msc/msc2020.html?t=15A42}{15A42}.

\ifDKAnonymous\else
\smallskip
\noindent\textbf{Corresponding author:} Jonathan Crall
(\texttt{jon.crall@kitware.com}).
\fi

\section{Introduction}
\label{sec:introduction}

An approximate eigenvector gives a candidate direction for the action of an
operator. If a unit vector $\sym{v}$ is paired with an approximate eigenvalue $\sym{mu}$,
the residual $\sym{A}\sym{v}-\sym{mu} \sym{v}$ measures how closely it satisfies the eigenvector
equation. A small residual can identify an accurate approximation to an
invariant subspace even when repeated or nearby eigenvalues make individual
eigenvectors difficult to distinguish. Davis--Kahan theory bounds the angles
between approximate and exact spectral subspaces in terms of a residual or
perturbation and a separation between the relevant spectra
\citep{davis1970rotation}.

In applications, one often works with a perturbed operator
  $\sym{Acurrent}=\sym{A}+\sym{H}$ while the reference operator $\sym{A}$ is
  unknown.
Taking an eigenspace of $\sym{Acurrent}$ as a trial subspace for $\sym{A}$, the norm of its
  residual relative to $\sym{A}$ is bounded by the norm of $\sym{H}$, allowing
  Davis--Kahan theory to bound the angle between the two subspaces without
  requiring $\sym{A}$ to be known or its residuals to be evaluated directly.

Our project began with a variant of this bound needed for
\ifDKAnonymous a spectral-estimation application\else the data kernel
perspective space (DKPS) \citep{helm2026query}\fi.

We extended the formalization to the results established in the original 1970
paper, including real and complex Hilbert spaces, unbounded self-adjoint
operators, and unitarily invariant operator-ideal norms. In particular, the
spectral-separation hypotheses allow half-infinite intervals
(e.g.\ $(-\infty,a]$). These settings require operator domains, spectral
calculus, and norm estimates beyond finite matrices.

We formalized the four headline theorems \citep[Section~2]{davis1970rotation} and every
named theorem, proposition, lemma, and corollary established in the 1970 paper.

This gives \DKResultCount{} targets, enumerated in
Appendix~\ref{app:dk-inventory}. We prove \DKProvedCount{} of them and refute
the printed statement \citep[Proposition~4.4]{davis1970rotation}.
That proposition asserts a full-displacement
optimality property for real direct rotations. We give a family of real
four-dimensional examples exhibiting trace-norm failure at every prescribed
positive acute angle and identify
the invalid singular-value comparison in its proof. The phase construction
behind the counterexample occurs over the complex field in the authors'
example \citep[Example~4.2]{davis1970rotation}. To the best of our knowledge,
its contradiction of the printed real, small-angle claim has not previously
been made explicit.
For the constructed pair, whose two principal angles both equal
$\sym{theta}$, the competing transporter attains the global trace-norm minimum
$4\sin\sym{theta}$ for the constructed subspaces. Section~\ref{sec:prop44}
compares these statements and repairs the optimality claim by restricting
the norm class to Q-norms (N is a Q-norm when $N(T)^2 = M(T^* T)$ for some
unitarily invariant norm $M$). This repair follows from the squared-displacement
comparison \citep[Proposition~4.3]{davis1970rotation};
earlier Q-norm minimization results appear in \citet{bolotnikov2001minimal}.

% TODO (lineage): Consolidate the scattered lineage records through the
% maintained dev/tauceti/dkps-upstream-contribution-census.md and the evidence
% index in AUTHORING_NOTES.md. Audit each attribution against an immutable
% donor commit/path, the accepted upstream commit/PR, and the corresponding
% declaration statements and proofs. Preserve Mathlib prerequisites, Spectra
% ancestry where applicable, and independent work as distinct relationships.
% Courant--Fischer/Weyl checked 2026-10-28: developed here in June 2026;
% August donor df78246f965881705c9301f96feddc270c98ceb1 predates the
% accepted Tau Ceti port 7ebf17ba04725c1a575ecf3d099ce0630cd949a0 (#13599).
% See COURANT_FISCHER_LINEAGE_AUDIT.md for the proof comparison and sources.

The proofs required us to formalize a significant amount of operator
theory, namely: principal angles and direct rotations, approximation numbers
and symmetric operator ideals, Borel spectral calculus, and Sylvester
estimates for unbounded operators. These developments build on Mathlib and
adapt Spectra's projection-valued-measure and Borel-calculus constructions
\citep{mathlib2020,bornemann2026spectra}.
These operator-theory foundations include our
Courant--Fischer and Weyl proofs, rectangular unitarily invariant seminorms,
and projection-valued measures, which were subsequently adapted into Tau
Ceti \citep{tauceti2026}. The measure constructions retain their earlier Spectra ancestry.
Section~\ref{sec:related} explains the mathematical history and formal
ancestry; Section~\ref{sec:foundations} describes the foundations used in
our proofs.

The same principal-angle and operator-ideal results support a further
optimization theorem. In Section~\ref{sec:example-transport}, we prove that
minimizing trace-norm displacement over all real orthogonal maps carrying
one subspace onto another reduces to weighted matching on principal angles.

\paragraph{Contributions.}
\begin{enumerate}
\item \textbf{Davis--Kahan formalization.} We prove 28 of the 29 source
results and refute the full-displacement assertion
\citep[Proposition~4.4]{davis1970rotation}. The sine-theta theorem covers real and
complex Hilbert spaces, unbounded self-adjoint operators, the source's
spectral-separation alternatives, and operator-ideal norms where finite
(\cref{sec:coverage,sec:sine-theta}).
\item \textbf{Formalized operator-theoretic foundations.} We formalized the
geometric and analytic results required by these theorems, adapted existing
formal mathematics, and contributed foundations later used in Tau Ceti
(\cref{sec:related-formal,sec:foundations}).
\item \textbf{Source correction.} We identify a false comparison
of Ky Fan sums in the printed claim \citep[Proposition~4.4]{davis1970rotation},
give a real counterexample at every positive acute angle,
and prove the Q-norm consequence of the valid squared-displacement theorem
(\cref{sec:prop44}).
\item \textbf{A reduction to weighted matching.} We reduce unrestricted
real trace-norm transport to a weighted-matching problem on principal angles
(\cref{sec:example-transport}).

\end{enumerate}

\section{Mathematical background}
\label{sec:related}

The Davis--Kahan angle estimates use three constructions. Principal angles
measure the error between subspaces; a residual measures how well a proposed
subspace represents the operator; and a Sylvester estimate bounds the first
quantity by the second under spectral separation. Operator ideals allow the
same argument to express several norms. We describe the development of these
constructions before stating the formalized sine-theta theorem.

Existing historical surveys cover these developments. Paige and Wei trace
the cosine--sine decomposition from Jordan's subspace geometry through the
principal ideas assembled by Davis and Kahan to Stewart's later completion
and explicit formulation \citep{paige1994history}.
Simon's historical survey of Kato's work in non-relativistic quantum mechanics
traces Sz.-Nagy's unitary construction for pairs of projections and Kato's
later algebraic formulation \citep[\S5]{simon2018kato}.
We use this account for the history of the projection transport described in
\cref{sec:related-projections}.
B\"ottcher and Spitkovsky survey pairs of projections and their operator algebras.
Their historical discussion relates Halmos's 1969 block representation to the
earlier formulations of Dixmier (1948) and Davis (1958), and to the related
projection constructions in Davis and Kahan's 1970 paper
\citep[Remark~1.3]{bottcher2010gentle}.
We use these accounts together with the primary
sources. Table~\ref{tab:historical-lineage} summarizes the mathematical
constructions used in our formalization.

\begin{table}[!htbp]
\centering
\small
\setlength{\tabcolsep}{5pt}
\renewcommand{\arraystretch}{1.18}
\begin{tabularx}{\linewidth}{@{}>{\raggedright\arraybackslash}p{0.26\linewidth}>{\raggedright\arraybackslash}X>{\raggedright\arraybackslash}p{0.28\linewidth}@{}}
\toprule
\textbf{Lineage} & \textbf{Mathematical role} & \textbf{Formal dependency}\\
\midrule
Jordan (1875); Davis (1958); Halmos (1969) & Principal angles, invariant subspace geometry,
  and pairs of projections & Angle operators; projection blocks\\
Rayleigh (1877); Ritz (1909); Weyl (1912) & Approximate eigenpairs,
  variational spectra, and eigenvalue stability & Trial coordinates;
  min--max and eigenvalue inequalities\\
Rellich (1937--1942); Sz.-Nagy (1947); Kato (1966) & Moving spectral projections and unitary
  transport between nearby subspaces & Spectral calculus; polar
  direct rotation\\
Sylvester (1884); Rosenblum (1956, 1969) & Spectral separation controls an
  off-diagonal operator & Domain-aware Sylvester estimates\\
Von Neumann (1937); Schatten (1950); Ky Fan (1951) & Comparison of spectral errors
  through singular-value gauges & Approximation numbers;
  symmetric operator ideals\\
Davis (1963, 1965); Davis--Kahan (1969, 1970) & Residual-to-angle bounds, tangent estimates,
  rotation extremality, and spectral branch selection & Source-level
  sine, tangent, and direct-rotation theorems\\
\bottomrule
\end{tabularx}
\caption{Mathematical sources used in the Davis--Kahan formalization.
Dates identify the publications discussed in the text.}
\label{tab:historical-lineage}
\end{table}
\FloatBarrier

\subsection{Principal angles}
\label{sec:related-classical}

Camille Jordan studied the geometry of linear subspaces in $n$ dimensions in
  1875 \citep{jordan1875essai}.

For two lines, a single acute angle describes their relative position.
For subspaces of higher dimension, one needs a sequence of \emph{principal
  angles}.
The successive principal vectors are selected under orthogonality constraints,
  so their angles are independent of the coordinates used to describe either
  subspace.
This separates the geometry of a subspace from the arbitrary choice of its
  orthonormal basis.

For finite-dimensional subspaces $\sym{Usource}$ and $\sym{Vtarget}$ of equal dimension, choose
isometries $\sym{Xframe}$ and $\sym{Yframe}$ with those ranges.
The singular values of the overlap matrix are the principal cosines:
\begin{equation}
  \sigma_j(\like{Xframe}{X^*}\sym{Yframe})=\cos\like{theta}{\theta_j},
  \qquad 0\leq\like{theta}{\theta_j}\leq\frac{\pi}{2}.
  \label{eq:history-principal-cosines}
\end{equation}

The corresponding principal vectors $\like{SourceVector}{x_j}\in \sym{Usource}$ and $\like{TargetVector}{y_j}\in \sym{Vtarget}$ satisfy
$\langle \like{SourceVector}{x_i},\like{TargetVector}{y_j}\rangle=\delta_{ij}\cos\like{theta}{\theta_j}$ in suitably oriented
coordinates.  For nonzero acute angles, each pair determines a real
two-dimensional plane in which the relative position is an ordinary
rotation.  An entire pair of subspaces is described by these planes,
together with their shared and orthogonal components.

The same geometric data appeared in statistics: Hotelling's 1936 canonical
  correlation analysis sought pairs of maximally correlated linear combinations
  of two sets of variables \citep{hotelling1936relations}.
After appropriate whitening, the canonical correlations are singular values of a
  cross-coordinate map.
In numerical linear algebra, Bj\"orck and Golub developed algorithms for
  computing principal angles and vectors in 1973 \citep{bjorck1973numerical}.
Their computational formulation makes the basis-invariant angles available from
  orthonormal matrix factorizations.

A further formulation uses the cosine--sine decomposition of a
partitioned orthogonal or unitary matrix.
The blocks of the decomposition contain compatible diagonal cosine and sine
  factors.
Davis and Kahan describe the geometry using orthogonal projections,
angle operators, and the direct rotation carrying one subspace onto the other
\citep[Section~3]{davis1970rotation}. Their work assembled the principal ideas
of the cosine--sine decomposition; Stewart completed and explicitly formulated
the matrix decomposition in 1977
\citep{stewart1977perturbation,paige1994history}.
Paige and Wei explain how this decomposition and direct rotation represent
the same principal-angle data
\citep{paige1994history}. Our formalization uses projections and functional
calculus for the Hilbert-space results and principal-plane coordinates for
the finite-dimensional counterexample.

\subsection{Eigenvalue approximation}
\label{sec:related-ritz-weyl}

Jordan's angles answered how to compare subspaces once they are given.
An independent question came from mechanics and numerical analysis:
how can one approximate an eigenvalue or eigenmode when the underlying
differential equation cannot be solved exactly?  In the nineteenth
century Rayleigh estimated vibration frequencies using trial shapes
and their energy ratios \citep{rayleigh1877sound}.  Ritz developed a systematic finite-dimensional
variational approximation in the work published in 1909
\citep{ritz1909variations,leissa2005historical}.  The method selects
a trial space and studies the action of the operator through that
space. A trial space is a chosen space of candidate approximate solutions.
For an eigenvalue problem, one searches for an approximate eigenvector among
its vectors. This reduces an infinite-dimensional problem to a smaller
problem whose residual can then be assessed.

For a Hermitian matrix $\sym{A}$, the Rayleigh quotient
$\langle \sym{A}\like{v}{x},\like{v}{x}\rangle/\langle \like{v}{x},\like{v}{x}\rangle$ expresses the scalar action
of $\sym{A}$ along a trial vector. A normalized trial vector $\sym{v}$ gives the
Rayleigh--Ritz approximate eigenvalue $\sym{mu}=\langle \sym{A}\sym{v},\sym{v}\rangle$ and the
residual $\sym{A}\sym{v}-\sym{mu} \sym{v}$. The Courant--Fischer min--max
principle organizes the eigenvalues through extremal quotients over
subspaces \citep{bhatia1997matrix}.  Weyl's eigenvalue comparison
inequalities then control the motion of ordered eigenvalues under a
Hermitian perturbation \citep{weyl1912eigenwerte,bhatia1997matrix}.  In modern notation,
if $\sym{A}$ and $\sym{Acurrent}=\sym{A}+\sym{H}$ are Hermitian $n$-by-$n$ matrices and their eigenvalues
are arranged in the same order, then
\begin{equation}
  \max_j |\like{SpectralValue}{\lambda_j}(\sym{A})-\like{SpectralValue}{\lambda_j}(\sym{Acurrent})|\leq\|\sym{H}\|_{\mathrm{op}}.
  \label{eq:history-weyl}
\end{equation}
This is an estimate on \emph{eigenvalues}; it does not tell us which
individual eigenvector of $\sym{A}$ corresponds to which eigenvector of $\sym{Acurrent}$.

The limitation already appears in two dimensions.  For any $\varepsilon>0$,
put
\begin{equation}
 \like{A}{A_\varepsilon}=\begin{pmatrix}0&0\\0&\varepsilon\end{pmatrix},
 \qquad
 \like{Acurrent}{\widehat A_\varepsilon}=\begin{pmatrix}\varepsilon&0\\0&0\end{pmatrix}.
 \label{eq:history-swap-example}
\end{equation}
The matrices have identical ordered eigenvalues and
$\|\like{A}{A_\varepsilon}-\like{Acurrent}{\widehat A_\varepsilon}\|_{\mathrm{op}}=\varepsilon$,
yet their lowest-eigenvalue lines are perpendicular. The gap between the
two eigenvalues of each matrix is $\varepsilon$, equal to the perturbation
norm. An eigenspace estimate therefore needs to compare the perturbation
size with the spectral separation.  For a cluster with repeated eigenvalues, an
entire invariant subspace is likewise more natural than a selected
orthonormal eigenbasis.

Rellich developed a rigorous perturbation theory of self-adjoint operators
in papers published from 1937 to 1942. Kato's 1966 book gives a systematic
account \citep{simon2018kato,kato1966perturbation}. The Davis--Kahan residual
bounds address a related approximation question: given a trial space and its
residual, spectral separation bounds its angle from an exact spectral
subspace. The estimate can therefore assess the accuracy of a
Rayleigh--Ritz approximation.

\subsection{Two orthogonal projections}
\label{sec:related-two-projections}

For a closed subspace of a Hilbert space $\mathcal H$, the orthogonal
projection remains a bounded operator even when the subspace is
infinite dimensional.  Denote the projections onto two closed
subspaces by $\sym{Psource}$ and $\sym{Qtarget}$. Compressing an operator to $\DKRange(\sym{Psource})$ means
restricting its input to that subspace and projecting its output back with
$\sym{Psource}$. The positive compression of $\sym{Qtarget}$ is therefore $\sym{Psource}\sym{Qtarget}\sym{Psource}$ restricted to
$\DKRange(\sym{Psource})$. It is the operator version of the squared overlap matrix: in finite dimensions its eigenvalues are
$\cos^2\like{theta}{\theta_j}$.  In Hilbert space it need not have a discrete
spectrum.  Functional calculus nevertheless defines its positive
square root and the associated cosine and sine of an angle operator.

Write $\sym{Theta0}$ for this directed angle operator on
$\DKRange(\sym{Psource})$, and $\sym{Theta1}$ for the corresponding
operator for the complementary pair $I-\sym{Psource},I-\sym{Qtarget}$.
The full angle operator is
\[
 \sym{Theta}=\operatorname{diag}(\sym{Theta0},\sym{Theta1})
 \quad\text{on }\DKRange(\sym{Psource})\oplus\DKKernel(\sym{Psource}).
\]
For $c>0$, the comparisons $\sym{Theta}\le c$ and $\sym{Theta}<c$ mean that its
operator norm is at most, or strictly less than, $c$, respectively.
In finite dimensions this norm is the largest eigenvalue of either directed
angle block, comparing the source with the target and their orthogonal
complements.
This includes right angles from unmatched directions: if a source vector is
orthogonal to the target, its cosine compression is zero and its directed
angle is $\pi/2$. The compression construction applies even when the subspace
dimensions differ. Target directions orthogonal to the source appear as
right angles in the complementary block. Later, isometric frames express these angle blocks on
coordinate Hilbert spaces.

The relative position of two closed subspaces was studied by Dixmier
in 1948 and by Davis in his 1958 paper \emph{Separation of Two Linear
Subspaces} \citep{dixmier1948position,davis1958separation}.  Davis's
closeness and separation operators place the cosine/sine data in an
operator-theoretic setting.  Halmos's 1969 \emph{Two Subspaces} gives a
block representation of the two projections
\citep{halmos1969two}.  One first splits off the four intersection
spaces corresponding to the combinations of membership in the ranges
and kernels of $\sym{Psource}$ and $\sym{Qtarget}$.  On their remaining, generic part, the pair
is unitarily equivalent to
\begin{equation}
 \like{Psource}{P_{\mathrm{gen}}}=\begin{pmatrix}I&0\\0&0\end{pmatrix},
 \qquad
 \like{Qtarget}{Q_{\mathrm{gen}}}=\begin{pmatrix}\like{CosAngle}{C^2}&\like{CosAngle}{C}\like{SinAngle}{S}\\\like{CosAngle}{C}\like{SinAngle}{S}&\like{SinAngle}{S^2}\end{pmatrix},
 \qquad \like{CosAngle}{C^2}+\like{SinAngle}{S^2}=I,
 \quad \like{CosAngle}{C}\like{SinAngle}{S}=\like{SinAngle}{S}\like{CosAngle}{C},
 \label{eq:history-halmos-block}
\end{equation}
where $\like{CosAngle}{C}$ and $\like{SinAngle}{S}$ are commuting positive contractions on a common
Hilbert space.  Thus an operator-valued angle replaces a finite list
of principal angles.  B\"ottcher and Spitkovsky survey this
representation and its consequences for the algebra generated by two
projections \citep{bottcher2010gentle}.

The four intersection spaces have a second role. The right-angle defect spaces
\[
 \sym{CrossedSourceDefect}=\DKRange(\sym{Psource})\cap\DKKernel(\sym{Qtarget}),
 \qquad
 \sym{CrossedTargetDefect}=\DKKernel(\sym{Psource})\cap\DKRange(\sym{Qtarget})
\]
contain source directions orthogonal to the target, and target directions
orthogonal to the source, respectively. At acute angles
both disappear.  For a general pair they affect whether a specified
intertwining operator exists. An operator $\sym{IntertwinerT}$ intertwines $\sym{Psource}$ and $\sym{Qtarget}$ when
$\sym{IntertwinerT}\sym{Psource}=\sym{Qtarget}\sym{IntertwinerT}$: applying $\sym{IntertwinerT}$ after the source projection agrees with applying the
target projection after $\sym{IntertwinerT}$. In particular, $\sym{IntertwinerT}$ maps the source range into
the target range and its orthogonal complement into the target complement.
The existence criteria for direct rotations include conditions on the
dimensions of these defect spaces \citep[Section~3]{davis1970rotation}.

A finite-dimensional orthogonal matrix decomposes into invariant real
lines and two-dimensional rotation planes.  In a general Hilbert space an
angle operator can instead have continuous spectrum.  The spectral
calculus of the pair of projections takes the place of a finite list
of angles, and the intersection and defect spaces cannot be discarded.
Our Hilbert-space direct-rotation proofs therefore use projection identities
and functional calculus. In finite dimensions these identities recover the
principal-plane construction.

\subsection{Transport between nearby spectral subspaces}
\label{sec:related-projections}

A perturbation can change the invariant subspace associated with an isolated
part of a self-adjoint operator's spectrum. Let $\sym{Psource}$ and
$\sym{Qtarget}$ be the orthogonal projections onto the corresponding spectral
subspaces before and after perturbation. To compare these subspaces, we seek
a unitary carrying $\DKRange(\sym{Psource})$ onto
$\DKRange(\sym{Qtarget})$ that depends continuously, or analytically, on the
projections.
Sz.-Nagy developed the nearby-projection construction in work
published in 1947 \citep{sznagy1947perturbations,simon2018kato}.  Kato gave a streamlined
algebraic construction, described in his 1955 report and later in
\emph{Perturbation Theory for Linear Operators}
\citep[Chapter~I, \S4.6, equations~(4.33)--(4.42)]{kato1966perturbation};
see also \citet[Theorem~5.1]{simon2018kato}.

For orthogonal projections the useful intertwiner is
\begin{equation}
 \sym{IntertwinerT}=\sym{Qtarget}\sym{Psource}+(I-\sym{Qtarget})(I-\sym{Psource}), \qquad \sym{IntertwinerT}\sym{Psource}=\sym{Qtarget}\sym{IntertwinerT}.
 \label{eq:history-kato-T}
\end{equation}
The identity
\begin{equation}
 \like{IntertwinerT}{T^*}\sym{IntertwinerT}=I-(\sym{Psource}-\sym{Qtarget})^2
 \label{eq:history-kato-square}
\end{equation}
shows why $\|\sym{Psource}-\sym{Qtarget}\|<1$ is the natural threshold: it makes $\like{IntertwinerT}{T^*}\sym{IntertwinerT}$
strictly positive and invertible.  Hence
\begin{equation}
 \sym{Rrotation}=\sym{IntertwinerT}\bigl(I-(\sym{Psource}-\sym{Qtarget})^2\bigr)^{-1/2},
 \qquad \sym{Rrotation}\sym{Psource}\like{Rrotation}{D^*}=\sym{Qtarget}
 \label{eq:history-kato-unitary}
\end{equation}
is unitary. This is the polar decomposition $\sym{IntertwinerT}=\sym{Rrotation}|\sym{IntertwinerT}|$, where
$|\sym{IntertwinerT}|=(\like{IntertwinerT}{T^*}\sym{IntertwinerT})^{1/2}$ is positive. The factor $|\sym{IntertwinerT}|$ describes stretching;
the unitary factor $\sym{Rrotation}$ preserves lengths. Here $\sym{IntertwinerT}$ is invertible, so
$\sym{Rrotation}=\sym{IntertwinerT}|\sym{IntertwinerT}|^{-1}$. In principal-plane coordinates it is the direct rotation through the
principal angle, acting trivially on common directions.

The transport equation
$\sym{Wcompetitor}\sym{Psource}\sym{Wcompetitor}^*=\sym{Qtarget}$
has many unitary solutions.
The polar-factor formula selects one determined by the two projections
and continuously connected to the identity in the acute regime.  That
choice makes the construction valuable in perturbation theory: the
transport identifies moving spectral subspaces while avoiding an
arbitrary choice of orthonormal bases. A pair is called acute here when
$\|\sym{Psource}-\sym{Qtarget}\|<1$ (equivalently, the largest principal angle is less than $\pi/2$
for equal-dimensional finite subspaces). For an orthogonal or unitary transporter $\sym{Wcompetitor}$, the operator
$I-\sym{Wcompetitor}$ records the displacement of every vector. Davis and Kahan ask which
norms of this displacement are minimized by the direct rotation
\citep[Section~4]{davis1970rotation}.

Kato used such transformations to identify nearby spectral
subspaces of analytic operator families, thereby reducing the study
of an isolated moving eigenvalue cluster to a fixed coordinate
space.  Simon's historical account explains the earlier Sz.-Nagy
construction, Kato's formula, and subsequent work on pairs of
projections \citep[\S5]{simon2018kato}. Our proofs of direct-rotation
existence, uniqueness, and symmetry use the square-root and
polar-decomposition identities in
\cref{eq:history-kato-square,eq:history-kato-unitary}. The condition
$\|\sym{Psource}-\sym{Qtarget}\|<1$ applies in arbitrary Hilbert dimension.

\subsection{Spectral projections and residuals}
\label{sec:related-spectral}

The spectral theorem represents a self-adjoint operator as
$\sym{A}=\int\like{SpectralValue}{\lambda}\,d\sym{SpectralMeasure}(\like{SpectralValue}{\lambda})$ for a projection-valued measure $\sym{SpectralMeasure}$.
For a Borel set $\sym{SpectralSet}$, the operator $\sym{SpectralMeasure}(\sym{SpectralSet})$ is the orthogonal projection onto
the corresponding spectral subspace. This gives a second way of obtaining
a subspace: select a region of the spectrum and take the range of its
spectral projection. The functional-calculus
framework developed from the work of Riesz, von Neumann, and others;
Kato made spectral projections central to the perturbation theory of
closed operators \citep{kato1966perturbation,simon2018kato}.
Unlike the finite-dimensional eigenspace picture, a spectral
subspace need not be spanned by a finite list of eigenvectors.
A self-adjoint operator may be unbounded, while each of its spectral
projections is bounded and everywhere defined. The Hellinger--Toeplitz
theorem states that an everywhere-defined symmetric operator on a Hilbert
space is bounded \citep{hellinger1910grundlagen}. Thus an unbounded
self-adjoint operator has a proper domain. Its action $\sym{A}x$ is defined only
for $x\in\DKDom(\sym{A})$, so residual and intertwining equations must include
the required domain inclusions.

In the complex case, a number $z$ belongs to the resolvent set when
$zI-\sym{A}$ has a bounded, everywhere-defined inverse. That inverse,
$(zI-\sym{A})^{-1}$, is the resolvent at $z$. For an isolated bounded spectral
cluster, choose a positively oriented contour $\Gamma$ in the resolvent
set enclosing the cluster and no other spectrum. Its projection is
\begin{equation}
 \like{Pexact}{P_\Gamma}=\frac{1}{2\pi i}\int_\Gamma(zI-\sym{A})^{-1}\,dz.
 \label{eq:history-riesz-projection}
\end{equation}
The resolvent is bounded at each point of $\Gamma$ even when $\sym{A}$ is unbounded.
This makes the spectral projection a stable object for perturbation analysis,
  while preserving the separate domain requirements whenever $\sym{A}$ itself appears
  in an identity.
More general Borel spectral projections cover finite and half-infinite intervals,
  including spectral configurations for which a single isolated
  finite-eigenvalue cluster is not the right model.
Our formalization uses Borel functional calculus to construct projections for
  these spectral sets, including for unbounded self-adjoint operators.

The Courant--Fischer min--max principle and Weyl-type eigenvalue
perturbation bounds provide complementary information about spectral
\emph{values} \citep{bhatia1997matrix,kato1966perturbation}.  They do
not by themselves identify a stable basis inside a multiple-eigenvalue
cluster, or quantify the movement of its invariant subspace.  The Courant--Fischer and Weyl results that we formalized describe extremal
Rayleigh quotients and eigenvalue stability. The angle estimates add control
of the associated subspaces.

Davis studied the rotation of eigenvectors under perturbation in
1963 and 1965 \citep{davis1963rotation,davis1965rotation}.  The
1969 Davis--Kahan announcement preceded their extended 1970 treatment
\citep{daviskahan1969bounds,davis1970rotation}.  Their problem has
two closely related formulations.  In one, the original self-adjoint
operator is perturbed and the movement of an exact spectral subspace
is estimated.  In the other, one has an approximate trial subspace and
a computable residual: how far can it be from the desired exact
spectral subspace?  Unlike a comparison of individual eigenvectors,
these are meaningful even when an eigenvalue has multiplicity and its
orthonormal eigenbasis is nonunique.

Write $\sym{A}$ for a reference operator and $\sym{Acurrent}=\sym{A}+\sym{H}$ for the perturbed
ambient operator. We use this convention in the sine-theta statement below.
Let $\sym{E0}$ embed a trial-coordinate Hilbert space into the ambient space
isometrically, and let $\sym{A0}$ be a self-adjoint operator on those coordinates.
The residual is the bounded operator agreeing, where defined, with
\begin{equation}
 \sym{R}=\sym{Acurrent}\sym{E0}-\sym{E0}\sym{A0}.
 \label{eq:history-residual}
\end{equation}
If $\sym{R}=0$, applying $\sym{Acurrent}$ after the embedding agrees with applying
$\sym{A0}$ before it: the trial subspace is invariant and $\sym{A0}$ represents its
operator. For $\sym{R}\ne0$, the residual measures the failure of this identity.  A second isometric coordinate map for the
desired exact spectral space defines its orthogonal projection $\sym{Pexact}$.
The positive operator
\begin{equation}
 \sym{sinTheta0}=\bigl(I-\like{E0}{E_0^*}\sym{Pexact}\sym{E0}\bigr)^{1/2}
 \label{eq:history-angle-sine}
\end{equation}
measures the component of a trial vector lying outside the exact
subspace.  The critical distinction for the formal proof is that this
positive sine acts on the trial-coordinate space, whereas
$(I-\sym{Pexact})\sym{E0}$ maps into the ambient space.  For an operator $B$, its positive modulus $|B|=(B^*B)^{1/2}$ acts on its
input space and records its stretching factors. The modulus of $(I-\sym{Pexact})\sym{E0}$
is $\sym{sinTheta0}$, so the two maps have equal unitarily invariant ideal
norms whenever those norms are finite.

\subsection{Sylvester estimates}
\label{sec:related-sylvester}

The matrix equation $\sym{SylvesterLeft}\sym{X}-\sym{X}\like{SylvesterRight}{A}=\sym{SylvesterRHS}$ is associated with Sylvester's 1884
work \citep{sylvester1884equation}.  Rosenblum studied separated
operator spectra in 1956 and the equation with self-adjoint coefficients
in 1969 \citep{rosenblum1956equation,rosenblum1969selfadjoint}.

The scalar prototype is $(\like{SylvesterLeft}{b}-\like{SylvesterRight}{a})\like{X}{x}=\like{SylvesterRHS}{y}$: if $|\like{SylvesterLeft}{b}-\like{SylvesterRight}{a}|\geq\sym{delta}>0$,
then $\sym{delta}|\like{X}{x}|\leq|\like{SylvesterRHS}{y}|$.  In the operator problem the scalar
multipliers are replaced by the spectral actions of two operators.
The operator estimate has the form $\sym{delta} \sym{NormN}(\sym{X})\leq \sym{NormN}(\sym{SylvesterLeft}\sym{X}-\sym{X}\like{SylvesterRight}{A})$.
Its hypotheses specify the spectral positions of $\like{SylvesterRight}{A}$ and $\sym{SylvesterLeft}$, their domains
when unbounded, and the unitarily invariant norm $\sym{NormN}$.

Let $\sym{F1}$ give isometric coordinates for the unwanted complementary
spectral subspace, with self-adjoint block $\sym{Lambda1}$.  Writing
$\sym{X}=\like{F1}{F_1^*}\sym{E0}$ converts the residual identity into a Sylvester equation:
\begin{equation}
 \sym{Lambda1}\sym{X}-\sym{X}\sym{A0}=\like{F1}{F_1^*}\sym{R}
 \label{eq:history-sylvester}
\end{equation}
on the appropriate domain.  Thus $\sym{X}$ describes the geometric leakage
of the trial space into the unwanted block, while the right side is
controlled by the computable residual.  Under the interval/exterior
or ordered spectral-gap hypotheses of the sine-theta theorem,
separated Sylvester estimates give
\begin{equation}
 \sym{delta}\,\sym{NormN}(\sym{X})\leq \sym{NormN}(\like{F1}{F_1^*}\sym{R})\leq \sym{NormN}(\sym{R}).
 \label{eq:history-separation-bound}
\end{equation}
This is the bridge from a spectral-gap assumption to the bound on the
sine of a subspace angle.

The coefficient of $\sym{NormN}(\sym{R})/\sym{delta}$ in this bound is one.
This coefficient uses the ordered or interval/exterior spectral
positions specified by Davis and Kahan. Arbitrarily interlaced spectra can
require a different constant. Our theorem includes both of the source's
spectral configurations, including ordered cases where both operators are
unbounded.  A formal implementation
therefore needs the domains of unbounded self-adjoint operators,
their spectral projections, strong spectral cutoffs, and the
ideal-norm inequalities used to pass through bounded approximations.

\subsection{Operator ideals and displacement norms}
\label{sec:related-norms}

The geometry of a pair of subspaces is fixed, but there are many ways
to assign a cost to a unitary transporting one onto the other.
Von Neumann's 1937 study of unitarily invariant matrix norms identified
the role of symmetric gauge functions of singular values
\citep{vonneumann1937matrix}. Schatten's 1950 monograph developed normed
operator ideals; Ky Fan's 1951 eigenvalue and singular-value inequalities
provided the partial-sum comparisons underlying majorization
arguments \citep{schatten1950crossspaces,fan1951maximum,bhatia1997matrix}.
The singular values $s_j(Z)$ are the nonnegative square roots of the
eigenvalues of $Z^*Z$, ordered decreasingly. They record the stretching
factors of the matrix $Z$. For $1\leq p<\infty$, the Schatten norms include
\begin{equation}
 \|Z\|_{S_p}=\left(\sum_j s_j(Z)^p\right)^{1/p},
 \qquad
 \|Z\|_{S_1}=\sum_j s_j(Z),
 \quad
 \|Z\|_{S_2}=\left(\sum_j s_j(Z)^2\right)^{1/2}.
 \label{eq:history-schatten}
\end{equation}
An operator ideal is a linear class of operators closed under multiplication
on either side by bounded operators. A normed ideal also controls the norm
of these products. For compact operators in infinite dimensions, the
Schatten $p$ ideal consists of those whose singular values have a finite
$p$-power sum. Thus a trace-norm inequality requires trace-class membership;
this is an instance of the source's qualification that its norms be defined.

The largest singular value gives the operator norm, the Euclidean
norm of the singular-value sequence gives the Hilbert--Schmidt norm,
and its sum gives the trace norm. These norms can rank transporters
differently. In Section~\ref{sec:prop44}, the two displacement singular-value
lists are $(2\sin(\sym{theta}/2),\ldots,2\sin(\sym{theta}/2))$ with four entries and
$(2\sin\sym{theta},2\sin\sym{theta},0,0)$. The second has a smaller sum and a larger
maximum. Ky Fan norms compare partial sums of the ordered singular values;
these comparisons determine which norm inequalities follow.

Davis established Hilbert--Schmidt and finite-dimensional Schatten
extremal properties for exponents at least two in his 1958 subspace
paper \citep[Theorems~7.3--7.4]{davis1958separation}.  Davis and
Kahan compare three objectives \citep[Section~4]{davis1970rotation} for a transporter $\sym{Wcompetitor}$ carrying
$\sym{Usource}$ onto $\sym{Vtarget}$: its restricted displacement $(I-\sym{Wcompetitor})\like{Psource}{P_U}$, its squared full
displacement $(I-\sym{Wcompetitor})^*(I-\sym{Wcompetitor})$, and its full displacement $I-\sym{Wcompetitor}$. The last acts
on every ambient vector, including vectors outside $\sym{Usource}$. They gave a
complex-phase modification that decreases its trace norm \citep[Example~4.2]{davis1970rotation}.
The full-displacement claim \citep[Proposition~4.4]{davis1970rotation}
asserts that direct rotation minimizes
full displacement for real subspaces if their principal angles are no
larger than $\pi/3$.  The four-dimensional example in
\cref{sec:prop44} gives a smaller trace norm at every positive acute
angle.  It exposes the failed step: adding singular-value sums of
restrictions to separate principal planes is not a valid global lower
bound for an arbitrary ambient transporter.

The subsequent minimal-distortion theory of Bolotnikov, Li, and Rodman
studied the same class of subspace-transport constraints, including
Frobenius and Q-norm optimality and explicit subquadratic
counterexamples \citep{bolotnikov2001minimal}.  Later work on
Grassmannian metrics and canonical-angle inequalities also separates
intrinsic paths between subspaces from minimizing the extrinsic
operator displacement \citep{qiu2006grassmann,zhang2010canonicalangles}.
These results give the norm-theoretic context for the source
correction and for the trace-norm matching problem in
\cref{sec:example-transport}.

\subsection{The Davis--Kahan angle estimates}
\label{sec:related-synthesis}

Davis's papers of 1963 and 1965 and the Davis--Kahan announcement of 1969
preceded the extended 1970 treatment
\citep{davis1963rotation,davis1965rotation,daviskahan1969bounds,davis1970rotation}.
For the sine-theta estimate, let $\sym{R}=\sym{Acurrent} \sym{E0}-\sym{E0}\sym{A0}$ be the residual
and $\sym{Pexact}$ the desired exact spectral projection. The operator $(I-\sym{Pexact})\sym{E0}$
extracts the part of a trial vector outside that spectral subspace. In
isometric coordinates $\sym{F1}$ for the complementary space, this component is
represented by $\sym{X}=\like{F1}{F_1^*}\sym{E0}$. The residual identity gives
\cref{eq:history-sylvester}, and a separated Sylvester estimate bounds $\sym{NormN}(\sym{X})$.
Since $\sym{S}=\sym{F1}\sym{X}$ and $\sym{F1}$ is an isometry,
$\sym{NormN}(\sym{X})=\sym{NormN}(\sym{S})$.
The modulus identity $|\sym{S}|=\sym{sinTheta0}$ gives
$\sym{NormN}(\sym{S})=\sym{NormN}(\sym{sinTheta0})$.

The resulting chain of equalities and inequalities is
\begin{equation}
 \underbrace{\sym{delta} \sym{NormN}(\sym{sinTheta0})}_{\text{subspace error}}
 =\sym{delta} \sym{NormN}(\sym{X})
 \leq \sym{NormN}(\like{F1}{F_1^*}\sym{R})
 \leq\underbrace{\sym{NormN}(\sym{R})}_{\text{computable residual}}.
 \label{eq:history-unified-chain}
\end{equation}
The first equality follows from the positive-modulus identity for the
complementary projection. The first inequality applies the Sylvester
estimate to $\sym{X}$; the second uses $\sym{NormN}(BC)\leq\|B\|_{\mathrm{op}}\sym{NormN}(C)$ and
$\|\like{F1}{F_1^*}\|_{\mathrm{op}}\leq1$. For unbounded operators the Sylvester
identity holds on the trial operator's domain, with a bounded residual.
These are the geometric, spectral, and norm estimates formalized in our
proof of the theorem.

The tangent and double-angle estimates use related projection blocks.
For example, an estimate for $\sin(2\sym{Theta})$ allows angles near zero and
near $\pi/2$. The spectral-selection results \citep[Section~8]{davis1970rotation}
identify which spectral subspaces give the small-angle alternative.
The direct-rotation analysis \citep[Section~4]{davis1970rotation}
studies the orthogonal or unitary map carrying one subspace onto another
and asks which displacement norms the direct rotation minimizes.

\subsection{Numerical and statistical applications}
\label{sec:related-numerical}

The 1970 theory influenced more than the analysis of exact
self-adjoint eigenvectors.  Wedin's 1972 singular-subspace
perturbation bounds address the corresponding problem for singular
vectors of general rectangular matrices
\citep{wedin1972perturbation}.  Bj\"orck and Golub's angle algorithms
made the subspace quantities computable
\citep{bjorck1973numerical}.  Stewart's work on pseudoinverses,
projections, and least-squares sensitivity helped connect these
geometric notions with matrix computations
\citep{stewart1977perturbation,stewart1990matrix}.
Paige and Wei's 1994 history shows how the cosine--sine decomposition
links the principal-angle geometry of Davis and Kahan with numerical
matrix factorizations \citep{paige1994history}.

The same subspace-versus-eigenvector distinction remains central in
statistical spectral estimation.  For example, Yu, Wang, and
Samworth give a statistical variant whose eigenspace error bound uses
separation in the population spectrum \citep{yu2015davis}, while Chen and
collaborators survey spectral methods across statistical learning
problems \citep{chen2021spectral}.  In that setting, an empirical covariance matrix approximates a population
operator and the perturbation norm controls the error of estimated
eigenspaces. Our general theorem supplies the deterministic angle estimate;
using it in a statistical result also requires the corresponding probabilistic
control of the perturbation.

\subsection{Foundations used in our formalization}
\label{sec:related-formal}

Each historical component above corresponds to a specific requirement
in the formalization.  Jordan's angle geometry and Halmos's
projection blocks lead to positive principal-angle operators rather
than chosen eigenbases.  The Sz.-Nagy--Kato intertwiner requires
polar decomposition, square roots, and real/complex transfer.
Spectral projections require Borel functional calculus, including for
unbounded self-adjoint operators.  The Sylvester proof requires a
representation carrying operator domains and a controlled passage
through spectral cutoffs.  Schatten and Ky Fan theory requires
approximation numbers and symmetric operator ideals to express
inequalities with the source's norm generality.

We formalized principal-angle geometry, symmetric operator ideals, and
unbounded Sylvester estimates using Lean~4 and Mathlib
\citep{demoura2021lean4,mathlib2020}. Mathlib supplies the algebraic,
topological, and Hilbert-space framework and continuous functional calculus
\citep{dedecker2025cfc}. Our projection-valued-measure and Borel-calculus
work also adapts Adam Bornemann's Spectra formalization
\citep{bornemann2026spectra}.

Our Courant--Fischer and Weyl development, rectangular unitarily invariant
seminorms, and projection-valued measures were subsequently adapted into
Tau Ceti \citep{tauceti2026}. The last of these carries the earlier Spectra ancestry. We also use
operator-theory results developed independently in Tau Ceti, including its
canonical polar decomposition.
% Evidence: dev/tauceti/dkps-upstream-contribution-census.md, PRs 13599,
% 14241, 17527 (adaptations), and 16127/16450 (independent polar proofs).
% AUTHORING_NOTES.md indexes accepted-source and campaign evidence.
% Mathematical origin and formal proof ancestry are separate attributions.

This history also separates the two mathematical questions that meet
in the project.  The sine-theta theorem estimates \emph{how far} an
approximate or perturbed spectral subspace lies from an exact one.
The direct-rotation problem asks \emph{which} ambient transformation
carries one given subspace onto another at minimum displacement cost.
Our proofs of the optimization results use the principal-angle,
projection, and operator-ideal foundations developed for the formalization.
The matching theorem has its own proof; the shared foundations explain the
relationship between the two parts of the work. The resulting counterexample
to the full-displacement claim \citep[Proposition~4.4]{davis1970rotation}
and the follow-on matching theorem are examined in
\cref{sec:prop44,sec:example-transport}.

\section{Davis--Kahan formalization and the sine-theta theorem}
\label{sec:coverage}

Our targets are the four headline theorems \citep[Section~2]{davis1970rotation} and every
named theorem, proposition, lemma, and corollary established by Davis and
Kahan in the 1970 paper. These \DKResultCount{} results include the sine and
tangent inequalities, the geometry and extremal properties of direct
rotations, bounded and unbounded Sylvester estimates, the generalized
perturbation bounds \citep[Section~6]{davis1970rotation}, and the spectral-selection
theorems \citep[Section~8]{davis1970rotation}. We prove \DKProvedCount{} results and refute the printed
full-displacement claim \citep[Proposition~4.4]{davis1970rotation}.
Appendix~\ref{app:dk-inventory} lists these targets.

We develop the sine-theta theorem and its proof below; the remaining source
results are summarized in Appendix~\ref{app:dk-inventory}, with hypotheses
and conclusions in Section~\ref{app:dk-result-summaries}. The main text also
explains the spectral-selection results needed to interpret the angle bounds.
Appendix~\ref{app:notation} collects the notation used here and in the
direct-rotation argument.

\subsection{The sine-theta theorem for self-adjoint operators}
\label{sec:sine-theta}

\subsubsection{Trial and exact subspaces}

We use $\sym{A}$ for the reference operator and
$\sym{Acurrent}=\sym{A}+\sym{H}$ for the perturbed ambient operator.
The symbol $\sym{Acurrent}$ corresponds to $\sym{A}+\sym{H}$ in the source.

For a unit trial vector $\sym{v}$ and an approximate eigenvalue $\sym{mu}$, the residual
is $\sym{Acurrent}\sym{v}-\sym{mu} \sym{v}$. An isometric coordinate map $\sym{E0}$ generalizes
this construction to a trial subspace: its range consists of the candidate
approximate vectors, and a self-adjoint operator $\sym{A0}$ represents their
approximate dynamics. For a finite-dimensional Rayleigh--Ritz choice,
$\sym{A0}=\like{E0}{E_0^*}\sym{Acurrent}\sym{E0}$. The theorem allows other self-adjoint trial
operators as well.

Let $\sym{Eambient}$ be the ambient Hilbert space.  Let $\sym{A0}$ be a self-adjoint operator on a
trial-coordinate space $\sym{Ftrial}$, and let
\[
  \sym{E0}:\sym{Ftrial}\to \sym{Eambient}
\]
be an isometric coordinate map for the trial subspace.  The residual compares
applying the ambient operator after embedding with applying the trial operator
before embedding:
\begin{equation}
  \sym{R} = \sym{Acurrent}\sym{E0}-\sym{E0}\sym{A0}.
  \label{eq:residual}
\end{equation}
When the operators are unbounded, \cref{eq:residual} is asserted on
$\DKDom(\sym{A0})$ and $\sym{R}$ is the bounded residual that agrees there.
The Lean predicate \texttt{IsTrialResidual} packages the isometry, the required
domain inclusion, and this residual identity.

Let $\sym{F0}:\sym{Hexact}\to \sym{Eambient}$ be isometric coordinates for the desired exact spectral
subspace and let $\sym{F1}:\sym{Gunwanted}\to \sym{Eambient}$ be isometric coordinates for its orthogonal
complement.  A self-adjoint operator $\sym{Lambda1}$ on $\sym{Gunwanted}$ describes the
complementary spectral block through
\[
  \sym{Acurrent}\sym{F1} = \sym{F1}\sym{Lambda1}
\]
on its domain. This intertwining identity says that $\sym{Lambda1}$ gives
the action of $\sym{Acurrent}$ in complementary coordinates. The predicate
\texttt{IsExactSpectralDecomposition} also requires the domain inclusion
$\sym{F1}(\DKDom(\sym{Lambda1}))\subseteq\DKDom(\sym{Acurrent})$
and states orthogonality and completeness through
\[
 \like{F0}{F_0^*}\sym{F1}=0,\qquad
 \sym{F0}\like{F0}{F_0^*}+\sym{F1}\like{F1}{F_1^*}=I.
\]

\subsubsection{The sine operator and its rectangular representative}

The source angle operator acts on the trial-coordinate space.  If
$\sym{Pexact}=\sym{F0}\like{F0}{F_0^*}$ is the orthogonal projection onto the desired exact subspace, then
\[
  \cos\sym{Theta0} = (\like{E0}{E_0^*}\sym{Pexact}\sym{E0})^{1/2},
  \qquad
  \sym{sinTheta0} = (I-\like{E0}{E_0^*}\sym{Pexact}\sym{E0})^{1/2}.
\]
The analytic proof uses the rectangular complementary block
\begin{equation}
  \sym{S}=(I-\sym{Pexact})\sym{E0}.
  \label{eq:rectangular-sine}
\end{equation}
The theorem uses the positive sine on the trial-coordinate space.  The rectangular block has codomain $\sym{Eambient}$, while the
positive sine has codomain $\sym{Ftrial}$.  Their relationship is
\[
  \sym{S}^*\sym{S}
  = \like{E0}{E_0^*}(I-\sym{Pexact})\sym{E0}
  = I-\like{E0}{E_0^*}\sym{Pexact}\sym{E0},
\]
so
\begin{equation}
  |\sym{S}|=(\sym{S}^*\sym{S})^{1/2}=\sym{sinTheta0}.
  \label{eq:modulus-sine}
\end{equation}
% It might be worth noting that our formalization actually defines \sym{Theta0}
% and proves eq:modulus-sine whereas most treatments of the subject seam to
% just start from |\sym{S}|. Maybe it's worth enough to put in the paper, but
% it at least deserves a latex comment!
The polar decomposition $\sym{S}=\sym{PolarU}|\sym{S}|$ separates the positive
  modulus from a partial isometry.
Unitarily invariant ideal norms assign the same value to an operator and its
  modulus, so \[ \sym{NormN}(\sym{S})=\sym{NormN}(\sym{sinTheta0}) \] whenever
  the norm is finite.
This equality transfers an estimate for the complementary block $\sym{S}$ to the
  positive sine operator.

\subsubsection{Gap geometry and the inequality}
\label{sec:gap-geometry}

Fix $\sym{delta}>0$. There are two geometric forms of separation.
\emph{Interval/exterior separation:} for some real
$\sym{betaEndpoint}\leq\sym{alphaEndpoint}$, one of the two spectra
$\sigma(\sym{A0})$ and $\sigma(\sym{Lambda1})$ lies in $[\sym{betaEndpoint},\sym{alphaEndpoint}]$ and the other
lies in
\[
 (-\infty,\sym{betaEndpoint}-\sym{delta}]\;\cup\;[\sym{alphaEndpoint}+\sym{delta},\infty).
\]
\emph{Ordered separation:} for some real $b$, the operator forms satisfy either
\[
 \sym{A0}\geq (b+\sym{delta})I,\quad \sym{Lambda1}\leq bI,
 \qquad\text{or}\qquad
 \sym{A0}\leq bI,\quad \sym{Lambda1}\geq (b+\sym{delta})I.
\]
Each form inequality is asserted on the domain of the corresponding operator;
for example, $\sym{A0}\geq aI$ means
$\langle \sym{A0}x,x\rangle\geq a\|x\|^2$ for every
$x\in\DKDom(\sym{A0})$. For self-adjoint operators, the source's ordered
hypotheses on half-infinite spectral intervals, such as $(-\infty,b]$ and
$[b+\sym{delta},\infty)$, imply these form bounds. These cases are included
in the source's unbounded sine theorem
\citep[Section~6, discussion of unbounded operators]{davis1970rotation}.
Write $\sym{SineGap}(\sym{A0},\sym{Lambda1})$ for their union.
The Lean predicate \texttt{FormBoundedSylvesterGap} has three cases:
the finite-interval alternative, which includes either spectral order,
and the two ordered form alternatives.

Both blocks may be unbounded in an ordered configuration. For instance,
on $\ell^2(\mathbb N_{\geq1})$, take diagonal operators
$\sym{A0}e_n=(n+\sym{delta})e_n$ and $\sym{Lambda1}e_n=-ne_n$, on the domains where their
images are square summable. They satisfy the first ordered configuration
with $b=0$, while neither spectrum lies in a finite interval. This example
explains why retaining the half-infinite alternatives changes the theorem's
applicability.

For the unbounded trial configuration, the source assumes that
$\sym{Acurrent}\sym{E0}$ and $\sym{E0}\sym{A0}$ have a common dense domain.
Since $\sym{E0}$ is everywhere defined, this gives
\[
 \DKDom(\sym{Acurrent}\sym{E0})=\DKDom(\sym{A0}),
 \qquad
 \sym{E0}(\DKDom(\sym{A0}))\subseteq\DKDom(\sym{Acurrent}).
\]
The residual agrees with the difference on this domain and extends to a
bounded operator. The formal predicate \texttt{IsTrialResidual} includes the
domain inclusion and the residual identity.

The norm $\sym{NormN}$ is a normalized symmetric operator-ideal norm.
Normalization means that the rank-one map $uv^*:x\mapsto u\langle v,x\rangle$
has norm
\[
 \sym{NormN}(uv^*)=\|u\|\,\|v\|;
\]
in particular, its norm is one when $u,v$ are unit vectors.
Its hypotheses also include Ky Fan dominance whenever both ideal norms
are finite: domination of all finite singular-value sums implies domination
of the ideal norms. This comparison does not assert ideal membership.
The norm is invariant under isometries and positive modulus, and satisfies
$\sym{NormN}(BT)\leq\|B\|_{\mathrm{op}}\sym{NormN}(T)$. Trace-class and
Hilbert--Schmidt operators with their respective norms are examples; the
operator norm on bounded operators is another. For an ideal norm, both the
residual and the positive sine must belong to its domain.

In two dimensions, take a one-dimensional trial frame $\sym{E0}(x)=x\sym{v}$ with
$\|\sym{v}\|=1$. A scalar trial operator $\sym{A0}=\sym{rho}$ gives
$\sym{r}=\sym{R}(1)=\sym{Acurrent}\sym{v}-\sym{rho}\sym{v}$.
Choosing $\sym{rho}=\langle\sym{v},\sym{Acurrent}\sym{v}\rangle$
makes $\sym{rho}\sym{v}$ the orthogonal projection of
$\sym{Acurrent}\sym{v}$ onto the trial line. The residual is therefore
perpendicular to that line. Figure~\ref{fig:sine-theta-residual} pairs this
residual triangle with the exact complementary projection that measures
the subspace error.

\begin{figure}[tb]
\centering
\resizebox{\linewidth}{!}{\input{figures/sine_theta_residual_geometry.pgf}}
\caption{Rayleigh residual geometry for a one-dimensional trial space in
$\mathbb R^2$ (the rank-one case). With $\sym{E0}(x)=x\sym{v}$, the residual
operator satisfies $\sym{R}(x)=x\sym{r}$, where
$\sym{r}=\sym{Acurrent}\sym{v}-\sym{rho}\sym{v}$.
For a higher-dimensional trial space,
$\sym{R}=\sym{Acurrent}\sym{E0}-\sym{E0}\sym{A0}$ assigns a residual vector
to each trial-coordinate vector. In a finite-dimensional eigenbasis of
$\sym{A0}$, its columns are the individual eigenpair residuals.
The exact frames in the drawing are
$\sym{F0}(h)=(h,0)$ and $\sym{F1}(g)=(0,g)$, with
$\sym{Pexact}=\sym{F0}\like{F0}{F_0^*}$ and $\sym{PexactPerp}=I-\sym{Pexact}$.
(a) The pink leg $\sym{S}(1)=\sym{PexactPerp}\sym{v}$ has length $\sin\sym{theta}$.
(b) The green residual joins $\sym{rho}\sym{v}$ to $\sym{Acurrent}\sym{v}$;
for the Rayleigh choice it is perpendicular to the trial line. The dashed
legs are its exact wanted and unwanted components. Their lengths give
$\sym{delta}\sin\sym{theta}=\|\sym{PexactPerp}\sym{r}\|_2
\leq\|\sym{r}\|_2=\|\sym{R}\|_{\mathrm{op}}$.
The drawing uses $\sym{Acurrent}=\operatorname{diag}(1.9,0.7)$,
$\sym{theta}=40^\circ$, $\sym{rho}=\langle\sym{v},\sym{Acurrent}\sym{v}\rangle$,
and $\sym{delta}=\sym{rho}-0.7$.}
\label{fig:sine-theta-residual}
\end{figure}

Appendix~\ref{app:lean-guide} explains the general Lean syntax in the
following theorem signature; the geometric and spectral hypotheses have been
defined above.

\begin{theorem}[Sine-theta estimate]
\label{thm:sine-theta}
Let the scalar field be $\mathbb R$ or $\mathbb C$, let $\sym{Eambient}$ be separable,
and let $\sym{Acurrent}$, $\sym{A0}$, and $\sym{Lambda1}$ be self-adjoint operators on the
Hilbert spaces above. Suppose the isometries, domain inclusions, residual
identity, exact complementary decomposition, and the separation condition
$\sym{SineGap}(\sym{A0},\sym{Lambda1})$ of
Section~\ref{sec:gap-geometry} hold. Let $\sym{NormN}$ be a normalized symmetric
operator-ideal norm. Whenever $\sym{sinTheta0}$ and $\sym{R}$ belong to its domain,
\begin{equation}
 \sym{delta}\,\sym{NormN}(\sym{sinTheta0})\leq \sym{NormN}(\sym{R}).
 \label{eq:dk-sine}
\end{equation}
\end{theorem}

\begin{lstlisting}[style=leanpaper,
  caption={Lean signature for the sine-theta estimate. The surrounding
  scalar-field and Hilbert-space context is omitted.
  Lean's \texttt{A} denotes $\sym{Acurrent}$;
  its carrier types \texttt{E}, \texttt{F}, \texttt{G}, \texttt{H} denote
  $\sym{Eambient}$, $\sym{Ftrial}$, $\sym{Gunwanted}$, $\sym{Hexact}$, respectively,
  and \texttt{K} is the scalar field. Thus \texttt{H} here is a coordinate
  space, while $\sym{H}$ in the prose is the perturbation. The positive trial-coordinate
  sine is \texttt{sourceDirectedSinThetaOperator E0 F0};
  \texttt{FormBoundedSylvesterGap} encodes the interval/exterior and ordered
  half-infinite separation hypotheses described above.},
  label={lst:dk-sintheta-headline}]
theorem sinTheta_unbounded_formGap_whereDefinedUIN_rclike
    [TopologicalSpace.SeparableSpace E]
    (N : NormalizedSymmetricOperatorIdealFamily K)
    (A : LinearPMap K E E)
    (A0 : LinearPMap K F F)
    (Lambda1 : LinearPMap K G G)
    (E0 : ContinuousLinearMap K F E)
    (F0 : ContinuousLinearMap K H E)
    (F1 : ContinuousLinearMap K G E)
    (R : ContinuousLinearMap K F E)
    (hA : IsSelfAdjoint A)
    (hA0 : IsSelfAdjoint A0)
    (hLambda1 : IsSelfAdjoint Lambda1)
    (htrial : IsTrialResidual A A0 E0 R)
    (hexact : IsExactSpectralDecomposition A Lambda1 F0 F1)
    (delta : Real) (hdelta : 0 < delta)
    (hgap : FormBoundedSylvesterGap A0 Lambda1 delta) :
    N.Mem (sourceDirectedSinThetaOperator E0 F0) ->
    N.Mem R ->
      delta * N.gaugeReal (sourceDirectedSinThetaOperator E0 F0) <=
        N.gaugeReal R
\end{lstlisting}

The two ideal-membership hypotheses ensure that the norm values in the
conclusion are finite. The gap predicate includes the finite
interval/exterior case and both ordered cases. In the ordered cases,
$\sym{A0}$ and $\sym{Lambda1}$ may both be unbounded.

\subsubsection{Why the Sylvester estimate gives the angle bound}

Set $\sym{X}=\like{F1}{F_1^*}\sym{E0}$. The exact complementary decomposition gives
$\sym{S}=\sym{F1}\sym{X}$, and multiplication by the isometry $\sym{F1}$ preserves the
singular-value gauges. Self-adjointness, the domain hypotheses, and the
residual identity yield the Sylvester relation
\[
 \sym{Lambda1} \sym{X}-\sym{X}\sym{A0}=\like{F1}{F_1^*}\sym{R}
\]
on $\DKDom(\sym{A0})$, with the left side represented by the bounded
right side. The separated Sylvester estimate gives
\[
 \sym{delta} \sym{NormN}(\sym{X})\leq \sym{NormN}(\like{F1}{F_1^*}\sym{R})\leq \sym{NormN}(\sym{R}).
\]
Combining this with $\sym{NormN}(\sym{X})=\sym{NormN}(\sym{S})=\sym{NormN}(\sym{sinTheta0})$ proves the conclusion.
For unbounded operators, deriving the domain compatibility for $\sym{X}$ is part of
the argument. The displayed relation holds on $\DKDom(\sym{A0})$ and
has a bounded right side. The specified interval/exterior or ordered
spectral positions give the factor $1/\sym{delta}$, with no additional
multiplicative constant.

For a one-dimensional trial space, take
$\sym{Acurrent}=\operatorname{diag}(0,\sym{delta})$, the trial scalar
$\sym{A0}=0$, and the unit vector
$\like{v}{v_\theta}=(\cos\sym{theta},\sin\sym{theta})$.
Choose the desired exact subspace $\operatorname{span}\{(1,0)\}$;
its orthogonal complement is $\operatorname{span}\{(0,1)\}$, on which
$\sym{Lambda1}=\sym{delta}$.
The complementary component is $(0,\sin\sym{theta})$ and the
residual is $(0,\sym{delta}\sin\sym{theta})$. Thus equality holds
in the sine-theta estimate for this finite-dimensional configuration.

\section{Formalized operator-theoretic foundations}
\label{sec:foundations}

We formalized the geometry of two subspaces, symmetric operator ideals and
majorization, and Sylvester estimates for unbounded self-adjoint operators.
These results supply the constructions and inequalities used in the
Davis--Kahan proofs.

\subsection{Direct rotations and two-subspace geometry}

Let $\sym{Psource}$ and $\sym{Qtarget}$ project onto closed subspaces
$\sym{Usource}$ and $\sym{Vtarget}$. For an acute pair,
$\|\sym{Psource}-\sym{Qtarget}\|<1$. The compression
$\sym{Psource}\sym{Qtarget}\sym{Psource}$ on $\sym{Usource}$ determines
the positive cosine operator. Its square root and the complementary
compression give the cosine and sine factors in the direct-rotation
construction. Polar decomposition gives the isometries between the
corresponding parts of the two subspaces.

The resulting orthogonal or unitary map $\sym{Rrotation}$ carries
$\sym{Usource}$ onto $\sym{Vtarget}$. In finite-dimensional principal-plane
coordinates it rotates each plane through its principal angle and fixes
common directions. The projection formula from
\cref{eq:history-kato-unitary} defines it in Hilbert space without requiring
a discrete list of principal vectors.

Our formalization proves acute existence and uniqueness, the
principal-square-root characterization, commutation with the angle
operators, and reversal symmetry: interchanging the two subspaces replaces
the rotation by its adjoint. For general pairs, the existence criteria
include the right-angle defect spaces described in
\cref{sec:related-two-projections}. When the angle operator is compact,
its spectral decomposition recovers a sequence of principal angles
\citep[Section~3]{davis1970rotation}.

The double-angle theorems for $\sin(2\sym{Theta})$ and $\tan(2\sym{Theta})$
  require selecting the branch with angles below a quarter turn.
Here $\sym{Theta}$ includes both the source and complementary angle blocks;
$\sym{Theta0}$ alone is the directed angle on source coordinates.
Because the function $\sin(2\sym{theta})$ is small both when $\sym{theta}$ is
  near $0$ and when $\sym{theta}$ is near $\pi/2$, a bound on
  $\sin(2\sym{Theta})$ alone cannot determine whether an approximate subspace is
  close to or nearly orthogonal to the target.
The branch-selection theorem \citep[Theorem~8.1]{davis1970rotation}
uses an ordered reference gap. Let $\sym{Psource}$ reduce a self-adjoint
reference operator $\sym{A}$, with
blocks $\sym{ReferenceBlock0}$ on $\DKRange(\sym{Psource})$ and
$\sym{ReferenceBlock1}$ on $\DKKernel(\sym{Psource})$, satisfying
\[
 \sym{ReferenceBlock0}\le\sym{alphaEndpoint}I,
 \qquad
 \sym{ReferenceBlock1}\ge(\sym{alphaEndpoint}+\sym{delta})I,
 \qquad \sym{delta}>0,
\]
in the quadratic-form sense. Suppose $\sym{H}$ is bounded and self-adjoint,
with both diagonal blocks zero:
$\sym{Psource}\sym{H}\sym{Psource}
 =(I-\sym{Psource})\sym{H}(I-\sym{Psource})=0$.
These hypotheses imply that $\sym{Acurrent}=\sym{A}+\sym{H}$ has no spectrum
in $(\sym{alphaEndpoint},\sym{alphaEndpoint}+\sym{delta})$.

Fix a reducing subspace $M$ of $\sym{Acurrent}$, and write
$\sym{Lambda0},\sym{Lambda1}$ for the blocks of $\sym{Acurrent}$ on $M,M^\perp$.
The full angle $\sym{Theta}$ comparing $\DKRange(\sym{Psource})$ with $M$
satisfies $\sym{Theta}\le\pi/4$ if and only if
$\sym{Lambda0}\le\sym{alphaEndpoint}I$ and
$\sym{Lambda1}\ge(\sym{alphaEndpoint}+\sym{delta})I$.
Such a subspace exists: take the spectral subspace of $\sym{Acurrent}$
for $(-\infty,\sym{alphaEndpoint}]$. Our formalization also proves that
it is the unique reducing subspace satisfying this angle bound.

To compare the reference and perturbed blocks on common coordinate spaces,
choose isometric frames for $\DKRange(\sym{Psource})$ and
$\DKKernel(\sym{Psource})$. Apply the direct rotation $\sym{Rrotation}$ to
these frames to obtain frames for $M$ and $M^\perp$. Express
$\sym{ReferenceBlock0},\sym{ReferenceBlock1}$ in the original frames and
$\sym{Lambda0},\sym{Lambda1}$ in the transported frames. In these coordinates
the branch satisfies the stronger quadratic-form inequality
\[
 \sym{ReferenceBlock1}-\sym{alphaEndpoint}I
 \le\sym{ComplementCosineBlock}
       (\sym{Lambda1}-\sym{alphaEndpoint}I)\sym{ComplementCosineBlock},
 \qquad \sym{ComplementCosineBlock}=\cos\sym{Theta1},
\]
with the lower-block analogue. These inequalities quantify spectral
repulsion in terms of the angles and yield eigenvalue and symmetric-gauge
comparisons in finite dimensions.

The smallness criterion \citep[Theorem~8.2]{davis1970rotation} uses a central
spectral interval instead of an ordered gap. Let $\sym{Psource}$ reduce
the self-adjoint operator $\sym{A}$ and $\sym{Qtarget}$ reduce the
self-adjoint operator $\sym{Acurrent}$. The perturbed block on
$\DKRange(\sym{Qtarget})$ has spectrum in
$[\sym{betaEndpoint},\sym{alphaEndpoint}]$, while its complementary spectrum
lies in
\[
 (-\infty,\sym{betaEndpoint}-\sym{delta}]
 \ \cup\ [\sym{alphaEndpoint}+\sym{delta},\infty),
 \qquad \sym{betaEndpoint}\le\sym{alphaEndpoint},\quad\sym{delta}>0.
\]
The reference block $\sym{ReferenceBlock0}$, representing
$\sym{A}|_{\DKRange(\sym{Psource})}$ in source coordinates, must have spectrum
in $[\sym{betaEndpoint}-\sym{delta}/2,\sym{alphaEndpoint}+\sym{delta}/2]$.
For an isometric frame $\sym{E0}$ of the reference subspace, its residual
is $\sym{R}=\sym{Acurrent}\sym{E0}-\sym{E0}\sym{ReferenceBlock0}$
on the operator domain, as in \cref{eq:history-residual}.
Assume
$\dim\sym{CrossedSourceDefect}=\dim\sym{CrossedTargetDefect}$: the space of source
vectors orthogonal to the target and the space of target vectors orthogonal
to the source have equal Hilbert dimension.

If $\|\sym{H}\|_{\mathrm{op}}<\sym{delta}/2$ for a
bounded self-adjoint perturbation, or $\|\sym{R}\|_{\mathrm{op}}<\sym{delta}/2$
for the bounded residual on the reference subspace, then
$\sym{Theta}<\pi/4$. For each ideal norm whose displayed arguments belong
to its domain, the corresponding double-angle estimate also holds:
\[
 \sym{delta}\sym{NormN}(\sin(2\sym{Theta}))\le2\sym{NormN}(\sym{H}),
 \qquad\text{or}\qquad
 \sym{delta}\sym{NormN}(\sin(2\sym{Theta0}))\le2\sym{NormN}(\sym{R}).
\]
The perturbation estimate uses the full angle, and the residual estimate
the directed angle on source coordinates. In the unbounded residual form,
$\sym{A}$ and $\sym{Acurrent}$ are self-adjoint on a common domain; only
the residual is required to be bounded.

We formalized both theorems in Hilbert space using spectral calculus.

\subsection{Approximation numbers and operator-ideal norms}

In finite dimensions, a unitarily invariant norm is a symmetric gauge
function of the singular values. Approximation numbers extend the
singular-value comparisons needed for bounded operators; for compact
operators they give the usual singular-value sequence. We formalized
symmetric norming functions and Ky Fan partial sums to pass from
majorization to norm inequalities.

For example, let $\sym{Rrotation}$ be the direct rotation and
$\sym{Wcompetitor}$ any unitary carrying $\sym{Usource}$ onto
$\sym{Vtarget}$. The restricted displacement of the direct rotation has
singular-value minimality, which implies minimality for unitarily invariant
norms. The squared-displacement result compares
$(I-\sym{Rrotation})^*(I-\sym{Rrotation})$ and
$(I-\sym{Wcompetitor})^*(I-\sym{Wcompetitor})$
\citep[Proposition~4.3]{davis1970rotation}. Their singular values are the
squares of those of the unsquared displacements. Section~\ref{sec:prop44}
shows why the claimed trace-norm comparison fails and derives the valid
Q-norm conclusion.
% Additional formalized quantitative results, omitted from this article:
% DavisKahan/FiniteDimensional/DirectRotation/SharpDisplacement.lean and
% SharpDisplacementLimit.lean. AUTHORING_NOTES.md indexes the evidence.

The same norm inequalities enter the sine-theta proof. The assertion
$\sym{delta}\sym{NormN}(\sym{S})\leq\sym{NormN}(\sym{R})$ applies when
$\sym{S}$ and $\sym{R}$ have finite ideal norm. Our ideal-membership
hypotheses state this condition explicitly in arbitrary Hilbert dimension.

\subsection{Sylvester estimates and unbounded operators}

For operators $\sym{SylvesterLeft}$ and $\sym{SylvesterRight}$ on the target and source spaces of $\sym{X}$, the
Sylvester map is $\sym{X}\mapsto \sym{SylvesterLeft}\sym{X}-\sym{X}\sym{SylvesterRight}$. Under the interval/exterior or ordered
spectral-separation hypotheses, the estimate has the form
\[
 \sym{delta} \sym{NormN}(\sym{X})\leq \sym{NormN}(\sym{SylvesterLeft}\sym{X}-\sym{X}\sym{SylvesterRight}).
\]
The sine-theta proof uses $\sym{SylvesterLeft}=\sym{Lambda1}$, $\sym{SylvesterRight}=\sym{A0}$, and
$\sym{X}=\like{F1}{F_1^*}\sym{E0}$. Other perturbation estimates choose related
projection blocks for $\sym{X}$.

For unbounded $\sym{SylvesterLeft}$ and $\sym{SylvesterRight}$, the identity is asserted for
$x\in\DKDom(\sym{SylvesterRight})$, with $\sym{X}x\in\DKDom(\sym{SylvesterLeft})$, and $\sym{SylvesterLeft}\sym{X}-\sym{X}\sym{SylvesterRight}$ agrees there with a
bounded residual. We formalized these domain conditions using linear
partial maps. Spectral cutoffs give bounded approximations; convergence
of approximation numbers passes the norm estimates to the unbounded
operators \citep[Section~5 and Appendix to Section~6]{davis1970rotation}.

The generalized sine and tangent theorems also compare block operators.
Reflection pinching replaces a block matrix $T$ by
$(T+\sym{Jreflect}T\sym{Jreflect})/2$, where $\sym{Jreflect}$ acts as $+I$ on one coordinate subspace and $-I$
on the other. This removes the off-diagonal blocks and contracts a unitarily
invariant norm. Direct-sum comparisons relate the norms of the component
blocks to the whole operator. Together with Sylvester separation these give
the generalized angle estimates
\citep[Lemmas~6.1--6.2 and Theorems~6.1--6.3]{davis1970rotation}.

\subsection{Real and complex Hilbert spaces}

The spectral constructions use complex functional calculus. We formalized
the scalar-transfer results needed to carry these arguments to real Hilbert
spaces through complexification. This permits the source theorems to be
stated over either field, including their domain and ideal-norm hypotheses.
The underlying algebra, topology, and continuous functional calculus come
from Mathlib; the adapted Spectra constructions and later Tau Ceti
contributions are described in Section~\ref{sec:related-formal}.

\input{counterexample.tex}

\input{formalization_enabled_extension.tex}

\section{Discussion}
\label{sec:discussion}

The formalization covers the original perturbation theory on real and
complex Hilbert spaces, including unbounded self-adjoint operators and
operator-ideal norms. The unbounded statements require explicit domain
conditions and bounded residuals. The operator-ideal statements require
finite norm values. These hypotheses determine which approximation problems
the estimates address.

The correction of the full-displacement claim
\citep[Proposition~4.4]{davis1970rotation} separates two extremal statements.
The valid comparison \citep[Proposition~4.3]{davis1970rotation}
gives squared-displacement minimality; its Q-norm consequence
is valid. Our real four-dimensional family exhibits full trace-norm failure
at every prescribed positive acute angle. In each example, the competitor has
two nonzero displacement singular values while the direct rotation has four.
The sums and squared sums of these values account for their different norm
comparisons. The earlier complex example supplied the construction; our
calculation makes its contradiction of the printed real claim explicit.

The matching theorem addresses the resulting optimization problem for any
equal-dimensional pair of finite-dimensional real subspaces. Its proof uses
principal-angle geometry, real orthogonal block decompositions, and
singular-value norms. These are also among the foundations required for the
Davis--Kahan formalization. The proof of the matching theorem uses additional
optimization arguments, including a passage to the trace-norm endpoint.

\section{Conclusion}
\label{sec:conclusion}

We formalized 28 of the 29 selected Davis--Kahan results and refuted the
remaining printed statement \citep[Proposition~4.4]{davis1970rotation}.
The sine-theta theorem includes
unbounded self-adjoint operators and symmetric operator-ideal norms with
their domain and finiteness assumptions. The direct-rotation calculation
identifies the failed singular-value comparison and gives an explicit real
counterexample and a valid Q-norm consequence.

We also formalized the operator theory used in these results: principal
angles and direct rotations, Borel spectral calculus, approximation numbers
and symmetric ideals, and Sylvester separation. Parts of this development
were subsequently adapted into Tau Ceti. Using these foundations, we proved
an exact weighted-matching formula for the real trace-norm transport
minimum.

\section*{AI assistance disclosure}
ChatGPT and Claude assisted with the formalization and manuscript preparation.
Formalizations were carried out by LLMs with human oversight.
The new matching result was derived via GPT 6.1 sol and reviewed by the human
  authors.

\section*{Code and data availability}
\ifDKAnonymous
Artifact information is omitted for anonymous review.
\else
The formalization, result-to-Lean mapping, provenance records, and manuscript
sources are available at
\url{https://github.com/AIQ-Kitware/aiq-dkps-formalization}.
\fi

\ifDKAnonymous\else
\section*{Acknowledgments}

This material is based upon work supported by the Defense Advanced Research
Project Agency (DARPA) under Contract No. HR001125CE017. Any opinions, findings
and conclusions or recommendations expressed in this material are those of the
author(s) and do not necessarily reflect the views of the Defense Advanced
Research Project Agency (DARPA).
\fi

\bibliographystyle{unsrtnat}
\bibliography{references}

\clearpage
\appendix

\input{appendix_notation.tex}

\section{Reading the Lean signatures}
\label{app:lean-guide}
\label{app:reading-lean}

A Lean theorem declaration has the form
\[
  \texttt{theorem name (parameters) (hypotheses) : conclusion := proof}.
\]
Parameters introduce the mathematical objects under discussion.  A hypothesis
such as \texttt{(hA : IsSelfAdjoint A)} supplies the corresponding fact;
bundled predicates such as \texttt{IsTrialResidual} and
\texttt{IsExactSpectralDecomposition} collect related assumptions.  Braces mark
arguments normally inferred by Lean, while square brackets request typeclass
instances supplying ambient algebraic, analytic, or topological structure.  The proposition after the final colon is the theorem's conclusion
\citep{demoura2021lean4}.

In \Cref{lst:dk-sintheta-headline}, \texttt{RCLike} is the Mathlib abstraction
covering $\mathbb{R}$ and $\mathbb{C}$.  A type
\texttt{ContinuousLinearMap K F E} is a continuous linear map.  A type
\texttt{LinearPMap K E E} is a linear partial map whose
domain is a submodule; this is the representation used for possibly unbounded
operators.  \texttt{N.Mem T} states that the ideal gauge associated with $\sym{NormN}$ is
finite on $T$, and \texttt{N.gaugeReal T} is the corresponding real value.  The
suffix \texttt{Real} refers to the codomain of the gauge and not to the scalar
field.

\input{appendix_process.tex}

\end{document}
, without editing.
% EVERY \sym, \like, and \dkrole obeys this switch, including the glossary.
% No color is required to identify a mathematical object: the glyphs, indices,
% definitions and adjacent prose also distinguish them.
%
% Families (base hue): exact/target=blue, trial/source=amber, exact unwanted=
% cyan, angles and their operators=magenta, residual/displacement=green,
% gaps=violet. Norm gauges=slate; provenance retains brown/steel/vermilion.
% P_ex is the exact-target projection in sine-theta; P is the source
% projection in transport. R denotes a residual and D a direct rotation.
% Semantic keys retain their roles even when the printed notation changes.
\newif\ifDKMathColors
\DKMathColorstrue
\ifdefined\DKMonochromeMath
  \DKMathColorsfalse
\fi

% Family anchors (print-compatible versions of the Manim light palette).
\definecolor{dkfg}{HTML}{1B1F24}
\definecolor{dkmuted}{HTML}{6B7280}
\definecolor{dkwanted}{HTML}{1F6FD1}
\definecolor{dktrial}{HTML}{A66705}
\definecolor{dkunwanted}{HTML}{007C91}
\definecolor{dksine}{HTML}{D12A7B}
\definecolor{dkresid}{HTML}{087D43}
\definecolor{dkgap}{HTML}{6E4BD8}
\definecolor{dkold}{HTML}{705519}
\definecolor{dkcurrent}{HTML}{054669}
\definecolor{dkperturb}{HTML}{BB491B}
\definecolor{dkrefuted}{HTML}{C62828}

% Related yet distinguishable shades. These are NOT new unrelated categories:
% objects inherit their parent family; a different shade signals a different
% mathematical object (e.g., a frame versus its subspace or its projection).
\definecolor{dkwantedframe}{HTML}{1254A9}
\definecolor{dkwantedprojection}{HTML}{08618D}
\definecolor{dkwantedrotation}{HTML}{4248B0}
\definecolor{dkwantedspectrum}{HTML}{345F8F}
\definecolor{dkwantedvector}{HTML}{1468AA}
\definecolor{dktrialframe}{HTML}{905600}
\definecolor{dktrialprojection}{HTML}{79511B}
\definecolor{dktrialcompetitor}{HTML}{A04D10}
\definecolor{dktrialspectrum}{HTML}{79591C}
\definecolor{dktrialvector}{HTML}{9B522A}
\definecolor{dkunwantedframe}{HTML}{006477}
\definecolor{dkunwantedspectrum}{HTML}{00546B}
\definecolor{dkunwantedprojection}{HTML}{145B6D}
\definecolor{dkunwantedvector}{HTML}{256879}
\definecolor{dkangleoperator}{HTML}{923478}
\definecolor{dksineoperator}{HTML}{AA275A}
\definecolor{dksineblock}{HTML}{B43F70}
\definecolor{dksylvester}{HTML}{804074}
\definecolor{dkanglemax}{HTML}{B32185}
\definecolor{dkanglecos}{HTML}{824276}
\definecolor{dkresidvector}{HTML}{086038}
\definecolor{dkdisplacement}{HTML}{226C50}
\definecolor{dkgaporiginal}{HTML}{4E4D89}
\definecolor{dkgapstat}{HTML}{7039A7}
\definecolor{dkgapinterval}{HTML}{574A82}
\definecolor{dkambient}{HTML}{36516A}
\definecolor{dknorm}{HTML}{475875}
\definecolor{dknormaux}{HTML}{66526E}
\definecolor{dkintertwiner}{HTML}{4A6280}
\definecolor{dkturn}{HTML}{47647D}
\definecolor{dkmix}{HTML}{505B65}

% Coloring primitives; all descendants (including indices and adjoints) can
% be colored as ONE semantic expression by writing \like{key}{expression}.
\newcommand{\dkrole}[2]{\ifDKMathColors\textcolor{dk#1}{#2}\else#2\fi}
% An optional plain-text reminder supplies definitions and types on hover,
% while the final argument stays short enough for the printed notation table.
\newcommand{\dkDeclareSymbol}[5][]{%
  \expandafter\def\csname dkgloss@#2\endcsname{#5}%
  \if\relax\detokenize{#1}\relax
    \expandafter\pdfstringdef\csname dkpdfgloss@#2\endcsname{#5}%
  \else
    \expandafter\pdfstringdef\csname dkpdfgloss@#2\endcsname{\detokenize{#1}}%
  \fi
  \expandafter\def\csname dkglyph@#2\endcsname{#3}%
  \expandafter\def\csname dkrole@#2\endcsname{#4}%
  \expandafter\def\csname dksym@#2\endcsname{\dkrole{#4}{#3}}%
}
\DeclareRobustCommand{\sym}[1]{%
  \ifcsname dksym@#1\endcsname
    \dkSymbolTooltip{#1}{\csname dksym@#1\endcsname}%
  \else
    \PackageError{dk-math-colors}{Unknown symbol key #1}{See math_symbol_colors.tex}%
  \fi
}
\DeclareRobustCommand{\like}[2]{%
  \ifcsname dkrole@#1\endcsname
    \dkSymbolTooltip{#1}{\dkrole{\csname dkrole@#1\endcsname}{#2}}%
  \else
    \PackageError{dk-math-colors}{Unknown role key #1}{See math_symbol_colors.tex}%
  \fi
}

% A short description belongs to the semantic key, not the bare letter.
% The same descriptions are available to the notation appendix.
\DeclareRobustCommand{\dkSymbolDescription}[1]{\csname dkgloss@#1\endcsname}
\newif\ifDKSymbolTooltips
\DKSymbolTooltipstrue
\ifdefined\DKNoSymbolTooltips\DKSymbolTooltipsfalse\fi
\newsavebox{\DKTooltipBox}
\newcount\DKTooltipCount
% Register only shipped widgets through hyperref's existing form-field
% machinery. Poppler/Okular need the catalog's field list, not just page
% annotations. The deferred aux write also discards measuring-only copies.
\makeatletter
\newcommand{\dkRegisterTooltipField}{%
  \expandafter\HyField@@AddToFields\expandafter{\number\pdffeedback lastannot}%
}
\makeatother
\AtBeginDocument{\ifDKSymbolTooltips\Form[NeedAppearances=false]\fi}
\AtEndDocument{\ifDKSymbolTooltips\endForm\fi}
% Native LuaLaTeX PDF widgets: no navigation action or JavaScript. Avoid
% pdfcomment's soulpos dependency, which previously broke local Vim builds.
% Measure in the CURRENT math style, then emit the original expression as-is:
% boxing the visible math would change spacing, scripts and line breaking.
\newcommand{\dkSymbolTooltip}[2]{%
  \ifDKSymbolTooltips
    \begingroup
      \DKSymbolTooltipsfalse % suppress nested annotations while measuring
      \ifmmode
        \edef\DKTooltipStyle{\ifcase\mathstyle
          \displaystyle\or\displaystyle\or\textstyle\or\textstyle\or
          \scriptstyle\or\scriptstyle\or\scriptscriptstyle\or\scriptscriptstyle\fi}%
        \sbox{\DKTooltipBox}{$\DKTooltipStyle#2$}%
      \else
        \leavevmode\sbox{\DKTooltipBox}{#2}%
      \fi
      \global\advance\DKTooltipCount by 1
      \pdfextension annot width\wd\DKTooltipBox
        height\ht\DKTooltipBox depth\dp\DKTooltipBox
        {/Subtype /Widget /FT /Btn /Ff 65536 /F 512
         /T (dk-symbol-\the\DKTooltipCount) /TU (\csname dkpdfgloss@#1\endcsname)
         /H /N /BS << /W 0 >> /Border [0 0 0]}%
      \dkRegisterTooltipField
      #2%
    \endgroup
  \else
    #2%
  \fi
}
\pdfstringdefDisableCommands{%
  \def\sym#1{\csname dkglyph@#1\endcsname}%
  \def\like#1#2{#2}%
}

% Declarations: [optional detailed hover reminder], key, glyph, role, short description.
% This registry includes local proof notation. The notation appendix selects
% the symbols useful for reference; registration does not add a printed row.
% Manim's original symbol keys, now split into semantic subshades.
% Ambient and perturbed operators (provenance is independent of geometry).
\dkDeclareSymbol[Reference self-adjoint operator A on the ambient Hilbert space; the perturbed operator is Ahat = A + H. In a single-operator eigenproblem, A itself is the operator being approximated. An unbounded A acts only on domain(A).]
  {A}{A}{old}{Reference self-adjoint operator before perturbation.}
\dkDeclareSymbol{A1}{A_1}{old}{Complementary block of the reference operator.}
\dkDeclareSymbol{a1}{a_1}{old}{Reference eigenvalue or complementary spectral value.}
\dkDeclareSymbol[Reference block A_0: the restriction of A to the range of the reducing source projection P, expressed in source coordinates. In Theorem 8.1 it lies at or below alpha; in Theorem 8.2 its spectrum lies in the enlarged central interval. The perturbed blocks are denoted Lambda_0 and Lambda_1.]
  {ReferenceBlock0}{A_0}{old}{Reference operator restricted to the source subspace in Section 8.}
\dkDeclareSymbol[Reference block A_1: the restriction of A to the kernel of the reducing source projection P, expressed in complementary coordinates. In Theorem 8.1 it lies at or above alpha + delta in quadratic-form order.]
  {ReferenceBlock1}{A_1}{old}{Reference operator restricted to the source complement in Section 8.}
\dkDeclareSymbol[Self-adjoint ambient operator Ahat being studied; in the perturbation setting Ahat = A + H. Its action on a trial vector is compared with the proposed eigenvalue times that vector.]
  {Ahat}{\widehat A}{current}{Operator under study; in the perturbation setting, the reference operator plus its perturbation.}
\dkDeclareSymbol[Self-adjoint ambient operator Ahat being studied. In the perturbation setting Ahat = A + H. The trial residual is R = Ahat E_0 - E_0 A_0, and the exact complementary block satisfies Ahat F_1 = F_1 Lambda_1 on its domain.]
  {Acurrent}{\widehat A}{current}{Operator under study; in the perturbation setting, the reference operator plus its perturbation.}
\dkDeclareSymbol[Additive self-adjoint perturbation H = Ahat - A. A bound on its norm can control the residual of an eigenspace of one operator relative to the other.]
  {H}{H}{perturb}{Additive self-adjoint perturbation of the reference operator.}
\dkDeclareSymbol{H0}{H_0}{perturb}{Wanted diagonal block of the perturbation.}
\dkDeclareSymbol{H1}{H_1}{perturb}{Unwanted diagonal block of the perturbation.}
\dkDeclareSymbol{B}{B}{perturb}{Off-diagonal perturbation or coupling block.}
\dkDeclareSymbol{b}{b}{perturb}{Scalar perturbation or coupling coefficient.}
% Trial/source family: subspace, frame, compressed operator, projector, vector.
\dkDeclareSymbol{U}{U}{wanted}{Desired exact subspace in the local geometric notation.}% Manim's U remains wanted in local slide notation.
\dkDeclareSymbol{V}{V}{trial}{Trial subspace in the local geometric notation.}% Manim's V remains trial in local slide notation.
\dkDeclareSymbol[Source subspace U to be carried onto the target V by an orthogonal or unitary map W. In the transport sections P projects onto U and Q onto V. The matching theorem also names the source subspace K.]
  {Usource}{U}{trial}{Source subspace to be transported onto the target subspace.}
\dkDeclareSymbol[Trial frame E_0: F -> E, a linear isometry from trial-coordinate space F into ambient space E. Its range is the trial subspace; E_0* returns trial coordinates and E_0* E_0 = I_F. In rank one, E_0(x) = x v.]
  {E0}{E_0}{trialframe}{Isometric embedding of the trial subspace; its adjoint returns trial coordinates.}
\dkDeclareSymbol[Self-adjoint trial operator A_0 on trial-coordinate space F. The residual compares Ahat E_0 with E_0 A_0. Rayleigh--Ritz chooses A_0 = E_0* Ahat E_0 when this compression is defined; the theorem allows other trial operators.]
  {A0}{A_0}{trialspectrum}{Self-adjoint operator on trial coordinates; its eigenvalues are the trial spectral values.}
\dkDeclareSymbol[Trial vector: a proposed approximate eigenvector. In the rank-one geometry v has unit length and E_0(x) = x v. Its residual for trial scalar rho is r = Ahat v - rho v.]
  {v}{v}{trialvector}{Trial vector or vector in the trial subspace.}
\dkDeclareSymbol[Scalar trial operator in rank one, A_0(x) = rho x. Figure 1 uses the Rayleigh value rho = <v, Ahat v> for unit v; then rho v is the projection of Ahat v onto the trial line and r is perpendicular to v.]
  {rho}{\rho}{trialspectrum}{Trial eigenvalue or scalar trial operator.}
\dkDeclareSymbol[Approximate eigenvalue assigned to a trial vector v. For the local operator A, its residual is A v - mu v. For unit v, the Rayleigh--Ritz choice is mu = <A v, v>.]
  {mu}{\mu}{trialspectrum}{Trial spectral value.}
\dkDeclareSymbol{a0}{a_0}{trialspectrum}{Wanted trial eigenvalue or trial spectral value.}
\dkDeclareSymbol[Ambient orthogonal or unitary transporter W with W(U) = V, equivalently W P W* = Q. Its full displacement is I - W on the whole ambient space. The real matching theorem minimizes its trace norm over every such orthogonal W.]
  {Wcompetitor}{W}{trialcompetitor}{Orthogonal or unitary competitor carrying the source subspace onto the target.}
\dkDeclareSymbol[Orthogonal projection P onto the source subspace U in the transport problem; P = P* = P squared. A feasible orthogonal or unitary W satisfies W P W* = Q. Restricted displacement (I - W)P tests only source vectors.]
  {Psource}{P}{trialprojection}{Orthogonal projection onto the source subspace in the transport problem.}
\dkDeclareSymbol{ProjSource}{P_U}{trialprojection}{Orthogonal projection onto the source subspace in the transport problem.}
\dkDeclareSymbol[Vector x in the source subspace U. For principal vectors, unit x_j in U and y_j in V are chosen with <x_i, y_j> = 0 for i different from j and <x_j, y_j> = cos(theta_j).]
  {SourceVector}{x}{trialvector}{Vector in the source subspace.}
\dkDeclareSymbol[Orthonormal coordinate frame X for the source subspace, so X* X = I. Together with a target frame Y, the singular values of X* Y are the principal cosines. In the four-dimensional counterexample X is the vertical stack of I_2 and 0.]
  {Xframe}{X}{trialframe}{Orthonormal coordinate matrix for the source subspace in the counterexample.}
\dkDeclareSymbol[Hilbert space F of trial coordinates. E_0 embeds F isometrically into ambient space E; A_0 and the positive sine operator sin(Theta_0) act on F, while the residual R maps F into E.]
  {Ftrial}{F}{trialframe}{Hilbert space of trial coordinates.}
\dkDeclareSymbol[Ambient Hilbert space E over the real or complex scalars. It contains the trial subspace range(E_0) and the orthogonal exact subspaces range(F_0) and range(F_1). The source theorem permits infinite-dimensional separable E.]
  {Eambient}{E}{ambient}{Ambient Hilbert space containing the trial and exact subspaces.}
% Exact/wanted family: the target subspace, its frame, projection and rotation.
\dkDeclareSymbol[Target subspace V in the transport problem: feasible maps satisfy W(U) = V, or equivalently W P W* = Q. The matching theorem also names this subspace L.]
  {Vtarget}{V}{wanted}{Target subspace in the transport problem.}
\dkDeclareSymbol[Exact wanted frame F_0: F_ex -> E, a linear isometry onto the desired exact spectral subspace. F_0* returns its coordinates. P_ex = F_0 F_0* is its ambient orthogonal projection; in Figure 1, F_0(h) = (h, 0).]
  {F0}{F_0}{wantedframe}{Isometric coordinates for the desired exact spectral subspace; star denotes the adjoint.}
\dkDeclareSymbol[Self-adjoint operator block on coordinates for the desired exact spectral subspace. Its spectrum consists of the wanted exact spectral values; in the corresponding decomposition Ahat F_0 = F_0 Lambda_0 on the operator domain.]
  {Lambda0}{\Lambda_0}{wantedspectrum}{Self-adjoint block on the desired exact spectral subspace.}
\dkDeclareSymbol{u}{u}{wantedvector}{Vector in the desired exact subspace.}
\dkDeclareSymbol{f0}{f_0}{wantedvector}{Component or coordinate vector in the desired exact subspace.}
\dkDeclareSymbol{lambda0}{\lambda_0}{wantedspectrum}{Wanted exact eigenvalue or wanted spectral value.}
\dkDeclareSymbol{lambdau}{\lambda_u}{wantedspectrum}{Eigenvalue associated with a desired exact vector.}
\dkDeclareSymbol[Exact wanted orthogonal projection P_ex: E -> E, with P_ex = F_0 F_0*. It retains the desired exact spectral component. Its complementary projection is P_ex_perp = I - P_ex = F_1 F_1*; the trial leakage is S = P_ex_perp E_0.]
  {Pexact}{P_{\mathrm{ex}}}{wantedprojection}{Orthogonal projection onto the desired exact subspace in the sine-theta theorem.}
\dkDeclareSymbol[Exact complementary orthogonal projection P_ex_perp = I - P_ex = F_1 F_1*: E -> E. It retains the unwanted exact component. In Figure 1, P_ex_perp v = S(1), whereas P_ex_perp r is the unwanted component of the residual.]
  {PexactPerp}{P_{\mathrm{ex}}^\perp}{unwantedprojection}{Orthogonal projection onto the exact complementary subspace.}
\dkDeclareSymbol{ReflectionSigma}{\Sigma}{wantedprojection}{Reflection associated with the desired exact spectral projection.}
\dkDeclareSymbol[Orthogonal projection Q onto the target subspace V in the transport problem. A feasible orthogonal or unitary W satisfies W P W* = Q, where P projects onto the source U.]
  {Qtarget}{Q}{wantedprojection}{Orthogonal projection onto the target subspace in the transport problem.}
\dkDeclareSymbol{ProjTarget}{P_V}{wantedprojection}{Orthogonal projection onto the target subspace in the transport problem.}
\dkDeclareSymbol[Canonical direct rotation D from the source U to the target V for an acute pair. It is the unitary polar factor of T = QP + (I - Q)(I - P), and acts by the principal angle on each principal plane. R is reserved for the trial residual.]
  {Rrotation}{D}{wantedrotation}{Canonical direct rotation carrying the source subspace onto the target.}
\dkDeclareSymbol[Direct rotation D between the ranges of orthogonal projections P and Q, including the nonacute case. Relative to range(P) and kernel(P), its diagonal blocks are positive and its off-diagonal blocks are opposite adjoints. It exists when the two crossed defect spaces have equal dimension, and is unique when both defects vanish.]
  {SourceDirectRotation}{D}{wantedrotation}{Direct rotation, including the nonacute case in the source-result table.}
\dkDeclareSymbol[Orthonormal coordinate frame Y for the target subspace, so Y* Y = I. In the counterexample Y is the vertical stack of cos(theta) I_2 and sin(theta) J_0; hence X* Y = cos(theta) I_2 and both principal angles equal theta.]
  {Yframe}{Y}{wantedframe}{Orthonormal coordinate matrix for the target subspace in the counterexample.}
\dkDeclareSymbol[Vector y in the target subspace V. For principal vectors, unit x_j in U and y_j in V are paired so that <x_j, y_j> = cos(theta_j); distinct pairs have zero cross inner product.]
  {TargetVector}{y}{wantedvector}{Vector in the target subspace.}
\dkDeclareSymbol[Hilbert space F_ex of exact wanted coordinates. The isometry F_0: F_ex -> E identifies it with the desired exact spectral subspace. This space differs from H, the additive perturbation.]
  {Hexact}{F_{\mathrm{ex}}}{wantedframe}{Hilbert space of coordinates for the desired exact spectral subspace.}% Coordinate Hilbert space (not perturbation).
% Exact/unwanted family; the trial error operators are MAGENTA, not cyan.
\dkDeclareSymbol{Uperp}{U^{\perp}}{unwanted}{Orthogonal complement of the desired exact subspace in the local notation.}
\dkDeclareSymbol[Exact complementary frame F_1: G -> E, a linear isometry onto the orthogonal complement of range(F_0). F_1* returns unwanted coordinates; F_1 F_1* = I - P_ex. In Figure 1, F_1(g) = (0, g).]
  {F1}{F_1}{unwantedframe}{Isometric coordinates for the exact complementary subspace; star denotes the adjoint.}
\dkDeclareSymbol[Self-adjoint unwanted block Lambda_1 on G, satisfying Ahat F_1 = F_1 Lambda_1 on its domain. Its spectrum is separated from the trial spectrum of A_0 by the specified positive delta. Figure 1 uses the scalar block 0.7.]
  {Lambda1}{\Lambda_1}{unwantedspectrum}{Self-adjoint operator on exact complementary coordinates: the unwanted spectral block.}
% Spectral measure and set (distinct from the ambient space E and sine block S).
\dkDeclareSymbol[Projection-valued spectral measure E associated with a self-adjoint operator A: A = integral lambda dE(lambda). For a Borel set S, E(S) is the orthogonal projection onto the corresponding spectral subspace.]
  {SpectralMeasure}{E}{wantedprojection}{Projection-valued spectral measure of a self-adjoint operator.}
\dkDeclareSymbol[Borel subset S of the real spectral axis. The spectral projection E(S) selects the part of a self-adjoint operator with spectral values in S; its range is a reducing subspace.]
  {SpectralSet}{S}{wantedspectrum}{Borel spectral set used to select a spectral subspace.}
% Generic Sylvester coefficients and right side, by source/target role.
\dkDeclareSymbol[Left coefficient B of the Sylvester map X -> B X - X C; B acts on the output space of X. For the sine bound, B = Lambda_1 on the unwanted-coordinate space G.]
  {SylvesterLeft}{B}{unwantedspectrum}{Left coefficient of the Sylvester equation; corresponds to the unwanted spectral block.}
\dkDeclareSymbol[Right coefficient C of the Sylvester map X -> B X - X C; C acts on the input space of X. For the sine bound, C = A_0 on trial space F. The historical equation also denotes this coefficient by A.]
  {SylvesterRight}{C}{trialspectrum}{Right coefficient of the Sylvester equation; corresponds to the trial operator.}
\dkDeclareSymbol[Right-hand side Y in B X - X C = Y, with the same input and output spaces as X. For the sine bound, Y = F_1* R: F -> G, the residual expressed in unwanted exact coordinates.]
  {SylvesterRHS}{Y}{resid}{Right-hand side of the Sylvester equation, obtained by projecting the residual.}
\dkDeclareSymbol{w}{w}{unwantedvector}{Vector in an unwanted or complementary spectral subspace.}
\dkDeclareSymbol{f1}{f_1}{unwantedvector}{Component or coordinate vector in the unwanted exact subspace.}
\dkDeclareSymbol{lambda1}{\lambda_1}{unwantedspectrum}{Unwanted exact eigenvalue or complementary spectral value.}
\dkDeclareSymbol{lambdaw}{\lambda_w}{unwantedspectrum}{Eigenvalue associated with an unwanted exact vector.}
\dkDeclareSymbol[Hilbert space G of exact complementary coordinates. F_1: G -> E embeds it onto the unwanted exact subspace; Lambda_1 acts on G and X = F_1* E_0 maps trial coordinates into G.]
  {Gunwanted}{G}{unwantedframe}{Hilbert space of exact complementary coordinates.}
\dkDeclareSymbol{E1}{E_1}{unwantedframe}{Isometric coordinates for the complementary trial subspace.}
\dkDeclareSymbol{Vperp}{V^{\perp}}{unwanted}{Orthogonal complement of the trial subspace in the local notation.}
% Angle and error geometry, including distinct rectangular and positive maps.
\dkDeclareSymbol[Individual principal angle, between 0 and pi/2. For orthonormal subspace frames X and Y, the singular values of X* Y are the principal cosines. In Figure 1, theta is the angle between the unit trial vector v and the exact wanted line.]
  {theta}{\theta}{sine}{Principal angle between the subspaces under comparison.}
\dkDeclareSymbol[Full positive angle operator Theta = diag(Theta_0, Theta_1) for the source/target decompositions. Theta_0 compares the source with the target; Theta_1 compares their orthogonal complements. Its operator norm is the largest spectral angle. For positive c, the conditions Theta <= c and Theta < c mean operator norm at most c and strictly below c, respectively.]
  {Theta}{\Theta}{angleoperator}{Full angle operator, the diagonal sum of the source and complementary angle blocks.}
\dkDeclareSymbol[Directed principal-angle operator Theta_0 on source or trial coordinates. For a trial frame E_0 and exact wanted projection P_ex, cos(Theta_0) = sqrt(E_0* P_ex E_0), and sin(Theta_0) = sqrt(I_F - E_0* P_ex E_0). In the two-projection geometry it comes from the positive cosine compression P Q P restricted to range(P).]
  {Theta0}{\Theta_0}{angleoperator}{Directed principal-angle operator acting on source or trial coordinates.}
\dkDeclareSymbol[Complementary directed angle operator Theta_1, obtained from the projections I-P and I-Q by the same cosine-compression construction as Theta_0. It acts on the source complement, or its isometric coordinates. The full angle operator is diag(Theta_0, Theta_1).]
  {Theta1}{\Theta_1}{angleoperator}{Directed angle operator comparing the orthogonal complements.}
\dkDeclareSymbol[Right-angle source defect D_0 = range(P) intersect kernel(Q). These are source vectors orthogonal to the target. The full-angle conclusion in the Section 8 smallness criterion assumes dim(D_0) = dim(D_1); in Lean their identification is CrossedDefectsEquivalent.]
  {CrossedSourceDefect}{D_0}{fg}{Source directions perpendicular to the target subspace.}
\dkDeclareSymbol[Right-angle target defect D_1 = kernel(P) intersect range(Q). These are target vectors orthogonal to the source. Equality of its Hilbert dimension with that of D_0 is assumed for the full-angle conclusion of Theorem 8.2.]
  {CrossedTargetDefect}{D_1}{fg}{Target directions perpendicular to the source subspace.}
\dkDeclareSymbol[Positive complementary cosine block C_1 = cos(Theta_1). The reference complement is identified with the perturbed complement by the direct rotation D, so A_1 and Lambda_1 act on common coordinates. Theorem 8.1 compares A_1 - alpha I with C_1 (Lambda_1 - alpha I) C_1 as quadratic forms.]
  {ComplementCosineBlock}{C_1}{anglecos}{Cosine of the complementary angle block in the spectral-repulsion inequality.}
\dkDeclareSymbol[Scalar sine of a principal angle. For a unit principal trial vector v and exact wanted projection P_ex, sin(theta) = ||(I - P_ex)v||. In Figure 1 it is the length of the pink leg S(1).]
  {sintheta}{\sin\theta}{sine}{Sine of an individual principal angle: the length of the complementary component of a unit principal trial vector.}
\dkDeclareSymbol[Positive sine operator sin(Theta_0) = sqrt(I_F - E_0* P_ex E_0): F -> F. It equals |S| for S = (I - P_ex)E_0: F -> E. Its eigenvalues in finite dimensions are the principal sines; the sine theorem bounds its ideal norm by N(R)/delta.]
  {sinTheta0}{\sin\Theta_0}{sineoperator}{Positive sine operator on trial coordinates; the modulus of the rectangular complementary error.}
\dkDeclareSymbol[Rectangular complementary error S = (I - P_ex) E_0: F -> E. It records how much each trial vector leaks outside the exact wanted subspace. Its modulus |S| = sqrt(S* S) is sin(Theta_0): F -> F. In rank one S(1) = P_ex_perp v.]
  {S}{S}{sineblock}{Rectangular complementary error: the trial frame projected onto the unwanted exact subspace.}
\dkDeclareSymbol[Complementary-coordinate map X = F_1* E_0: F -> G, with S = F_1 X. It satisfies Lambda_1 X - X A_0 = F_1* R on domain(A_0). Isometric multiplication gives N(X) = N(S); the separate modulus identity abs(S) = sin(Theta_0) gives N(S) = N(sin(Theta_0)), whenever these norms are finite.]
  {X}{X}{sylvester}{Solution of the Sylvester equation; for the sine bound, the trial frame in unwanted exact coordinates.}
\dkDeclareSymbol{alphaAngle}{\alpha}{anglemax}{Largest principal angle.}
\dkDeclareSymbol[Scalar c = cos(theta) in the counterexample. In the general two-projection decomposition, the same role denotes a positive cosine operator C; the cosine and sine blocks commute and satisfy C squared + S squared = I.]
  {CosAngle}{c}{anglecos}{Cosine of a principal angle or the corresponding positive cosine block.}
\dkDeclareSymbol[Scalar s = sin(theta) in the counterexample. In the general two-projection decomposition, the same role denotes a positive sine operator S; it commutes with the positive cosine block C and satisfies C squared + S squared = I.]
  {SinAngle}{s}{sineoperator}{Sine of a principal angle or the corresponding positive sine block.}
\dkDeclareSymbol[Positive small-angle parameter in the local subspace comparison. It controls the allowed angular deviation; it is an angle parameter rather than a perturbation-operator norm.]
  {EpsilonAngle}{\varepsilon}{anglemax}{Angle parameter controlling the maximum subspace deviation.}
\dkDeclareSymbol{PhiAngle}{\varphi}{angleoperator}{Auxiliary angle or angle operator in the two-subspace comparison.}
\dkDeclareSymbol{Bblock}{B}{sineblock}{Off-diagonal block in the two-subspace decomposition.}
\dkDeclareSymbol{Cblock}{C}{sylvester}{Complementary or Sylvester block in the two-subspace decomposition.}
\dkDeclareSymbol{Ablock}{A}{anglecos}{Positive cosine-related block in the two-subspace decomposition.}
% Residual/error-displacement family.
\dkDeclareSymbol[Residual vector r = Ahat v - rho v = R(1) for the rank-one frame E_0(x) = x v. Figure 1 draws it from rho v to Ahat v. For unit v and the Rayleigh value rho = <v, Ahat v>, <v, r> = 0; a general trial scalar need not give orthogonality.]
  {r}{r}{residvector}{Residual vector of an approximate eigenpair.}
\dkDeclareSymbol[Bounded residual R: F -> E, agreeing with Ahat E_0 - E_0 A_0 on domain(A_0). The domain condition requires E_0(domain(A_0)) inside domain(Ahat). R = 0 means that E_0 intertwines A_0 with Ahat; its norm measures the defect.]
  {R}{R}{resid}{Bounded residual of the trial eigenproblem, agreeing with the operator difference on the trial domain.}
\dkDeclareSymbol[Full displacement Delta_W = I - W on the ambient space. In the failed Ky Fan comparison W = W_(pi/4); Delta_W has singular values sqrt(2), sqrt(2), 0, 0. Multiplying by a principal-plane projection restricts its input to that plane.]
  {Kdisplacement}{\Delta_W}{displacement}{Full displacement of a transporter from the identity.}
\dkDeclareSymbol[Squared full displacement Dcal(W) = (I-W)* (I-W), a positive operator. Proposition 4.3 compares its unitarily invariant norm for a direct rotation D with that for every feasible transporter W.]
  {SquaredDisplacement}{\mathcal D}{displacement}{Squared full displacement as a function of the transporter.}
% Matching proof: the gradient and its spectral level are not a space or gap.
\dkDeclareSymbol[Skew-symmetric spectral gradient G of Phi_p. It commutes with W and represents first variations by d/dt Phi_p(W exp(tB)) at t = 0 equal to -tr(G B) for skew-symmetric B. Feasible stationarity gives J_src G J_src = -G.]
  {SpectralGradient}{G}{displacement}{Skew-symmetric spectral gradient of the displacement power sum in the matching proof.}
\dkDeclareSymbol[Positive spectral level g of the skew-symmetric gradient G: on this invariant subspace, -G squared = g squared I. Consequently G/g squares to -I and supplies the complex structure used in the real matching proof.]
  {GradientLevel}{g}{displacement}{Positive gradient level; the gradient squared is minus this level squared times the identity.}
\dkDeclareSymbol[Complex structure I_g = G/g on the positive gradient-level space V_g. It is a real isometry with square -I. Multiplication by a+ib means a times identity plus b times I_g; source reflection J_src reverses I_g and acts as conjugation.]
  {GradientComplexStructure}{\mathcal I_g}{displacement}{Complex multiplication constructed from a positive gradient level.}
\dkDeclareSymbol[Real skew-symmetric matrix B, with B transpose = -B, giving the orthogonal variation W exp(tB). If B commutes with the source projection P, this curve preserves the constraint W(U) = V.]
  {SkewDirection}{B}{ambient}{Skew-symmetric direction for a first variation of an orthogonal transporter.}
% Spectral separation family: gap and original gap are NOT identical.
\dkDeclareSymbol[Positive spectral separation in the stated configuration. In the sine theorem it separates the trial operator from the exact complementary block; Theorem 8.1 uses it for the ordered reference gap, and Theorem 8.2 for the perturbed central-band separation. Figure 1 uses delta = abs(rho - 0.7).]
  {delta}{\delta}{gap}{Positive spectral separation in the stated gap configuration.}
\dkDeclareSymbol{g}{g}{gaporiginal}{Original spectral gap of the reference operator; distinct from the certified sine-theta separation.}
\dkDeclareSymbol{Delta}{\Delta}{gapstat}{Population or statistical spectral gap.}
\dkDeclareSymbol[Spectral endpoint alpha. In the interval/exterior configuration, one spectrum lies between beta and alpha, the other at or below beta - delta or at or above alpha + delta. Theorem 8.2 places the selected perturbed block in the central interval. In Theorem 8.1, alpha bounds the lower reference block and alpha + delta bounds the upper block from below.]
  {alphaEndpoint}{\alpha}{gapinterval}{Upper endpoint of the central band, or lower-block cut in the ordered gap.}
\dkDeclareSymbol[Lower endpoint beta of the finite interval from beta to alpha in the gap hypothesis. One spectrum lies in that interval; the other lies at or below beta - delta or at or above alpha + delta.]
  {betaEndpoint}{\beta}{gapinterval}{Lower endpoint of the finite interval in the interval/exterior gap hypothesis.}
% Spectral/projector and auxiliary matrix objects, context-scoped.
\dkDeclareSymbol{P}{P}{fg}{Orthogonal projection; its source or exact-target role is specified locally.}% Neutral historical Manim key; prefer Pexact/Psource.
\dkDeclareSymbol{Q}{Q}{fg}{Orthogonal projection; its role is specified locally.}% Neutral historical Manim key; prefer Qtarget.
\dkDeclareSymbol[Two-projection intertwiner T = QP + (I - Q)(I - P), satisfying TP = QT and T* T = I - (P - Q) squared. For ||P - Q|| < 1 it is invertible; its unitary polar factor T |T| inverse is the canonical direct rotation.]
  {IntertwinerT}{T}{intertwiner}{Two-projection intertwiner; its unitary polar factor is the acute direct rotation.}
\dkDeclareSymbol[Planar quarter-turn J_0 = ((0, -1), (1, 0)), with J_0 squared = -I_2 and J_0 transpose = -J_0. It defines the rotation C = cos(2 theta) I_2 + sin(2 theta) J_0 used to construct the real competitor.]
  {Jturn}{J_0}{turn}{Planar quarter-turn, with square minus the identity.}
\dkDeclareSymbol[Partial isometry V_pol in the polar decomposition S = V_pol |S|, where |S| = sqrt(S* S). It is isometric on the orthogonal complement of kernel(S) and zero on kernel(S). For an invertible operator its polar factor is unitary.]
  {PolarU}{V_{\mathrm{pol}}}{intertwiner}{Partial isometry in the polar decomposition; unitary when the operator is invertible.}
\dkDeclareSymbol[Planar rotation C = cos(2 theta) I_2 + sin(2 theta) J_0 through twice the principal angle. The competitor W_theta is orthogonally conjugate to C direct-sum I_2, so it rotates one two-plane and fixes the other.]
  {Cturn}{C}{turn}{Planar rotation through twice the principal angle in the counterexample.}
\dkDeclareSymbol[Orthogonal mixing matrix H_mix = (1/sqrt(2)) times the block matrix ((I_2, I_2), (I_2, -I_2)). It constructs W_theta = H_mix (C direct-sum I_2) H_mix transpose.]
  {Hmix}{H_{\mathrm{mix}}}{mix}{Orthogonal mixing matrix in the counterexample.}
\dkDeclareSymbol[Source reflection J_src = 2P - I: it acts as +I on the source subspace and -I on its orthogonal complement, so its square is I. In the matching proof stationarity gives J_src G J_src = -G for the skew-symmetric spectral gradient G.]
  {Jreflect}{J_{\mathrm{src}}}{trialprojection}{Reflection equal to twice the source projection minus the identity.}
\dkDeclareSymbol[Skew angle generator J_ang for a direct rotation D = exp(J_ang Theta). On each principal plane it is the oriented quarter-turn; it vanishes on the common fixed directions. It commutes with the full angle operator Theta.]
  {AngleGenerator}{J_{\mathrm{ang}}}{turn}{Skew angle generator in the direct-rotation result summaries.}
\dkDeclareSymbol{Dblock}{D}{mix}{Auxiliary diagonal or block matrix in the counterexample.}
\dkDeclareSymbol[Orthogonal projection onto the first principal plane span(e_0, e_3) in the four-dimensional counterexample. For Delta_W = I - W_(pi/4), the Ky Fan 2-norm of Delta_W Omega_1 equals 2.]
  {OmegaOne}{\Omega_1}{turn}{First principal-plane projection in the failed Ky Fan comparison.}
\dkDeclareSymbol[Orthogonal projection onto the second principal plane span(e_1, e_2) in the four-dimensional counterexample. For Delta_W = I - W_(pi/4), the Ky Fan 2-norm of Delta_W Omega_2 equals 2, although the trace norm of Delta_W is only 2 sqrt(2).]
  {OmegaTwo}{\Omega_2}{intertwiner}{Second principal-plane projection in the failed Ky Fan comparison.}
\dkDeclareSymbol{Dplanar}{D}{wantedrotation}{Planar direct-rotation block.}
\dkDeclareSymbol{Hreflection}{H}{trialcompetitor}{Planar reflection carrying the source line to the target.}
% Unitarily invariant gauge family; N, M are NOT a norm versus its input.
\dkDeclareSymbol[Unitarily invariant norm N, or a symmetric operator-ideal norm on its domain of finiteness. It satisfies N(U T V) = N(T) for unitaries and N(B T) <= ||B||op N(T). In the sine theorem it is normalized by N(u v*) = ||u|| ||v|| and has Ky Fan dominance where both norms are finite. Both R and sin(Theta_0) must belong to its ideal.]
  {NormN}{N}{norm}{Unitarily invariant norm on its domain of finiteness; in the sine theorem, normalized by $N(uv^*)=\|u\|\|v\|$.}
\dkDeclareSymbol[Sine-theta separation G_delta(B,C): finite interval/exterior separation in either spectral order, or ordered separation B <= a I and C >= (a+delta) I, or the reverse. The ordered alternatives allow both spectra to be unbounded.]
  {SineGap}{\mathcal G_\delta}{gap}{Interval/exterior separation or either ordered separation, as in the sine-theta theorem.}
\dkDeclareSymbol[Auxiliary unitarily invariant norm M witnessing the Q-norm property N(T) squared = M(T* T). For the Frobenius norm N, one can take M to be the trace norm.]
  {NormM}{M}{normaux}{Auxiliary unitarily invariant norm used to define a Q-norm through the squared modulus.}
\dkDeclareSymbol[Q-norm property: a unitarily invariant norm N is a Q-norm when N(T) squared = M(T* T) for some unitarily invariant norm M. The Frobenius norm is an example, with M the trace norm; the repaired full-displacement theorem uses this norm class.]
  {NormQ}{Q}{normaux}{Q-norm: its square is a unitarily invariant norm of the squared modulus.}
\dkDeclareSymbol[Singleton trace cost s_i = min(4 sin(theta_i/2), 2) for angle label i. It transports one real principal plane independently, choosing a rotation of cost 4 sin(theta_i/2) or a reflection of cost 2.]
  {SingletonCost}{s}{norm}{Minimum trace displacement for one real principal plane: choose a rotation or a reflection.}
\dkDeclareSymbol[Paired trace cost e_ij = 4 sin((theta_i + theta_j)/2) for distinct angle labels i,j. It is attained by coupling their principal planes in a real four-dimensional block. The edge weight s_i + s_j - e_ij is the difference between the cost of two singletons and this paired cost.]
  {PairCost}{e}{norm}{Trace displacement cost of coupling two principal planes in a four-dimensional block.}
\dkDeclareSymbol[Scalar phase cost f_p(t) = (2 sin(t/2)) to the power p, for 0 <= t <= pi. A real planar rotation through t contributes two copies, so its displacement power cost is 2 f_p(t).]
  {MatchingPhaseCost}{f_p}{norm}{Displacement power cost for one scalar phase.}
\dkDeclareSymbol[Singleton power cost a_i(p) = min(2 f_p(theta_i), 2 to the power p). These are the rotation and reflection choices for one principal plane. At p=1 it equals the trace-cost weight s_i.]
  {SingletonPowerCost}{a}{norm}{Minimum displacement power cost for one angle.}
\dkDeclareSymbol[Paired power cost b_ij(p) = min over d <= tau <= h of 2(f_p(h+tau) + f_p(h-tau)), where h=(theta_i+theta_j)/2 and d=abs(theta_i-theta_j)/2. This is the positive-determinant pair minimum; unrestricted block optimization also compares two singletons. At p=1 it equals e_ij.]
  {PairPowerCost}{b}{norm}{Positive-determinant displacement power cost for a pair of angles.}
\dkDeclareSymbol[Matching M on the finite general graph of positive principal-angle labels. Each label belongs to at most one edge; unmatched labels use singleton transporters. Minimizing paired and singleton costs is equivalent to maximum-weight matching with edge weights g_ij = s_i + s_j - e_ij. The empty matching is allowed.]
  {Matching}{M}{normaux}{Matching of positive principal-angle labels: each label is used in at most one pair.}
\dkDeclareSymbol[Displacement power sum Phi_p(W) = sum_j s_j(I - W)^p. The matching proof first minimizes it for 1 < p < 2, then lets p decrease to 1, where it becomes the trace norm ||I - W||_1.]
  {PowerCost}{\Phi}{norm}{Sum of the p-th powers of displacement singular values.}
% Auxiliary matching-proof notation: neutral ink, with the same hover registry.
\dkDeclareSymbol[Off-diagonal gradient block G_01 mapping the source complement K perpendicular into K. Stationarity makes the diagonal blocks of G zero, so G = ((0,G_01),(-G_01 transpose,0)) in the source/complement decomposition.]
  {GradientOffDiagonal}{G_{01}}{fg}{Complement-to-source block of the skew-symmetric gradient.}
\dkDeclareSymbol[Orthogonal group O(n): real n by n matrices W satisfying W transpose times W = I. Its elements preserve Euclidean inner products and lengths. The notation O(2) refers to the same group in dimension two.]
  {OrthogonalGroup}{O(n)}{fg}{Group of real orthogonal matrices.}
\dkDeclareSymbol[Feasible transporters F = {W in O(n) : W(K) = L}. Equivalently W P_K W transpose = P_L. This nonempty compact set is the optimization domain throughout the matching proof.]
  {FeasibleTransport}{\mathcal F}{fg}{Orthogonal maps carrying the source subspace onto the target.}
\dkDeclareSymbol[Singular value s_j(T), the nonnegative square root of an eigenvalue of T transpose times T, counted with multiplicity. Summing s_j(I-W) gives the trace displacement; summing its p-th powers gives Phi_p(W). These singular values differ from the singleton weights s_i.]
  {TransportSingularValues}{s_j}{fg}{A displacement singular value in the matching proof.}
\dkDeclareSymbol[Scalar singular-value magnitude s between zero and two for the displacement I-W of an orthogonal matrix. Uniform convergence of s to the power p to s on this interval permits the trace-norm limit. This variable differs from the rotation parameter s used earlier in the proof.]
  {DisplacementSize}{s}{fg}{Scalar displacement singular-value magnitude in the uniform limit.}
\dkDeclareSymbol[All-singleton cost C_0 = sum_i s_i, before any two angle labels are paired. A matching has cost C_0 minus the sum of its selected edge weights.]
  {SingletonBaseline}{C_0}{fg}{Trace displacement when every angle uses its singleton transporter.}
\dkDeclareSymbol[Matching power cost C_p(M): sum of b_ij(p) over selected pairs plus a_i(p) over unmatched labels. Its minimum equals the minimum of Phi_p over feasible transporters. At p=1 it is C(M).]
  {MatchingPowerCost}{C_p}{fg}{Total paired and singleton power cost of a matching.}
\dkDeclareSymbol[Invariant orthogonal block V_a preserved by the source projection, target projection, and minimizing transporter. The displacement cost adds over these blocks, and their principal-angle lists combine into the original list.]
  {MatchingBlock}{V_a}{fg}{Simultaneous invariant block in the matching proof.}
\dkDeclareSymbol[Gradient-level space V_g: the eigenspace of -G squared for a positive eigenvalue g squared. It is preserved by G, W, and the source reflection J_src. Multiplication by i on it is G/g.]
  {GradientLevelSpace}{V_g}{fg}{Positive gradient-level eigenspace.}
\dkDeclareSymbol[Eigenphase varphi of W on V_g viewed as a complex space with multiplication by i = G/g. It lies between zero and pi and satisfies f_p prime(varphi) = g. There are at most two such phases.]
  {GradientEigenphase}{\varphi}{fg}{One of the at most two phases at a positive gradient level.}
\dkDeclareSymbol[Orthogonal projection E onto the first phase space of W on V_g. Here W = exp(i varphi_1) E + exp(i varphi_2) (I-E). In coordinates where the source reflection is conjugation, E = X + iY.]
  {GradientPhaseProjection}{E}{fg}{Projection onto the first phase space in a gradient-level space.}
\dkDeclareSymbol[Real part X of the phase projection E = X + iY. It is a real symmetric matrix with spectrum in the closed interval from zero to one, satisfying X squared - Y squared = X and XY + YX = Y. This X differs from the source-coordinate and Sylvester matrices elsewhere in the paper.]
  {ProjectionRealPart}{X}{fg}{Real symmetric part of the phase projection.}
\dkDeclareSymbol[Imaginary part Y of the phase projection E = X + iY. It is a real skew-symmetric matrix. It maps the lambda eigenspace of X to the 1-lambda eigenspace and constructs the small invariant planes.]
  {ProjectionImaginaryPart}{Y}{fg}{Real skew-symmetric imaginary part of the phase projection.}
\dkDeclareSymbol[Eigenvalue lambda in the closed interval from zero to one of the real part X of the phase projection. It is local to the block construction and differs from the original operator's spectral parameter.]
  {ProjectionEigenvalue}{\lambda}{fg}{Eigenvalue of the real part of the phase projection.}
\dkDeclareSymbol[Plane coupling eta = sqrt(lambda (1-lambda)) for 0 < lambda < 1. If u is a unit lambda eigenvector of X, then v = Yu/eta completes its invariant real plane.]
  {ProjectionCoupling}{\eta}{fg}{Normalization used to construct an invariant real plane.}
\dkDeclareSymbol[Pair mean angle h = (theta_i + theta_j)/2. Together with d = (theta_i - theta_j)/2 after ordering theta_i >= theta_j, it parametrizes the two-angle block cost.]
  {PairMeanAngle}{h}{fg}{Mean of the two ordered principal angles.}
\dkDeclareSymbol[Pair half-difference d = (theta_i - theta_j)/2 after ordering theta_i >= theta_j. It satisfies 0 <= d <= h and h+d <= pi/2.]
  {PairHalfDifference}{d}{fg}{Half the difference of the two ordered principal angles.}
\dkDeclareSymbol[Scalar pair parameter tau in the closed interval from d to h. The paired power cost is the minimum of 2(f_p(h+tau) + f_p(h-tau)) over this interval. At p=1 a minimum is attained at tau=h.]
  {PairPhaseParameter}{\tau}{fg}{Scalar parameter in the two-angle power-cost minimum.}
\dkDeclareSymbol[Attaining rotation angle gamma = arccos(cos(tau)/cos(d)). Choosing U_0 = U_1 = R_gamma realizes phase parameters x=h and y=tau in the four-dimensional block.]
  {PairAttainingAngle}{\gamma}{fg}{Rotation angle realizing a chosen scalar pair parameter.}
\dkDeclareSymbol[Four-dimensional direct rotation D = ((C,-S),(S,C)) in source and complementary principal coordinates, with C the diagonal principal cosines and S the diagonal principal sines. Every transporter on this block is D diag(U_0,U_1).]
  {PairDirectRotation}{D}{fg}{Direct rotation on the two principal planes.}
\dkDeclareSymbol[Diagonal cosine block C = diag(cos(theta_i), cos(theta_j)) in the four-dimensional direct rotation D. This matrix differs from the scalar displacement cost C(M).]
  {PairCosineBlock}{C}{fg}{Diagonal principal-cosine block.}
\dkDeclareSymbol[Diagonal sine block S = diag(sin(theta_i), sin(theta_j)) in the four-dimensional direct rotation D. This matrix is local to the matching proof.]
  {PairSineBlock}{S}{fg}{Diagonal principal-sine block.}
\dkDeclareSymbol[Free planar orthogonal matrices U_0 and U_1 acting on the two-dimensional source and its two-dimensional complement, respectively. Multiplying D by diag(U_0,U_1) parametrizes every transporter on this four-dimensional block.]
  {PairFreeOrthogonal}{U}{fg}{Free source or complementary planar orthogonal factor.}
\dkDeclareSymbol[Phase parameter x in the closed interval from h to pi-h for a positive-determinant four-dimensional transporter. Its eigenvalues are exp(plus or minus i(x+y)) and exp(plus or minus i(x-y)).]
  {PairPhaseX}{x}{fg}{First phase parameter of a four-dimensional transporter.}
\dkDeclareSymbol[Phase parameter y in the closed interval from d to pi-d for a positive-determinant four-dimensional transporter. Together with x it gives rotational phases x+y and x-y, folded into the closed interval from zero to pi for the displacement cost.]
  {PairPhaseY}{y}{fg}{Second phase parameter of a four-dimensional transporter.}
\dkDeclareSymbol[Scalar a = cos(x) in the four-dimensional phase calculation. Together with b = cos(y), it expresses the trace as 4ab and determines the reciprocal characteristic polynomial.]
  {PairPhaseCosineX}{a}{fg}{Cosine of the first phase parameter in the trace calculation.}
\dkDeclareSymbol[Scalar b = cos(y) in the four-dimensional phase calculation, paired with a = cos(x). These scalars are local to the characteristic-polynomial calculation.]
  {PairPhaseCosineY}{b}{fg}{Cosine of the second phase parameter in the trace calculation.}
\dkDeclareSymbol[Rotation parameter r: the two free orthogonal factors use planar rotation angles r+s and r-s. It appears in cos(y) = cos(d) cos(r) for two rotations.]
  {PairRotationR}{r}{fg}{Half the sum of the two free planar rotation angles.}
\dkDeclareSymbol[Rotation parameter s: the two free orthogonal factors use planar rotation angles r+s and r-s. It appears in cos(x) = cos(h) cos(s) for two rotations. This scalar differs from the singular values and singleton costs.]
  {PairRotationS}{s}{fg}{Half the difference of the two free planar rotation angles.}
\dkDeclareSymbol[Planar reflection J_2 = diag(1,-1). Every two-dimensional orthogonal matrix of determinant -1 can be written as a planar rotation times J_2.]
  {PairPlanarReflection}{J_2}{fg}{Fixed planar reflection used to parametrize determinant -1 factors.}
\dkDeclareSymbol[Real planar rotation R_t = ((cos(t),-sin(t)),(sin(t),cos(t))). Its two displacement singular values are both 2 times abs(sin(t/2)). In the matching proof a rotation through t in the closed interval from zero to pi contributes 2 f_p(t).]
  {MatchingPlanarRotation}{R_t}{fg}{Planar rotation through the indicated angle.}
% Statistical application / operator provenance.
\dkDeclareSymbol{Sigma}{\Sigma}{old}{Population covariance matrix.}
\dkDeclareSymbol{SigmaHat}{\widehat\Sigma}{current}{Estimated or perturbed covariance matrix.}
\dkDeclareSymbol{YWSV}{V}{wanted}{Orthonormal coordinates for the population eigenspace in the statistical application.}
\dkDeclareSymbol{YWSVhat}{\widehat V}{trialframe}{Estimated eigenspace frame in the statistical application.}

% Figure-specific objects retain their own meanings when bare letters overlap.
\dkDeclareSymbol[Signed rotation angle phi of the displayed map on its invariant real two-plane. At theta = pi/4 the direct rotation turns its two principal planes by opposite signed angles of magnitude pi/4; the competitor turns its moving plane by pi/2 and fixes its other plane.]
  {PlaneRotationAngle}{\varphi}{sine}{Signed rotation angle of the orthogonal map on the displayed invariant plane.}
\dkDeclareSymbol[Unit vector q in the displayed invariant plane. The map T is the direct rotation D or the competitor W in that panel; Tq is the image and (I - T)q = q - Tq is the displacement.]
  {FigureInput}{q}{ambient}{Unit input to the orthogonal map shown in this invariant-plane panel.}
\dkDeclareSymbol[Image Tq of the panel unit input q under T = D or W. Both q and Tq lie on the unit circle; the green displacement arrow joins Tq to q, so it represents q - Tq.]
  {FigureImage}{Tq}{wantedvector}{Image of the unit input under the orthogonal map shown in this panel.}
\dkDeclareSymbol[Displacement (I - T)q = q - Tq, drawn from the image Tq to the unit input q. On a plane rotated through phi its length is 2 |sin(phi/2)|; it vanishes on a fixed plane.]
  {DisplacementVector}{(I-T)q}{displacement}{Difference between a unit input and its image under the displayed orthogonal map.}
\dkDeclareSymbol[First principal plane Pi_1 = span(e_0, e_3) for the source-target pair in the real four-dimensional counterexample. At theta = pi/4 the direct rotation acts on this plane by a quarter of pi, with displacement singular value sqrt(2 - sqrt(2)).]
  {CounterPrincipalOne}{\Pi_1}{fg}{First principal plane in the real counterexample: the span of e0 and e3.}
\dkDeclareSymbol[Second principal plane Pi_2 = span(e_1, e_2) for the source-target pair in the real four-dimensional counterexample. The direct rotation acts by the same angle magnitude as on Pi_1, with the opposite sign in the displayed plane coordinates.]
  {CounterPrincipalTwo}{\Pi_2}{fg}{Second principal plane in the real counterexample: the span of e1 and e2.}
\dkDeclareSymbol[Moving plane M = range(W - I) = span(e_0 + e_2, e_1 + e_3) for the competitor. W rotates it by 2 theta; at theta = pi/4 the two displacement singular values contributed by this plane are both sqrt(2).]
  {CounterMovingPlane}{M}{ambient}{Moving plane of the competing transporter: the range of W minus the identity.}
\dkDeclareSymbol[Fixed plane Z = kernel(W - I) = span(e_0 - e_2, e_1 - e_3) for the competitor. W acts as the identity here, so this plane contributes two zero displacement singular values.]
  {CounterFixedPlane}{Z}{ambient}{Fixed plane of the competing transporter: the kernel of W minus the identity.}
\dkDeclareSymbol[Direct-rotation displacement singular value a = 2 sin(pi/8) = sqrt(2 - sqrt(2)) at theta = pi/4. Each of its two principal planes contributes two copies, giving trace displacement 4a.]
  {CounterChordA}{a}{displacement}{Direct-rotation displacement singular value in the counterexample: sqrt(2 - sqrt(2)).}
\dkDeclareSymbol[Competitor displacement singular value b = 2 sin(pi/4) = sqrt(2) at theta = pi/4. Its moving plane contributes two copies; the fixed plane contributes two zeros, giving trace displacement 2b.]
  {CounterChordB}{b}{displacement}{Competing-transporter displacement singular value in the counterexample: square root of 2.}
\dkDeclareSymbol[Singular value s_j(I - T) of the full ambient displacement, ordered from largest to smallest. Each rotating real two-plane contributes two values 2 |sin(phi/2)|; fixed directions contribute zeros. Their sum is the trace cost.]
  {DisplacementSingular}{s_j}{displacement}{An individual singular value of the full displacement from the identity.}
\dkDeclareSymbol[Edge weight g_ij = s_i + s_j - e_ij, the difference between the cost of two singletons and the paired cost. A matching has total cost C_0 minus the sum of its edge weights, with C_0 = sum_i s_i. Maximizing this sum gives the minimum transport cost; edges of negative weight may be omitted.]
  {MatchingEdgeWeight}{g_{ij}}{current}{Edge weight: the cost of two singletons minus their paired cost.}
\dkDeclareSymbol[Total trace-norm displacement cost. For a transporter W, C(W) = ||I - W||_1, the sum of its displacement singular values. For a matching M, C(M) is the sum of its paired and unmatched singleton costs.]
  {TraceCost}{C}{norm}{Trace-norm displacement cost of the indicated orthogonal transporter or matching.}

% Spectral values have their own key, distinct from an operator block.
\dkDeclareSymbol[Spectral parameter lambda of a self-adjoint operator. In a finite-dimensional eigenvalue list lambda_j(A) it denotes an individual eigenvalue; in A = integral lambda dE(lambda), it is the real integration variable.]{SpectralValue}{\lambda}{wantedspectrum}{Real spectral parameter or individual eigenvalue of a self-adjoint operator.}
 in paper.tex)
% or pass \def\DKMonochromeMath{1} before \documentclass[11pt]{article}
% PDF version used for the manuscript and its symbol annotations.
\pdfvariable minorversion=7\relax

\usepackage[margin=1in]{geometry}
\usepackage{amsmath,amssymb,amsthm,mathtools}
\usepackage{booktabs,tabularx,array,longtable}
\usepackage[section]{placeins}
\usepackage{float}
\usepackage{microtype}
\usepackage{graphicx}
\usepackage{pgf} % Live vector figures with LaTeX symbol labels and tooltips
% Editing citations: switch to the venue's numeric style at submission.
\usepackage[authoryear,round]{natbib}
\usepackage{xurl}
\usepackage[colorlinks=true,allcolors=blue]{hyperref}
\usepackage[nameinlink]{cleveref}
\usepackage[T1]{fontenc}
\usepackage{lmodern} % Scalable LaTeX lettering, shared with the Manim slides
\usepackage{xcolor}
% Draft controls. Set these independently for review, print, or arXiv versions.
\newif\ifDKAnonymous
\DKAnonymousfalse
\ifdefined\DKAnonymousReview\DKAnonymoustrue\fi
\newif\ifDKColoredSymbols
\DKColoredSymbolstrue
\ifdefined\DKMonochromeMath\DKColoredSymbolsfalse\fi
% The colors retain the Manim glossary's light-theme mathematical roles.
% Semantic, multi-shade mathematical notation for the Davis--Kahan paper.
% The six core families follow visualizations/dkvis/palette.py (LIGHT).
% Subshades distinguish objects that are RELATED but not IDENTICAL.
% Values were chosen for legibility against white manuscript paper; some
% amber, green, and vermilion anchors are a little darker than the slides.
%
% Color control (independent of hover reminders):
%   \DKMathColorsfalse   (after \input{math_symbol_colors.tex} in paper.tex)
% or pass \def\DKMonochromeMath{1} before \input{paper.tex}, without editing.
% EVERY \sym, \like, and \dkrole obeys this switch, including the glossary.
% No color is required to identify a mathematical object: the glyphs, indices,
% definitions and adjacent prose also distinguish them.
%
% Families (base hue): exact/target=blue, trial/source=amber, exact unwanted=
% cyan, angles and their operators=magenta, residual/displacement=green,
% gaps=violet. Norm gauges=slate; provenance retains brown/steel/vermilion.
% P_ex is the exact-target projection in sine-theta; P is the source
% projection in transport. R denotes a residual and D a direct rotation.
% Semantic keys retain their roles even when the printed notation changes.
\newif\ifDKMathColors
\DKMathColorstrue
\ifdefined\DKMonochromeMath
  \DKMathColorsfalse
\fi

% Family anchors (print-compatible versions of the Manim light palette).
\definecolor{dkfg}{HTML}{1B1F24}
\definecolor{dkmuted}{HTML}{6B7280}
\definecolor{dkwanted}{HTML}{1F6FD1}
\definecolor{dktrial}{HTML}{A66705}
\definecolor{dkunwanted}{HTML}{007C91}
\definecolor{dksine}{HTML}{D12A7B}
\definecolor{dkresid}{HTML}{087D43}
\definecolor{dkgap}{HTML}{6E4BD8}
\definecolor{dkold}{HTML}{705519}
\definecolor{dkcurrent}{HTML}{054669}
\definecolor{dkperturb}{HTML}{BB491B}
\definecolor{dkrefuted}{HTML}{C62828}

% Related yet distinguishable shades. These are NOT new unrelated categories:
% objects inherit their parent family; a different shade signals a different
% mathematical object (e.g., a frame versus its subspace or its projection).
\definecolor{dkwantedframe}{HTML}{1254A9}
\definecolor{dkwantedprojection}{HTML}{08618D}
\definecolor{dkwantedrotation}{HTML}{4248B0}
\definecolor{dkwantedspectrum}{HTML}{345F8F}
\definecolor{dkwantedvector}{HTML}{1468AA}
\definecolor{dktrialframe}{HTML}{905600}
\definecolor{dktrialprojection}{HTML}{79511B}
\definecolor{dktrialcompetitor}{HTML}{A04D10}
\definecolor{dktrialspectrum}{HTML}{79591C}
\definecolor{dktrialvector}{HTML}{9B522A}
\definecolor{dkunwantedframe}{HTML}{006477}
\definecolor{dkunwantedspectrum}{HTML}{00546B}
\definecolor{dkunwantedprojection}{HTML}{145B6D}
\definecolor{dkunwantedvector}{HTML}{256879}
\definecolor{dkangleoperator}{HTML}{923478}
\definecolor{dksineoperator}{HTML}{AA275A}
\definecolor{dksineblock}{HTML}{B43F70}
\definecolor{dksylvester}{HTML}{804074}
\definecolor{dkanglemax}{HTML}{B32185}
\definecolor{dkanglecos}{HTML}{824276}
\definecolor{dkresidvector}{HTML}{086038}
\definecolor{dkdisplacement}{HTML}{226C50}
\definecolor{dkgaporiginal}{HTML}{4E4D89}
\definecolor{dkgapstat}{HTML}{7039A7}
\definecolor{dkgapinterval}{HTML}{574A82}
\definecolor{dkambient}{HTML}{36516A}
\definecolor{dknorm}{HTML}{475875}
\definecolor{dknormaux}{HTML}{66526E}
\definecolor{dkintertwiner}{HTML}{4A6280}
\definecolor{dkturn}{HTML}{47647D}
\definecolor{dkmix}{HTML}{505B65}

% Coloring primitives; all descendants (including indices and adjoints) can
% be colored as ONE semantic expression by writing \like{key}{expression}.
\newcommand{\dkrole}[2]{\ifDKMathColors\textcolor{dk#1}{#2}\else#2\fi}
% An optional plain-text reminder supplies definitions and types on hover,
% while the final argument stays short enough for the printed notation table.
\newcommand{\dkDeclareSymbol}[5][]{%
  \expandafter\def\csname dkgloss@#2\endcsname{#5}%
  \if\relax\detokenize{#1}\relax
    \expandafter\pdfstringdef\csname dkpdfgloss@#2\endcsname{#5}%
  \else
    \expandafter\pdfstringdef\csname dkpdfgloss@#2\endcsname{\detokenize{#1}}%
  \fi
  \expandafter\def\csname dkglyph@#2\endcsname{#3}%
  \expandafter\def\csname dkrole@#2\endcsname{#4}%
  \expandafter\def\csname dksym@#2\endcsname{\dkrole{#4}{#3}}%
}
\DeclareRobustCommand{\sym}[1]{%
  \ifcsname dksym@#1\endcsname
    \dkSymbolTooltip{#1}{\csname dksym@#1\endcsname}%
  \else
    \PackageError{dk-math-colors}{Unknown symbol key #1}{See math_symbol_colors.tex}%
  \fi
}
\DeclareRobustCommand{\like}[2]{%
  \ifcsname dkrole@#1\endcsname
    \dkSymbolTooltip{#1}{\dkrole{\csname dkrole@#1\endcsname}{#2}}%
  \else
    \PackageError{dk-math-colors}{Unknown role key #1}{See math_symbol_colors.tex}%
  \fi
}

% A short description belongs to the semantic key, not the bare letter.
% The same descriptions are available to the notation appendix.
\DeclareRobustCommand{\dkSymbolDescription}[1]{\csname dkgloss@#1\endcsname}
\newif\ifDKSymbolTooltips
\DKSymbolTooltipstrue
\ifdefined\DKNoSymbolTooltips\DKSymbolTooltipsfalse\fi
\newsavebox{\DKTooltipBox}
\newcount\DKTooltipCount
% Register only shipped widgets through hyperref's existing form-field
% machinery. Poppler/Okular need the catalog's field list, not just page
% annotations. The deferred aux write also discards measuring-only copies.
\makeatletter
\newcommand{\dkRegisterTooltipField}{%
  \expandafter\HyField@@AddToFields\expandafter{\number\pdffeedback lastannot}%
}
\makeatother
\AtBeginDocument{\ifDKSymbolTooltips\Form[NeedAppearances=false]\fi}
\AtEndDocument{\ifDKSymbolTooltips\endForm\fi}
% Native LuaLaTeX PDF widgets: no navigation action or JavaScript. Avoid
% pdfcomment's soulpos dependency, which previously broke local Vim builds.
% Measure in the CURRENT math style, then emit the original expression as-is:
% boxing the visible math would change spacing, scripts and line breaking.
\newcommand{\dkSymbolTooltip}[2]{%
  \ifDKSymbolTooltips
    \begingroup
      \DKSymbolTooltipsfalse % suppress nested annotations while measuring
      \ifmmode
        \edef\DKTooltipStyle{\ifcase\mathstyle
          \displaystyle\or\displaystyle\or\textstyle\or\textstyle\or
          \scriptstyle\or\scriptstyle\or\scriptscriptstyle\or\scriptscriptstyle\fi}%
        \sbox{\DKTooltipBox}{$\DKTooltipStyle#2$}%
      \else
        \leavevmode\sbox{\DKTooltipBox}{#2}%
      \fi
      \global\advance\DKTooltipCount by 1
      \pdfextension annot width\wd\DKTooltipBox
        height\ht\DKTooltipBox depth\dp\DKTooltipBox
        {/Subtype /Widget /FT /Btn /Ff 65536 /F 512
         /T (dk-symbol-\the\DKTooltipCount) /TU (\csname dkpdfgloss@#1\endcsname)
         /H /N /BS << /W 0 >> /Border [0 0 0]}%
      \dkRegisterTooltipField
      #2%
    \endgroup
  \else
    #2%
  \fi
}
\pdfstringdefDisableCommands{%
  \def\sym#1{\csname dkglyph@#1\endcsname}%
  \def\like#1#2{#2}%
}

% Declarations: [optional detailed hover reminder], key, glyph, role, short description.
% This registry includes local proof notation. The notation appendix selects
% the symbols useful for reference; registration does not add a printed row.
% Manim's original symbol keys, now split into semantic subshades.
% Ambient and perturbed operators (provenance is independent of geometry).
\dkDeclareSymbol[Reference self-adjoint operator A on the ambient Hilbert space; the perturbed operator is Ahat = A + H. In a single-operator eigenproblem, A itself is the operator being approximated. An unbounded A acts only on domain(A).]
  {A}{A}{old}{Reference self-adjoint operator before perturbation.}
\dkDeclareSymbol{A1}{A_1}{old}{Complementary block of the reference operator.}
\dkDeclareSymbol{a1}{a_1}{old}{Reference eigenvalue or complementary spectral value.}
\dkDeclareSymbol[Reference block A_0: the restriction of A to the range of the reducing source projection P, expressed in source coordinates. In Theorem 8.1 it lies at or below alpha; in Theorem 8.2 its spectrum lies in the enlarged central interval. The perturbed blocks are denoted Lambda_0 and Lambda_1.]
  {ReferenceBlock0}{A_0}{old}{Reference operator restricted to the source subspace in Section 8.}
\dkDeclareSymbol[Reference block A_1: the restriction of A to the kernel of the reducing source projection P, expressed in complementary coordinates. In Theorem 8.1 it lies at or above alpha + delta in quadratic-form order.]
  {ReferenceBlock1}{A_1}{old}{Reference operator restricted to the source complement in Section 8.}
\dkDeclareSymbol[Self-adjoint ambient operator Ahat being studied; in the perturbation setting Ahat = A + H. Its action on a trial vector is compared with the proposed eigenvalue times that vector.]
  {Ahat}{\widehat A}{current}{Operator under study; in the perturbation setting, the reference operator plus its perturbation.}
\dkDeclareSymbol[Self-adjoint ambient operator Ahat being studied. In the perturbation setting Ahat = A + H. The trial residual is R = Ahat E_0 - E_0 A_0, and the exact complementary block satisfies Ahat F_1 = F_1 Lambda_1 on its domain.]
  {Acurrent}{\widehat A}{current}{Operator under study; in the perturbation setting, the reference operator plus its perturbation.}
\dkDeclareSymbol[Additive self-adjoint perturbation H = Ahat - A. A bound on its norm can control the residual of an eigenspace of one operator relative to the other.]
  {H}{H}{perturb}{Additive self-adjoint perturbation of the reference operator.}
\dkDeclareSymbol{H0}{H_0}{perturb}{Wanted diagonal block of the perturbation.}
\dkDeclareSymbol{H1}{H_1}{perturb}{Unwanted diagonal block of the perturbation.}
\dkDeclareSymbol{B}{B}{perturb}{Off-diagonal perturbation or coupling block.}
\dkDeclareSymbol{b}{b}{perturb}{Scalar perturbation or coupling coefficient.}
% Trial/source family: subspace, frame, compressed operator, projector, vector.
\dkDeclareSymbol{U}{U}{wanted}{Desired exact subspace in the local geometric notation.}% Manim's U remains wanted in local slide notation.
\dkDeclareSymbol{V}{V}{trial}{Trial subspace in the local geometric notation.}% Manim's V remains trial in local slide notation.
\dkDeclareSymbol[Source subspace U to be carried onto the target V by an orthogonal or unitary map W. In the transport sections P projects onto U and Q onto V. The matching theorem also names the source subspace K.]
  {Usource}{U}{trial}{Source subspace to be transported onto the target subspace.}
\dkDeclareSymbol[Trial frame E_0: F -> E, a linear isometry from trial-coordinate space F into ambient space E. Its range is the trial subspace; E_0* returns trial coordinates and E_0* E_0 = I_F. In rank one, E_0(x) = x v.]
  {E0}{E_0}{trialframe}{Isometric embedding of the trial subspace; its adjoint returns trial coordinates.}
\dkDeclareSymbol[Self-adjoint trial operator A_0 on trial-coordinate space F. The residual compares Ahat E_0 with E_0 A_0. Rayleigh--Ritz chooses A_0 = E_0* Ahat E_0 when this compression is defined; the theorem allows other trial operators.]
  {A0}{A_0}{trialspectrum}{Self-adjoint operator on trial coordinates; its eigenvalues are the trial spectral values.}
\dkDeclareSymbol[Trial vector: a proposed approximate eigenvector. In the rank-one geometry v has unit length and E_0(x) = x v. Its residual for trial scalar rho is r = Ahat v - rho v.]
  {v}{v}{trialvector}{Trial vector or vector in the trial subspace.}
\dkDeclareSymbol[Scalar trial operator in rank one, A_0(x) = rho x. Figure 1 uses the Rayleigh value rho = <v, Ahat v> for unit v; then rho v is the projection of Ahat v onto the trial line and r is perpendicular to v.]
  {rho}{\rho}{trialspectrum}{Trial eigenvalue or scalar trial operator.}
\dkDeclareSymbol[Approximate eigenvalue assigned to a trial vector v. For the local operator A, its residual is A v - mu v. For unit v, the Rayleigh--Ritz choice is mu = <A v, v>.]
  {mu}{\mu}{trialspectrum}{Trial spectral value.}
\dkDeclareSymbol{a0}{a_0}{trialspectrum}{Wanted trial eigenvalue or trial spectral value.}
\dkDeclareSymbol[Ambient orthogonal or unitary transporter W with W(U) = V, equivalently W P W* = Q. Its full displacement is I - W on the whole ambient space. The real matching theorem minimizes its trace norm over every such orthogonal W.]
  {Wcompetitor}{W}{trialcompetitor}{Orthogonal or unitary competitor carrying the source subspace onto the target.}
\dkDeclareSymbol[Orthogonal projection P onto the source subspace U in the transport problem; P = P* = P squared. A feasible orthogonal or unitary W satisfies W P W* = Q. Restricted displacement (I - W)P tests only source vectors.]
  {Psource}{P}{trialprojection}{Orthogonal projection onto the source subspace in the transport problem.}
\dkDeclareSymbol{ProjSource}{P_U}{trialprojection}{Orthogonal projection onto the source subspace in the transport problem.}
\dkDeclareSymbol[Vector x in the source subspace U. For principal vectors, unit x_j in U and y_j in V are chosen with <x_i, y_j> = 0 for i different from j and <x_j, y_j> = cos(theta_j).]
  {SourceVector}{x}{trialvector}{Vector in the source subspace.}
\dkDeclareSymbol[Orthonormal coordinate frame X for the source subspace, so X* X = I. Together with a target frame Y, the singular values of X* Y are the principal cosines. In the four-dimensional counterexample X is the vertical stack of I_2 and 0.]
  {Xframe}{X}{trialframe}{Orthonormal coordinate matrix for the source subspace in the counterexample.}
\dkDeclareSymbol[Hilbert space F of trial coordinates. E_0 embeds F isometrically into ambient space E; A_0 and the positive sine operator sin(Theta_0) act on F, while the residual R maps F into E.]
  {Ftrial}{F}{trialframe}{Hilbert space of trial coordinates.}
\dkDeclareSymbol[Ambient Hilbert space E over the real or complex scalars. It contains the trial subspace range(E_0) and the orthogonal exact subspaces range(F_0) and range(F_1). The source theorem permits infinite-dimensional separable E.]
  {Eambient}{E}{ambient}{Ambient Hilbert space containing the trial and exact subspaces.}
% Exact/wanted family: the target subspace, its frame, projection and rotation.
\dkDeclareSymbol[Target subspace V in the transport problem: feasible maps satisfy W(U) = V, or equivalently W P W* = Q. The matching theorem also names this subspace L.]
  {Vtarget}{V}{wanted}{Target subspace in the transport problem.}
\dkDeclareSymbol[Exact wanted frame F_0: F_ex -> E, a linear isometry onto the desired exact spectral subspace. F_0* returns its coordinates. P_ex = F_0 F_0* is its ambient orthogonal projection; in Figure 1, F_0(h) = (h, 0).]
  {F0}{F_0}{wantedframe}{Isometric coordinates for the desired exact spectral subspace; star denotes the adjoint.}
\dkDeclareSymbol[Self-adjoint operator block on coordinates for the desired exact spectral subspace. Its spectrum consists of the wanted exact spectral values; in the corresponding decomposition Ahat F_0 = F_0 Lambda_0 on the operator domain.]
  {Lambda0}{\Lambda_0}{wantedspectrum}{Self-adjoint block on the desired exact spectral subspace.}
\dkDeclareSymbol{u}{u}{wantedvector}{Vector in the desired exact subspace.}
\dkDeclareSymbol{f0}{f_0}{wantedvector}{Component or coordinate vector in the desired exact subspace.}
\dkDeclareSymbol{lambda0}{\lambda_0}{wantedspectrum}{Wanted exact eigenvalue or wanted spectral value.}
\dkDeclareSymbol{lambdau}{\lambda_u}{wantedspectrum}{Eigenvalue associated with a desired exact vector.}
\dkDeclareSymbol[Exact wanted orthogonal projection P_ex: E -> E, with P_ex = F_0 F_0*. It retains the desired exact spectral component. Its complementary projection is P_ex_perp = I - P_ex = F_1 F_1*; the trial leakage is S = P_ex_perp E_0.]
  {Pexact}{P_{\mathrm{ex}}}{wantedprojection}{Orthogonal projection onto the desired exact subspace in the sine-theta theorem.}
\dkDeclareSymbol[Exact complementary orthogonal projection P_ex_perp = I - P_ex = F_1 F_1*: E -> E. It retains the unwanted exact component. In Figure 1, P_ex_perp v = S(1), whereas P_ex_perp r is the unwanted component of the residual.]
  {PexactPerp}{P_{\mathrm{ex}}^\perp}{unwantedprojection}{Orthogonal projection onto the exact complementary subspace.}
\dkDeclareSymbol{ReflectionSigma}{\Sigma}{wantedprojection}{Reflection associated with the desired exact spectral projection.}
\dkDeclareSymbol[Orthogonal projection Q onto the target subspace V in the transport problem. A feasible orthogonal or unitary W satisfies W P W* = Q, where P projects onto the source U.]
  {Qtarget}{Q}{wantedprojection}{Orthogonal projection onto the target subspace in the transport problem.}
\dkDeclareSymbol{ProjTarget}{P_V}{wantedprojection}{Orthogonal projection onto the target subspace in the transport problem.}
\dkDeclareSymbol[Canonical direct rotation D from the source U to the target V for an acute pair. It is the unitary polar factor of T = QP + (I - Q)(I - P), and acts by the principal angle on each principal plane. R is reserved for the trial residual.]
  {Rrotation}{D}{wantedrotation}{Canonical direct rotation carrying the source subspace onto the target.}
\dkDeclareSymbol[Direct rotation D between the ranges of orthogonal projections P and Q, including the nonacute case. Relative to range(P) and kernel(P), its diagonal blocks are positive and its off-diagonal blocks are opposite adjoints. It exists when the two crossed defect spaces have equal dimension, and is unique when both defects vanish.]
  {SourceDirectRotation}{D}{wantedrotation}{Direct rotation, including the nonacute case in the source-result table.}
\dkDeclareSymbol[Orthonormal coordinate frame Y for the target subspace, so Y* Y = I. In the counterexample Y is the vertical stack of cos(theta) I_2 and sin(theta) J_0; hence X* Y = cos(theta) I_2 and both principal angles equal theta.]
  {Yframe}{Y}{wantedframe}{Orthonormal coordinate matrix for the target subspace in the counterexample.}
\dkDeclareSymbol[Vector y in the target subspace V. For principal vectors, unit x_j in U and y_j in V are paired so that <x_j, y_j> = cos(theta_j); distinct pairs have zero cross inner product.]
  {TargetVector}{y}{wantedvector}{Vector in the target subspace.}
\dkDeclareSymbol[Hilbert space F_ex of exact wanted coordinates. The isometry F_0: F_ex -> E identifies it with the desired exact spectral subspace. This space differs from H, the additive perturbation.]
  {Hexact}{F_{\mathrm{ex}}}{wantedframe}{Hilbert space of coordinates for the desired exact spectral subspace.}% Coordinate Hilbert space (not perturbation).
% Exact/unwanted family; the trial error operators are MAGENTA, not cyan.
\dkDeclareSymbol{Uperp}{U^{\perp}}{unwanted}{Orthogonal complement of the desired exact subspace in the local notation.}
\dkDeclareSymbol[Exact complementary frame F_1: G -> E, a linear isometry onto the orthogonal complement of range(F_0). F_1* returns unwanted coordinates; F_1 F_1* = I - P_ex. In Figure 1, F_1(g) = (0, g).]
  {F1}{F_1}{unwantedframe}{Isometric coordinates for the exact complementary subspace; star denotes the adjoint.}
\dkDeclareSymbol[Self-adjoint unwanted block Lambda_1 on G, satisfying Ahat F_1 = F_1 Lambda_1 on its domain. Its spectrum is separated from the trial spectrum of A_0 by the specified positive delta. Figure 1 uses the scalar block 0.7.]
  {Lambda1}{\Lambda_1}{unwantedspectrum}{Self-adjoint operator on exact complementary coordinates: the unwanted spectral block.}
% Spectral measure and set (distinct from the ambient space E and sine block S).
\dkDeclareSymbol[Projection-valued spectral measure E associated with a self-adjoint operator A: A = integral lambda dE(lambda). For a Borel set S, E(S) is the orthogonal projection onto the corresponding spectral subspace.]
  {SpectralMeasure}{E}{wantedprojection}{Projection-valued spectral measure of a self-adjoint operator.}
\dkDeclareSymbol[Borel subset S of the real spectral axis. The spectral projection E(S) selects the part of a self-adjoint operator with spectral values in S; its range is a reducing subspace.]
  {SpectralSet}{S}{wantedspectrum}{Borel spectral set used to select a spectral subspace.}
% Generic Sylvester coefficients and right side, by source/target role.
\dkDeclareSymbol[Left coefficient B of the Sylvester map X -> B X - X C; B acts on the output space of X. For the sine bound, B = Lambda_1 on the unwanted-coordinate space G.]
  {SylvesterLeft}{B}{unwantedspectrum}{Left coefficient of the Sylvester equation; corresponds to the unwanted spectral block.}
\dkDeclareSymbol[Right coefficient C of the Sylvester map X -> B X - X C; C acts on the input space of X. For the sine bound, C = A_0 on trial space F. The historical equation also denotes this coefficient by A.]
  {SylvesterRight}{C}{trialspectrum}{Right coefficient of the Sylvester equation; corresponds to the trial operator.}
\dkDeclareSymbol[Right-hand side Y in B X - X C = Y, with the same input and output spaces as X. For the sine bound, Y = F_1* R: F -> G, the residual expressed in unwanted exact coordinates.]
  {SylvesterRHS}{Y}{resid}{Right-hand side of the Sylvester equation, obtained by projecting the residual.}
\dkDeclareSymbol{w}{w}{unwantedvector}{Vector in an unwanted or complementary spectral subspace.}
\dkDeclareSymbol{f1}{f_1}{unwantedvector}{Component or coordinate vector in the unwanted exact subspace.}
\dkDeclareSymbol{lambda1}{\lambda_1}{unwantedspectrum}{Unwanted exact eigenvalue or complementary spectral value.}
\dkDeclareSymbol{lambdaw}{\lambda_w}{unwantedspectrum}{Eigenvalue associated with an unwanted exact vector.}
\dkDeclareSymbol[Hilbert space G of exact complementary coordinates. F_1: G -> E embeds it onto the unwanted exact subspace; Lambda_1 acts on G and X = F_1* E_0 maps trial coordinates into G.]
  {Gunwanted}{G}{unwantedframe}{Hilbert space of exact complementary coordinates.}
\dkDeclareSymbol{E1}{E_1}{unwantedframe}{Isometric coordinates for the complementary trial subspace.}
\dkDeclareSymbol{Vperp}{V^{\perp}}{unwanted}{Orthogonal complement of the trial subspace in the local notation.}
% Angle and error geometry, including distinct rectangular and positive maps.
\dkDeclareSymbol[Individual principal angle, between 0 and pi/2. For orthonormal subspace frames X and Y, the singular values of X* Y are the principal cosines. In Figure 1, theta is the angle between the unit trial vector v and the exact wanted line.]
  {theta}{\theta}{sine}{Principal angle between the subspaces under comparison.}
\dkDeclareSymbol[Full positive angle operator Theta = diag(Theta_0, Theta_1) for the source/target decompositions. Theta_0 compares the source with the target; Theta_1 compares their orthogonal complements. Its operator norm is the largest spectral angle. For positive c, the conditions Theta <= c and Theta < c mean operator norm at most c and strictly below c, respectively.]
  {Theta}{\Theta}{angleoperator}{Full angle operator, the diagonal sum of the source and complementary angle blocks.}
\dkDeclareSymbol[Directed principal-angle operator Theta_0 on source or trial coordinates. For a trial frame E_0 and exact wanted projection P_ex, cos(Theta_0) = sqrt(E_0* P_ex E_0), and sin(Theta_0) = sqrt(I_F - E_0* P_ex E_0). In the two-projection geometry it comes from the positive cosine compression P Q P restricted to range(P).]
  {Theta0}{\Theta_0}{angleoperator}{Directed principal-angle operator acting on source or trial coordinates.}
\dkDeclareSymbol[Complementary directed angle operator Theta_1, obtained from the projections I-P and I-Q by the same cosine-compression construction as Theta_0. It acts on the source complement, or its isometric coordinates. The full angle operator is diag(Theta_0, Theta_1).]
  {Theta1}{\Theta_1}{angleoperator}{Directed angle operator comparing the orthogonal complements.}
\dkDeclareSymbol[Right-angle source defect D_0 = range(P) intersect kernel(Q). These are source vectors orthogonal to the target. The full-angle conclusion in the Section 8 smallness criterion assumes dim(D_0) = dim(D_1); in Lean their identification is CrossedDefectsEquivalent.]
  {CrossedSourceDefect}{D_0}{fg}{Source directions perpendicular to the target subspace.}
\dkDeclareSymbol[Right-angle target defect D_1 = kernel(P) intersect range(Q). These are target vectors orthogonal to the source. Equality of its Hilbert dimension with that of D_0 is assumed for the full-angle conclusion of Theorem 8.2.]
  {CrossedTargetDefect}{D_1}{fg}{Target directions perpendicular to the source subspace.}
\dkDeclareSymbol[Positive complementary cosine block C_1 = cos(Theta_1). The reference complement is identified with the perturbed complement by the direct rotation D, so A_1 and Lambda_1 act on common coordinates. Theorem 8.1 compares A_1 - alpha I with C_1 (Lambda_1 - alpha I) C_1 as quadratic forms.]
  {ComplementCosineBlock}{C_1}{anglecos}{Cosine of the complementary angle block in the spectral-repulsion inequality.}
\dkDeclareSymbol[Scalar sine of a principal angle. For a unit principal trial vector v and exact wanted projection P_ex, sin(theta) = ||(I - P_ex)v||. In Figure 1 it is the length of the pink leg S(1).]
  {sintheta}{\sin\theta}{sine}{Sine of an individual principal angle: the length of the complementary component of a unit principal trial vector.}
\dkDeclareSymbol[Positive sine operator sin(Theta_0) = sqrt(I_F - E_0* P_ex E_0): F -> F. It equals |S| for S = (I - P_ex)E_0: F -> E. Its eigenvalues in finite dimensions are the principal sines; the sine theorem bounds its ideal norm by N(R)/delta.]
  {sinTheta0}{\sin\Theta_0}{sineoperator}{Positive sine operator on trial coordinates; the modulus of the rectangular complementary error.}
\dkDeclareSymbol[Rectangular complementary error S = (I - P_ex) E_0: F -> E. It records how much each trial vector leaks outside the exact wanted subspace. Its modulus |S| = sqrt(S* S) is sin(Theta_0): F -> F. In rank one S(1) = P_ex_perp v.]
  {S}{S}{sineblock}{Rectangular complementary error: the trial frame projected onto the unwanted exact subspace.}
\dkDeclareSymbol[Complementary-coordinate map X = F_1* E_0: F -> G, with S = F_1 X. It satisfies Lambda_1 X - X A_0 = F_1* R on domain(A_0). Isometric multiplication gives N(X) = N(S); the separate modulus identity abs(S) = sin(Theta_0) gives N(S) = N(sin(Theta_0)), whenever these norms are finite.]
  {X}{X}{sylvester}{Solution of the Sylvester equation; for the sine bound, the trial frame in unwanted exact coordinates.}
\dkDeclareSymbol{alphaAngle}{\alpha}{anglemax}{Largest principal angle.}
\dkDeclareSymbol[Scalar c = cos(theta) in the counterexample. In the general two-projection decomposition, the same role denotes a positive cosine operator C; the cosine and sine blocks commute and satisfy C squared + S squared = I.]
  {CosAngle}{c}{anglecos}{Cosine of a principal angle or the corresponding positive cosine block.}
\dkDeclareSymbol[Scalar s = sin(theta) in the counterexample. In the general two-projection decomposition, the same role denotes a positive sine operator S; it commutes with the positive cosine block C and satisfies C squared + S squared = I.]
  {SinAngle}{s}{sineoperator}{Sine of a principal angle or the corresponding positive sine block.}
\dkDeclareSymbol[Positive small-angle parameter in the local subspace comparison. It controls the allowed angular deviation; it is an angle parameter rather than a perturbation-operator norm.]
  {EpsilonAngle}{\varepsilon}{anglemax}{Angle parameter controlling the maximum subspace deviation.}
\dkDeclareSymbol{PhiAngle}{\varphi}{angleoperator}{Auxiliary angle or angle operator in the two-subspace comparison.}
\dkDeclareSymbol{Bblock}{B}{sineblock}{Off-diagonal block in the two-subspace decomposition.}
\dkDeclareSymbol{Cblock}{C}{sylvester}{Complementary or Sylvester block in the two-subspace decomposition.}
\dkDeclareSymbol{Ablock}{A}{anglecos}{Positive cosine-related block in the two-subspace decomposition.}
% Residual/error-displacement family.
\dkDeclareSymbol[Residual vector r = Ahat v - rho v = R(1) for the rank-one frame E_0(x) = x v. Figure 1 draws it from rho v to Ahat v. For unit v and the Rayleigh value rho = <v, Ahat v>, <v, r> = 0; a general trial scalar need not give orthogonality.]
  {r}{r}{residvector}{Residual vector of an approximate eigenpair.}
\dkDeclareSymbol[Bounded residual R: F -> E, agreeing with Ahat E_0 - E_0 A_0 on domain(A_0). The domain condition requires E_0(domain(A_0)) inside domain(Ahat). R = 0 means that E_0 intertwines A_0 with Ahat; its norm measures the defect.]
  {R}{R}{resid}{Bounded residual of the trial eigenproblem, agreeing with the operator difference on the trial domain.}
\dkDeclareSymbol[Full displacement Delta_W = I - W on the ambient space. In the failed Ky Fan comparison W = W_(pi/4); Delta_W has singular values sqrt(2), sqrt(2), 0, 0. Multiplying by a principal-plane projection restricts its input to that plane.]
  {Kdisplacement}{\Delta_W}{displacement}{Full displacement of a transporter from the identity.}
\dkDeclareSymbol[Squared full displacement Dcal(W) = (I-W)* (I-W), a positive operator. Proposition 4.3 compares its unitarily invariant norm for a direct rotation D with that for every feasible transporter W.]
  {SquaredDisplacement}{\mathcal D}{displacement}{Squared full displacement as a function of the transporter.}
% Matching proof: the gradient and its spectral level are not a space or gap.
\dkDeclareSymbol[Skew-symmetric spectral gradient G of Phi_p. It commutes with W and represents first variations by d/dt Phi_p(W exp(tB)) at t = 0 equal to -tr(G B) for skew-symmetric B. Feasible stationarity gives J_src G J_src = -G.]
  {SpectralGradient}{G}{displacement}{Skew-symmetric spectral gradient of the displacement power sum in the matching proof.}
\dkDeclareSymbol[Positive spectral level g of the skew-symmetric gradient G: on this invariant subspace, -G squared = g squared I. Consequently G/g squares to -I and supplies the complex structure used in the real matching proof.]
  {GradientLevel}{g}{displacement}{Positive gradient level; the gradient squared is minus this level squared times the identity.}
\dkDeclareSymbol[Complex structure I_g = G/g on the positive gradient-level space V_g. It is a real isometry with square -I. Multiplication by a+ib means a times identity plus b times I_g; source reflection J_src reverses I_g and acts as conjugation.]
  {GradientComplexStructure}{\mathcal I_g}{displacement}{Complex multiplication constructed from a positive gradient level.}
\dkDeclareSymbol[Real skew-symmetric matrix B, with B transpose = -B, giving the orthogonal variation W exp(tB). If B commutes with the source projection P, this curve preserves the constraint W(U) = V.]
  {SkewDirection}{B}{ambient}{Skew-symmetric direction for a first variation of an orthogonal transporter.}
% Spectral separation family: gap and original gap are NOT identical.
\dkDeclareSymbol[Positive spectral separation in the stated configuration. In the sine theorem it separates the trial operator from the exact complementary block; Theorem 8.1 uses it for the ordered reference gap, and Theorem 8.2 for the perturbed central-band separation. Figure 1 uses delta = abs(rho - 0.7).]
  {delta}{\delta}{gap}{Positive spectral separation in the stated gap configuration.}
\dkDeclareSymbol{g}{g}{gaporiginal}{Original spectral gap of the reference operator; distinct from the certified sine-theta separation.}
\dkDeclareSymbol{Delta}{\Delta}{gapstat}{Population or statistical spectral gap.}
\dkDeclareSymbol[Spectral endpoint alpha. In the interval/exterior configuration, one spectrum lies between beta and alpha, the other at or below beta - delta or at or above alpha + delta. Theorem 8.2 places the selected perturbed block in the central interval. In Theorem 8.1, alpha bounds the lower reference block and alpha + delta bounds the upper block from below.]
  {alphaEndpoint}{\alpha}{gapinterval}{Upper endpoint of the central band, or lower-block cut in the ordered gap.}
\dkDeclareSymbol[Lower endpoint beta of the finite interval from beta to alpha in the gap hypothesis. One spectrum lies in that interval; the other lies at or below beta - delta or at or above alpha + delta.]
  {betaEndpoint}{\beta}{gapinterval}{Lower endpoint of the finite interval in the interval/exterior gap hypothesis.}
% Spectral/projector and auxiliary matrix objects, context-scoped.
\dkDeclareSymbol{P}{P}{fg}{Orthogonal projection; its source or exact-target role is specified locally.}% Neutral historical Manim key; prefer Pexact/Psource.
\dkDeclareSymbol{Q}{Q}{fg}{Orthogonal projection; its role is specified locally.}% Neutral historical Manim key; prefer Qtarget.
\dkDeclareSymbol[Two-projection intertwiner T = QP + (I - Q)(I - P), satisfying TP = QT and T* T = I - (P - Q) squared. For ||P - Q|| < 1 it is invertible; its unitary polar factor T |T| inverse is the canonical direct rotation.]
  {IntertwinerT}{T}{intertwiner}{Two-projection intertwiner; its unitary polar factor is the acute direct rotation.}
\dkDeclareSymbol[Planar quarter-turn J_0 = ((0, -1), (1, 0)), with J_0 squared = -I_2 and J_0 transpose = -J_0. It defines the rotation C = cos(2 theta) I_2 + sin(2 theta) J_0 used to construct the real competitor.]
  {Jturn}{J_0}{turn}{Planar quarter-turn, with square minus the identity.}
\dkDeclareSymbol[Partial isometry V_pol in the polar decomposition S = V_pol |S|, where |S| = sqrt(S* S). It is isometric on the orthogonal complement of kernel(S) and zero on kernel(S). For an invertible operator its polar factor is unitary.]
  {PolarU}{V_{\mathrm{pol}}}{intertwiner}{Partial isometry in the polar decomposition; unitary when the operator is invertible.}
\dkDeclareSymbol[Planar rotation C = cos(2 theta) I_2 + sin(2 theta) J_0 through twice the principal angle. The competitor W_theta is orthogonally conjugate to C direct-sum I_2, so it rotates one two-plane and fixes the other.]
  {Cturn}{C}{turn}{Planar rotation through twice the principal angle in the counterexample.}
\dkDeclareSymbol[Orthogonal mixing matrix H_mix = (1/sqrt(2)) times the block matrix ((I_2, I_2), (I_2, -I_2)). It constructs W_theta = H_mix (C direct-sum I_2) H_mix transpose.]
  {Hmix}{H_{\mathrm{mix}}}{mix}{Orthogonal mixing matrix in the counterexample.}
\dkDeclareSymbol[Source reflection J_src = 2P - I: it acts as +I on the source subspace and -I on its orthogonal complement, so its square is I. In the matching proof stationarity gives J_src G J_src = -G for the skew-symmetric spectral gradient G.]
  {Jreflect}{J_{\mathrm{src}}}{trialprojection}{Reflection equal to twice the source projection minus the identity.}
\dkDeclareSymbol[Skew angle generator J_ang for a direct rotation D = exp(J_ang Theta). On each principal plane it is the oriented quarter-turn; it vanishes on the common fixed directions. It commutes with the full angle operator Theta.]
  {AngleGenerator}{J_{\mathrm{ang}}}{turn}{Skew angle generator in the direct-rotation result summaries.}
\dkDeclareSymbol{Dblock}{D}{mix}{Auxiliary diagonal or block matrix in the counterexample.}
\dkDeclareSymbol[Orthogonal projection onto the first principal plane span(e_0, e_3) in the four-dimensional counterexample. For Delta_W = I - W_(pi/4), the Ky Fan 2-norm of Delta_W Omega_1 equals 2.]
  {OmegaOne}{\Omega_1}{turn}{First principal-plane projection in the failed Ky Fan comparison.}
\dkDeclareSymbol[Orthogonal projection onto the second principal plane span(e_1, e_2) in the four-dimensional counterexample. For Delta_W = I - W_(pi/4), the Ky Fan 2-norm of Delta_W Omega_2 equals 2, although the trace norm of Delta_W is only 2 sqrt(2).]
  {OmegaTwo}{\Omega_2}{intertwiner}{Second principal-plane projection in the failed Ky Fan comparison.}
\dkDeclareSymbol{Dplanar}{D}{wantedrotation}{Planar direct-rotation block.}
\dkDeclareSymbol{Hreflection}{H}{trialcompetitor}{Planar reflection carrying the source line to the target.}
% Unitarily invariant gauge family; N, M are NOT a norm versus its input.
\dkDeclareSymbol[Unitarily invariant norm N, or a symmetric operator-ideal norm on its domain of finiteness. It satisfies N(U T V) = N(T) for unitaries and N(B T) <= ||B||op N(T). In the sine theorem it is normalized by N(u v*) = ||u|| ||v|| and has Ky Fan dominance where both norms are finite. Both R and sin(Theta_0) must belong to its ideal.]
  {NormN}{N}{norm}{Unitarily invariant norm on its domain of finiteness; in the sine theorem, normalized by $N(uv^*)=\|u\|\|v\|$.}
\dkDeclareSymbol[Sine-theta separation G_delta(B,C): finite interval/exterior separation in either spectral order, or ordered separation B <= a I and C >= (a+delta) I, or the reverse. The ordered alternatives allow both spectra to be unbounded.]
  {SineGap}{\mathcal G_\delta}{gap}{Interval/exterior separation or either ordered separation, as in the sine-theta theorem.}
\dkDeclareSymbol[Auxiliary unitarily invariant norm M witnessing the Q-norm property N(T) squared = M(T* T). For the Frobenius norm N, one can take M to be the trace norm.]
  {NormM}{M}{normaux}{Auxiliary unitarily invariant norm used to define a Q-norm through the squared modulus.}
\dkDeclareSymbol[Q-norm property: a unitarily invariant norm N is a Q-norm when N(T) squared = M(T* T) for some unitarily invariant norm M. The Frobenius norm is an example, with M the trace norm; the repaired full-displacement theorem uses this norm class.]
  {NormQ}{Q}{normaux}{Q-norm: its square is a unitarily invariant norm of the squared modulus.}
\dkDeclareSymbol[Singleton trace cost s_i = min(4 sin(theta_i/2), 2) for angle label i. It transports one real principal plane independently, choosing a rotation of cost 4 sin(theta_i/2) or a reflection of cost 2.]
  {SingletonCost}{s}{norm}{Minimum trace displacement for one real principal plane: choose a rotation or a reflection.}
\dkDeclareSymbol[Paired trace cost e_ij = 4 sin((theta_i + theta_j)/2) for distinct angle labels i,j. It is attained by coupling their principal planes in a real four-dimensional block. The edge weight s_i + s_j - e_ij is the difference between the cost of two singletons and this paired cost.]
  {PairCost}{e}{norm}{Trace displacement cost of coupling two principal planes in a four-dimensional block.}
\dkDeclareSymbol[Scalar phase cost f_p(t) = (2 sin(t/2)) to the power p, for 0 <= t <= pi. A real planar rotation through t contributes two copies, so its displacement power cost is 2 f_p(t).]
  {MatchingPhaseCost}{f_p}{norm}{Displacement power cost for one scalar phase.}
\dkDeclareSymbol[Singleton power cost a_i(p) = min(2 f_p(theta_i), 2 to the power p). These are the rotation and reflection choices for one principal plane. At p=1 it equals the trace-cost weight s_i.]
  {SingletonPowerCost}{a}{norm}{Minimum displacement power cost for one angle.}
\dkDeclareSymbol[Paired power cost b_ij(p) = min over d <= tau <= h of 2(f_p(h+tau) + f_p(h-tau)), where h=(theta_i+theta_j)/2 and d=abs(theta_i-theta_j)/2. This is the positive-determinant pair minimum; unrestricted block optimization also compares two singletons. At p=1 it equals e_ij.]
  {PairPowerCost}{b}{norm}{Positive-determinant displacement power cost for a pair of angles.}
\dkDeclareSymbol[Matching M on the finite general graph of positive principal-angle labels. Each label belongs to at most one edge; unmatched labels use singleton transporters. Minimizing paired and singleton costs is equivalent to maximum-weight matching with edge weights g_ij = s_i + s_j - e_ij. The empty matching is allowed.]
  {Matching}{M}{normaux}{Matching of positive principal-angle labels: each label is used in at most one pair.}
\dkDeclareSymbol[Displacement power sum Phi_p(W) = sum_j s_j(I - W)^p. The matching proof first minimizes it for 1 < p < 2, then lets p decrease to 1, where it becomes the trace norm ||I - W||_1.]
  {PowerCost}{\Phi}{norm}{Sum of the p-th powers of displacement singular values.}
% Auxiliary matching-proof notation: neutral ink, with the same hover registry.
\dkDeclareSymbol[Off-diagonal gradient block G_01 mapping the source complement K perpendicular into K. Stationarity makes the diagonal blocks of G zero, so G = ((0,G_01),(-G_01 transpose,0)) in the source/complement decomposition.]
  {GradientOffDiagonal}{G_{01}}{fg}{Complement-to-source block of the skew-symmetric gradient.}
\dkDeclareSymbol[Orthogonal group O(n): real n by n matrices W satisfying W transpose times W = I. Its elements preserve Euclidean inner products and lengths. The notation O(2) refers to the same group in dimension two.]
  {OrthogonalGroup}{O(n)}{fg}{Group of real orthogonal matrices.}
\dkDeclareSymbol[Feasible transporters F = {W in O(n) : W(K) = L}. Equivalently W P_K W transpose = P_L. This nonempty compact set is the optimization domain throughout the matching proof.]
  {FeasibleTransport}{\mathcal F}{fg}{Orthogonal maps carrying the source subspace onto the target.}
\dkDeclareSymbol[Singular value s_j(T), the nonnegative square root of an eigenvalue of T transpose times T, counted with multiplicity. Summing s_j(I-W) gives the trace displacement; summing its p-th powers gives Phi_p(W). These singular values differ from the singleton weights s_i.]
  {TransportSingularValues}{s_j}{fg}{A displacement singular value in the matching proof.}
\dkDeclareSymbol[Scalar singular-value magnitude s between zero and two for the displacement I-W of an orthogonal matrix. Uniform convergence of s to the power p to s on this interval permits the trace-norm limit. This variable differs from the rotation parameter s used earlier in the proof.]
  {DisplacementSize}{s}{fg}{Scalar displacement singular-value magnitude in the uniform limit.}
\dkDeclareSymbol[All-singleton cost C_0 = sum_i s_i, before any two angle labels are paired. A matching has cost C_0 minus the sum of its selected edge weights.]
  {SingletonBaseline}{C_0}{fg}{Trace displacement when every angle uses its singleton transporter.}
\dkDeclareSymbol[Matching power cost C_p(M): sum of b_ij(p) over selected pairs plus a_i(p) over unmatched labels. Its minimum equals the minimum of Phi_p over feasible transporters. At p=1 it is C(M).]
  {MatchingPowerCost}{C_p}{fg}{Total paired and singleton power cost of a matching.}
\dkDeclareSymbol[Invariant orthogonal block V_a preserved by the source projection, target projection, and minimizing transporter. The displacement cost adds over these blocks, and their principal-angle lists combine into the original list.]
  {MatchingBlock}{V_a}{fg}{Simultaneous invariant block in the matching proof.}
\dkDeclareSymbol[Gradient-level space V_g: the eigenspace of -G squared for a positive eigenvalue g squared. It is preserved by G, W, and the source reflection J_src. Multiplication by i on it is G/g.]
  {GradientLevelSpace}{V_g}{fg}{Positive gradient-level eigenspace.}
\dkDeclareSymbol[Eigenphase varphi of W on V_g viewed as a complex space with multiplication by i = G/g. It lies between zero and pi and satisfies f_p prime(varphi) = g. There are at most two such phases.]
  {GradientEigenphase}{\varphi}{fg}{One of the at most two phases at a positive gradient level.}
\dkDeclareSymbol[Orthogonal projection E onto the first phase space of W on V_g. Here W = exp(i varphi_1) E + exp(i varphi_2) (I-E). In coordinates where the source reflection is conjugation, E = X + iY.]
  {GradientPhaseProjection}{E}{fg}{Projection onto the first phase space in a gradient-level space.}
\dkDeclareSymbol[Real part X of the phase projection E = X + iY. It is a real symmetric matrix with spectrum in the closed interval from zero to one, satisfying X squared - Y squared = X and XY + YX = Y. This X differs from the source-coordinate and Sylvester matrices elsewhere in the paper.]
  {ProjectionRealPart}{X}{fg}{Real symmetric part of the phase projection.}
\dkDeclareSymbol[Imaginary part Y of the phase projection E = X + iY. It is a real skew-symmetric matrix. It maps the lambda eigenspace of X to the 1-lambda eigenspace and constructs the small invariant planes.]
  {ProjectionImaginaryPart}{Y}{fg}{Real skew-symmetric imaginary part of the phase projection.}
\dkDeclareSymbol[Eigenvalue lambda in the closed interval from zero to one of the real part X of the phase projection. It is local to the block construction and differs from the original operator's spectral parameter.]
  {ProjectionEigenvalue}{\lambda}{fg}{Eigenvalue of the real part of the phase projection.}
\dkDeclareSymbol[Plane coupling eta = sqrt(lambda (1-lambda)) for 0 < lambda < 1. If u is a unit lambda eigenvector of X, then v = Yu/eta completes its invariant real plane.]
  {ProjectionCoupling}{\eta}{fg}{Normalization used to construct an invariant real plane.}
\dkDeclareSymbol[Pair mean angle h = (theta_i + theta_j)/2. Together with d = (theta_i - theta_j)/2 after ordering theta_i >= theta_j, it parametrizes the two-angle block cost.]
  {PairMeanAngle}{h}{fg}{Mean of the two ordered principal angles.}
\dkDeclareSymbol[Pair half-difference d = (theta_i - theta_j)/2 after ordering theta_i >= theta_j. It satisfies 0 <= d <= h and h+d <= pi/2.]
  {PairHalfDifference}{d}{fg}{Half the difference of the two ordered principal angles.}
\dkDeclareSymbol[Scalar pair parameter tau in the closed interval from d to h. The paired power cost is the minimum of 2(f_p(h+tau) + f_p(h-tau)) over this interval. At p=1 a minimum is attained at tau=h.]
  {PairPhaseParameter}{\tau}{fg}{Scalar parameter in the two-angle power-cost minimum.}
\dkDeclareSymbol[Attaining rotation angle gamma = arccos(cos(tau)/cos(d)). Choosing U_0 = U_1 = R_gamma realizes phase parameters x=h and y=tau in the four-dimensional block.]
  {PairAttainingAngle}{\gamma}{fg}{Rotation angle realizing a chosen scalar pair parameter.}
\dkDeclareSymbol[Four-dimensional direct rotation D = ((C,-S),(S,C)) in source and complementary principal coordinates, with C the diagonal principal cosines and S the diagonal principal sines. Every transporter on this block is D diag(U_0,U_1).]
  {PairDirectRotation}{D}{fg}{Direct rotation on the two principal planes.}
\dkDeclareSymbol[Diagonal cosine block C = diag(cos(theta_i), cos(theta_j)) in the four-dimensional direct rotation D. This matrix differs from the scalar displacement cost C(M).]
  {PairCosineBlock}{C}{fg}{Diagonal principal-cosine block.}
\dkDeclareSymbol[Diagonal sine block S = diag(sin(theta_i), sin(theta_j)) in the four-dimensional direct rotation D. This matrix is local to the matching proof.]
  {PairSineBlock}{S}{fg}{Diagonal principal-sine block.}
\dkDeclareSymbol[Free planar orthogonal matrices U_0 and U_1 acting on the two-dimensional source and its two-dimensional complement, respectively. Multiplying D by diag(U_0,U_1) parametrizes every transporter on this four-dimensional block.]
  {PairFreeOrthogonal}{U}{fg}{Free source or complementary planar orthogonal factor.}
\dkDeclareSymbol[Phase parameter x in the closed interval from h to pi-h for a positive-determinant four-dimensional transporter. Its eigenvalues are exp(plus or minus i(x+y)) and exp(plus or minus i(x-y)).]
  {PairPhaseX}{x}{fg}{First phase parameter of a four-dimensional transporter.}
\dkDeclareSymbol[Phase parameter y in the closed interval from d to pi-d for a positive-determinant four-dimensional transporter. Together with x it gives rotational phases x+y and x-y, folded into the closed interval from zero to pi for the displacement cost.]
  {PairPhaseY}{y}{fg}{Second phase parameter of a four-dimensional transporter.}
\dkDeclareSymbol[Scalar a = cos(x) in the four-dimensional phase calculation. Together with b = cos(y), it expresses the trace as 4ab and determines the reciprocal characteristic polynomial.]
  {PairPhaseCosineX}{a}{fg}{Cosine of the first phase parameter in the trace calculation.}
\dkDeclareSymbol[Scalar b = cos(y) in the four-dimensional phase calculation, paired with a = cos(x). These scalars are local to the characteristic-polynomial calculation.]
  {PairPhaseCosineY}{b}{fg}{Cosine of the second phase parameter in the trace calculation.}
\dkDeclareSymbol[Rotation parameter r: the two free orthogonal factors use planar rotation angles r+s and r-s. It appears in cos(y) = cos(d) cos(r) for two rotations.]
  {PairRotationR}{r}{fg}{Half the sum of the two free planar rotation angles.}
\dkDeclareSymbol[Rotation parameter s: the two free orthogonal factors use planar rotation angles r+s and r-s. It appears in cos(x) = cos(h) cos(s) for two rotations. This scalar differs from the singular values and singleton costs.]
  {PairRotationS}{s}{fg}{Half the difference of the two free planar rotation angles.}
\dkDeclareSymbol[Planar reflection J_2 = diag(1,-1). Every two-dimensional orthogonal matrix of determinant -1 can be written as a planar rotation times J_2.]
  {PairPlanarReflection}{J_2}{fg}{Fixed planar reflection used to parametrize determinant -1 factors.}
\dkDeclareSymbol[Real planar rotation R_t = ((cos(t),-sin(t)),(sin(t),cos(t))). Its two displacement singular values are both 2 times abs(sin(t/2)). In the matching proof a rotation through t in the closed interval from zero to pi contributes 2 f_p(t).]
  {MatchingPlanarRotation}{R_t}{fg}{Planar rotation through the indicated angle.}
% Statistical application / operator provenance.
\dkDeclareSymbol{Sigma}{\Sigma}{old}{Population covariance matrix.}
\dkDeclareSymbol{SigmaHat}{\widehat\Sigma}{current}{Estimated or perturbed covariance matrix.}
\dkDeclareSymbol{YWSV}{V}{wanted}{Orthonormal coordinates for the population eigenspace in the statistical application.}
\dkDeclareSymbol{YWSVhat}{\widehat V}{trialframe}{Estimated eigenspace frame in the statistical application.}

% Figure-specific objects retain their own meanings when bare letters overlap.
\dkDeclareSymbol[Signed rotation angle phi of the displayed map on its invariant real two-plane. At theta = pi/4 the direct rotation turns its two principal planes by opposite signed angles of magnitude pi/4; the competitor turns its moving plane by pi/2 and fixes its other plane.]
  {PlaneRotationAngle}{\varphi}{sine}{Signed rotation angle of the orthogonal map on the displayed invariant plane.}
\dkDeclareSymbol[Unit vector q in the displayed invariant plane. The map T is the direct rotation D or the competitor W in that panel; Tq is the image and (I - T)q = q - Tq is the displacement.]
  {FigureInput}{q}{ambient}{Unit input to the orthogonal map shown in this invariant-plane panel.}
\dkDeclareSymbol[Image Tq of the panel unit input q under T = D or W. Both q and Tq lie on the unit circle; the green displacement arrow joins Tq to q, so it represents q - Tq.]
  {FigureImage}{Tq}{wantedvector}{Image of the unit input under the orthogonal map shown in this panel.}
\dkDeclareSymbol[Displacement (I - T)q = q - Tq, drawn from the image Tq to the unit input q. On a plane rotated through phi its length is 2 |sin(phi/2)|; it vanishes on a fixed plane.]
  {DisplacementVector}{(I-T)q}{displacement}{Difference between a unit input and its image under the displayed orthogonal map.}
\dkDeclareSymbol[First principal plane Pi_1 = span(e_0, e_3) for the source-target pair in the real four-dimensional counterexample. At theta = pi/4 the direct rotation acts on this plane by a quarter of pi, with displacement singular value sqrt(2 - sqrt(2)).]
  {CounterPrincipalOne}{\Pi_1}{fg}{First principal plane in the real counterexample: the span of e0 and e3.}
\dkDeclareSymbol[Second principal plane Pi_2 = span(e_1, e_2) for the source-target pair in the real four-dimensional counterexample. The direct rotation acts by the same angle magnitude as on Pi_1, with the opposite sign in the displayed plane coordinates.]
  {CounterPrincipalTwo}{\Pi_2}{fg}{Second principal plane in the real counterexample: the span of e1 and e2.}
\dkDeclareSymbol[Moving plane M = range(W - I) = span(e_0 + e_2, e_1 + e_3) for the competitor. W rotates it by 2 theta; at theta = pi/4 the two displacement singular values contributed by this plane are both sqrt(2).]
  {CounterMovingPlane}{M}{ambient}{Moving plane of the competing transporter: the range of W minus the identity.}
\dkDeclareSymbol[Fixed plane Z = kernel(W - I) = span(e_0 - e_2, e_1 - e_3) for the competitor. W acts as the identity here, so this plane contributes two zero displacement singular values.]
  {CounterFixedPlane}{Z}{ambient}{Fixed plane of the competing transporter: the kernel of W minus the identity.}
\dkDeclareSymbol[Direct-rotation displacement singular value a = 2 sin(pi/8) = sqrt(2 - sqrt(2)) at theta = pi/4. Each of its two principal planes contributes two copies, giving trace displacement 4a.]
  {CounterChordA}{a}{displacement}{Direct-rotation displacement singular value in the counterexample: sqrt(2 - sqrt(2)).}
\dkDeclareSymbol[Competitor displacement singular value b = 2 sin(pi/4) = sqrt(2) at theta = pi/4. Its moving plane contributes two copies; the fixed plane contributes two zeros, giving trace displacement 2b.]
  {CounterChordB}{b}{displacement}{Competing-transporter displacement singular value in the counterexample: square root of 2.}
\dkDeclareSymbol[Singular value s_j(I - T) of the full ambient displacement, ordered from largest to smallest. Each rotating real two-plane contributes two values 2 |sin(phi/2)|; fixed directions contribute zeros. Their sum is the trace cost.]
  {DisplacementSingular}{s_j}{displacement}{An individual singular value of the full displacement from the identity.}
\dkDeclareSymbol[Edge weight g_ij = s_i + s_j - e_ij, the difference between the cost of two singletons and the paired cost. A matching has total cost C_0 minus the sum of its edge weights, with C_0 = sum_i s_i. Maximizing this sum gives the minimum transport cost; edges of negative weight may be omitted.]
  {MatchingEdgeWeight}{g_{ij}}{current}{Edge weight: the cost of two singletons minus their paired cost.}
\dkDeclareSymbol[Total trace-norm displacement cost. For a transporter W, C(W) = ||I - W||_1, the sum of its displacement singular values. For a matching M, C(M) is the sum of its paired and unmatched singleton costs.]
  {TraceCost}{C}{norm}{Trace-norm displacement cost of the indicated orthogonal transporter or matching.}

% Spectral values have their own key, distinct from an operator block.
\dkDeclareSymbol[Spectral parameter lambda of a self-adjoint operator. In a finite-dimensional eigenvalue list lambda_j(A) it denotes an individual eigenvalue; in A = integral lambda dE(lambda), it is the real integration variable.]{SpectralValue}{\lambda}{wantedspectrum}{Real spectral parameter or individual eigenvalue of a self-adjoint operator.}

\ifDKColoredSymbols\else\DKMathColorsfalse\fi

%\DeclareMathOperator{\DKRange}{range}
%\DeclareMathOperator{\DKKernel}{kernel}
%\DeclareMathOperator{\DKDom}{domain}
\DeclareMathOperator{\DKRange}{ran}
\DeclareMathOperator{\DKKernel}{ker}
\DeclareMathOperator{\DKDom}{dom}
\usepackage[most]{tcolorbox}
\usepackage{listings}
\usepackage{listingsutf8}
\usepackage{authblk}
\usepackage{emoji} % Requires LuaLaTeX compiler

\setlength{\emergencystretch}{3em}

% ArXiv-friendly Lean code: no shell escape, minted, or external highlighter.
\definecolor{leankeyword}{HTML}{243B6B}
\definecolor{leantactic}{HTML}{70477A}
\definecolor{leantype}{HTML}{1F6473}
\definecolor{leancomment}{HTML}{39734D}
\definecolor{leanstring}{HTML}{8A4B2A}
\definecolor{leannumber}{HTML}{777777}
\definecolor{leanbackground}{HTML}{F7F7F3}
\lstdefinelanguage{lean}{%
  sensitive=true,
  morecomment=[l]{--},
  morecomment=[s]{/-}{-/},
  morestring=[b]",
  morekeywords=[1]{import,open,namespace,section,end,noncomputable,variable,variables,
    universe,universes,set_option,attribute,axiom,opaque,theorem,lemma,def,abbrev,
    example,instance,structure,class,inductive,notation,local,deriving,extends,
    by,where,let,have,show,fun,match,with,if,then,else,do,return,in,forall},
  morekeywords=[2]{simp,simpa,rw,erw,dsimp,unfold,ring,ring_nf,nlinarith,linarith,
    positivity,norm_num,field_simp,omega,exact,refine,apply,assumption,intro,intros,
    ext,funext,cases,rcases,obtain,constructor,left,right,classical,by_contra,
    contrapose,convert,using,suffices,calc},
  morekeywords=[3]{Prop,Type,Nat,Int,Real,Fin,Set,Submodule,Matrix,EuclideanSpace,
    LinearMap,ContinuousLinearMap,LinearPMap,RCLike,IsTrialResidual,
    IsExactSpectralDecomposition,FormBoundedSylvesterGap,
    NormalizedSymmetricOperatorIdealFamily},
}
\lstdefinestyle{leanpaper}{%
  language=lean,
  columns=fullflexible,
  backgroundcolor=\color{leanbackground},
  basicstyle=\ttfamily\scriptsize,
  commentstyle=\color{leancomment}\itshape,
  keywordstyle=[1]\color{leankeyword}\bfseries,
  keywordstyle=[2]\color{leantactic},
  keywordstyle=[3]\color{leantype},
  stringstyle=\color{leanstring},
  numberstyle=\tiny\color{leannumber},
  breakatwhitespace=false,
  breaklines=true,
  captionpos=b,
  keepspaces=true,
  numbers=left,
  numbersep=6pt,
  showspaces=false,
  showstringspaces=false,
  showtabs=false,
  tabsize=2,
  aboveskip=0.8em,
  belowskip=0.8em,
  literate=
    {\\KField}{{$\Bbbk$}}1
    {∃}{{$\exists$}}1
    {∧}{{$\land$}}1
    {∀}{{$\forall$}}1
    {≠}{{$\neq$}}1
    {⊤}{{$\top$}}1
}
\crefname{lstlisting}{formalization}{formalizations}
\Crefname{lstlisting}{Formalization}{Formalizations}
\renewcommand{\lstlistingname}{Formalization}

\newtheorem{theorem}{Theorem}
\newtheorem{proposition}[theorem]{Proposition}
\newtheorem{lemma}[theorem]{Lemma}
\newtheorem{remark}[theorem]{Remark}
\newtheorem{corollary}[theorem]{Corollary}

\newcommand{\Lean}{\textsc{Lean}}
\newcommand{\DK}{Davis--Kahan}
\newcommand{\TauCeti}{Tau Ceti}
% Use \ie{} and \eg{}; supply surrounding punctuation at the call site.
\newcommand{\ie}{i.e.\@}
\newcommand{\eg}{e.g.\@}
\newcommand{\DKResultCount}{29}
\newcommand{\DKProvedCount}{28}
\newcommand{\DKRefutedCount}{1}

% Usage: \margincomment{Check this formula!}
%\newcommand{\jonsays}[1]{%
%  \marginpar{%
%    \raggedright\color{red}\tiny% Change to \footnotesize if text is too small
%    \emoji{person-raising-hand} \textbf{Comment:}\\#1%
%  }%
%}

\geometry{
  % Keeps your main text layout intact while expanding the side margin box
  marginparwidth=2.2cm,
  marginparsep=0.3cm
}

% Usage: \margincomment{Your comment that now wraps perfectly!}
\newcommand{\jonsays}[1]{%
  \marginpar{%
    \raggedright% Forces text to wrap beautifully without ugly hyphenation
    \color{red}\scriptsize% Using \scriptsize makes the text tighter
    \emoji{person-raising-hand}: \\#1%
  }%
}
%\newcommand{\jonsays}[1]{{\color{red}\small\textbf{[\emoji{person-raising-hand}: #1]}}}

% Usage: \margincomment{Your comment that now wraps perfectly!}
%\newcommand{\agentsays}[1]{%
%  \marginpar{%
%    \raggedright% Forces text to wrap beautifully without ugly hyphenation
%    \color{red}\scriptsize% Using \scriptsize makes the text tighter
%    \emoji{robot}: \\#1%
%  }%
%}
\newcommand{\agentsays}[1]{{\color{red}\small\textbf{[\emoji{robot}: #1]}}}

\newcommand{\jonrsp}[1]{{\color{red}\small\textbf{[\emoji{person-raising-hand}: #1]}}}

\title{Formalizing Davis--Kahan Perturbation Theory in Lean 4}
% Additional authors are expected; keep the confirmed ordering below stable.

\ifDKAnonymous
\author{Anonymous authors}
\else
\author[1]{Jonathan Crall}
\author[2]{Edward Wang}
\author[1]{Brian Hu}
\author[2]{Carey E. Priebe}

\affil[1]{%
  Kitware, Inc.\\
  \texttt{jon.crall@kitware.com},
  \texttt{brian.hu@kitware.com}
}

\affil[2]{%
  Johns Hopkins University\\
  \texttt{ewang43@jhu.edu},
  \texttt{cep@jhu.edu}
}
\fi

\date{28 October 2026}

\begin{document}
\maketitle

\begin{abstract}

We formalize the Davis--Kahan theory of spectral-subspace perturbations.
From Davis and Kahan's landmark 1970 paper we identify \DKResultCount{}
  results as the targets of the formalization: of these
  \DKProvedCount{} are proved in Lean~4 under their stated hypotheses, and the remaining
  statement --- as originally written --- is refuted.
The formal development includes the sine and tangent theorems, direct rotations,
  principal-angle operators, Sylvester separation, unitarily invariant
  operator-ideal norms, and unbounded self-adjoint operators.
We place these constructions in their mathematical lineage, from the geometry of
  principal angles through the theory of projections and the spectral methods
  underlying Davis and Kahan's estimates.


The refuted statement \citep[Proposition~4.4]{davis1970rotation} concerns
  minimization of full ambient displacement by direct rotation.
We construct real four-dimensional examples in which direct rotation fails
  to minimize the trace norm at every prescribed positive acute angle.
We locate the failed singular-value comparison and prove a corrected optimality
  statement for Q-norms, including the Frobenius norm.

The formalization builds on Mathlib and adapts the Spectra project's
  projection-valued-measure and Borel-calculus constructions.
Portions of the reusable operator theory developed here were subsequently
  upstreamed into Tau Ceti.
The resulting infrastructure also supports a theorem reducing trace-norm
  optimal real subspace transport to weighted matching on principal angles.
These outcomes show how formalizing the foundations of a classical mathematical
  theory can support validating results, identifying historical errors, and
  developing new formal research.

\end{abstract}

\noindent\textbf{Keywords:} Davis--Kahan theory; spectral perturbation;
formal verification; Lean 4; operator theory; direct rotation; principal angles.

\smallskip

\noindent\textbf{Mathematics Subject Classification (2020):}
\href{https://mathscinet.ams.org/msc/msc2020.html?t=68V20}{68V20},
\href{https://mathscinet.ams.org/msc/msc2020.html?t=47A55}{47A55},
\href{https://mathscinet.ams.org/msc/msc2020.html?t=15A42}{15A42}.

\ifDKAnonymous\else
\smallskip
\noindent\textbf{Corresponding author:} Jonathan Crall
(\texttt{jon.crall@kitware.com}).
\fi

\section{Introduction}
\label{sec:introduction}

An approximate eigenvector gives a candidate direction for the action of an
operator. If a unit vector $\sym{v}$ is paired with an approximate eigenvalue $\sym{mu}$,
the residual $\sym{A}\sym{v}-\sym{mu} \sym{v}$ measures how closely it satisfies the eigenvector
equation. A small residual can identify an accurate approximation to an
invariant subspace even when repeated or nearby eigenvalues make individual
eigenvectors difficult to distinguish. Davis--Kahan theory bounds the angles
between approximate and exact spectral subspaces in terms of a residual or
perturbation and a separation between the relevant spectra
\citep{davis1970rotation}.

In applications, one often works with a perturbed operator
  $\sym{Acurrent}=\sym{A}+\sym{H}$ while the reference operator $\sym{A}$ is
  unknown.
Taking an eigenspace of $\sym{Acurrent}$ as a trial subspace for $\sym{A}$, the norm of its
  residual relative to $\sym{A}$ is bounded by the norm of $\sym{H}$, allowing
  Davis--Kahan theory to bound the angle between the two subspaces without
  requiring $\sym{A}$ to be known or its residuals to be evaluated directly.

Our project began with a variant of this bound needed for
\ifDKAnonymous a spectral-estimation application\else the data kernel
perspective space (DKPS) \citep{helm2026query}\fi.

We extended the formalization to the results established in the original 1970
paper, including real and complex Hilbert spaces, unbounded self-adjoint
operators, and unitarily invariant operator-ideal norms. In particular, the
spectral-separation hypotheses allow half-infinite intervals
(e.g.\ $(-\infty,a]$). These settings require operator domains, spectral
calculus, and norm estimates beyond finite matrices.

We formalized the four headline theorems \citep[Section~2]{davis1970rotation} and every
named theorem, proposition, lemma, and corollary established in the 1970 paper.

This gives \DKResultCount{} targets, enumerated in
Appendix~\ref{app:dk-inventory}. We prove \DKProvedCount{} of them and refute
the printed statement \citep[Proposition~4.4]{davis1970rotation}.
That proposition asserts a full-displacement
optimality property for real direct rotations. We give a family of real
four-dimensional examples exhibiting trace-norm failure at every prescribed
positive acute angle and identify
the invalid singular-value comparison in its proof. The phase construction
behind the counterexample occurs over the complex field in the authors'
example \citep[Example~4.2]{davis1970rotation}. To the best of our knowledge,
its contradiction of the printed real, small-angle claim has not previously
been made explicit.
For the constructed pair, whose two principal angles both equal
$\sym{theta}$, the competing transporter attains the global trace-norm minimum
$4\sin\sym{theta}$ for the constructed subspaces. Section~\ref{sec:prop44}
compares these statements and repairs the optimality claim by restricting
the norm class to Q-norms (N is a Q-norm when $N(T)^2 = M(T^* T)$ for some
unitarily invariant norm $M$). This repair follows from the squared-displacement
comparison \citep[Proposition~4.3]{davis1970rotation};
earlier Q-norm minimization results appear in \citet{bolotnikov2001minimal}.

% TODO (lineage): Consolidate the scattered lineage records through the
% maintained dev/tauceti/dkps-upstream-contribution-census.md and the evidence
% index in AUTHORING_NOTES.md. Audit each attribution against an immutable
% donor commit/path, the accepted upstream commit/PR, and the corresponding
% declaration statements and proofs. Preserve Mathlib prerequisites, Spectra
% ancestry where applicable, and independent work as distinct relationships.
% Courant--Fischer/Weyl checked 2026-10-28: developed here in June 2026;
% August donor df78246f965881705c9301f96feddc270c98ceb1 predates the
% accepted Tau Ceti port 7ebf17ba04725c1a575ecf3d099ce0630cd949a0 (#13599).
% See COURANT_FISCHER_LINEAGE_AUDIT.md for the proof comparison and sources.

The proofs required us to formalize a significant amount of operator
theory, namely: principal angles and direct rotations, approximation numbers
and symmetric operator ideals, Borel spectral calculus, and Sylvester
estimates for unbounded operators. These developments build on Mathlib and
adapt Spectra's projection-valued-measure and Borel-calculus constructions
\citep{mathlib2020,bornemann2026spectra}.
These operator-theory foundations include our
Courant--Fischer and Weyl proofs, rectangular unitarily invariant seminorms,
and projection-valued measures, which were subsequently adapted into Tau
Ceti \citep{tauceti2026}. The measure constructions retain their earlier Spectra ancestry.
Section~\ref{sec:related} explains the mathematical history and formal
ancestry; Section~\ref{sec:foundations} describes the foundations used in
our proofs.

The same principal-angle and operator-ideal results support a further
optimization theorem. In Section~\ref{sec:example-transport}, we prove that
minimizing trace-norm displacement over all real orthogonal maps carrying
one subspace onto another reduces to weighted matching on principal angles.

\paragraph{Contributions.}
\begin{enumerate}
\item \textbf{Davis--Kahan formalization.} We prove 28 of the 29 source
results and refute the full-displacement assertion
\citep[Proposition~4.4]{davis1970rotation}. The sine-theta theorem covers real and
complex Hilbert spaces, unbounded self-adjoint operators, the source's
spectral-separation alternatives, and operator-ideal norms where finite
(\cref{sec:coverage,sec:sine-theta}).
\item \textbf{Formalized operator-theoretic foundations.} We formalized the
geometric and analytic results required by these theorems, adapted existing
formal mathematics, and contributed foundations later used in Tau Ceti
(\cref{sec:related-formal,sec:foundations}).
\item \textbf{Source correction.} We identify a false comparison
of Ky Fan sums in the printed claim \citep[Proposition~4.4]{davis1970rotation},
give a real counterexample at every positive acute angle,
and prove the Q-norm consequence of the valid squared-displacement theorem
(\cref{sec:prop44}).
\item \textbf{A reduction to weighted matching.} We reduce unrestricted
real trace-norm transport to a weighted-matching problem on principal angles
(\cref{sec:example-transport}).

\end{enumerate}

\section{Mathematical background}
\label{sec:related}

The Davis--Kahan angle estimates use three constructions. Principal angles
measure the error between subspaces; a residual measures how well a proposed
subspace represents the operator; and a Sylvester estimate bounds the first
quantity by the second under spectral separation. Operator ideals allow the
same argument to express several norms. We describe the development of these
constructions before stating the formalized sine-theta theorem.

Existing historical surveys cover these developments. Paige and Wei trace
the cosine--sine decomposition from Jordan's subspace geometry through the
principal ideas assembled by Davis and Kahan to Stewart's later completion
and explicit formulation \citep{paige1994history}.
Simon's historical survey of Kato's work in non-relativistic quantum mechanics
traces Sz.-Nagy's unitary construction for pairs of projections and Kato's
later algebraic formulation \citep[\S5]{simon2018kato}.
We use this account for the history of the projection transport described in
\cref{sec:related-projections}.
B\"ottcher and Spitkovsky survey pairs of projections and their operator algebras.
Their historical discussion relates Halmos's 1969 block representation to the
earlier formulations of Dixmier (1948) and Davis (1958), and to the related
projection constructions in Davis and Kahan's 1970 paper
\citep[Remark~1.3]{bottcher2010gentle}.
We use these accounts together with the primary
sources. Table~\ref{tab:historical-lineage} summarizes the mathematical
constructions used in our formalization.

\begin{table}[!htbp]
\centering
\small
\setlength{\tabcolsep}{5pt}
\renewcommand{\arraystretch}{1.18}
\begin{tabularx}{\linewidth}{@{}>{\raggedright\arraybackslash}p{0.26\linewidth}>{\raggedright\arraybackslash}X>{\raggedright\arraybackslash}p{0.28\linewidth}@{}}
\toprule
\textbf{Lineage} & \textbf{Mathematical role} & \textbf{Formal dependency}\\
\midrule
Jordan (1875); Davis (1958); Halmos (1969) & Principal angles, invariant subspace geometry,
  and pairs of projections & Angle operators; projection blocks\\
Rayleigh (1877); Ritz (1909); Weyl (1912) & Approximate eigenpairs,
  variational spectra, and eigenvalue stability & Trial coordinates;
  min--max and eigenvalue inequalities\\
Rellich (1937--1942); Sz.-Nagy (1947); Kato (1966) & Moving spectral projections and unitary
  transport between nearby subspaces & Spectral calculus; polar
  direct rotation\\
Sylvester (1884); Rosenblum (1956, 1969) & Spectral separation controls an
  off-diagonal operator & Domain-aware Sylvester estimates\\
Von Neumann (1937); Schatten (1950); Ky Fan (1951) & Comparison of spectral errors
  through singular-value gauges & Approximation numbers;
  symmetric operator ideals\\
Davis (1963, 1965); Davis--Kahan (1969, 1970) & Residual-to-angle bounds, tangent estimates,
  rotation extremality, and spectral branch selection & Source-level
  sine, tangent, and direct-rotation theorems\\
\bottomrule
\end{tabularx}
\caption{Mathematical sources used in the Davis--Kahan formalization.
Dates identify the publications discussed in the text.}
\label{tab:historical-lineage}
\end{table}
\FloatBarrier

\subsection{Principal angles}
\label{sec:related-classical}

Camille Jordan studied the geometry of linear subspaces in $n$ dimensions in
  1875 \citep{jordan1875essai}.

For two lines, a single acute angle describes their relative position.
For subspaces of higher dimension, one needs a sequence of \emph{principal
  angles}.
The successive principal vectors are selected under orthogonality constraints,
  so their angles are independent of the coordinates used to describe either
  subspace.
This separates the geometry of a subspace from the arbitrary choice of its
  orthonormal basis.

For finite-dimensional subspaces $\sym{Usource}$ and $\sym{Vtarget}$ of equal dimension, choose
isometries $\sym{Xframe}$ and $\sym{Yframe}$ with those ranges.
The singular values of the overlap matrix are the principal cosines:
\begin{equation}
  \sigma_j(\like{Xframe}{X^*}\sym{Yframe})=\cos\like{theta}{\theta_j},
  \qquad 0\leq\like{theta}{\theta_j}\leq\frac{\pi}{2}.
  \label{eq:history-principal-cosines}
\end{equation}

The corresponding principal vectors $\like{SourceVector}{x_j}\in \sym{Usource}$ and $\like{TargetVector}{y_j}\in \sym{Vtarget}$ satisfy
$\langle \like{SourceVector}{x_i},\like{TargetVector}{y_j}\rangle=\delta_{ij}\cos\like{theta}{\theta_j}$ in suitably oriented
coordinates.  For nonzero acute angles, each pair determines a real
two-dimensional plane in which the relative position is an ordinary
rotation.  An entire pair of subspaces is described by these planes,
together with their shared and orthogonal components.

The same geometric data appeared in statistics: Hotelling's 1936 canonical
  correlation analysis sought pairs of maximally correlated linear combinations
  of two sets of variables \citep{hotelling1936relations}.
After appropriate whitening, the canonical correlations are singular values of a
  cross-coordinate map.
In numerical linear algebra, Bj\"orck and Golub developed algorithms for
  computing principal angles and vectors in 1973 \citep{bjorck1973numerical}.
Their computational formulation makes the basis-invariant angles available from
  orthonormal matrix factorizations.

A further formulation uses the cosine--sine decomposition of a
partitioned orthogonal or unitary matrix.
The blocks of the decomposition contain compatible diagonal cosine and sine
  factors.
Davis and Kahan describe the geometry using orthogonal projections,
angle operators, and the direct rotation carrying one subspace onto the other
\citep[Section~3]{davis1970rotation}. Their work assembled the principal ideas
of the cosine--sine decomposition; Stewart completed and explicitly formulated
the matrix decomposition in 1977
\citep{stewart1977perturbation,paige1994history}.
Paige and Wei explain how this decomposition and direct rotation represent
the same principal-angle data
\citep{paige1994history}. Our formalization uses projections and functional
calculus for the Hilbert-space results and principal-plane coordinates for
the finite-dimensional counterexample.

\subsection{Eigenvalue approximation}
\label{sec:related-ritz-weyl}

Jordan's angles answered how to compare subspaces once they are given.
An independent question came from mechanics and numerical analysis:
how can one approximate an eigenvalue or eigenmode when the underlying
differential equation cannot be solved exactly?  In the nineteenth
century Rayleigh estimated vibration frequencies using trial shapes
and their energy ratios \citep{rayleigh1877sound}.  Ritz developed a systematic finite-dimensional
variational approximation in the work published in 1909
\citep{ritz1909variations,leissa2005historical}.  The method selects
a trial space and studies the action of the operator through that
space. A trial space is a chosen space of candidate approximate solutions.
For an eigenvalue problem, one searches for an approximate eigenvector among
its vectors. This reduces an infinite-dimensional problem to a smaller
problem whose residual can then be assessed.

For a Hermitian matrix $\sym{A}$, the Rayleigh quotient
$\langle \sym{A}\like{v}{x},\like{v}{x}\rangle/\langle \like{v}{x},\like{v}{x}\rangle$ expresses the scalar action
of $\sym{A}$ along a trial vector. A normalized trial vector $\sym{v}$ gives the
Rayleigh--Ritz approximate eigenvalue $\sym{mu}=\langle \sym{A}\sym{v},\sym{v}\rangle$ and the
residual $\sym{A}\sym{v}-\sym{mu} \sym{v}$. The Courant--Fischer min--max
principle organizes the eigenvalues through extremal quotients over
subspaces \citep{bhatia1997matrix}.  Weyl's eigenvalue comparison
inequalities then control the motion of ordered eigenvalues under a
Hermitian perturbation \citep{weyl1912eigenwerte,bhatia1997matrix}.  In modern notation,
if $\sym{A}$ and $\sym{Acurrent}=\sym{A}+\sym{H}$ are Hermitian $n$-by-$n$ matrices and their eigenvalues
are arranged in the same order, then
\begin{equation}
  \max_j |\like{SpectralValue}{\lambda_j}(\sym{A})-\like{SpectralValue}{\lambda_j}(\sym{Acurrent})|\leq\|\sym{H}\|_{\mathrm{op}}.
  \label{eq:history-weyl}
\end{equation}
This is an estimate on \emph{eigenvalues}; it does not tell us which
individual eigenvector of $\sym{A}$ corresponds to which eigenvector of $\sym{Acurrent}$.

The limitation already appears in two dimensions.  For any $\varepsilon>0$,
put
\begin{equation}
 \like{A}{A_\varepsilon}=\begin{pmatrix}0&0\\0&\varepsilon\end{pmatrix},
 \qquad
 \like{Acurrent}{\widehat A_\varepsilon}=\begin{pmatrix}\varepsilon&0\\0&0\end{pmatrix}.
 \label{eq:history-swap-example}
\end{equation}
The matrices have identical ordered eigenvalues and
$\|\like{A}{A_\varepsilon}-\like{Acurrent}{\widehat A_\varepsilon}\|_{\mathrm{op}}=\varepsilon$,
yet their lowest-eigenvalue lines are perpendicular. The gap between the
two eigenvalues of each matrix is $\varepsilon$, equal to the perturbation
norm. An eigenspace estimate therefore needs to compare the perturbation
size with the spectral separation.  For a cluster with repeated eigenvalues, an
entire invariant subspace is likewise more natural than a selected
orthonormal eigenbasis.

Rellich developed a rigorous perturbation theory of self-adjoint operators
in papers published from 1937 to 1942. Kato's 1966 book gives a systematic
account \citep{simon2018kato,kato1966perturbation}. The Davis--Kahan residual
bounds address a related approximation question: given a trial space and its
residual, spectral separation bounds its angle from an exact spectral
subspace. The estimate can therefore assess the accuracy of a
Rayleigh--Ritz approximation.

\subsection{Two orthogonal projections}
\label{sec:related-two-projections}

For a closed subspace of a Hilbert space $\mathcal H$, the orthogonal
projection remains a bounded operator even when the subspace is
infinite dimensional.  Denote the projections onto two closed
subspaces by $\sym{Psource}$ and $\sym{Qtarget}$. Compressing an operator to $\DKRange(\sym{Psource})$ means
restricting its input to that subspace and projecting its output back with
$\sym{Psource}$. The positive compression of $\sym{Qtarget}$ is therefore $\sym{Psource}\sym{Qtarget}\sym{Psource}$ restricted to
$\DKRange(\sym{Psource})$. It is the operator version of the squared overlap matrix: in finite dimensions its eigenvalues are
$\cos^2\like{theta}{\theta_j}$.  In Hilbert space it need not have a discrete
spectrum.  Functional calculus nevertheless defines its positive
square root and the associated cosine and sine of an angle operator.

Write $\sym{Theta0}$ for this directed angle operator on
$\DKRange(\sym{Psource})$, and $\sym{Theta1}$ for the corresponding
operator for the complementary pair $I-\sym{Psource},I-\sym{Qtarget}$.
The full angle operator is
\[
 \sym{Theta}=\operatorname{diag}(\sym{Theta0},\sym{Theta1})
 \quad\text{on }\DKRange(\sym{Psource})\oplus\DKKernel(\sym{Psource}).
\]
For $c>0$, the comparisons $\sym{Theta}\le c$ and $\sym{Theta}<c$ mean that its
operator norm is at most, or strictly less than, $c$, respectively.
In finite dimensions this norm is the largest eigenvalue of either directed
angle block, comparing the source with the target and their orthogonal
complements.
This includes right angles from unmatched directions: if a source vector is
orthogonal to the target, its cosine compression is zero and its directed
angle is $\pi/2$. The compression construction applies even when the subspace
dimensions differ. Target directions orthogonal to the source appear as
right angles in the complementary block. Later, isometric frames express these angle blocks on
coordinate Hilbert spaces.

The relative position of two closed subspaces was studied by Dixmier
in 1948 and by Davis in his 1958 paper \emph{Separation of Two Linear
Subspaces} \citep{dixmier1948position,davis1958separation}.  Davis's
closeness and separation operators place the cosine/sine data in an
operator-theoretic setting.  Halmos's 1969 \emph{Two Subspaces} gives a
block representation of the two projections
\citep{halmos1969two}.  One first splits off the four intersection
spaces corresponding to the combinations of membership in the ranges
and kernels of $\sym{Psource}$ and $\sym{Qtarget}$.  On their remaining, generic part, the pair
is unitarily equivalent to
\begin{equation}
 \like{Psource}{P_{\mathrm{gen}}}=\begin{pmatrix}I&0\\0&0\end{pmatrix},
 \qquad
 \like{Qtarget}{Q_{\mathrm{gen}}}=\begin{pmatrix}\like{CosAngle}{C^2}&\like{CosAngle}{C}\like{SinAngle}{S}\\\like{CosAngle}{C}\like{SinAngle}{S}&\like{SinAngle}{S^2}\end{pmatrix},
 \qquad \like{CosAngle}{C^2}+\like{SinAngle}{S^2}=I,
 \quad \like{CosAngle}{C}\like{SinAngle}{S}=\like{SinAngle}{S}\like{CosAngle}{C},
 \label{eq:history-halmos-block}
\end{equation}
where $\like{CosAngle}{C}$ and $\like{SinAngle}{S}$ are commuting positive contractions on a common
Hilbert space.  Thus an operator-valued angle replaces a finite list
of principal angles.  B\"ottcher and Spitkovsky survey this
representation and its consequences for the algebra generated by two
projections \citep{bottcher2010gentle}.

The four intersection spaces have a second role. The right-angle defect spaces
\[
 \sym{CrossedSourceDefect}=\DKRange(\sym{Psource})\cap\DKKernel(\sym{Qtarget}),
 \qquad
 \sym{CrossedTargetDefect}=\DKKernel(\sym{Psource})\cap\DKRange(\sym{Qtarget})
\]
contain source directions orthogonal to the target, and target directions
orthogonal to the source, respectively. At acute angles
both disappear.  For a general pair they affect whether a specified
intertwining operator exists. An operator $\sym{IntertwinerT}$ intertwines $\sym{Psource}$ and $\sym{Qtarget}$ when
$\sym{IntertwinerT}\sym{Psource}=\sym{Qtarget}\sym{IntertwinerT}$: applying $\sym{IntertwinerT}$ after the source projection agrees with applying the
target projection after $\sym{IntertwinerT}$. In particular, $\sym{IntertwinerT}$ maps the source range into
the target range and its orthogonal complement into the target complement.
The existence criteria for direct rotations include conditions on the
dimensions of these defect spaces \citep[Section~3]{davis1970rotation}.

A finite-dimensional orthogonal matrix decomposes into invariant real
lines and two-dimensional rotation planes.  In a general Hilbert space an
angle operator can instead have continuous spectrum.  The spectral
calculus of the pair of projections takes the place of a finite list
of angles, and the intersection and defect spaces cannot be discarded.
Our Hilbert-space direct-rotation proofs therefore use projection identities
and functional calculus. In finite dimensions these identities recover the
principal-plane construction.

\subsection{Transport between nearby spectral subspaces}
\label{sec:related-projections}

A perturbation can change the invariant subspace associated with an isolated
part of a self-adjoint operator's spectrum. Let $\sym{Psource}$ and
$\sym{Qtarget}$ be the orthogonal projections onto the corresponding spectral
subspaces before and after perturbation. To compare these subspaces, we seek
a unitary carrying $\DKRange(\sym{Psource})$ onto
$\DKRange(\sym{Qtarget})$ that depends continuously, or analytically, on the
projections.
Sz.-Nagy developed the nearby-projection construction in work
published in 1947 \citep{sznagy1947perturbations,simon2018kato}.  Kato gave a streamlined
algebraic construction, described in his 1955 report and later in
\emph{Perturbation Theory for Linear Operators}
\citep[Chapter~I, \S4.6, equations~(4.33)--(4.42)]{kato1966perturbation};
see also \citet[Theorem~5.1]{simon2018kato}.

For orthogonal projections the useful intertwiner is
\begin{equation}
 \sym{IntertwinerT}=\sym{Qtarget}\sym{Psource}+(I-\sym{Qtarget})(I-\sym{Psource}), \qquad \sym{IntertwinerT}\sym{Psource}=\sym{Qtarget}\sym{IntertwinerT}.
 \label{eq:history-kato-T}
\end{equation}
The identity
\begin{equation}
 \like{IntertwinerT}{T^*}\sym{IntertwinerT}=I-(\sym{Psource}-\sym{Qtarget})^2
 \label{eq:history-kato-square}
\end{equation}
shows why $\|\sym{Psource}-\sym{Qtarget}\|<1$ is the natural threshold: it makes $\like{IntertwinerT}{T^*}\sym{IntertwinerT}$
strictly positive and invertible.  Hence
\begin{equation}
 \sym{Rrotation}=\sym{IntertwinerT}\bigl(I-(\sym{Psource}-\sym{Qtarget})^2\bigr)^{-1/2},
 \qquad \sym{Rrotation}\sym{Psource}\like{Rrotation}{D^*}=\sym{Qtarget}
 \label{eq:history-kato-unitary}
\end{equation}
is unitary. This is the polar decomposition $\sym{IntertwinerT}=\sym{Rrotation}|\sym{IntertwinerT}|$, where
$|\sym{IntertwinerT}|=(\like{IntertwinerT}{T^*}\sym{IntertwinerT})^{1/2}$ is positive. The factor $|\sym{IntertwinerT}|$ describes stretching;
the unitary factor $\sym{Rrotation}$ preserves lengths. Here $\sym{IntertwinerT}$ is invertible, so
$\sym{Rrotation}=\sym{IntertwinerT}|\sym{IntertwinerT}|^{-1}$. In principal-plane coordinates it is the direct rotation through the
principal angle, acting trivially on common directions.

The transport equation
$\sym{Wcompetitor}\sym{Psource}\sym{Wcompetitor}^*=\sym{Qtarget}$
has many unitary solutions.
The polar-factor formula selects one determined by the two projections
and continuously connected to the identity in the acute regime.  That
choice makes the construction valuable in perturbation theory: the
transport identifies moving spectral subspaces while avoiding an
arbitrary choice of orthonormal bases. A pair is called acute here when
$\|\sym{Psource}-\sym{Qtarget}\|<1$ (equivalently, the largest principal angle is less than $\pi/2$
for equal-dimensional finite subspaces). For an orthogonal or unitary transporter $\sym{Wcompetitor}$, the operator
$I-\sym{Wcompetitor}$ records the displacement of every vector. Davis and Kahan ask which
norms of this displacement are minimized by the direct rotation
\citep[Section~4]{davis1970rotation}.

Kato used such transformations to identify nearby spectral
subspaces of analytic operator families, thereby reducing the study
of an isolated moving eigenvalue cluster to a fixed coordinate
space.  Simon's historical account explains the earlier Sz.-Nagy
construction, Kato's formula, and subsequent work on pairs of
projections \citep[\S5]{simon2018kato}. Our proofs of direct-rotation
existence, uniqueness, and symmetry use the square-root and
polar-decomposition identities in
\cref{eq:history-kato-square,eq:history-kato-unitary}. The condition
$\|\sym{Psource}-\sym{Qtarget}\|<1$ applies in arbitrary Hilbert dimension.

\subsection{Spectral projections and residuals}
\label{sec:related-spectral}

The spectral theorem represents a self-adjoint operator as
$\sym{A}=\int\like{SpectralValue}{\lambda}\,d\sym{SpectralMeasure}(\like{SpectralValue}{\lambda})$ for a projection-valued measure $\sym{SpectralMeasure}$.
For a Borel set $\sym{SpectralSet}$, the operator $\sym{SpectralMeasure}(\sym{SpectralSet})$ is the orthogonal projection onto
the corresponding spectral subspace. This gives a second way of obtaining
a subspace: select a region of the spectrum and take the range of its
spectral projection. The functional-calculus
framework developed from the work of Riesz, von Neumann, and others;
Kato made spectral projections central to the perturbation theory of
closed operators \citep{kato1966perturbation,simon2018kato}.
Unlike the finite-dimensional eigenspace picture, a spectral
subspace need not be spanned by a finite list of eigenvectors.
A self-adjoint operator may be unbounded, while each of its spectral
projections is bounded and everywhere defined. The Hellinger--Toeplitz
theorem states that an everywhere-defined symmetric operator on a Hilbert
space is bounded \citep{hellinger1910grundlagen}. Thus an unbounded
self-adjoint operator has a proper domain. Its action $\sym{A}x$ is defined only
for $x\in\DKDom(\sym{A})$, so residual and intertwining equations must include
the required domain inclusions.

In the complex case, a number $z$ belongs to the resolvent set when
$zI-\sym{A}$ has a bounded, everywhere-defined inverse. That inverse,
$(zI-\sym{A})^{-1}$, is the resolvent at $z$. For an isolated bounded spectral
cluster, choose a positively oriented contour $\Gamma$ in the resolvent
set enclosing the cluster and no other spectrum. Its projection is
\begin{equation}
 \like{Pexact}{P_\Gamma}=\frac{1}{2\pi i}\int_\Gamma(zI-\sym{A})^{-1}\,dz.
 \label{eq:history-riesz-projection}
\end{equation}
The resolvent is bounded at each point of $\Gamma$ even when $\sym{A}$ is unbounded.
This makes the spectral projection a stable object for perturbation analysis,
  while preserving the separate domain requirements whenever $\sym{A}$ itself appears
  in an identity.
More general Borel spectral projections cover finite and half-infinite intervals,
  including spectral configurations for which a single isolated
  finite-eigenvalue cluster is not the right model.
Our formalization uses Borel functional calculus to construct projections for
  these spectral sets, including for unbounded self-adjoint operators.

The Courant--Fischer min--max principle and Weyl-type eigenvalue
perturbation bounds provide complementary information about spectral
\emph{values} \citep{bhatia1997matrix,kato1966perturbation}.  They do
not by themselves identify a stable basis inside a multiple-eigenvalue
cluster, or quantify the movement of its invariant subspace.  The Courant--Fischer and Weyl results that we formalized describe extremal
Rayleigh quotients and eigenvalue stability. The angle estimates add control
of the associated subspaces.

Davis studied the rotation of eigenvectors under perturbation in
1963 and 1965 \citep{davis1963rotation,davis1965rotation}.  The
1969 Davis--Kahan announcement preceded their extended 1970 treatment
\citep{daviskahan1969bounds,davis1970rotation}.  Their problem has
two closely related formulations.  In one, the original self-adjoint
operator is perturbed and the movement of an exact spectral subspace
is estimated.  In the other, one has an approximate trial subspace and
a computable residual: how far can it be from the desired exact
spectral subspace?  Unlike a comparison of individual eigenvectors,
these are meaningful even when an eigenvalue has multiplicity and its
orthonormal eigenbasis is nonunique.

Write $\sym{A}$ for a reference operator and $\sym{Acurrent}=\sym{A}+\sym{H}$ for the perturbed
ambient operator. We use this convention in the sine-theta statement below.
Let $\sym{E0}$ embed a trial-coordinate Hilbert space into the ambient space
isometrically, and let $\sym{A0}$ be a self-adjoint operator on those coordinates.
The residual is the bounded operator agreeing, where defined, with
\begin{equation}
 \sym{R}=\sym{Acurrent}\sym{E0}-\sym{E0}\sym{A0}.
 \label{eq:history-residual}
\end{equation}
If $\sym{R}=0$, applying $\sym{Acurrent}$ after the embedding agrees with applying
$\sym{A0}$ before it: the trial subspace is invariant and $\sym{A0}$ represents its
operator. For $\sym{R}\ne0$, the residual measures the failure of this identity.  A second isometric coordinate map for the
desired exact spectral space defines its orthogonal projection $\sym{Pexact}$.
The positive operator
\begin{equation}
 \sym{sinTheta0}=\bigl(I-\like{E0}{E_0^*}\sym{Pexact}\sym{E0}\bigr)^{1/2}
 \label{eq:history-angle-sine}
\end{equation}
measures the component of a trial vector lying outside the exact
subspace.  The critical distinction for the formal proof is that this
positive sine acts on the trial-coordinate space, whereas
$(I-\sym{Pexact})\sym{E0}$ maps into the ambient space.  For an operator $B$, its positive modulus $|B|=(B^*B)^{1/2}$ acts on its
input space and records its stretching factors. The modulus of $(I-\sym{Pexact})\sym{E0}$
is $\sym{sinTheta0}$, so the two maps have equal unitarily invariant ideal
norms whenever those norms are finite.

\subsection{Sylvester estimates}
\label{sec:related-sylvester}

The matrix equation $\sym{SylvesterLeft}\sym{X}-\sym{X}\like{SylvesterRight}{A}=\sym{SylvesterRHS}$ is associated with Sylvester's 1884
work \citep{sylvester1884equation}.  Rosenblum studied separated
operator spectra in 1956 and the equation with self-adjoint coefficients
in 1969 \citep{rosenblum1956equation,rosenblum1969selfadjoint}.

The scalar prototype is $(\like{SylvesterLeft}{b}-\like{SylvesterRight}{a})\like{X}{x}=\like{SylvesterRHS}{y}$: if $|\like{SylvesterLeft}{b}-\like{SylvesterRight}{a}|\geq\sym{delta}>0$,
then $\sym{delta}|\like{X}{x}|\leq|\like{SylvesterRHS}{y}|$.  In the operator problem the scalar
multipliers are replaced by the spectral actions of two operators.
The operator estimate has the form $\sym{delta} \sym{NormN}(\sym{X})\leq \sym{NormN}(\sym{SylvesterLeft}\sym{X}-\sym{X}\like{SylvesterRight}{A})$.
Its hypotheses specify the spectral positions of $\like{SylvesterRight}{A}$ and $\sym{SylvesterLeft}$, their domains
when unbounded, and the unitarily invariant norm $\sym{NormN}$.

Let $\sym{F1}$ give isometric coordinates for the unwanted complementary
spectral subspace, with self-adjoint block $\sym{Lambda1}$.  Writing
$\sym{X}=\like{F1}{F_1^*}\sym{E0}$ converts the residual identity into a Sylvester equation:
\begin{equation}
 \sym{Lambda1}\sym{X}-\sym{X}\sym{A0}=\like{F1}{F_1^*}\sym{R}
 \label{eq:history-sylvester}
\end{equation}
on the appropriate domain.  Thus $\sym{X}$ describes the geometric leakage
of the trial space into the unwanted block, while the right side is
controlled by the computable residual.  Under the interval/exterior
or ordered spectral-gap hypotheses of the sine-theta theorem,
separated Sylvester estimates give
\begin{equation}
 \sym{delta}\,\sym{NormN}(\sym{X})\leq \sym{NormN}(\like{F1}{F_1^*}\sym{R})\leq \sym{NormN}(\sym{R}).
 \label{eq:history-separation-bound}
\end{equation}
This is the bridge from a spectral-gap assumption to the bound on the
sine of a subspace angle.

The coefficient of $\sym{NormN}(\sym{R})/\sym{delta}$ in this bound is one.
This coefficient uses the ordered or interval/exterior spectral
positions specified by Davis and Kahan. Arbitrarily interlaced spectra can
require a different constant. Our theorem includes both of the source's
spectral configurations, including ordered cases where both operators are
unbounded.  A formal implementation
therefore needs the domains of unbounded self-adjoint operators,
their spectral projections, strong spectral cutoffs, and the
ideal-norm inequalities used to pass through bounded approximations.

\subsection{Operator ideals and displacement norms}
\label{sec:related-norms}

The geometry of a pair of subspaces is fixed, but there are many ways
to assign a cost to a unitary transporting one onto the other.
Von Neumann's 1937 study of unitarily invariant matrix norms identified
the role of symmetric gauge functions of singular values
\citep{vonneumann1937matrix}. Schatten's 1950 monograph developed normed
operator ideals; Ky Fan's 1951 eigenvalue and singular-value inequalities
provided the partial-sum comparisons underlying majorization
arguments \citep{schatten1950crossspaces,fan1951maximum,bhatia1997matrix}.
The singular values $s_j(Z)$ are the nonnegative square roots of the
eigenvalues of $Z^*Z$, ordered decreasingly. They record the stretching
factors of the matrix $Z$. For $1\leq p<\infty$, the Schatten norms include
\begin{equation}
 \|Z\|_{S_p}=\left(\sum_j s_j(Z)^p\right)^{1/p},
 \qquad
 \|Z\|_{S_1}=\sum_j s_j(Z),
 \quad
 \|Z\|_{S_2}=\left(\sum_j s_j(Z)^2\right)^{1/2}.
 \label{eq:history-schatten}
\end{equation}
An operator ideal is a linear class of operators closed under multiplication
on either side by bounded operators. A normed ideal also controls the norm
of these products. For compact operators in infinite dimensions, the
Schatten $p$ ideal consists of those whose singular values have a finite
$p$-power sum. Thus a trace-norm inequality requires trace-class membership;
this is an instance of the source's qualification that its norms be defined.

The largest singular value gives the operator norm, the Euclidean
norm of the singular-value sequence gives the Hilbert--Schmidt norm,
and its sum gives the trace norm. These norms can rank transporters
differently. In Section~\ref{sec:prop44}, the two displacement singular-value
lists are $(2\sin(\sym{theta}/2),\ldots,2\sin(\sym{theta}/2))$ with four entries and
$(2\sin\sym{theta},2\sin\sym{theta},0,0)$. The second has a smaller sum and a larger
maximum. Ky Fan norms compare partial sums of the ordered singular values;
these comparisons determine which norm inequalities follow.

Davis established Hilbert--Schmidt and finite-dimensional Schatten
extremal properties for exponents at least two in his 1958 subspace
paper \citep[Theorems~7.3--7.4]{davis1958separation}.  Davis and
Kahan compare three objectives \citep[Section~4]{davis1970rotation} for a transporter $\sym{Wcompetitor}$ carrying
$\sym{Usource}$ onto $\sym{Vtarget}$: its restricted displacement $(I-\sym{Wcompetitor})\like{Psource}{P_U}$, its squared full
displacement $(I-\sym{Wcompetitor})^*(I-\sym{Wcompetitor})$, and its full displacement $I-\sym{Wcompetitor}$. The last acts
on every ambient vector, including vectors outside $\sym{Usource}$. They gave a
complex-phase modification that decreases its trace norm \citep[Example~4.2]{davis1970rotation}.
The full-displacement claim \citep[Proposition~4.4]{davis1970rotation}
asserts that direct rotation minimizes
full displacement for real subspaces if their principal angles are no
larger than $\pi/3$.  The four-dimensional example in
\cref{sec:prop44} gives a smaller trace norm at every positive acute
angle.  It exposes the failed step: adding singular-value sums of
restrictions to separate principal planes is not a valid global lower
bound for an arbitrary ambient transporter.

The subsequent minimal-distortion theory of Bolotnikov, Li, and Rodman
studied the same class of subspace-transport constraints, including
Frobenius and Q-norm optimality and explicit subquadratic
counterexamples \citep{bolotnikov2001minimal}.  Later work on
Grassmannian metrics and canonical-angle inequalities also separates
intrinsic paths between subspaces from minimizing the extrinsic
operator displacement \citep{qiu2006grassmann,zhang2010canonicalangles}.
These results give the norm-theoretic context for the source
correction and for the trace-norm matching problem in
\cref{sec:example-transport}.

\subsection{The Davis--Kahan angle estimates}
\label{sec:related-synthesis}

Davis's papers of 1963 and 1965 and the Davis--Kahan announcement of 1969
preceded the extended 1970 treatment
\citep{davis1963rotation,davis1965rotation,daviskahan1969bounds,davis1970rotation}.
For the sine-theta estimate, let $\sym{R}=\sym{Acurrent} \sym{E0}-\sym{E0}\sym{A0}$ be the residual
and $\sym{Pexact}$ the desired exact spectral projection. The operator $(I-\sym{Pexact})\sym{E0}$
extracts the part of a trial vector outside that spectral subspace. In
isometric coordinates $\sym{F1}$ for the complementary space, this component is
represented by $\sym{X}=\like{F1}{F_1^*}\sym{E0}$. The residual identity gives
\cref{eq:history-sylvester}, and a separated Sylvester estimate bounds $\sym{NormN}(\sym{X})$.
Since $\sym{S}=\sym{F1}\sym{X}$ and $\sym{F1}$ is an isometry,
$\sym{NormN}(\sym{X})=\sym{NormN}(\sym{S})$.
The modulus identity $|\sym{S}|=\sym{sinTheta0}$ gives
$\sym{NormN}(\sym{S})=\sym{NormN}(\sym{sinTheta0})$.

The resulting chain of equalities and inequalities is
\begin{equation}
 \underbrace{\sym{delta} \sym{NormN}(\sym{sinTheta0})}_{\text{subspace error}}
 =\sym{delta} \sym{NormN}(\sym{X})
 \leq \sym{NormN}(\like{F1}{F_1^*}\sym{R})
 \leq\underbrace{\sym{NormN}(\sym{R})}_{\text{computable residual}}.
 \label{eq:history-unified-chain}
\end{equation}
The first equality follows from the positive-modulus identity for the
complementary projection. The first inequality applies the Sylvester
estimate to $\sym{X}$; the second uses $\sym{NormN}(BC)\leq\|B\|_{\mathrm{op}}\sym{NormN}(C)$ and
$\|\like{F1}{F_1^*}\|_{\mathrm{op}}\leq1$. For unbounded operators the Sylvester
identity holds on the trial operator's domain, with a bounded residual.
These are the geometric, spectral, and norm estimates formalized in our
proof of the theorem.

The tangent and double-angle estimates use related projection blocks.
For example, an estimate for $\sin(2\sym{Theta})$ allows angles near zero and
near $\pi/2$. The spectral-selection results \citep[Section~8]{davis1970rotation}
identify which spectral subspaces give the small-angle alternative.
The direct-rotation analysis \citep[Section~4]{davis1970rotation}
studies the orthogonal or unitary map carrying one subspace onto another
and asks which displacement norms the direct rotation minimizes.

\subsection{Numerical and statistical applications}
\label{sec:related-numerical}

The 1970 theory influenced more than the analysis of exact
self-adjoint eigenvectors.  Wedin's 1972 singular-subspace
perturbation bounds address the corresponding problem for singular
vectors of general rectangular matrices
\citep{wedin1972perturbation}.  Bj\"orck and Golub's angle algorithms
made the subspace quantities computable
\citep{bjorck1973numerical}.  Stewart's work on pseudoinverses,
projections, and least-squares sensitivity helped connect these
geometric notions with matrix computations
\citep{stewart1977perturbation,stewart1990matrix}.
Paige and Wei's 1994 history shows how the cosine--sine decomposition
links the principal-angle geometry of Davis and Kahan with numerical
matrix factorizations \citep{paige1994history}.

The same subspace-versus-eigenvector distinction remains central in
statistical spectral estimation.  For example, Yu, Wang, and
Samworth give a statistical variant whose eigenspace error bound uses
separation in the population spectrum \citep{yu2015davis}, while Chen and
collaborators survey spectral methods across statistical learning
problems \citep{chen2021spectral}.  In that setting, an empirical covariance matrix approximates a population
operator and the perturbation norm controls the error of estimated
eigenspaces. Our general theorem supplies the deterministic angle estimate;
using it in a statistical result also requires the corresponding probabilistic
control of the perturbation.

\subsection{Foundations used in our formalization}
\label{sec:related-formal}

Each historical component above corresponds to a specific requirement
in the formalization.  Jordan's angle geometry and Halmos's
projection blocks lead to positive principal-angle operators rather
than chosen eigenbases.  The Sz.-Nagy--Kato intertwiner requires
polar decomposition, square roots, and real/complex transfer.
Spectral projections require Borel functional calculus, including for
unbounded self-adjoint operators.  The Sylvester proof requires a
representation carrying operator domains and a controlled passage
through spectral cutoffs.  Schatten and Ky Fan theory requires
approximation numbers and symmetric operator ideals to express
inequalities with the source's norm generality.

We formalized principal-angle geometry, symmetric operator ideals, and
unbounded Sylvester estimates using Lean~4 and Mathlib
\citep{demoura2021lean4,mathlib2020}. Mathlib supplies the algebraic,
topological, and Hilbert-space framework and continuous functional calculus
\citep{dedecker2025cfc}. Our projection-valued-measure and Borel-calculus
work also adapts Adam Bornemann's Spectra formalization
\citep{bornemann2026spectra}.

Our Courant--Fischer and Weyl development, rectangular unitarily invariant
seminorms, and projection-valued measures were subsequently adapted into
Tau Ceti \citep{tauceti2026}. The last of these carries the earlier Spectra ancestry. We also use
operator-theory results developed independently in Tau Ceti, including its
canonical polar decomposition.
% Evidence: dev/tauceti/dkps-upstream-contribution-census.md, PRs 13599,
% 14241, 17527 (adaptations), and 16127/16450 (independent polar proofs).
% AUTHORING_NOTES.md indexes accepted-source and campaign evidence.
% Mathematical origin and formal proof ancestry are separate attributions.

This history also separates the two mathematical questions that meet
in the project.  The sine-theta theorem estimates \emph{how far} an
approximate or perturbed spectral subspace lies from an exact one.
The direct-rotation problem asks \emph{which} ambient transformation
carries one given subspace onto another at minimum displacement cost.
Our proofs of the optimization results use the principal-angle,
projection, and operator-ideal foundations developed for the formalization.
The matching theorem has its own proof; the shared foundations explain the
relationship between the two parts of the work. The resulting counterexample
to the full-displacement claim \citep[Proposition~4.4]{davis1970rotation}
and the follow-on matching theorem are examined in
\cref{sec:prop44,sec:example-transport}.

\section{Davis--Kahan formalization and the sine-theta theorem}
\label{sec:coverage}

Our targets are the four headline theorems \citep[Section~2]{davis1970rotation} and every
named theorem, proposition, lemma, and corollary established by Davis and
Kahan in the 1970 paper. These \DKResultCount{} results include the sine and
tangent inequalities, the geometry and extremal properties of direct
rotations, bounded and unbounded Sylvester estimates, the generalized
perturbation bounds \citep[Section~6]{davis1970rotation}, and the spectral-selection
theorems \citep[Section~8]{davis1970rotation}. We prove \DKProvedCount{} results and refute the printed
full-displacement claim \citep[Proposition~4.4]{davis1970rotation}.
Appendix~\ref{app:dk-inventory} lists these targets.

We develop the sine-theta theorem and its proof below; the remaining source
results are summarized in Appendix~\ref{app:dk-inventory}, with hypotheses
and conclusions in Section~\ref{app:dk-result-summaries}. The main text also
explains the spectral-selection results needed to interpret the angle bounds.
Appendix~\ref{app:notation} collects the notation used here and in the
direct-rotation argument.

\subsection{The sine-theta theorem for self-adjoint operators}
\label{sec:sine-theta}

\subsubsection{Trial and exact subspaces}

We use $\sym{A}$ for the reference operator and
$\sym{Acurrent}=\sym{A}+\sym{H}$ for the perturbed ambient operator.
The symbol $\sym{Acurrent}$ corresponds to $\sym{A}+\sym{H}$ in the source.

For a unit trial vector $\sym{v}$ and an approximate eigenvalue $\sym{mu}$, the residual
is $\sym{Acurrent}\sym{v}-\sym{mu} \sym{v}$. An isometric coordinate map $\sym{E0}$ generalizes
this construction to a trial subspace: its range consists of the candidate
approximate vectors, and a self-adjoint operator $\sym{A0}$ represents their
approximate dynamics. For a finite-dimensional Rayleigh--Ritz choice,
$\sym{A0}=\like{E0}{E_0^*}\sym{Acurrent}\sym{E0}$. The theorem allows other self-adjoint trial
operators as well.

Let $\sym{Eambient}$ be the ambient Hilbert space.  Let $\sym{A0}$ be a self-adjoint operator on a
trial-coordinate space $\sym{Ftrial}$, and let
\[
  \sym{E0}:\sym{Ftrial}\to \sym{Eambient}
\]
be an isometric coordinate map for the trial subspace.  The residual compares
applying the ambient operator after embedding with applying the trial operator
before embedding:
\begin{equation}
  \sym{R} = \sym{Acurrent}\sym{E0}-\sym{E0}\sym{A0}.
  \label{eq:residual}
\end{equation}
When the operators are unbounded, \cref{eq:residual} is asserted on
$\DKDom(\sym{A0})$ and $\sym{R}$ is the bounded residual that agrees there.
The Lean predicate \texttt{IsTrialResidual} packages the isometry, the required
domain inclusion, and this residual identity.

Let $\sym{F0}:\sym{Hexact}\to \sym{Eambient}$ be isometric coordinates for the desired exact spectral
subspace and let $\sym{F1}:\sym{Gunwanted}\to \sym{Eambient}$ be isometric coordinates for its orthogonal
complement.  A self-adjoint operator $\sym{Lambda1}$ on $\sym{Gunwanted}$ describes the
complementary spectral block through
\[
  \sym{Acurrent}\sym{F1} = \sym{F1}\sym{Lambda1}
\]
on its domain. This intertwining identity says that $\sym{Lambda1}$ gives
the action of $\sym{Acurrent}$ in complementary coordinates. The predicate
\texttt{IsExactSpectralDecomposition} also requires the domain inclusion
$\sym{F1}(\DKDom(\sym{Lambda1}))\subseteq\DKDom(\sym{Acurrent})$
and states orthogonality and completeness through
\[
 \like{F0}{F_0^*}\sym{F1}=0,\qquad
 \sym{F0}\like{F0}{F_0^*}+\sym{F1}\like{F1}{F_1^*}=I.
\]

\subsubsection{The sine operator and its rectangular representative}

The source angle operator acts on the trial-coordinate space.  If
$\sym{Pexact}=\sym{F0}\like{F0}{F_0^*}$ is the orthogonal projection onto the desired exact subspace, then
\[
  \cos\sym{Theta0} = (\like{E0}{E_0^*}\sym{Pexact}\sym{E0})^{1/2},
  \qquad
  \sym{sinTheta0} = (I-\like{E0}{E_0^*}\sym{Pexact}\sym{E0})^{1/2}.
\]
The analytic proof uses the rectangular complementary block
\begin{equation}
  \sym{S}=(I-\sym{Pexact})\sym{E0}.
  \label{eq:rectangular-sine}
\end{equation}
The theorem uses the positive sine on the trial-coordinate space.  The rectangular block has codomain $\sym{Eambient}$, while the
positive sine has codomain $\sym{Ftrial}$.  Their relationship is
\[
  \sym{S}^*\sym{S}
  = \like{E0}{E_0^*}(I-\sym{Pexact})\sym{E0}
  = I-\like{E0}{E_0^*}\sym{Pexact}\sym{E0},
\]
so
\begin{equation}
  |\sym{S}|=(\sym{S}^*\sym{S})^{1/2}=\sym{sinTheta0}.
  \label{eq:modulus-sine}
\end{equation}
% It might be worth noting that our formalization actually defines \sym{Theta0}
% and proves eq:modulus-sine whereas most treatments of the subject seam to
% just start from |\sym{S}|. Maybe it's worth enough to put in the paper, but
% it at least deserves a latex comment!
The polar decomposition $\sym{S}=\sym{PolarU}|\sym{S}|$ separates the positive
  modulus from a partial isometry.
Unitarily invariant ideal norms assign the same value to an operator and its
  modulus, so \[ \sym{NormN}(\sym{S})=\sym{NormN}(\sym{sinTheta0}) \] whenever
  the norm is finite.
This equality transfers an estimate for the complementary block $\sym{S}$ to the
  positive sine operator.

\subsubsection{Gap geometry and the inequality}
\label{sec:gap-geometry}

Fix $\sym{delta}>0$. There are two geometric forms of separation.
\emph{Interval/exterior separation:} for some real
$\sym{betaEndpoint}\leq\sym{alphaEndpoint}$, one of the two spectra
$\sigma(\sym{A0})$ and $\sigma(\sym{Lambda1})$ lies in $[\sym{betaEndpoint},\sym{alphaEndpoint}]$ and the other
lies in
\[
 (-\infty,\sym{betaEndpoint}-\sym{delta}]\;\cup\;[\sym{alphaEndpoint}+\sym{delta},\infty).
\]
\emph{Ordered separation:} for some real $b$, the operator forms satisfy either
\[
 \sym{A0}\geq (b+\sym{delta})I,\quad \sym{Lambda1}\leq bI,
 \qquad\text{or}\qquad
 \sym{A0}\leq bI,\quad \sym{Lambda1}\geq (b+\sym{delta})I.
\]
Each form inequality is asserted on the domain of the corresponding operator;
for example, $\sym{A0}\geq aI$ means
$\langle \sym{A0}x,x\rangle\geq a\|x\|^2$ for every
$x\in\DKDom(\sym{A0})$. For self-adjoint operators, the source's ordered
hypotheses on half-infinite spectral intervals, such as $(-\infty,b]$ and
$[b+\sym{delta},\infty)$, imply these form bounds. These cases are included
in the source's unbounded sine theorem
\citep[Section~6, discussion of unbounded operators]{davis1970rotation}.
Write $\sym{SineGap}(\sym{A0},\sym{Lambda1})$ for their union.
The Lean predicate \texttt{FormBoundedSylvesterGap} has three cases:
the finite-interval alternative, which includes either spectral order,
and the two ordered form alternatives.

Both blocks may be unbounded in an ordered configuration. For instance,
on $\ell^2(\mathbb N_{\geq1})$, take diagonal operators
$\sym{A0}e_n=(n+\sym{delta})e_n$ and $\sym{Lambda1}e_n=-ne_n$, on the domains where their
images are square summable. They satisfy the first ordered configuration
with $b=0$, while neither spectrum lies in a finite interval. This example
explains why retaining the half-infinite alternatives changes the theorem's
applicability.

For the unbounded trial configuration, the source assumes that
$\sym{Acurrent}\sym{E0}$ and $\sym{E0}\sym{A0}$ have a common dense domain.
Since $\sym{E0}$ is everywhere defined, this gives
\[
 \DKDom(\sym{Acurrent}\sym{E0})=\DKDom(\sym{A0}),
 \qquad
 \sym{E0}(\DKDom(\sym{A0}))\subseteq\DKDom(\sym{Acurrent}).
\]
The residual agrees with the difference on this domain and extends to a
bounded operator. The formal predicate \texttt{IsTrialResidual} includes the
domain inclusion and the residual identity.

The norm $\sym{NormN}$ is a normalized symmetric operator-ideal norm.
Normalization means that the rank-one map $uv^*:x\mapsto u\langle v,x\rangle$
has norm
\[
 \sym{NormN}(uv^*)=\|u\|\,\|v\|;
\]
in particular, its norm is one when $u,v$ are unit vectors.
Its hypotheses also include Ky Fan dominance whenever both ideal norms
are finite: domination of all finite singular-value sums implies domination
of the ideal norms. This comparison does not assert ideal membership.
The norm is invariant under isometries and positive modulus, and satisfies
$\sym{NormN}(BT)\leq\|B\|_{\mathrm{op}}\sym{NormN}(T)$. Trace-class and
Hilbert--Schmidt operators with their respective norms are examples; the
operator norm on bounded operators is another. For an ideal norm, both the
residual and the positive sine must belong to its domain.

In two dimensions, take a one-dimensional trial frame $\sym{E0}(x)=x\sym{v}$ with
$\|\sym{v}\|=1$. A scalar trial operator $\sym{A0}=\sym{rho}$ gives
$\sym{r}=\sym{R}(1)=\sym{Acurrent}\sym{v}-\sym{rho}\sym{v}$.
Choosing $\sym{rho}=\langle\sym{v},\sym{Acurrent}\sym{v}\rangle$
makes $\sym{rho}\sym{v}$ the orthogonal projection of
$\sym{Acurrent}\sym{v}$ onto the trial line. The residual is therefore
perpendicular to that line. Figure~\ref{fig:sine-theta-residual} pairs this
residual triangle with the exact complementary projection that measures
the subspace error.

\begin{figure}[tb]
\centering
\resizebox{\linewidth}{!}{%% Creator: Matplotlib, PGF backend
%%
%% To include the figure in your LaTeX document, write
%%   \input{<filename>.pgf}
%%
%% Make sure the required packages are loaded in your preamble
%%   \usepackage{pgf}
%%
%% Also ensure that all the required font packages are loaded; for instance,
%% the lmodern package is sometimes necessary when using math font.
%%   \usepackage{lmodern}
%%
%% Figures using additional raster images can only be included by \input if
%% they are in the same directory as the main LaTeX file. For loading figures
%% from other directories you can use the `import` package
%%   \usepackage{import}
%%
%% and then include the figures with
%%   \import{<path to file>}{<filename>.pgf}
%%
%% Matplotlib used the following preamble
%%   \def\mathdefault#1{#1}
%%   \everymath=\expandafter{\the\everymath\displaystyle}
%%   \IfFileExists{scrextend.sty}{
%%     \usepackage[fontsize=18.000000pt]{scrextend}
%%   }{
%%     \renewcommand{\normalsize}{\fontsize{18.000000}{21.600000}\selectfont}
%%     \normalsize
%%   }
%%   \usepackage{lmodern,amsmath,amssymb,xcolor,hyperref}\def\DKNoSymbolTooltips{1}\input{math_symbol_colors.tex}
%%   \makeatletter\@ifpackageloaded{underscore}{}{\usepackage[strings]{underscore}}\makeatother
%%
\begingroup%
\makeatletter%
\begin{pgfpicture}%
\pgfpathrectangle{\pgfpointorigin}{\pgfqpoint{12.400000in}{5.700000in}}%
\pgfusepath{use as bounding box, clip}%
\begin{pgfscope}%
\pgfsetbuttcap%
\pgfsetmiterjoin%
\definecolor{currentfill}{rgb}{1.000000,1.000000,1.000000}%
\pgfsetfillcolor{currentfill}%
\pgfsetlinewidth{0.000000pt}%
\definecolor{currentstroke}{rgb}{1.000000,1.000000,1.000000}%
\pgfsetstrokecolor{currentstroke}%
\pgfsetdash{}{0pt}%
\pgfpathmoveto{\pgfqpoint{0.000000in}{0.000000in}}%
\pgfpathlineto{\pgfqpoint{12.400000in}{0.000000in}}%
\pgfpathlineto{\pgfqpoint{12.400000in}{5.700000in}}%
\pgfpathlineto{\pgfqpoint{0.000000in}{5.700000in}}%
\pgfpathlineto{\pgfqpoint{0.000000in}{0.000000in}}%
\pgfpathclose%
\pgfusepath{fill}%
\end{pgfscope}%
\begin{pgfscope}%
\pgfpathrectangle{\pgfqpoint{1.232860in}{1.653000in}}{\pgfqpoint{3.920281in}{3.249000in}}%
\pgfusepath{clip}%
\pgfsetbuttcap%
\pgfsetmiterjoin%
\pgfsetlinewidth{1.003750pt}%
\definecolor{currentstroke}{rgb}{0.835294,0.850980,0.870588}%
\pgfsetstrokecolor{currentstroke}%
\pgfsetdash{}{0pt}%
\pgfpathmoveto{\pgfqpoint{4.347603in}{2.163174in}}%
\pgfpathcurveto{\pgfqpoint{4.347603in}{2.875040in}}{\pgfqpoint{4.064514in}{3.558478in}}{\pgfqpoint{3.561149in}{4.061843in}}%
\pgfpathcurveto{\pgfqpoint{3.057783in}{4.565208in}}{\pgfqpoint{2.374345in}{4.848298in}}{\pgfqpoint{1.662479in}{4.848298in}}%
\pgfusepath{stroke}%
\end{pgfscope}%
\begin{pgfscope}%
\pgfpathrectangle{\pgfqpoint{1.232860in}{1.653000in}}{\pgfqpoint{3.920281in}{3.249000in}}%
\pgfusepath{clip}%
\pgfsetroundcap%
\pgfsetroundjoin%
\pgfsetlinewidth{2.509375pt}%
\definecolor{currentstroke}{rgb}{0.780392,0.486275,0.000000}%
\pgfsetstrokecolor{currentstroke}%
\pgfsetdash{}{0pt}%
\pgfpathmoveto{\pgfqpoint{1.662479in}{2.163174in}}%
\pgfpathquadraticcurveto{\pgfqpoint{2.762242in}{2.933235in}}{\pgfqpoint{3.830204in}{3.681031in}}%
\pgfusepath{stroke}%
\end{pgfscope}%
\begin{pgfscope}%
\pgfpathrectangle{\pgfqpoint{1.232860in}{1.653000in}}{\pgfqpoint{3.920281in}{3.249000in}}%
\pgfusepath{clip}%
\pgfsetroundcap%
\pgfsetroundjoin%
\definecolor{currentfill}{rgb}{0.780392,0.486275,0.000000}%
\pgfsetfillcolor{currentfill}%
\pgfsetlinewidth{2.509375pt}%
\definecolor{currentstroke}{rgb}{0.780392,0.486275,0.000000}%
\pgfsetstrokecolor{currentstroke}%
\pgfsetdash{}{0pt}%
\pgfpathmoveto{\pgfqpoint{3.744187in}{3.668275in}}%
\pgfpathlineto{\pgfqpoint{3.830204in}{3.681031in}}%
\pgfpathlineto{\pgfqpoint{3.788798in}{3.604563in}}%
\pgfpathlineto{\pgfqpoint{3.744187in}{3.668275in}}%
\pgfpathclose%
\pgfusepath{stroke,fill}%
\end{pgfscope}%
\begin{pgfscope}%
\pgfpathrectangle{\pgfqpoint{1.232860in}{1.653000in}}{\pgfqpoint{3.920281in}{3.249000in}}%
\pgfusepath{clip}%
\pgfsetroundcap%
\pgfsetroundjoin%
\pgfsetlinewidth{2.509375pt}%
\definecolor{currentstroke}{rgb}{0.121569,0.435294,0.819608}%
\pgfsetstrokecolor{currentstroke}%
\pgfsetdash{}{0pt}%
\pgfpathmoveto{\pgfqpoint{1.662479in}{2.163174in}}%
\pgfpathquadraticcurveto{\pgfqpoint{2.762242in}{2.163174in}}{\pgfqpoint{3.823184in}{2.163174in}}%
\pgfusepath{stroke}%
\end{pgfscope}%
\begin{pgfscope}%
\pgfpathrectangle{\pgfqpoint{1.232860in}{1.653000in}}{\pgfqpoint{3.920281in}{3.249000in}}%
\pgfusepath{clip}%
\pgfsetroundcap%
\pgfsetroundjoin%
\definecolor{currentfill}{rgb}{0.121569,0.435294,0.819608}%
\pgfsetfillcolor{currentfill}%
\pgfsetlinewidth{2.509375pt}%
\definecolor{currentstroke}{rgb}{0.121569,0.435294,0.819608}%
\pgfsetstrokecolor{currentstroke}%
\pgfsetdash{}{0pt}%
\pgfpathmoveto{\pgfqpoint{3.745406in}{2.202062in}}%
\pgfpathlineto{\pgfqpoint{3.823184in}{2.163174in}}%
\pgfpathlineto{\pgfqpoint{3.745406in}{2.124285in}}%
\pgfpathlineto{\pgfqpoint{3.745406in}{2.202062in}}%
\pgfpathclose%
\pgfusepath{stroke,fill}%
\end{pgfscope}%
\begin{pgfscope}%
\pgfpathrectangle{\pgfqpoint{1.232860in}{1.653000in}}{\pgfqpoint{3.920281in}{3.249000in}}%
\pgfusepath{clip}%
\pgfsetroundcap%
\pgfsetroundjoin%
\pgfsetlinewidth{3.011250pt}%
\definecolor{currentstroke}{rgb}{0.819608,0.164706,0.482353}%
\pgfsetstrokecolor{currentstroke}%
\pgfsetdash{}{0pt}%
\pgfpathmoveto{\pgfqpoint{3.862004in}{2.163174in}}%
\pgfpathquadraticcurveto{\pgfqpoint{3.862004in}{2.933235in}}{\pgfqpoint{3.862004in}{3.656713in}}%
\pgfusepath{stroke}%
\end{pgfscope}%
\begin{pgfscope}%
\pgfpathrectangle{\pgfqpoint{1.232860in}{1.653000in}}{\pgfqpoint{3.920281in}{3.249000in}}%
\pgfusepath{clip}%
\pgfsetroundcap%
\pgfsetroundjoin%
\definecolor{currentfill}{rgb}{0.819608,0.164706,0.482353}%
\pgfsetfillcolor{currentfill}%
\pgfsetlinewidth{3.011250pt}%
\definecolor{currentstroke}{rgb}{0.819608,0.164706,0.482353}%
\pgfsetstrokecolor{currentstroke}%
\pgfsetdash{}{0pt}%
\pgfpathmoveto{\pgfqpoint{3.823115in}{3.578935in}}%
\pgfpathlineto{\pgfqpoint{3.862004in}{3.656713in}}%
\pgfpathlineto{\pgfqpoint{3.900893in}{3.578935in}}%
\pgfpathlineto{\pgfqpoint{3.823115in}{3.578935in}}%
\pgfpathclose%
\pgfusepath{stroke,fill}%
\end{pgfscope}%
\begin{pgfscope}%
\pgfpathrectangle{\pgfqpoint{1.232860in}{1.653000in}}{\pgfqpoint{3.920281in}{3.249000in}}%
\pgfusepath{clip}%
\pgfsetbuttcap%
\pgfsetmiterjoin%
\pgfsetlinewidth{1.606000pt}%
\definecolor{currentstroke}{rgb}{0.819608,0.164706,0.482353}%
\pgfsetstrokecolor{currentstroke}%
\pgfsetdash{}{0pt}%
\pgfpathmoveto{\pgfqpoint{2.360612in}{2.163174in}}%
\pgfpathcurveto{\pgfqpoint{2.360612in}{2.234389in}}{\pgfqpoint{2.349715in}{2.305187in}}{\pgfqpoint{2.328300in}{2.373106in}}%
\pgfpathcurveto{\pgfqpoint{2.306885in}{2.441025in}}{\pgfqpoint{2.275203in}{2.505270in}}{\pgfqpoint{2.234356in}{2.563606in}}%
\pgfusepath{stroke}%
\end{pgfscope}%
\begin{pgfscope}%
\pgfpathrectangle{\pgfqpoint{1.232860in}{1.653000in}}{\pgfqpoint{3.920281in}{3.249000in}}%
\pgfusepath{clip}%
\pgfsetrectcap%
\pgfsetroundjoin%
\pgfsetlinewidth{1.304875pt}%
\definecolor{currentstroke}{rgb}{0.121569,0.435294,0.819608}%
\pgfsetstrokecolor{currentstroke}%
\pgfsetdash{}{0pt}%
\pgfpathmoveto{\pgfqpoint{1.393967in}{2.163174in}}%
\pgfpathlineto{\pgfqpoint{4.884628in}{2.163174in}}%
\pgfusepath{stroke}%
\end{pgfscope}%
\begin{pgfscope}%
\pgfpathrectangle{\pgfqpoint{1.232860in}{1.653000in}}{\pgfqpoint{3.920281in}{3.249000in}}%
\pgfusepath{clip}%
\pgfsetrectcap%
\pgfsetroundjoin%
\pgfsetlinewidth{1.304875pt}%
\definecolor{currentstroke}{rgb}{0.000000,0.486275,0.568627}%
\pgfsetstrokecolor{currentstroke}%
\pgfsetdash{}{0pt}%
\pgfpathmoveto{\pgfqpoint{1.662479in}{2.163174in}}%
\pgfpathlineto{\pgfqpoint{1.662479in}{4.660339in}}%
\pgfusepath{stroke}%
\end{pgfscope}%
\begin{pgfscope}%
\pgfpathrectangle{\pgfqpoint{1.232860in}{1.653000in}}{\pgfqpoint{3.920281in}{3.249000in}}%
\pgfusepath{clip}%
\pgfsetrectcap%
\pgfsetroundjoin%
\pgfsetlinewidth{1.003750pt}%
\definecolor{currentstroke}{rgb}{0.835294,0.850980,0.870588}%
\pgfsetstrokecolor{currentstroke}%
\pgfsetdash{}{0pt}%
\pgfpathmoveto{\pgfqpoint{3.714322in}{2.163174in}}%
\pgfpathlineto{\pgfqpoint{3.714322in}{2.310855in}}%
\pgfpathlineto{\pgfqpoint{3.862004in}{2.310855in}}%
\pgfusepath{stroke}%
\end{pgfscope}%
\begin{pgfscope}%
\definecolor{textcolor}{rgb}{0.105882,0.121569,0.141176}%
\pgfsetstrokecolor{textcolor}%
\pgfsetfillcolor{textcolor}%
\pgftext[x=2.441165in,y=2.364558in,left,base]{\color{textcolor}{\rmfamily\fontsize{21.000000}{25.200000}\selectfont\catcode`\^=\active\def^{\ifmmode\sp\else\^{}\fi}\catcode`\%=\active\def%{\%}$\sym{theta}$}}%
\end{pgfscope}%
\begin{pgfscope}%
\definecolor{textcolor}{rgb}{0.105882,0.121569,0.141176}%
\pgfsetstrokecolor{textcolor}%
\pgfsetfillcolor{textcolor}%
\pgftext[x=2.736529in,y=3.559438in,,base]{\color{textcolor}{\rmfamily\fontsize{18.000000}{21.600000}\selectfont\catcode`\^=\active\def^{\ifmmode\sp\else\^{}\fi}\catcode`\%=\active\def%{\%}$\like{v}{v_\theta}=\sym{E0}(1)$}}%
\end{pgfscope}%
\begin{pgfscope}%
\definecolor{textcolor}{rgb}{0.105882,0.121569,0.141176}%
\pgfsetstrokecolor{textcolor}%
\pgfsetfillcolor{textcolor}%
\pgftext[x=2.951339in,y=1.840959in,,base]{\color{textcolor}{\rmfamily\fontsize{18.000000}{21.600000}\selectfont\catcode`\^=\active\def^{\ifmmode\sp\else\^{}\fi}\catcode`\%=\active\def%{\%}$\sym{Pexact}\like{v}{v_\theta}$}}%
\end{pgfscope}%
\begin{pgfscope}%
\definecolor{textcolor}{rgb}{0.105882,0.121569,0.141176}%
\pgfsetstrokecolor{textcolor}%
\pgfsetfillcolor{textcolor}%
\pgftext[x=3.971686in,y=2.915008in,left,base]{\color{textcolor}{\rmfamily\fontsize{19.000000}{22.800000}\selectfont\catcode`\^=\active\def^{\ifmmode\sp\else\^{}\fi}\catcode`\%=\active\def%{\%}$\sym{S}(1)$}}%
\end{pgfscope}%
\begin{pgfscope}%
\definecolor{textcolor}{rgb}{0.105882,0.121569,0.141176}%
\pgfsetstrokecolor{textcolor}%
\pgfsetfillcolor{textcolor}%
\pgftext[x=3.971686in,y=2.619645in,left,base]{\color{textcolor}{\rmfamily\fontsize{16.000000}{19.200000}\selectfont\catcode`\^=\active\def^{\ifmmode\sp\else\^{}\fi}\catcode`\%=\active\def%{\%}$\Vert\sym{S}(1)\Vert=\sym{sintheta}$}}%
\end{pgfscope}%
\begin{pgfscope}%
\definecolor{textcolor}{rgb}{0.105882,0.121569,0.141176}%
\pgfsetstrokecolor{textcolor}%
\pgfsetfillcolor{textcolor}%
\pgftext[x=1.796736in,y=4.579785in,left,base]{\color{textcolor}{\rmfamily\fontsize{16.000000}{19.200000}\selectfont\catcode`\^=\active\def^{\ifmmode\sp\else\^{}\fi}\catcode`\%=\active\def%{\%}$\mathrm{range}(\sym{F1})$}}%
\end{pgfscope}%
\begin{pgfscope}%
\definecolor{textcolor}{rgb}{0.105882,0.121569,0.141176}%
\pgfsetstrokecolor{textcolor}%
\pgfsetfillcolor{textcolor}%
\pgftext[x=5.045736in,y=1.840959in,right,base]{\color{textcolor}{\rmfamily\fontsize{16.000000}{19.200000}\selectfont\catcode`\^=\active\def^{\ifmmode\sp\else\^{}\fi}\catcode`\%=\active\def%{\%}$\mathrm{range}(\sym{F0})$}}%
\end{pgfscope}%
\begin{pgfscope}%
\pgfpathrectangle{\pgfqpoint{7.440000in}{3.876000in}}{\pgfqpoint{3.720000in}{0.570000in}}%
\pgfusepath{clip}%
\pgfsetrectcap%
\pgfsetroundjoin%
\pgfsetlinewidth{1.606000pt}%
\definecolor{currentstroke}{rgb}{0.835294,0.850980,0.870588}%
\pgfsetstrokecolor{currentstroke}%
\pgfsetdash{}{0pt}%
\pgfpathmoveto{\pgfqpoint{7.615059in}{4.161000in}}%
\pgfpathlineto{\pgfqpoint{10.984941in}{4.161000in}}%
\pgfusepath{stroke}%
\end{pgfscope}%
\begin{pgfscope}%
\pgfpathrectangle{\pgfqpoint{7.440000in}{3.876000in}}{\pgfqpoint{3.720000in}{0.570000in}}%
\pgfusepath{clip}%
\pgfsetbuttcap%
\pgfsetroundjoin%
\definecolor{currentfill}{rgb}{0.780392,0.486275,0.000000}%
\pgfsetfillcolor{currentfill}%
\pgfsetlinewidth{1.003750pt}%
\definecolor{currentstroke}{rgb}{0.780392,0.486275,0.000000}%
\pgfsetstrokecolor{currentstroke}%
\pgfsetdash{}{0pt}%
\pgfsys@defobject{currentmarker}{\pgfqpoint{-0.048611in}{-0.048611in}}{\pgfqpoint{0.048611in}{0.048611in}}{%
\pgfpathmoveto{\pgfqpoint{0.000000in}{-0.048611in}}%
\pgfpathcurveto{\pgfqpoint{0.012892in}{-0.048611in}}{\pgfqpoint{0.025257in}{-0.043489in}}{\pgfqpoint{0.034373in}{-0.034373in}}%
\pgfpathcurveto{\pgfqpoint{0.043489in}{-0.025257in}}{\pgfqpoint{0.048611in}{-0.012892in}}{\pgfqpoint{0.048611in}{0.000000in}}%
\pgfpathcurveto{\pgfqpoint{0.048611in}{0.012892in}}{\pgfqpoint{0.043489in}{0.025257in}}{\pgfqpoint{0.034373in}{0.034373in}}%
\pgfpathcurveto{\pgfqpoint{0.025257in}{0.043489in}}{\pgfqpoint{0.012892in}{0.048611in}}{\pgfqpoint{0.000000in}{0.048611in}}%
\pgfpathcurveto{\pgfqpoint{-0.012892in}{0.048611in}}{\pgfqpoint{-0.025257in}{0.043489in}}{\pgfqpoint{-0.034373in}{0.034373in}}%
\pgfpathcurveto{\pgfqpoint{-0.043489in}{0.025257in}}{\pgfqpoint{-0.048611in}{0.012892in}}{\pgfqpoint{-0.048611in}{0.000000in}}%
\pgfpathcurveto{\pgfqpoint{-0.048611in}{-0.012892in}}{\pgfqpoint{-0.043489in}{-0.025257in}}{\pgfqpoint{-0.034373in}{-0.034373in}}%
\pgfpathcurveto{\pgfqpoint{-0.025257in}{-0.043489in}}{\pgfqpoint{-0.012892in}{-0.048611in}}{\pgfqpoint{0.000000in}{-0.048611in}}%
\pgfpathlineto{\pgfqpoint{0.000000in}{-0.048611in}}%
\pgfpathclose%
\pgfusepath{stroke,fill}%
}%
\begin{pgfscope}%
\pgfsys@transformshift{7.877647in}{4.161000in}%
\pgfsys@useobject{currentmarker}{}%
\end{pgfscope}%
\end{pgfscope}%
\begin{pgfscope}%
\pgfpathrectangle{\pgfqpoint{7.440000in}{3.876000in}}{\pgfqpoint{3.720000in}{0.570000in}}%
\pgfusepath{clip}%
\pgfsetbuttcap%
\pgfsetroundjoin%
\definecolor{currentfill}{rgb}{0.431373,0.294118,0.847059}%
\pgfsetfillcolor{currentfill}%
\pgfsetlinewidth{1.003750pt}%
\definecolor{currentstroke}{rgb}{0.431373,0.294118,0.847059}%
\pgfsetstrokecolor{currentstroke}%
\pgfsetdash{}{0pt}%
\pgfsys@defobject{currentmarker}{\pgfqpoint{-0.048611in}{-0.048611in}}{\pgfqpoint{0.048611in}{0.048611in}}{%
\pgfpathmoveto{\pgfqpoint{0.000000in}{-0.048611in}}%
\pgfpathcurveto{\pgfqpoint{0.012892in}{-0.048611in}}{\pgfqpoint{0.025257in}{-0.043489in}}{\pgfqpoint{0.034373in}{-0.034373in}}%
\pgfpathcurveto{\pgfqpoint{0.043489in}{-0.025257in}}{\pgfqpoint{0.048611in}{-0.012892in}}{\pgfqpoint{0.048611in}{0.000000in}}%
\pgfpathcurveto{\pgfqpoint{0.048611in}{0.012892in}}{\pgfqpoint{0.043489in}{0.025257in}}{\pgfqpoint{0.034373in}{0.034373in}}%
\pgfpathcurveto{\pgfqpoint{0.025257in}{0.043489in}}{\pgfqpoint{0.012892in}{0.048611in}}{\pgfqpoint{0.000000in}{0.048611in}}%
\pgfpathcurveto{\pgfqpoint{-0.012892in}{0.048611in}}{\pgfqpoint{-0.025257in}{0.043489in}}{\pgfqpoint{-0.034373in}{0.034373in}}%
\pgfpathcurveto{\pgfqpoint{-0.043489in}{0.025257in}}{\pgfqpoint{-0.048611in}{0.012892in}}{\pgfqpoint{-0.048611in}{0.000000in}}%
\pgfpathcurveto{\pgfqpoint{-0.048611in}{-0.012892in}}{\pgfqpoint{-0.043489in}{-0.025257in}}{\pgfqpoint{-0.034373in}{-0.034373in}}%
\pgfpathcurveto{\pgfqpoint{-0.025257in}{-0.043489in}}{\pgfqpoint{-0.012892in}{-0.048611in}}{\pgfqpoint{0.000000in}{-0.048611in}}%
\pgfpathlineto{\pgfqpoint{0.000000in}{-0.048611in}}%
\pgfpathclose%
\pgfusepath{stroke,fill}%
}%
\begin{pgfscope}%
\pgfsys@transformshift{10.722353in}{4.161000in}%
\pgfsys@useobject{currentmarker}{}%
\end{pgfscope}%
\end{pgfscope}%
\begin{pgfscope}%
\definecolor{textcolor}{rgb}{0.105882,0.121569,0.141176}%
\pgfsetstrokecolor{textcolor}%
\pgfsetfillcolor{textcolor}%
\pgftext[x=7.877647in,y=4.047000in,,top]{\color{textcolor}{\rmfamily\fontsize{16.000000}{19.200000}\selectfont\catcode`\^=\active\def^{\ifmmode\sp\else\^{}\fi}\catcode`\%=\active\def%{\%}$\sym{A0}:0$}}%
\end{pgfscope}%
\begin{pgfscope}%
\definecolor{textcolor}{rgb}{0.105882,0.121569,0.141176}%
\pgfsetstrokecolor{textcolor}%
\pgfsetfillcolor{textcolor}%
\pgftext[x=10.722353in,y=4.047000in,,top]{\color{textcolor}{\rmfamily\fontsize{16.000000}{19.200000}\selectfont\catcode`\^=\active\def^{\ifmmode\sp\else\^{}\fi}\catcode`\%=\active\def%{\%}$\sym{Lambda1}:\sym{delta}$}}%
\end{pgfscope}%
\begin{pgfscope}%
\pgfsetroundcap%
\pgfsetroundjoin%
\pgfsetlinewidth{1.003750pt}%
\definecolor{currentstroke}{rgb}{0.431373,0.294118,0.847059}%
\pgfsetstrokecolor{currentstroke}%
\pgfsetdash{}{0pt}%
\pgfpathmoveto{\pgfqpoint{10.694616in}{4.339125in}}%
\pgfpathquadraticcurveto{\pgfqpoint{9.300037in}{4.339125in}}{\pgfqpoint{7.905457in}{4.339125in}}%
\pgfusepath{stroke}%
\end{pgfscope}%
\begin{pgfscope}%
\pgfsetroundcap%
\pgfsetroundjoin%
\pgfsetlinewidth{1.003750pt}%
\definecolor{currentstroke}{rgb}{0.431373,0.294118,0.847059}%
\pgfsetstrokecolor{currentstroke}%
\pgfsetdash{}{0pt}%
\pgfpathmoveto{\pgfqpoint{10.694616in}{4.089125in}}%
\pgfpathlineto{\pgfqpoint{10.694616in}{4.089125in}}%
\pgfpathlineto{\pgfqpoint{10.694616in}{4.589125in}}%
\pgfpathlineto{\pgfqpoint{10.694616in}{4.589125in}}%
\pgfusepath{stroke}%
\end{pgfscope}%
\begin{pgfscope}%
\pgfsetroundcap%
\pgfsetroundjoin%
\pgfsetlinewidth{1.003750pt}%
\definecolor{currentstroke}{rgb}{0.431373,0.294118,0.847059}%
\pgfsetstrokecolor{currentstroke}%
\pgfsetdash{}{0pt}%
\pgfpathmoveto{\pgfqpoint{7.905457in}{4.589125in}}%
\pgfpathlineto{\pgfqpoint{7.905457in}{4.589125in}}%
\pgfpathlineto{\pgfqpoint{7.905457in}{4.089125in}}%
\pgfpathlineto{\pgfqpoint{7.905457in}{4.089125in}}%
\pgfusepath{stroke}%
\end{pgfscope}%
\begin{pgfscope}%
\definecolor{textcolor}{rgb}{0.105882,0.121569,0.141176}%
\pgfsetstrokecolor{textcolor}%
\pgfsetfillcolor{textcolor}%
\pgftext[x=9.300000in,y=4.360500in,,bottom]{\color{textcolor}{\rmfamily\fontsize{16.000000}{19.200000}\selectfont\catcode`\^=\active\def^{\ifmmode\sp\else\^{}\fi}\catcode`\%=\active\def%{\%}gap $\sym{delta}$}}%
\end{pgfscope}%
\begin{pgfscope}%
\pgfpathrectangle{\pgfqpoint{7.710598in}{1.653000in}}{\pgfqpoint{3.302804in}{1.767000in}}%
\pgfusepath{clip}%
\pgfsetroundcap%
\pgfsetroundjoin%
\pgfsetlinewidth{3.011250pt}%
\definecolor{currentstroke}{rgb}{0.819608,0.164706,0.482353}%
\pgfsetstrokecolor{currentstroke}%
\pgfsetdash{}{0pt}%
\pgfpathmoveto{\pgfqpoint{8.123449in}{1.917224in}}%
\pgfpathquadraticcurveto{\pgfqpoint{8.123449in}{2.390827in}}{\pgfqpoint{8.123449in}{2.817845in}}%
\pgfusepath{stroke}%
\end{pgfscope}%
\begin{pgfscope}%
\pgfpathrectangle{\pgfqpoint{7.710598in}{1.653000in}}{\pgfqpoint{3.302804in}{1.767000in}}%
\pgfusepath{clip}%
\pgfsetroundcap%
\pgfsetroundjoin%
\definecolor{currentfill}{rgb}{0.819608,0.164706,0.482353}%
\pgfsetfillcolor{currentfill}%
\pgfsetlinewidth{3.011250pt}%
\definecolor{currentstroke}{rgb}{0.819608,0.164706,0.482353}%
\pgfsetstrokecolor{currentstroke}%
\pgfsetdash{}{0pt}%
\pgfpathmoveto{\pgfqpoint{8.084560in}{2.740067in}}%
\pgfpathlineto{\pgfqpoint{8.123449in}{2.817845in}}%
\pgfpathlineto{\pgfqpoint{8.162337in}{2.740067in}}%
\pgfpathlineto{\pgfqpoint{8.084560in}{2.740067in}}%
\pgfpathclose%
\pgfusepath{stroke,fill}%
\end{pgfscope}%
\begin{pgfscope}%
\pgfpathrectangle{\pgfqpoint{7.710598in}{1.653000in}}{\pgfqpoint{3.302804in}{1.767000in}}%
\pgfusepath{clip}%
\pgfsetroundcap%
\pgfsetroundjoin%
\pgfsetlinewidth{3.011250pt}%
\definecolor{currentstroke}{rgb}{0.074510,0.541176,0.294118}%
\pgfsetstrokecolor{currentstroke}%
\pgfsetdash{}{0pt}%
\pgfpathmoveto{\pgfqpoint{9.692280in}{1.917224in}}%
\pgfpathquadraticcurveto{\pgfqpoint{9.692280in}{2.532908in}}{\pgfqpoint{9.692280in}{3.102006in}}%
\pgfusepath{stroke}%
\end{pgfscope}%
\begin{pgfscope}%
\pgfpathrectangle{\pgfqpoint{7.710598in}{1.653000in}}{\pgfqpoint{3.302804in}{1.767000in}}%
\pgfusepath{clip}%
\pgfsetroundcap%
\pgfsetroundjoin%
\definecolor{currentfill}{rgb}{0.074510,0.541176,0.294118}%
\pgfsetfillcolor{currentfill}%
\pgfsetlinewidth{3.011250pt}%
\definecolor{currentstroke}{rgb}{0.074510,0.541176,0.294118}%
\pgfsetstrokecolor{currentstroke}%
\pgfsetdash{}{0pt}%
\pgfpathmoveto{\pgfqpoint{9.653391in}{3.024229in}}%
\pgfpathlineto{\pgfqpoint{9.692280in}{3.102006in}}%
\pgfpathlineto{\pgfqpoint{9.731169in}{3.024229in}}%
\pgfpathlineto{\pgfqpoint{9.653391in}{3.024229in}}%
\pgfpathclose%
\pgfusepath{stroke,fill}%
\end{pgfscope}%
\begin{pgfscope}%
\definecolor{textcolor}{rgb}{0.105882,0.121569,0.141176}%
\pgfsetstrokecolor{textcolor}%
\pgfsetfillcolor{textcolor}%
\pgftext[x=8.024364in,y=1.752084in,,base]{\color{textcolor}{\rmfamily\fontsize{18.000000}{21.600000}\selectfont\catcode`\^=\active\def^{\ifmmode\sp\else\^{}\fi}\catcode`\%=\active\def%{\%}$\sym{S}(1)$}}%
\end{pgfscope}%
\begin{pgfscope}%
\definecolor{textcolor}{rgb}{0.105882,0.121569,0.141176}%
\pgfsetstrokecolor{textcolor}%
\pgfsetfillcolor{textcolor}%
\pgftext[x=9.774850in,y=1.752084in,,base]{\color{textcolor}{\rmfamily\fontsize{18.000000}{21.600000}\selectfont\catcode`\^=\active\def^{\ifmmode\sp\else\^{}\fi}\catcode`\%=\active\def%{\%}$\sym{R}(1)$}}%
\end{pgfscope}%
\begin{pgfscope}%
\definecolor{textcolor}{rgb}{0.105882,0.121569,0.141176}%
\pgfsetstrokecolor{textcolor}%
\pgfsetfillcolor{textcolor}%
\pgftext[x=8.255561in,y=2.396131in,left,base]{\color{textcolor}{\rmfamily\fontsize{17.000000}{20.400000}\selectfont\catcode`\^=\active\def^{\ifmmode\sp\else\^{}\fi}\catcode`\%=\active\def%{\%}$\sym{sintheta}$}}%
\end{pgfscope}%
\begin{pgfscope}%
\definecolor{textcolor}{rgb}{0.105882,0.121569,0.141176}%
\pgfsetstrokecolor{textcolor}%
\pgfsetfillcolor{textcolor}%
\pgftext[x=9.824393in,y=2.577785in,left,base]{\color{textcolor}{\rmfamily\fontsize{17.000000}{20.400000}\selectfont\catcode`\^=\active\def^{\ifmmode\sp\else\^{}\fi}\catcode`\%=\active\def%{\%}$\sym{delta}\sym{sintheta}$}}%
\end{pgfscope}%
\begin{pgfscope}%
\pgfsetroundcap%
\pgfsetroundjoin%
\pgfsetlinewidth{1.003750pt}%
\definecolor{currentstroke}{rgb}{0.431373,0.294118,0.847059}%
\pgfsetstrokecolor{currentstroke}%
\pgfsetdash{}{0pt}%
\pgfpathmoveto{\pgfqpoint{8.754691in}{3.122460in}}%
\pgfpathquadraticcurveto{\pgfqpoint{9.154665in}{3.215195in}}{\pgfqpoint{9.522417in}{3.011828in}}%
\pgfusepath{stroke}%
\end{pgfscope}%
\begin{pgfscope}%
\pgfsetroundcap%
\pgfsetroundjoin%
\pgfsetlinewidth{1.003750pt}%
\definecolor{currentstroke}{rgb}{0.431373,0.294118,0.847059}%
\pgfsetstrokecolor{currentstroke}%
\pgfsetdash{}{0pt}%
\pgfpathmoveto{\pgfqpoint{9.459103in}{3.103976in}}%
\pgfpathlineto{\pgfqpoint{9.522417in}{3.011828in}}%
\pgfpathlineto{\pgfqpoint{9.410710in}{3.016466in}}%
\pgfusepath{stroke}%
\end{pgfscope}%
\begin{pgfscope}%
\definecolor{textcolor}{rgb}{0.105882,0.121569,0.141176}%
\pgfsetstrokecolor{textcolor}%
\pgfsetfillcolor{textcolor}%
\pgftext[x=8.387673in,y=2.990636in,left,base]{\color{textcolor}{\rmfamily\fontsize{18.000000}{21.600000}\selectfont\catcode`\^=\active\def^{\ifmmode\sp\else\^{}\fi}\catcode`\%=\active\def%{\%}$\times\sym{delta}$}}%
\end{pgfscope}%
\begin{pgfscope}%
\pgfsetrectcap%
\pgfsetroundjoin%
\pgfsetlinewidth{1.003750pt}%
\definecolor{currentstroke}{rgb}{0.835294,0.850980,0.870588}%
\pgfsetstrokecolor{currentstroke}%
\pgfsetdash{}{0pt}%
\pgfpathmoveto{\pgfqpoint{6.200000in}{1.482000in}}%
\pgfpathlineto{\pgfqpoint{6.200000in}{5.073000in}}%
\pgfusepath{stroke}%
\end{pgfscope}%
\begin{pgfscope}%
\definecolor{textcolor}{rgb}{0.105882,0.121569,0.141176}%
\pgfsetstrokecolor{textcolor}%
\pgfsetfillcolor{textcolor}%
\pgftext[x=3.038000in,y=5.358000in,,base]{\color{textcolor}{\rmfamily\fontsize{20.000000}{24.000000}\selectfont\catcode`\^=\active\def^{\ifmmode\sp\else\^{}\fi}\catcode`\%=\active\def%{\%}(a) Project the trial vector}}%
\end{pgfscope}%
\begin{pgfscope}%
\definecolor{textcolor}{rgb}{0.105882,0.121569,0.141176}%
\pgfsetstrokecolor{textcolor}%
\pgfsetfillcolor{textcolor}%
\pgftext[x=9.300000in,y=5.358000in,,base]{\color{textcolor}{\rmfamily\fontsize{20.000000}{24.000000}\selectfont\catcode`\^=\active\def^{\ifmmode\sp\else\^{}\fi}\catcode`\%=\active\def%{\%}(b) Scale the unwanted component}}%
\end{pgfscope}%
\begin{pgfscope}%
\definecolor{textcolor}{rgb}{0.105882,0.121569,0.141176}%
\pgfsetstrokecolor{textcolor}%
\pgfsetfillcolor{textcolor}%
\pgftext[x=9.300000in,y=4.845000in,,base]{\color{textcolor}{\rmfamily\fontsize{19.000000}{22.800000}\selectfont\catcode`\^=\active\def^{\ifmmode\sp\else\^{}\fi}\catcode`\%=\active\def%{\%}$\sym{Acurrent}=\mathrm{diag}(0,\sym{delta}),\quad\sym{A0}=0$}}%
\end{pgfscope}%
\begin{pgfscope}%
\definecolor{textcolor}{rgb}{0.105882,0.121569,0.141176}%
\pgfsetstrokecolor{textcolor}%
\pgfsetfillcolor{textcolor}%
\pgftext[x=6.200000in,y=1.026000in,,base]{\color{textcolor}{\rmfamily\fontsize{18.000000}{21.600000}\selectfont\catcode`\^=\active\def^{\ifmmode\sp\else\^{}\fi}\catcode`\%=\active\def%{\%}$\sym{E0}:\mathbb{R}\to\mathbb{R}^2,\ x\mapsto x\like{v}{v_\theta}\qquad\sym{S}:\mathbb{R}\to\mathbb{R}^2,\ x\mapsto(0,x\sym{sintheta})$}}%
\end{pgfscope}%
\begin{pgfscope}%
\definecolor{textcolor}{rgb}{0.105882,0.121569,0.141176}%
\pgfsetstrokecolor{textcolor}%
\pgfsetfillcolor{textcolor}%
\pgftext[x=6.200000in,y=0.484500in,,base]{\color{textcolor}{\rmfamily\fontsize{18.000000}{21.600000}\selectfont\catcode`\^=\active\def^{\ifmmode\sp\else\^{}\fi}\catcode`\%=\active\def%{\%}$\sym{sinTheta0}:\mathbb{R}\to\mathbb{R},\ x\mapsto x\sym{sintheta}\qquad\sym{delta}\Vert\sym{sinTheta0}\Vert_{\mathrm{op}}=\Vert\sym{R}\Vert_{\mathrm{op}}$}}%
\end{pgfscope}%
\end{pgfpicture}%
\makeatother%
\endgroup%
}
\caption{Rayleigh residual geometry for a one-dimensional trial space in
$\mathbb R^2$ (the rank-one case). With $\sym{E0}(x)=x\sym{v}$, the residual
operator satisfies $\sym{R}(x)=x\sym{r}$, where
$\sym{r}=\sym{Acurrent}\sym{v}-\sym{rho}\sym{v}$.
For a higher-dimensional trial space,
$\sym{R}=\sym{Acurrent}\sym{E0}-\sym{E0}\sym{A0}$ assigns a residual vector
to each trial-coordinate vector. In a finite-dimensional eigenbasis of
$\sym{A0}$, its columns are the individual eigenpair residuals.
The exact frames in the drawing are
$\sym{F0}(h)=(h,0)$ and $\sym{F1}(g)=(0,g)$, with
$\sym{Pexact}=\sym{F0}\like{F0}{F_0^*}$ and $\sym{PexactPerp}=I-\sym{Pexact}$.
(a) The pink leg $\sym{S}(1)=\sym{PexactPerp}\sym{v}$ has length $\sin\sym{theta}$.
(b) The green residual joins $\sym{rho}\sym{v}$ to $\sym{Acurrent}\sym{v}$;
for the Rayleigh choice it is perpendicular to the trial line. The dashed
legs are its exact wanted and unwanted components. Their lengths give
$\sym{delta}\sin\sym{theta}=\|\sym{PexactPerp}\sym{r}\|_2
\leq\|\sym{r}\|_2=\|\sym{R}\|_{\mathrm{op}}$.
The drawing uses $\sym{Acurrent}=\operatorname{diag}(1.9,0.7)$,
$\sym{theta}=40^\circ$, $\sym{rho}=\langle\sym{v},\sym{Acurrent}\sym{v}\rangle$,
and $\sym{delta}=\sym{rho}-0.7$.}
\label{fig:sine-theta-residual}
\end{figure}

Appendix~\ref{app:lean-guide} explains the general Lean syntax in the
following theorem signature; the geometric and spectral hypotheses have been
defined above.

\begin{theorem}[Sine-theta estimate]
\label{thm:sine-theta}
Let the scalar field be $\mathbb R$ or $\mathbb C$, let $\sym{Eambient}$ be separable,
and let $\sym{Acurrent}$, $\sym{A0}$, and $\sym{Lambda1}$ be self-adjoint operators on the
Hilbert spaces above. Suppose the isometries, domain inclusions, residual
identity, exact complementary decomposition, and the separation condition
$\sym{SineGap}(\sym{A0},\sym{Lambda1})$ of
Section~\ref{sec:gap-geometry} hold. Let $\sym{NormN}$ be a normalized symmetric
operator-ideal norm. Whenever $\sym{sinTheta0}$ and $\sym{R}$ belong to its domain,
\begin{equation}
 \sym{delta}\,\sym{NormN}(\sym{sinTheta0})\leq \sym{NormN}(\sym{R}).
 \label{eq:dk-sine}
\end{equation}
\end{theorem}

\begin{lstlisting}[style=leanpaper,
  caption={Lean signature for the sine-theta estimate. The surrounding
  scalar-field and Hilbert-space context is omitted.
  Lean's \texttt{A} denotes $\sym{Acurrent}$;
  its carrier types \texttt{E}, \texttt{F}, \texttt{G}, \texttt{H} denote
  $\sym{Eambient}$, $\sym{Ftrial}$, $\sym{Gunwanted}$, $\sym{Hexact}$, respectively,
  and \texttt{K} is the scalar field. Thus \texttt{H} here is a coordinate
  space, while $\sym{H}$ in the prose is the perturbation. The positive trial-coordinate
  sine is \texttt{sourceDirectedSinThetaOperator E0 F0};
  \texttt{FormBoundedSylvesterGap} encodes the interval/exterior and ordered
  half-infinite separation hypotheses described above.},
  label={lst:dk-sintheta-headline}]
theorem sinTheta_unbounded_formGap_whereDefinedUIN_rclike
    [TopologicalSpace.SeparableSpace E]
    (N : NormalizedSymmetricOperatorIdealFamily K)
    (A : LinearPMap K E E)
    (A0 : LinearPMap K F F)
    (Lambda1 : LinearPMap K G G)
    (E0 : ContinuousLinearMap K F E)
    (F0 : ContinuousLinearMap K H E)
    (F1 : ContinuousLinearMap K G E)
    (R : ContinuousLinearMap K F E)
    (hA : IsSelfAdjoint A)
    (hA0 : IsSelfAdjoint A0)
    (hLambda1 : IsSelfAdjoint Lambda1)
    (htrial : IsTrialResidual A A0 E0 R)
    (hexact : IsExactSpectralDecomposition A Lambda1 F0 F1)
    (delta : Real) (hdelta : 0 < delta)
    (hgap : FormBoundedSylvesterGap A0 Lambda1 delta) :
    N.Mem (sourceDirectedSinThetaOperator E0 F0) ->
    N.Mem R ->
      delta * N.gaugeReal (sourceDirectedSinThetaOperator E0 F0) <=
        N.gaugeReal R
\end{lstlisting}

The two ideal-membership hypotheses ensure that the norm values in the
conclusion are finite. The gap predicate includes the finite
interval/exterior case and both ordered cases. In the ordered cases,
$\sym{A0}$ and $\sym{Lambda1}$ may both be unbounded.

\subsubsection{Why the Sylvester estimate gives the angle bound}

Set $\sym{X}=\like{F1}{F_1^*}\sym{E0}$. The exact complementary decomposition gives
$\sym{S}=\sym{F1}\sym{X}$, and multiplication by the isometry $\sym{F1}$ preserves the
singular-value gauges. Self-adjointness, the domain hypotheses, and the
residual identity yield the Sylvester relation
\[
 \sym{Lambda1} \sym{X}-\sym{X}\sym{A0}=\like{F1}{F_1^*}\sym{R}
\]
on $\DKDom(\sym{A0})$, with the left side represented by the bounded
right side. The separated Sylvester estimate gives
\[
 \sym{delta} \sym{NormN}(\sym{X})\leq \sym{NormN}(\like{F1}{F_1^*}\sym{R})\leq \sym{NormN}(\sym{R}).
\]
Combining this with $\sym{NormN}(\sym{X})=\sym{NormN}(\sym{S})=\sym{NormN}(\sym{sinTheta0})$ proves the conclusion.
For unbounded operators, deriving the domain compatibility for $\sym{X}$ is part of
the argument. The displayed relation holds on $\DKDom(\sym{A0})$ and
has a bounded right side. The specified interval/exterior or ordered
spectral positions give the factor $1/\sym{delta}$, with no additional
multiplicative constant.

For a one-dimensional trial space, take
$\sym{Acurrent}=\operatorname{diag}(0,\sym{delta})$, the trial scalar
$\sym{A0}=0$, and the unit vector
$\like{v}{v_\theta}=(\cos\sym{theta},\sin\sym{theta})$.
Choose the desired exact subspace $\operatorname{span}\{(1,0)\}$;
its orthogonal complement is $\operatorname{span}\{(0,1)\}$, on which
$\sym{Lambda1}=\sym{delta}$.
The complementary component is $(0,\sin\sym{theta})$ and the
residual is $(0,\sym{delta}\sin\sym{theta})$. Thus equality holds
in the sine-theta estimate for this finite-dimensional configuration.

\section{Formalized operator-theoretic foundations}
\label{sec:foundations}

We formalized the geometry of two subspaces, symmetric operator ideals and
majorization, and Sylvester estimates for unbounded self-adjoint operators.
These results supply the constructions and inequalities used in the
Davis--Kahan proofs.

\subsection{Direct rotations and two-subspace geometry}

Let $\sym{Psource}$ and $\sym{Qtarget}$ project onto closed subspaces
$\sym{Usource}$ and $\sym{Vtarget}$. For an acute pair,
$\|\sym{Psource}-\sym{Qtarget}\|<1$. The compression
$\sym{Psource}\sym{Qtarget}\sym{Psource}$ on $\sym{Usource}$ determines
the positive cosine operator. Its square root and the complementary
compression give the cosine and sine factors in the direct-rotation
construction. Polar decomposition gives the isometries between the
corresponding parts of the two subspaces.

The resulting orthogonal or unitary map $\sym{Rrotation}$ carries
$\sym{Usource}$ onto $\sym{Vtarget}$. In finite-dimensional principal-plane
coordinates it rotates each plane through its principal angle and fixes
common directions. The projection formula from
\cref{eq:history-kato-unitary} defines it in Hilbert space without requiring
a discrete list of principal vectors.

Our formalization proves acute existence and uniqueness, the
principal-square-root characterization, commutation with the angle
operators, and reversal symmetry: interchanging the two subspaces replaces
the rotation by its adjoint. For general pairs, the existence criteria
include the right-angle defect spaces described in
\cref{sec:related-two-projections}. When the angle operator is compact,
its spectral decomposition recovers a sequence of principal angles
\citep[Section~3]{davis1970rotation}.

The double-angle theorems for $\sin(2\sym{Theta})$ and $\tan(2\sym{Theta})$
  require selecting the branch with angles below a quarter turn.
Here $\sym{Theta}$ includes both the source and complementary angle blocks;
$\sym{Theta0}$ alone is the directed angle on source coordinates.
Because the function $\sin(2\sym{theta})$ is small both when $\sym{theta}$ is
  near $0$ and when $\sym{theta}$ is near $\pi/2$, a bound on
  $\sin(2\sym{Theta})$ alone cannot determine whether an approximate subspace is
  close to or nearly orthogonal to the target.
The branch-selection theorem \citep[Theorem~8.1]{davis1970rotation}
uses an ordered reference gap. Let $\sym{Psource}$ reduce a self-adjoint
reference operator $\sym{A}$, with
blocks $\sym{ReferenceBlock0}$ on $\DKRange(\sym{Psource})$ and
$\sym{ReferenceBlock1}$ on $\DKKernel(\sym{Psource})$, satisfying
\[
 \sym{ReferenceBlock0}\le\sym{alphaEndpoint}I,
 \qquad
 \sym{ReferenceBlock1}\ge(\sym{alphaEndpoint}+\sym{delta})I,
 \qquad \sym{delta}>0,
\]
in the quadratic-form sense. Suppose $\sym{H}$ is bounded and self-adjoint,
with both diagonal blocks zero:
$\sym{Psource}\sym{H}\sym{Psource}
 =(I-\sym{Psource})\sym{H}(I-\sym{Psource})=0$.
These hypotheses imply that $\sym{Acurrent}=\sym{A}+\sym{H}$ has no spectrum
in $(\sym{alphaEndpoint},\sym{alphaEndpoint}+\sym{delta})$.

Fix a reducing subspace $M$ of $\sym{Acurrent}$, and write
$\sym{Lambda0},\sym{Lambda1}$ for the blocks of $\sym{Acurrent}$ on $M,M^\perp$.
The full angle $\sym{Theta}$ comparing $\DKRange(\sym{Psource})$ with $M$
satisfies $\sym{Theta}\le\pi/4$ if and only if
$\sym{Lambda0}\le\sym{alphaEndpoint}I$ and
$\sym{Lambda1}\ge(\sym{alphaEndpoint}+\sym{delta})I$.
Such a subspace exists: take the spectral subspace of $\sym{Acurrent}$
for $(-\infty,\sym{alphaEndpoint}]$. Our formalization also proves that
it is the unique reducing subspace satisfying this angle bound.

To compare the reference and perturbed blocks on common coordinate spaces,
choose isometric frames for $\DKRange(\sym{Psource})$ and
$\DKKernel(\sym{Psource})$. Apply the direct rotation $\sym{Rrotation}$ to
these frames to obtain frames for $M$ and $M^\perp$. Express
$\sym{ReferenceBlock0},\sym{ReferenceBlock1}$ in the original frames and
$\sym{Lambda0},\sym{Lambda1}$ in the transported frames. In these coordinates
the branch satisfies the stronger quadratic-form inequality
\[
 \sym{ReferenceBlock1}-\sym{alphaEndpoint}I
 \le\sym{ComplementCosineBlock}
       (\sym{Lambda1}-\sym{alphaEndpoint}I)\sym{ComplementCosineBlock},
 \qquad \sym{ComplementCosineBlock}=\cos\sym{Theta1},
\]
with the lower-block analogue. These inequalities quantify spectral
repulsion in terms of the angles and yield eigenvalue and symmetric-gauge
comparisons in finite dimensions.

The smallness criterion \citep[Theorem~8.2]{davis1970rotation} uses a central
spectral interval instead of an ordered gap. Let $\sym{Psource}$ reduce
the self-adjoint operator $\sym{A}$ and $\sym{Qtarget}$ reduce the
self-adjoint operator $\sym{Acurrent}$. The perturbed block on
$\DKRange(\sym{Qtarget})$ has spectrum in
$[\sym{betaEndpoint},\sym{alphaEndpoint}]$, while its complementary spectrum
lies in
\[
 (-\infty,\sym{betaEndpoint}-\sym{delta}]
 \ \cup\ [\sym{alphaEndpoint}+\sym{delta},\infty),
 \qquad \sym{betaEndpoint}\le\sym{alphaEndpoint},\quad\sym{delta}>0.
\]
The reference block $\sym{ReferenceBlock0}$, representing
$\sym{A}|_{\DKRange(\sym{Psource})}$ in source coordinates, must have spectrum
in $[\sym{betaEndpoint}-\sym{delta}/2,\sym{alphaEndpoint}+\sym{delta}/2]$.
For an isometric frame $\sym{E0}$ of the reference subspace, its residual
is $\sym{R}=\sym{Acurrent}\sym{E0}-\sym{E0}\sym{ReferenceBlock0}$
on the operator domain, as in \cref{eq:history-residual}.
Assume
$\dim\sym{CrossedSourceDefect}=\dim\sym{CrossedTargetDefect}$: the space of source
vectors orthogonal to the target and the space of target vectors orthogonal
to the source have equal Hilbert dimension.

If $\|\sym{H}\|_{\mathrm{op}}<\sym{delta}/2$ for a
bounded self-adjoint perturbation, or $\|\sym{R}\|_{\mathrm{op}}<\sym{delta}/2$
for the bounded residual on the reference subspace, then
$\sym{Theta}<\pi/4$. For each ideal norm whose displayed arguments belong
to its domain, the corresponding double-angle estimate also holds:
\[
 \sym{delta}\sym{NormN}(\sin(2\sym{Theta}))\le2\sym{NormN}(\sym{H}),
 \qquad\text{or}\qquad
 \sym{delta}\sym{NormN}(\sin(2\sym{Theta0}))\le2\sym{NormN}(\sym{R}).
\]
The perturbation estimate uses the full angle, and the residual estimate
the directed angle on source coordinates. In the unbounded residual form,
$\sym{A}$ and $\sym{Acurrent}$ are self-adjoint on a common domain; only
the residual is required to be bounded.

We formalized both theorems in Hilbert space using spectral calculus.

\subsection{Approximation numbers and operator-ideal norms}

In finite dimensions, a unitarily invariant norm is a symmetric gauge
function of the singular values. Approximation numbers extend the
singular-value comparisons needed for bounded operators; for compact
operators they give the usual singular-value sequence. We formalized
symmetric norming functions and Ky Fan partial sums to pass from
majorization to norm inequalities.

For example, let $\sym{Rrotation}$ be the direct rotation and
$\sym{Wcompetitor}$ any unitary carrying $\sym{Usource}$ onto
$\sym{Vtarget}$. The restricted displacement of the direct rotation has
singular-value minimality, which implies minimality for unitarily invariant
norms. The squared-displacement result compares
$(I-\sym{Rrotation})^*(I-\sym{Rrotation})$ and
$(I-\sym{Wcompetitor})^*(I-\sym{Wcompetitor})$
\citep[Proposition~4.3]{davis1970rotation}. Their singular values are the
squares of those of the unsquared displacements. Section~\ref{sec:prop44}
shows why the claimed trace-norm comparison fails and derives the valid
Q-norm conclusion.
% Additional formalized quantitative results, omitted from this article:
% DavisKahan/FiniteDimensional/DirectRotation/SharpDisplacement.lean and
% SharpDisplacementLimit.lean. AUTHORING_NOTES.md indexes the evidence.

The same norm inequalities enter the sine-theta proof. The assertion
$\sym{delta}\sym{NormN}(\sym{S})\leq\sym{NormN}(\sym{R})$ applies when
$\sym{S}$ and $\sym{R}$ have finite ideal norm. Our ideal-membership
hypotheses state this condition explicitly in arbitrary Hilbert dimension.

\subsection{Sylvester estimates and unbounded operators}

For operators $\sym{SylvesterLeft}$ and $\sym{SylvesterRight}$ on the target and source spaces of $\sym{X}$, the
Sylvester map is $\sym{X}\mapsto \sym{SylvesterLeft}\sym{X}-\sym{X}\sym{SylvesterRight}$. Under the interval/exterior or ordered
spectral-separation hypotheses, the estimate has the form
\[
 \sym{delta} \sym{NormN}(\sym{X})\leq \sym{NormN}(\sym{SylvesterLeft}\sym{X}-\sym{X}\sym{SylvesterRight}).
\]
The sine-theta proof uses $\sym{SylvesterLeft}=\sym{Lambda1}$, $\sym{SylvesterRight}=\sym{A0}$, and
$\sym{X}=\like{F1}{F_1^*}\sym{E0}$. Other perturbation estimates choose related
projection blocks for $\sym{X}$.

For unbounded $\sym{SylvesterLeft}$ and $\sym{SylvesterRight}$, the identity is asserted for
$x\in\DKDom(\sym{SylvesterRight})$, with $\sym{X}x\in\DKDom(\sym{SylvesterLeft})$, and $\sym{SylvesterLeft}\sym{X}-\sym{X}\sym{SylvesterRight}$ agrees there with a
bounded residual. We formalized these domain conditions using linear
partial maps. Spectral cutoffs give bounded approximations; convergence
of approximation numbers passes the norm estimates to the unbounded
operators \citep[Section~5 and Appendix to Section~6]{davis1970rotation}.

The generalized sine and tangent theorems also compare block operators.
Reflection pinching replaces a block matrix $T$ by
$(T+\sym{Jreflect}T\sym{Jreflect})/2$, where $\sym{Jreflect}$ acts as $+I$ on one coordinate subspace and $-I$
on the other. This removes the off-diagonal blocks and contracts a unitarily
invariant norm. Direct-sum comparisons relate the norms of the component
blocks to the whole operator. Together with Sylvester separation these give
the generalized angle estimates
\citep[Lemmas~6.1--6.2 and Theorems~6.1--6.3]{davis1970rotation}.

\subsection{Real and complex Hilbert spaces}

The spectral constructions use complex functional calculus. We formalized
the scalar-transfer results needed to carry these arguments to real Hilbert
spaces through complexification. This permits the source theorems to be
stated over either field, including their domain and ideal-norm hypotheses.
The underlying algebra, topology, and continuous functional calculus come
from Mathlib; the adapted Spectra constructions and later Tau Ceti
contributions are described in Section~\ref{sec:related-formal}.

\section{A counterexample to full-displacement optimality}
\label{sec:prop44}

This section gives an explicit real counterexample to the small-angle
optimality assertion \citep[Proposition~4.4]{davis1970rotation} and
identifies the invalid inequality in its proof. To the best of our knowledge,
this contradiction of the printed real-space claim has not previously been
published. Its complex-phase construction is already in the authors'
example \citep[Example~4.2]{davis1970rotation}; the distinction between trace-norm failure and Q-norm optimality
also has earlier treatments \citep{bolotnikov2001minimal}. We give the
counterexample, the failed comparison, and a Q-norm repair below.

Write $\sym{Psource}$ and $\sym{Qtarget}$ for the orthogonal projections
onto subspaces $\sym{Usource}$ and $\sym{Vtarget}$. For an orthogonal or
unitary map $\sym{Wcompetitor}$ with
$\sym{Wcompetitor}(\sym{Usource})=\sym{Vtarget}$, the vector
$(I-\sym{Wcompetitor})x$ is the displacement of $x$ up to sign. Full
displacement measures this operator on the whole ambient space; restricted
displacement measures $(I-\sym{Wcompetitor})\sym{Psource}$, which tests only
vectors in the source subspace. We use the trace norm
$\|Z\|_1=\sum_j s_j(Z)$.

For a finite-dimensional acute pair, the direct rotation $\sym{Rrotation}$
is the orthogonal or unitary polar factor of
\[
 \sym{IntertwinerT}=\sym{Qtarget}\sym{Psource}
 +(I-\sym{Qtarget})(I-\sym{Psource}).
\]
The printed claim \citep[Proposition~4.4]{davis1970rotation} asserts that,
for real subspaces with largest principal angle
at most $\pi/3$, it minimizes $\sym{NormN}(I-\sym{Wcompetitor})$ for every
unitarily invariant norm $\sym{NormN}$.

\subsection{The real four-dimensional counterexample}

The competitor below rotates one real two-plane and fixes its orthogonal
complement. The direct rotation moves both principal planes. This permits a
smaller trace norm even though the competitor's largest displacement is
larger.

\begin{theorem}[Trace-norm failure throughout the acute angle range]
\label{thm:small-angle}
For every $0<\sym{theta}<\pi/2$, there are two-dimensional subspaces
$\sym{Usource},\like{Vtarget}{V_\theta}\subset\mathbb R^4$ and an orthogonal map $\like{Wcompetitor}{W_\theta}$ with
$\like{Wcompetitor}{W_\theta} \sym{Usource}=\like{Vtarget}{V_\theta}$, such that both principal angles equal $\sym{theta}$ and
the canonical direct rotation $\like{Rrotation}{D_\theta}$ satisfies
\begin{align}
 \|I-\like{Wcompetitor}{W_\theta}\|_1&=4\sin\sym{theta},\label{eq:competitor-trace}\\
 \|I-\like{Rrotation}{D_\theta}\|_1&=8\sin(\sym{theta}/2)>4\sin\sym{theta}.
 \label{eq:direct-trace}
\end{align}
\end{theorem}

\begin{proof}
Put $\sym{CosAngle}=\cos\sym{theta}$, $\sym{SinAngle}=\sin\sym{theta}$, and let
\[
 \sym{Jturn}=\begin{pmatrix}0&-1\\1&0\end{pmatrix},\quad
 \sym{Cturn}=\cos(2\sym{theta})I_2+\sin(2\sym{theta})\sym{Jturn},\quad
 \sym{Hmix}=\frac{1}{\sqrt2}\begin{pmatrix}I_2&I_2\\I_2&-I_2\end{pmatrix}.
\]
Define
\begin{equation}
 \like{Wcompetitor}{W_\theta}=\sym{Hmix}\begin{pmatrix}\sym{Cturn}&0\\0&I_2\end{pmatrix}\sym{Hmix}^T
 =\begin{pmatrix}\like{CosAngle}{c^2}I_2+\sym{SinAngle}\sym{CosAngle}\sym{Jturn}&-\like{SinAngle}{s^2}I_2+\sym{SinAngle}\sym{CosAngle}\sym{Jturn}\\
                  -\like{SinAngle}{s^2}I_2+\sym{SinAngle}\sym{CosAngle}\sym{Jturn}&\like{CosAngle}{c^2}I_2+\sym{SinAngle}\sym{CosAngle}\sym{Jturn}\end{pmatrix}.
 \label{eq:wtheta}
\end{equation}
Both $\sym{Hmix}$ and $\sym{Cturn}$ are orthogonal, so $\like{Wcompetitor}{W_\theta}$ is orthogonal. Let
$\sym{Usource}=\mathbb R^2\oplus\{0\}$ and $\like{Vtarget}{V_\theta}=\like{Wcompetitor}{W_\theta} \sym{Usource}$.
The first two columns of $\like{Wcompetitor}{W_\theta}$ factor as
\[
 \begin{pmatrix}\like{CosAngle}{c^2}I_2+\sym{SinAngle}\sym{CosAngle}\sym{Jturn}\\-\like{SinAngle}{s^2}I_2+\sym{SinAngle}\sym{CosAngle}\sym{Jturn}\end{pmatrix}
 =\begin{pmatrix}\sym{CosAngle}I_2\\\sym{SinAngle}\sym{Jturn}\end{pmatrix}(\sym{CosAngle}I_2+\sym{SinAngle}\sym{Jturn}).
\]
Since $\sym{CosAngle}I_2+\sym{SinAngle}\sym{Jturn}$ is orthogonal, the columns of
$\sym{Yframe}=\binom{\sym{CosAngle}I_2}{\sym{SinAngle}\sym{Jturn}}$ form an orthonormal basis of $\like{Vtarget}{V_\theta}$.
With $\sym{Xframe}=\binom{I_2}{0}$, one has $\sym{Xframe}^T\sym{Yframe}=\sym{CosAngle}I_2$.
Thus both principal cosines are $\sym{CosAngle}>0$, and both principal angles are
$\sym{theta}$.

The corresponding projections and canonical intertwiner are
\[
 \sym{Psource}=\begin{pmatrix}I_2&0\\0&0\end{pmatrix},\qquad
 \sym{Qtarget}=\sym{Yframe}\sym{Yframe}^T=\begin{pmatrix}\like{CosAngle}{c^2}I_2&-\sym{SinAngle}\sym{CosAngle}\sym{Jturn}\\\sym{SinAngle}\sym{CosAngle}\sym{Jturn}&\like{SinAngle}{s^2}I_2\end{pmatrix},
\]
\begin{equation}
 \sym{IntertwinerT}=\sym{Qtarget}\sym{Psource}+(I-\sym{Qtarget})(I-\sym{Psource})
   =\sym{CosAngle}\begin{pmatrix}\sym{CosAngle}I_2&\sym{SinAngle}\sym{Jturn}\\\sym{SinAngle}\sym{Jturn}&\sym{CosAngle}I_2\end{pmatrix}.
 \label{eq:intertwiner-theta}
\end{equation}
The last block matrix is orthogonal and $\sym{CosAngle}>0$, so $|\sym{IntertwinerT}|=\sym{CosAngle}I_4$ and
its polar factor is exactly
\[
 \like{Rrotation}{D_\theta}=\begin{pmatrix}\sym{CosAngle}I_2&\sym{SinAngle}\sym{Jturn}\\\sym{SinAngle}\sym{Jturn}&\sym{CosAngle}I_2\end{pmatrix}.
\]
In particular, this calculation identifies the canonical direct rotation
of the constructed pair. Since $\like{Rrotation}{D_\theta}+\like{Rrotation}{D_\theta}^T=2\sym{CosAngle}I_4$,
\[
 (I-\like{Rrotation}{D_\theta})^T(I-\like{Rrotation}{D_\theta})=(2-2\sym{CosAngle})I_4.
\]
Consequently all four displacement singular values equal
$\sqrt{2-2\sym{CosAngle}}=2\sin(\sym{theta}/2)$.

On the other hand, \cref{eq:wtheta} shows that the singular values of
$I-\like{Wcompetitor}{W_\theta}$ are those of $(I_2-\sym{Cturn})\oplus0_2$.
The identity
$(I_2-\sym{Cturn})^T(I_2-\sym{Cturn})=4\like{SinAngle}{s^2}I_2$ gives the list
$(2\sym{SinAngle},2\sym{SinAngle},0,0)$. The two trace formulas follow. Finally,
\[
 4\sin\sym{theta}
 =8\sin(\sym{theta}/2)\cos(\sym{theta}/2)
 <8\sin(\sym{theta}/2),
\]
because $0<\sym{theta}/2<\pi/4$.
\end{proof}

\begin{corollary}[No positive small-angle threshold]
\label{cor:no-threshold}
For every $\sym{EpsilonAngle}>0$, there are real two-dimensional subspaces of
$\mathbb R^4$ whose largest principal angle is strictly between zero and
$\sym{EpsilonAngle}$ and is at most $\pi/3$, together with an orthogonal competitor
having strictly smaller trace-norm displacement than the direct rotation.
\end{corollary}

\begin{proof}
Apply \cref{thm:small-angle} with
$\sym{theta}=\min\{\sym{EpsilonAngle}/2,\pi/4\}$.
\end{proof}

The competitor rotates a diagonal real two-plane through $2\sym{theta}$ and
fixes its orthogonal complement; its displacement has rank two, compared
with rank four for $\like{Rrotation}{D_\theta}$. The operator norm favors the canonical
rotation instead:
$\|I-\like{Rrotation}{D_\theta}\|_{\mathrm{op}}=2\sin(\sym{theta}/2)
 <2\sin\sym{theta}=\|I-\like{Wcompetitor}{W_\theta}\|_{\mathrm{op}}$.
The competitor attains the global trace-norm minimum $4\sin\sym{theta}$.
Indeed, the projection difference has trace norm
$\|\sym{Psource}-\sym{Qtarget}\|_1=4\sin\sym{theta}$, and every orthogonal
transporter $\sym{Wcompetitor}$ satisfies
$\|\sym{Psource}-\sym{Qtarget}\|_1\leq\|I-\sym{Wcompetitor}\|_1$.
Equation~\eqref{eq:competitor-trace} attains this lower bound.
% Formal proof: DavisKahan/FiniteDimensional/DirectRotation/TraceOptimality.lean,
% SmallAngleCounterexample.traceDisplacement_lower_bound and Wequiv_traceMinimizer.
% Dimension minimality is now proved in DirectRotation/DimensionMinimality.lean:
% nuclear_directRotation_le_of_finrank_le_three compares the actual direct
% rotation with every ambient real orthogonal transporter in dimension <= 3,
% under the acute-pair and largest-principal-angle <= pi/3 hypotheses.
% isLeast_real_shortAngle_nuclear_counterexample_dimension proves the least
% dimension is four; the same file retains arbitrarily small positive witnesses.
% Proposed corollary for the next prose review: Four is the least real ambient
% dimension for trace-norm failure under the printed small-angle hypotheses.

At $\sym{theta}=\pi/4$, \cref{fig:prop44-displacement-planes}
shows the different displacement patterns of the canonical direct rotation
and the competing transporter. Both carry the same source subspace onto the
same target, with both principal angles equal to $45^\circ$. The competitor
fixes half of the ambient space; the direct rotation moves both principal planes.

\begin{figure}[tbp]
\centering
\resizebox{119mm}{!}{\input{figures/Fig2.pgf}}
\caption{A counterexample to the full-displacement claim
\citep[Proposition~4.4]{davis1970rotation} at
$\sym{theta}=\pi/4$. Both orthogonal maps carry $\sym{Usource}$ onto
$\sym{Vtarget}$. Each circle shows the action on one invariant real
two-plane: (a,b) the principal planes $\Pi_1,\Pi_2$ for $\sym{Rrotation}$,
and (c,d) the moving and fixed planes $M,Z$ for $\sym{Wcompetitor}$.
The unit inputs are $u_1=e_0$, $u_2=e_1$,
$x=(e_0+e_2)/\sqrt2$, and $z=(e_0-e_2)/\sqrt2$.
For $T=\sym{Rrotation}$ or $\sym{Wcompetitor}$, the green arrow from
the image $Tq$ to the input $q$ is the displacement $(I-T)q$.
The bars list the four singular values of $I-T$, grouped by contributing
plane, on the same vertical scale. Their sums give
$\|I-\sym{Wcompetitor}\|_1=2\sqrt2<4\sqrt{2-\sqrt2}=\|I-\sym{Rrotation}\|_1$.
The two maps have different invariant-plane decompositions}
\label{fig:prop44-displacement-planes}
\end{figure}

\subsection{Failure of the published proof}

At $\sym{theta}=\pi/4$, the family gives
\[
 \like{Wcompetitor}{W_{\pi/4}}=\frac12
 \begin{pmatrix}
 1&-1&-1&-1\\1&1&1&-1\\-1&-1&1&-1\\1&-1&1&1
 \end{pmatrix}.
\]
The two trace norms are
\[
 \|I-\like{Wcompetitor}{W_{\pi/4}}\|_1=2\sqrt2
 <4\sqrt{2-\sqrt2}=\|I-\like{Rrotation}{D_{\pi/4}}\|_1.
\]

For $\sym{Kdisplacement}=I-\like{Wcompetitor}{W_{\pi/4}}$, the even-index part
of the published comparison \citep[equation~(4.3)]{davis1970rotation}
reads, in our notation,
\[
 \|\sym{Kdisplacement}\|_{(\nu)}\geq
 \sum_{k=1}^{\nu/2}\|\sym{Kdisplacement}\Omega_k\|_{(2)},
 \qquad \nu\ \text{even},
\]
where $\|\cdot\|_{(j)}$ is the Ky Fan $j$-norm and $\Omega_k$ projects onto
the $k$th principal plane. Choose $\nu=4$ and let
$\sym{OmegaOne}$ and $\sym{OmegaTwo}$ project onto
$\operatorname{span}\{e_0,e_3\}$ and
$\operatorname{span}\{e_1,e_2\}$. These are principal planes for the pair.
The comparison then requires
\begin{equation}
 \|\sym{Kdisplacement}\|_{(4)}\geq
 \|\sym{Kdisplacement}\sym{OmegaOne}\|_{(2)}+\|\sym{Kdisplacement}\sym{OmegaTwo}\|_{(2)},
 \label{eq:dk43-false}
\end{equation}
The checked calculations give
\[
 \|\sym{Kdisplacement}\|_{(4)}=2\sqrt2,\qquad
 \|\sym{Kdisplacement}\sym{OmegaOne}\|_{(2)}=\|\sym{Kdisplacement}\sym{OmegaTwo}\|_{(2)}=2.
\]
Thus \cref{eq:dk43-false} would imply $2\sqrt2\geq4$.
The displayed matrix also gives
\[
 \DKRange(\sym{Kdisplacement}\sym{OmegaOne})
 =\DKRange(\sym{Kdisplacement}\sym{OmegaTwo})
 =\sym{CounterMovingPlane}
 =\operatorname{span}\{e_0+e_2,e_1+e_3\}.
\]
Both restrictions map onto the same two-dimensional output plane. Each
has singular values $1,1$, whereas the full displacement has just two
nonzero singular values, $\sqrt2,\sqrt2$. The proposed sum treats the
restrictions as independent output blocks. A direct-sum argument would
require orthogonal images as well as orthogonal domain planes; here the
images coincide.

\subsection{A Q-norm repair from squared-displacement optimality}

Call a unitarily invariant norm $\sym{NormN}$ a $\sym{NormQ}$-norm when there is a unitarily
invariant norm $\sym{NormM}$ satisfying
\begin{equation}
 \sym{NormN}(Z)^2=\sym{NormM}(Z^*Z)\qquad\text{for every }Z.
 \label{eq:qnorm-def}
\end{equation}

For example, the Frobenius norm is a $\sym{NormQ}$-norm because
$\|Z\|_F^2=\|Z^*Z\|_1$: its squared norm is the sum of the eigenvalues of
the positive matrix $Z^*Z$. This identity uses the trace norm for
$\sym{NormM}$. The squared-displacement comparison
\citep[Proposition~4.3]{davis1970rotation} gives the following Q-norm consequence.
Earlier Q-norm minimization results appear in
\citet[Theorem~3.1]{bolotnikov2001minimal}.

\begin{proposition}[$\sym{NormQ}$-norm consequence]
\label{prop:qnorm-repair}
Let $\sym{Usource},\sym{Vtarget}$ be an acute pair of subspaces of a finite-dimensional real or
complex Hilbert space, and let $\sym{Rrotation}$ be their canonical direct rotation.
For every $\sym{NormQ}$-norm $\sym{NormN}$ and every orthogonal or unitary map $\sym{Wcompetitor}$ with $\sym{Wcompetitor}\sym{Usource}=\sym{Vtarget}$,
\begin{equation}
 \sym{NormN}(I-\sym{Rrotation})\leq \sym{NormN}(I-\sym{Wcompetitor}).
 \label{eq:qnorm-repair}
\end{equation}
No $\pi/3$ restriction is required.
\end{proposition}

\begin{proof}
The valid squared-displacement theorem gives
\[
 \sym{NormM}((I-\sym{Rrotation})^*(I-\sym{Rrotation}))\leq \sym{NormM}((I-\sym{Wcompetitor})^*(I-\sym{Wcompetitor})).
\]
By \cref{eq:qnorm-def}, these quantities are $\sym{NormN}(I-\sym{Rrotation})^2$ and $\sym{NormN}(I-\sym{Wcompetitor})^2$.
Both norm values are nonnegative, so taking square roots proves the claim.
\end{proof}

The trace norm on the counterexample space does not satisfy the
$\sym{NormQ}$-norm condition. Its smaller competing displacement in
\cref{thm:small-angle} contradicts the minimality that this condition
would imply.

\subsection{Historical context}

Davis and Kahan already exhibited a complex two-dimensional phase
modification that decreases the trace norm of full displacement
\citep[Example~4.2]{davis1970rotation}. Our real family is its
realification. Indeed, under the identification $\mathbb C^2\simeq
\mathbb R^2\oplus\mathbb R^2$ in which multiplication by $i$ is
represented by $\sym{Jturn}$ from \cref{eq:wtheta}, the matrices
\[
 \like{Rrotation}{D_\theta}^{\mathbb C}=
 \begin{pmatrix}\sym{CosAngle}&i\sym{SinAngle}\\i\sym{SinAngle}&\sym{CosAngle}\end{pmatrix},\qquad
 \like{Wcompetitor}{W_\theta}^{\mathbb C}=e^{i\sym{theta}}\like{Rrotation}{D_\theta}^{\mathbb C}
\]
become precisely $\like{Rrotation}{D_\theta}$ and $\like{Wcompetitor}{W_\theta}$ above. Realification doubles
the complex singular-value multiplicities. Thus the real small-angle
restriction in the full-displacement claim
\citep[Proposition~4.4]{davis1970rotation} does not eliminate the mechanism of the
authors' preceding example.

\citet{bolotnikov2001minimal} studied related unitary completion and
subspace-mapping problems. Their results give a $\sym{NormQ}$-norm
minimization theorem and a counterexample for Schatten $p$-norms with
$1\leq p<2$ \citep[Theorem~3.1 and Lemma~3.3]{bolotnikov2001minimal}.
Those results establish the earlier norm-class
distinction. \citet[Section~5]{qiu2006grassmann} subsequently repeated
the $\pi/3$ assertion, whereas
\citet[Section~4, p.~1615]{zhang2010canonicalangles} explicitly noted
general trace-norm nonminimality without addressing that particular
real-space restriction. The present refutation identifies the invalid
inequality in the printed full-displacement claim
\citep[Proposition~4.4]{davis1970rotation}. The Q-norm conclusion above follows
from the earlier squared-displacement extremality statement.


\section{A follow-on theorem enabled by the formal development}
\label{sec:example-transport}

We prove an exact formula for the smallest trace-norm displacement of a real
orthogonal map carrying one given subspace onto another in finite dimension.
The formula reduces this optimization to maximum-weight matching on a finite
general graph whose vertices label the positive principal angles.
Its proof uses the principal-angle geometry, real and complex operator
  constructions, and singular-value norms developed for the formalization,
  together with additional optimization arguments.

Let $\like{Usource}{K},\like{Vtarget}{L}\subseteq\mathbb R^n$ have the same dimension and let
$\like{theta}{\theta_1},\ldots,\like{theta}{\theta_m}\in(0,\pi/2]$ be their positive principal angles,
counted with multiplicity. A matching $\sym{Matching}$ chooses unordered
pairs of labels from $\{1,\ldots,m\}$, using each label at most once.
Labels that are not used in a pair are unmatched.
For three labels, $\sym{Matching}=\{\{1,2\}\}$ pairs labels $1,2$ and
leaves $3$ unmatched. The choice $\{\{1,2\},\{2,3\}\}$ is disallowed:
label $2$ would occur in two pairs.

Each positive principal angle describes a source line and a target line
in a two-dimensional principal plane. Transporting that plane independently
allows a rotation through $\like{theta}{\theta_i}$, with trace displacement
$4\sin(\like{theta}{\theta_i}/2)$, or a reflection carrying the source line
to the target, with trace displacement $2$. A pair of labels $\{i,j\}$
allows the two principal planes to interact on their four-dimensional sum.
For each possible pair of distinct labels $i,j$, define its paired cost
$\like{PairCost}{e_{ij}}$; for each unmatched label $i$, define its singleton
cost $\like{SingletonCost}{s_i}$:
\begin{equation}\label{eq:nuclear-matching-costs}
 \like{PairCost}{e_{ij}}=4\sin\bigl((\like{theta}{\theta_i}+\like{theta}{\theta_j})/2\bigr),\qquad
 \like{SingletonCost}{s_i}=\min\{4\sin(\like{theta}{\theta_i}/2),2\}.
\end{equation}
The paired cost is attained by an orthogonal transporter on the two-plane sum.
Write $\sym{OrthogonalGroup}$ for the orthogonal group of real
$n\times n$ matrices, defined by $\sym{Wcompetitor}^T\sym{Wcompetitor}=I$.

\begin{theorem}[Real trace-norm matching reduction]
\label{thm:real-matching-example}
For equal-dimensional real subspaces $\like{Usource}{K},\like{Vtarget}{L}$ as above,
\begin{equation}\label{eq:real-nuclear-match}
\min_{\substack{\sym{Wcompetitor}\in \sym{OrthogonalGroup}\\\sym{Wcompetitor}(\like{Usource}{K})=\like{Vtarget}{L}}}\|I-\sym{Wcompetitor}\|_1
=
\min_{\sym{Matching}\text{ a matching}}
\left(
\sum_{\{i,j\}\in \sym{Matching}}\like{PairCost}{e_{ij}}
+
\sum_{i\text{ unmatched}}\like{SingletonCost}{s_i}
\right).
\end{equation}
Both minima are attained.  The equality includes repeated angles,
right angles, and common ambient directions; the left side has no
principal-plane or determinant constraint.
\end{theorem}

To express the formula as maximum-weight matching, start from the
all-singleton cost $\sym{SingletonBaseline}=\sum_i \like{SingletonCost}{s_i}$ and assign each edge $\{i,j\}$ the weight
\[
 \sym{MatchingEdgeWeight}=\like{SingletonCost}{s_i}+\like{SingletonCost}{s_j}-\like{PairCost}{e_{ij}}.
\]
This weight is the difference between the cost of two singletons and the
paired cost. Thus a matching $\sym{Matching}$ has cost
\[
 \like{TraceCost}{C(M)}=\sym{SingletonBaseline}-
       \sum_{\{i,j\}\in\sym{Matching}}\sym{MatchingEdgeWeight}.
\]
Minimizing this cost is equivalent to maximizing the sum of edge weights.
The empty matching is allowed, and edges of negative weight can be omitted.
Figure~\ref{fig:real-nuclear-matching} shows a three-angle example.

\begin{figure}[tb]
\centering
\resizebox{\linewidth}{!}{\input{figures/real_nuclear_matching.pgf}}
\caption{
Graph of principal-angle pairing choices for two three-dimensional subspaces
of $\mathbb R^6$ whose intersection is $\{0\}$ and whose principal angles are
$30^\circ$, $45^\circ$, and $75^\circ$.
Vertex positions are arbitrary; each vertex represents an angle label and
each edge a possible pairing. Vertices carry singleton costs $\like{SingletonCost}{s_i}$;
edges carry weights $\sym{MatchingEdgeWeight}=\like{SingletonCost}{s_i}+\like{SingletonCost}{s_j}-\like{PairCost}{e_{ij}}$. The solid edge
pairs labels $1,2$ and leaves label $3$ unmatched. The table lists all four
matchings; costs and edge weights are rounded to four decimals. Geometrically,
the selected pair is realized by a four-dimensional transporter and the
unmatched angle by a two-dimensional reflection. Their orthogonal sum
attains the minimum on $\mathbb R^6$, with trace displacement
$4\sin(37.5^\circ)+2\approx4.4350$.}
\label{fig:real-nuclear-matching}
\end{figure}

\begin{proof}[Proof and explanation]
The trace norm, corresponding to $p=1$, need not be differentiable.
We first replace it by a differentiable displacement cost with $1<p<2$.
At a global minimizer, derivatives along constraint-preserving variations
vanish. We use these equations to decompose the minimizer into orthogonal
blocks carrying one or two principal angles. The blocks determine a matching;
conversely, every matching constructs a transporter with the same cost.
This proves equality of the two minimum values for $1<p<2$.
We then let $p$ approach $1$ from above to recover the trace norm.

\emph{Step 1: choose a differentiable cost and a global minimizer.}
For $1<p<2$, set
\[
 \sym{FeasibleTransport}=\{\sym{Wcompetitor}\in \sym{OrthogonalGroup}:
                 \sym{Wcompetitor}(\like{Usource}{K})=\like{Vtarget}{L}\},
 \qquad
 \like{PowerCost}{\Phi_p}(\sym{Wcompetitor})
       =\sum_{j=1}^n \like{TransportSingularValues}{s_j}(I-\sym{Wcompetitor})^p.
\]
Here $\like{TransportSingularValues}{s_j}(T)$ denotes a singular value of $T$. For $p>1$ this cost is
differentiable, including at zero displacement singular values.

Write $\like{Psource}{P_K}$ and $\like{Qtarget}{P_L}$ for the orthogonal
projections onto the source and target. An orthogonal map carrying an
orthonormal basis of $\like{Usource}{K}$ to one of $\like{Vtarget}{L}$,
extended to their orthogonal complements, belongs to $\sym{FeasibleTransport}$.
Thus $\sym{FeasibleTransport}$ is nonempty. The group $\sym{OrthogonalGroup}$ is closed and bounded
in $\mathbb R^{n\times n}$, hence compact by Heine--Borel.
The equivalent transport equation
$\sym{Wcompetitor}\like{Psource}{P_K}\sym{Wcompetitor}^T
 =\like{Qtarget}{P_L}$ is a closed constraint.
Since $\like{PowerCost}{\Phi_p}$ is continuous, the extreme value theorem
gives a global minimizer $\sym{Wcompetitor}\in\sym{FeasibleTransport}$.

The decomposition we seek is an orthogonal sum
$\mathbb R^n=\bigoplus_a\sym{MatchingBlock}$ in which each $\sym{MatchingBlock}$ is preserved by
$\like{Psource}{P_K}$, $\like{Qtarget}{P_L}$, and $\sym{Wcompetitor}$.
In a basis adapted to this sum all three matrices are block diagonal.
Consequently, the source and target split as the sums of their
intersections with the $\sym{MatchingBlock}$, and
\begin{equation}\label{eq:matching-block-additivity}
 \like{PowerCost}{\Phi_p}(\sym{Wcompetitor})
 =\sum_a\like{PowerCost}{\Phi_p}(\sym{Wcompetitor}|_{\sym{MatchingBlock}}).
\end{equation}
Indeed, the displacement is block diagonal too, and its singular-value
list is the union of the block singular-value lists. The principal-angle
lists also combine in this way: the compression
$\like{Psource}{P_K}\like{Qtarget}{P_L}\like{Psource}{P_K}$
is block diagonal, and its restriction to $\like{Usource}{K}$ has the
squared principal cosines as eigenvalues.
These two additivity facts will turn the block decomposition into a
matching of the original angle labels.

\emph{Step 2: use derivatives along maps that still transport the source.}
Consider smooth curves $\sym{Wcompetitor}(t)\in\sym{FeasibleTransport}$, for real $t$
near zero, with $\sym{Wcompetitor}(0)=\sym{Wcompetitor}$.
The cost derivative at zero vanishes along every such curve because
$\sym{Wcompetitor}$ is a minimizer. This is the stationarity condition.
We use it as a necessary condition on the global minimizer already
obtained in Step~1.

Let $\sym{SkewDirection}$ be skew-symmetric (\ie{}
$\sym{SkewDirection}^T=-\sym{SkewDirection}$).
For real $t$, its exponential is orthogonal, since
\[
 \exp(t\sym{SkewDirection})^T
 =\exp(-t\sym{SkewDirection})
 =\exp(t\sym{SkewDirection})^{-1}.
\]
If $\sym{SkewDirection}$ commutes with $\like{Psource}{P_K}$, so does its
exponential.\jonsays{Why? Can use parentheticals for obvious things} It therefore maps the source onto itself, and
$\sym{Wcompetitor}\exp(t\sym{SkewDirection})$ remains in $\sym{FeasibleTransport}$.
This gives explicit curves on which to test stationarity.

To describe the derivative, define the scalar phase cost
\[
 \sym{MatchingPhaseCost}(t)=(2\sin(t/2))^p,\qquad 0\le t\le\pi.
\]
A planar rotation $\like{MatchingPlanarRotation}{R_t}$ through $t$ satisfies
\[
 (I-\like{MatchingPlanarRotation}{R_t})^T(I-\like{MatchingPlanarRotation}{R_t})
   =2I-\like{MatchingPlanarRotation}{R_t}-\like{MatchingPlanarRotation}{R_t}^T
   =4\sin^2(t/2)I.
\]
Its two displacement singular values are therefore both $2\sin(t/2)$,
and its contribution is $2\sym{MatchingPhaseCost}(t)$.
On such a rotation plane, define the skew-symmetric spectral gradient by
\[
 \sym{SpectralGradient}
 =\sym{MatchingPhaseCost}'(t)
       \begin{pmatrix}0&-1\\1&0\end{pmatrix}.
\]
On the $+1$ and $-1$ eigenspaces of the transporter it is \jonsays{Say what ``it'' is} zero.
Every real orthogonal matrix splits into these scalar blocks and planar
rotations: diagonalize over $\mathbb C$ and pair conjugate eigenvectors
to obtain the real planes. This defines $\sym{SpectralGradient}$ on the
whole space and shows that it commutes with $\sym{Wcompetitor}$.
These rotation planes may mix source and complementary directions;
stationarity will relate them to $\like{Psource}{P_K}$.
Differentiating the spectral cost gives
\begin{equation}\label{eq:matching-first-variation}
 \left.\frac{d}{dt}
   \like{PowerCost}{\Phi_p}(\sym{Wcompetitor}\exp(t\sym{SkewDirection}))
 \right|_{t=0}
 =-\operatorname{tr}(\sym{SpectralGradient}\sym{SkewDirection}).
\end{equation}
For example, on a rotation plane with angle $\vartheta$, a direction
$\sym{SkewDirection}=b\bigl(\begin{smallmatrix}0&-1\\1&0\end{smallmatrix}\bigr)$
changes the angle at speed $b$. Both sides of
\cref{eq:matching-first-variation} are then
$2b\sym{MatchingPhaseCost}'(\vartheta)$.

Now take the particular direction
$\sym{SkewDirection}=\like{Psource}{P_K}\sym{SpectralGradient}\like{Psource}{P_K}$.
It is skew-symmetric and commutes with $\like{Psource}{P_K}$, so stationarity
and cyclicity of trace give
\[
 0=-\operatorname{tr}\bigl(
      (\like{Psource}{P_K}\sym{SpectralGradient}\like{Psource}{P_K})^2
     \bigr)
   =\|\like{Psource}{P_K}\sym{SpectralGradient}\like{Psource}{P_K}\|_F^2.
\]
The source diagonal block of the gradient is therefore zero.
Using $I-\like{Psource}{P_K}$ instead gives the same conclusion for the
complementary diagonal block. Relative to
$\like{Usource}{K}\oplus\like{Usource}{K}^{\perp}$, the gradient and source
reflection thus have the forms
\[
 \sym{SpectralGradient}
  =\begin{pmatrix}0&\sym{GradientOffDiagonal}\\-\sym{GradientOffDiagonal}^T&0\end{pmatrix},
 \qquad
 \sym{Jreflect}=2\like{Psource}{P_K}-I
  =\begin{pmatrix}I&0\\0&-I\end{pmatrix}.
\]
Multiplication gives
$\sym{Jreflect}\sym{SpectralGradient}\sym{Jreflect}
 =-\sym{SpectralGradient}$.
This identity is what makes the source projection compatible with the
decomposition constructed next.

\emph{Step 3: obtain blocks with one or two source directions.}
The matrix $-\like{SpectralGradient}{G^2}$ is symmetric and nonnegative,
since
$\langle-\like{SpectralGradient}{G^2}x,x\rangle
  =\|\sym{SpectralGradient}x\|^2$.
Decompose the space into its eigenspaces. The gradient, transporter,
and source reflection all preserve them. On an eigenspace $\sym{GradientLevelSpace}$
with positive eigenvalue $\like{GradientLevel}{g^2}$, put
$\sym{GradientComplexStructure}=\sym{SpectralGradient}/\sym{GradientLevel}$.
It satisfies $\sym{GradientComplexStructure}^2=-I$ and preserves lengths.
We can therefore treat $\sym{GradientLevelSpace}$ as a complex space by defining
\[
 (a+ib)x=ax+b\sym{GradientComplexStructure}x.
\]
A real plane spanned by $x,\sym{GradientComplexStructure}x$ becomes one
complex line. Commutation with the gradient makes the transporter
complex-linear and unitary in these coordinates.
The source reflection reverses the complex structure:
\[
 \sym{Jreflect}\sym{GradientComplexStructure}
   =-\sym{GradientComplexStructure}\sym{Jreflect}.
\]
Consequently
$\sym{Jreflect}((a+ib)x)=(a-ib)\sym{Jreflect}x$.
Together with $\sym{Jreflect}^2=I$, this says that it acts as conjugation.

The definition of the gradient says that an eigenphase $t$ of the
transporter on this complex space satisfies
$\sym{MatchingPhaseCost}'(t)=\sym{GradientLevel}$, with $0<t<\pi$.
There are at most two such phases. To check this directly,
\[
 \sym{MatchingPhaseCost}'(t)
   =p(2\sin(t/2))^{p-1}\cos(t/2),
 \qquad
 \sym{MatchingPhaseCost}''(t)
   =p(2\sin(t/2))^{p-2}
       \bigl((p-1)\cos^2(t/2)-\sin^2(t/2)\bigr).
\]
The second derivative changes sign once, where
$\tan^2(t/2)=p-1$. The first derivative therefore increases and then
decreases, so a positive horizontal level has at most two intersections.
Writing the two phases as $\like{GradientEigenphase}{\varphi_1},\like{GradientEigenphase}{\varphi_2}$, with equality allowed,
we obtain
\[
 \sym{Wcompetitor}|_{\sym{GradientLevelSpace}}
   =e^{i\like{GradientEigenphase}{\varphi_1}}\sym{GradientPhaseProjection}+e^{i\like{GradientEigenphase}{\varphi_2}}(I-\sym{GradientPhaseProjection}),
\]
where $\sym{GradientPhaseProjection}$ is the orthogonal projection onto the first phase space.
A subspace preserved by $\sym{GradientPhaseProjection}$ is preserved by this transporter.
To preserve the source projection as well, we require invariance under
the conjugation $\sym{Jreflect}$.

Choose complex coordinates in which $\sym{Jreflect}$ is coordinate
conjugation and write $\sym{GradientPhaseProjection}=\sym{ProjectionRealPart}+i\sym{ProjectionImaginaryPart}$ with real matrices $\sym{ProjectionRealPart},\sym{ProjectionImaginaryPart}$.
The equations $\sym{GradientPhaseProjection}^*=\sym{GradientPhaseProjection}$ and $\sym{GradientPhaseProjection}^2=\sym{GradientPhaseProjection}$ give
\[
 \sym{ProjectionRealPart}^T=\sym{ProjectionRealPart},\qquad \sym{ProjectionImaginaryPart}^T=-\sym{ProjectionImaginaryPart},\qquad
 \sym{ProjectionRealPart}^2-\sym{ProjectionImaginaryPart}^2=\sym{ProjectionRealPart},\qquad \sym{ProjectionRealPart}\sym{ProjectionImaginaryPart}+\sym{ProjectionImaginaryPart}\sym{ProjectionRealPart}=\sym{ProjectionImaginaryPart}.
\]
For real coordinate vectors $x$,
$\langle \sym{ProjectionRealPart}x,x\rangle=\operatorname{Re}\langle \sym{GradientPhaseProjection}x,x\rangle$
lies between zero and $\|x\|^2$, so the eigenvalues of $\sym{ProjectionRealPart}$ lie in $[0,1]$.
If $\sym{ProjectionRealPart}u=\sym{ProjectionEigenvalue} u$, the last two identities imply
\[
 \sym{ProjectionRealPart}(\sym{ProjectionImaginaryPart}u)=(1-\sym{ProjectionEigenvalue})\sym{ProjectionImaginaryPart}u,\qquad
 \sym{ProjectionImaginaryPart}^2u=-\sym{ProjectionEigenvalue}(1-\sym{ProjectionEigenvalue})u.
\]
For a unit $u$ with $0<\sym{ProjectionEigenvalue}<1$, set
$\sym{ProjectionCoupling}=\sqrt{\sym{ProjectionEigenvalue}(1-\sym{ProjectionEigenvalue})}$ and $v=\sym{ProjectionImaginaryPart}u/\sym{ProjectionCoupling}$.
Skew-symmetry shows that $u,v$ are orthonormal, and
\[
 \sym{ProjectionRealPart}u=\sym{ProjectionEigenvalue} u,\quad \sym{ProjectionRealPart}v=(1-\sym{ProjectionEigenvalue})v,\quad
 \sym{ProjectionImaginaryPart}u=\sym{ProjectionCoupling} v,\quad \sym{ProjectionImaginaryPart}v=-\sym{ProjectionCoupling} u.
\]
Thus their real span is preserved by both $\sym{ProjectionRealPart}$ and $\sym{ProjectionImaginaryPart}$.
This includes $\sym{ProjectionEigenvalue}=1/2$. At $\sym{ProjectionEigenvalue}=0$ or $1$, $\sym{ProjectionImaginaryPart}u=0$ and a
single eigenvector gives an invariant real line.
By symmetry of $\sym{ProjectionRealPart}$ and skew-symmetry of $\sym{ProjectionImaginaryPart}$, the orthogonal complement
of each chosen line or plane is again preserved by $\sym{ProjectionRealPart},\sym{ProjectionImaginaryPart}$.
The construction can be repeated until the real
coordinate space is exhausted.

The complex spans of these lines and planes are invariant under
$\sym{GradientPhaseProjection},\sym{Jreflect}$, and $\sym{Wcompetitor}$. Returning to the original
real space gives blocks of dimension two or four. Each is also
preserved by
$\like{Psource}{P_K}=(\sym{Jreflect}+I)/2$, hence by
$\like{Qtarget}{P_L}
 =\sym{Wcompetitor}\like{Psource}{P_K}\sym{Wcompetitor}^T$.
The $+1$ eigenspace of conjugation has half the real dimension.
Accordingly, a two-dimensional block contains one source direction
and a four-dimensional block contains two. Their principal-angle lists
have one and two entries, respectively.

On $\ker\sym{SpectralGradient}$ the transporter has eigenvalues $+1,-1$,
because these are exactly the phases where the scalar derivative vanishes.
It is therefore a reflection there. Apply the two-projection decomposition
to $\like{Psource}{P_K}$ and the projection onto its $-1$ eigenspace:
this gives real lines and planes preserved by both projections and the
transporter. On an invariant line the source and target are either both
zero or both the whole line. Replacing a one-dimensional $-1$ block by
$+1$ preserves them and decreases the cost by $2^p$. Such a block cannot
occur at a global minimum. The remaining blocks are fixed directions
or planes with one source direction.
This completes the decomposition needed for
\cref{eq:matching-block-additivity}.

\emph{Step 4: calculate the costs on a plane and on two planes.}
For a principal angle $\like{theta}{\theta_i}$, the least rotation cost
is $2\sym{MatchingPhaseCost}(\like{theta}{\theta_i})$.
The alternative rotation sends a source unit vector to the opposite
target unit vector and turns through $\pi-\like{theta}{\theta_i}$,
which is at least $\like{theta}{\theta_i}$.
A planar reflection has eigenvalues $+1,-1$, hence displacement singular values $0,2$,
so its cost is $2^p$. The one-angle minimum is
\[
 \like{SingletonPowerCost}{a_i(p)}
   =\min\{2\sym{MatchingPhaseCost}(\like{theta}{\theta_i}),\,2^p\}.
\]

For two angles, relabel within the pair so that
$\like{theta}{\theta_i}\ge\like{theta}{\theta_j}$, and set
\[
 \sym{PairMeanAngle}=(\like{theta}{\theta_i}+\like{theta}{\theta_j})/2,\qquad
 \sym{PairHalfDifference}=(\like{theta}{\theta_i}-\like{theta}{\theta_j})/2.
\]
Then $0\le \sym{PairHalfDifference}\le \sym{PairMeanAngle}$ and $\sym{PairMeanAngle}+\sym{PairHalfDifference}\le\pi/2$.
In principal coordinates the direct rotation is
\[
 \sym{PairDirectRotation}=\begin{pmatrix}\sym{PairCosineBlock}&-\sym{PairSineBlock}\\\sym{PairSineBlock}&\sym{PairCosineBlock}\end{pmatrix},\qquad
 \sym{PairCosineBlock}=\operatorname{diag}(\cos\like{theta}{\theta_i},
                      \cos\like{theta}{\theta_j}),\quad
 \sym{PairSineBlock}=\operatorname{diag}(\sin\like{theta}{\theta_i},
                      \sin\like{theta}{\theta_j}).
\]
Every orthogonal transporter on this four-dimensional space is
$\sym{PairDirectRotation}\operatorname{diag}(\like{PairFreeOrthogonal}{U_0},\like{PairFreeOrthogonal}{U_1})$ with $\like{PairFreeOrthogonal}{U_0},\like{PairFreeOrthogonal}{U_1}\in \like{OrthogonalGroup}{O(2)}$:
multiplying by $\sym{PairDirectRotation}^T$ gives an orthogonal map preserving the source
and its complement. This parametrization allows mixing between the
principal planes.

For determinant $+1$, the two matrices $\like{PairFreeOrthogonal}{U_0},\like{PairFreeOrthogonal}{U_1}$ are either both
rotations or both reflections. The planar matrix calculation gives
phase parameters $\sym{PairPhaseX},\sym{PairPhaseY}$ with
\[
 \sym{PairMeanAngle}\le \sym{PairPhaseX}\le\pi-\sym{PairMeanAngle},\qquad \sym{PairHalfDifference}\le \sym{PairPhaseY}\le\pi-\sym{PairHalfDifference},
\]
and eigenvalues $e^{\pm i(\sym{PairPhaseX}+\sym{PairPhaseY})},e^{\pm i(\sym{PairPhaseX}-\sym{PairPhaseY})}$.
After folding the angles into $[0,\pi]$, the displacement cost is
$2\bigl(\sym{MatchingPhaseCost}(\min\{\sym{PairPhaseX}+\sym{PairPhaseY},2\pi-\sym{PairPhaseX}-\sym{PairPhaseY}\})
       +\sym{MatchingPhaseCost}(|\sym{PairPhaseX}-\sym{PairPhaseY}|)\bigr)$.
For example, if $\like{PairFreeOrthogonal}{U_0},\like{PairFreeOrthogonal}{U_1}$ rotate through $\sym{PairRotationR}+\sym{PairRotationS},\sym{PairRotationR}-\sym{PairRotationS}$, respectively,
these parameters satisfy
$\cos \sym{PairPhaseX}=\cos \sym{PairMeanAngle}\cos \sym{PairRotationS}$ and $\cos \sym{PairPhaseY}=\cos \sym{PairHalfDifference}\cos \sym{PairRotationR}$.
For two reflections, write
$\like{PairFreeOrthogonal}{U_0}=\like{MatchingPlanarRotation}{R_{r+s}}\sym{PairPlanarReflection}$ and $\like{PairFreeOrthogonal}{U_1}=\like{MatchingPlanarRotation}{R_{r-s}}\sym{PairPlanarReflection}$ with
$\sym{PairPlanarReflection}=\operatorname{diag}(1,-1)$. The calculation gives
$\cos \sym{PairPhaseX}=-\sin \sym{PairHalfDifference}\cos \sym{PairRotationR}$ and $\cos \sym{PairPhaseY}=\sin \sym{PairMeanAngle}\cos \sym{PairRotationS}$,
which lie in the same ranges because $\sym{PairMeanAngle}+\sym{PairHalfDifference}\le\pi/2$.

We can replace $(\sym{PairPhaseX},\sym{PairPhaseY})$ by $(\pi-\sym{PairPhaseX},\pi-\sym{PairPhaseY})$ when $\sym{PairPhaseX}+\sym{PairPhaseY}>\pi$ without
changing this cost. Suppose therefore that $\sym{PairPhaseX}+\sym{PairPhaseY}\le\pi$.
To replace the first parameter by $\sym{PairMeanAngle}$ without increasing the smaller
phase, we need $\sym{PairMeanAngle}-\sym{PairPhaseParameter}\le|\sym{PairPhaseX}-\sym{PairPhaseY}|$, as well as $\sym{PairPhaseParameter}\ge \sym{PairHalfDifference}$.
Choose the smallest such parameter,
$\sym{PairPhaseParameter}=\max\{\sym{PairHalfDifference},\sym{PairMeanAngle}-|\sym{PairPhaseX}-\sym{PairPhaseY}|\}$.
The identities $\sym{PairPhaseX}+\sym{PairPhaseY}+|\sym{PairPhaseX}-\sym{PairPhaseY}|=2\max\{\sym{PairPhaseX},\sym{PairPhaseY}\}\ge2\sym{PairMeanAngle}$ and $\sym{PairPhaseX}+\sym{PairPhaseY}\ge \sym{PairMeanAngle}+\sym{PairHalfDifference}$
also give $\sym{PairMeanAngle}+\sym{PairPhaseParameter}\le \sym{PairPhaseX}+\sym{PairPhaseY}$. Thus
\[
 \sym{PairHalfDifference}\le\sym{PairPhaseParameter}\le \sym{PairMeanAngle},\qquad
 \sym{PairMeanAngle}+\sym{PairPhaseParameter}\le \sym{PairPhaseX}+\sym{PairPhaseY},\qquad \sym{PairMeanAngle}-\sym{PairPhaseParameter}\le|\sym{PairPhaseX}-\sym{PairPhaseY}|.
\]
Since $\sym{MatchingPhaseCost}$ is increasing on $[0,\pi]$, this
transporter costs at least
$2\bigl(\sym{MatchingPhaseCost}(\sym{PairMeanAngle}+\sym{PairPhaseParameter})+
        \sym{MatchingPhaseCost}(\sym{PairMeanAngle}-\sym{PairPhaseParameter})\bigr)$.
Every $\sym{PairPhaseParameter}\in[\sym{PairHalfDifference},\sym{PairMeanAngle}]$ is attained: take
$\like{PairFreeOrthogonal}{U_0}=\like{PairFreeOrthogonal}{U_1}=\like{MatchingPlanarRotation}{R_\gamma}$, where
$\sym{PairAttainingAngle}=\arccos(\cos\sym{PairPhaseParameter}/\cos \sym{PairHalfDifference})$.
The denominator is positive since $\sym{PairHalfDifference}\le\pi/4$.
The preceding phase equations give $\sym{PairPhaseX}=\sym{PairMeanAngle},\sym{PairPhaseY}=\sym{PairPhaseParameter}$.
Thus the positive-determinant pair minimum is
\begin{equation}\label{eq:matching-pair-power-cost}
 \like{PairPowerCost}{b_{ij}(p)}
 =\min_{\sym{PairPhaseParameter}\in[\sym{PairHalfDifference},\sym{PairMeanAngle}]}
     2\bigl(\sym{MatchingPhaseCost}(\sym{PairMeanAngle}+\sym{PairPhaseParameter})+
             \sym{MatchingPhaseCost}(\sym{PairMeanAngle}-\sym{PairPhaseParameter})\bigr).
\end{equation}

For determinant $-1$, the eigenvalues are $+1,-1$ and a conjugate
rotational pair $e^{\pm it}$, with $0\le t\le\pi$.
The planar calculation gives
$|\operatorname{tr}\sym{Wcompetitor}|\le2\cos\like{theta}{\theta_j}$.
Since this trace is $2\cos t$, the rotational phase satisfies
$t\ge\like{theta}{\theta_j}$.
Its least cost is therefore
$2^p+2\sym{MatchingPhaseCost}(\like{theta}{\theta_j})$.
It is attained by reflecting on the larger-angle plane and rotating
on the smaller-angle plane.
Taking $\sym{PairPhaseParameter}=\sym{PairHalfDifference}$ bounds the positive-determinant minimum by the cost
of two rotations. Taking $\sym{PairPhaseParameter}=\sym{PairMeanAngle}$ bounds it by
$2\sym{MatchingPhaseCost}(\like{theta}{\theta_i}+\like{theta}{\theta_j})
 \le2\cdot2^p$, the cost of two reflections.
The less expensive mixed rotation/reflection choice is the
negative-determinant minimum just computed.
Comparing these four independent choices shows that the unrestricted
four-dimensional minimum is
\[
 \min\{\like{PairPowerCost}{b_{ij}(p)},\,
         \like{SingletonPowerCost}{a_i(p)}
          +\like{SingletonPowerCost}{a_j(p)}\}.
\]
A four-dimensional block therefore offers either a paired choice or
two independent singleton choices.

\emph{Step 5: prove equality of the minimum costs and recover a transporter.}
We construct a transporter from every matching and extract a matching
from a minimizing transporter, preserving the cost in both constructions.
Thus solving the matching problem gives the minimum cost and an attaining
transporter, although individual transporters and matchings need not
correspond one-to-one.
For a matching $\sym{Matching}$ define
\[
 \sym{MatchingPowerCost}(\sym{Matching})
   =\sum_{\{i,j\}\in\sym{Matching}}\like{PairPowerCost}{b_{ij}(p)}
     +\sum_{i\text{ unmatched}}\like{SingletonPowerCost}{a_i(p)}.
\]
The principal-angle decomposition gives mutually orthogonal planes
for the positive angle labels. Choose an attaining two-angle transporter for each selected pair,
an attaining planar rotation or reflection for each unmatched label,
and identity on the remaining common directions.
Their orthogonal sum is an element of $\sym{FeasibleTransport}$.
By \cref{eq:matching-block-additivity}, its cost is
$\sym{MatchingPowerCost}(\sym{Matching})$. Hence
$\min_{\sym{FeasibleTransport}}\like{PowerCost}{\Phi_p}
 \le\min_{\sym{Matching}}\sym{MatchingPowerCost}(\sym{Matching})$.

For the reverse bound, use the global minimizer from Step~1 and its
decomposition from Step~3. Its restriction to each block must minimize
the corresponding block problem: otherwise replacing that restriction
by a cheaper transporter would lower the global cost while retaining
the same source and target. By Step~4, a plane supplies a singleton,
and a four-dimensional block supplies a pair or two singletons.
Each occurrence of an angle belongs to one block, so no label can be used twice.
This produces a matching with the same total cost as the minimizer.

Zero angles do not supply cheaper extra partners. Their singleton cost
is zero; pairing an angle $\theta$ with zero forces $\sym{PairHalfDifference}=\sym{PairMeanAngle}=\theta/2$
in \cref{eq:matching-pair-power-cost}, giving $2\sym{MatchingPhaseCost}(\theta)$,
which is at least its singleton cost.
Thus we can leave zero-angle directions fixed and use only the original
positive labels. We have proved
\begin{equation}\label{eq:matching-power-equality}
 \min_{\sym{FeasibleTransport}}\like{PowerCost}{\Phi_p}
   =\min_{\sym{Matching}}\sym{MatchingPowerCost}(\sym{Matching}),
 \qquad 1<p<2.
\end{equation}
For the matching in \cref{fig:real-nuclear-matching}, this construction
uses one four-dimensional block for labels $1,2$ and one plane for label~$3$.

\emph{Step 6: pass to the trace norm and evaluate the weights.}
Every displacement singular value lies in $[0,2]$.
Since $\sym{DisplacementSize}^p$ converges uniformly to $\sym{DisplacementSize}$ on this interval as $p$ approaches $1$ from above,
\[
 \sup_{\sym{Wcompetitor}\in\sym{FeasibleTransport}}
   \left|\like{PowerCost}{\Phi_p}(\sym{Wcompetitor})
                  -\|I-\sym{Wcompetitor}\|_1\right|
 \le n\sup_{0\le\sym{DisplacementSize}\le2}|\sym{DisplacementSize}^p-\sym{DisplacementSize}|\ \longrightarrow\ 0.
\]
The minimum values on the left of \cref{eq:matching-power-equality}
therefore converge to the minimum trace norm.
The scalar block minima also converge: the one-angle minimum is a
minimum of two continuous quantities, and the pair minimum is over
the fixed compact interval $[\sym{PairHalfDifference},\sym{PairMeanAngle}]$ with a uniformly convergent cost.
There are finitely many matchings, so their minimum values converge too.

At $p=1$ the singleton formula gives
$\like{SingletonPowerCost}{a_i(1)}
 =\min\{4\sin(\like{theta}{\theta_i}/2),2\}
 =\like{SingletonCost}{s_i}$.
For the paired cost, the scalar expression becomes
\[
 2\bigl(\sym{MatchingPhaseCost}(\sym{PairMeanAngle}+\sym{PairPhaseParameter})
          +\sym{MatchingPhaseCost}(\sym{PairMeanAngle}-\sym{PairPhaseParameter})\bigr)\big|_{p=1}
 =8\sin(\sym{PairMeanAngle}/2)\cos(\sym{PairPhaseParameter}/2).
\]
It is nonincreasing as $\sym{PairPhaseParameter}$ runs from $\sym{PairHalfDifference}$ to $\sym{PairMeanAngle}$, so its minimum is at $\sym{PairPhaseParameter}=\sym{PairMeanAngle}$,
with value $4\sin \sym{PairMeanAngle}=\like{PairCost}{e_{ij}}$.
At this endpoint one rotational phase is $2\sym{PairMeanAngle}$ and the other is zero:
the paired transporter fixes two real directions and rotates the other
two through $\like{theta}{\theta_i}+\like{theta}{\theta_j}$.
Taking the limit in \cref{eq:matching-power-equality} now gives
\cref{eq:real-nuclear-match}. Compactness supplies an attaining transporter
for the trace norm as well.
\end{proof}
% Formal proof locations: RealNuclearMatching.lean (endpoint continuation),
% SchattenRealDecomposition.lean / SchattenRealRestrictions.lean (block
% reduction), SchattenRealMatchingWeights.lean / RealMatchingAttainment.lean
% (matching and attainment), SchattenRealAngleCosts.lean (angle costs),
% RealFourCoordinates.lean / RealFourPositive.lean / RealFourOptimization.lean
% (all two-angle competitors and the scalar edge reduction).
% All paths are under DavisKahan/FiniteDimensional/DirectRotation/.
% See AUTHORING_NOTES.md for the mathematical and provenance evidence.

For $\like{theta}{\theta_1}=\like{theta}{\theta_2}=\pi/6$, the direct rotation costs
$8\sin(\pi/12)>2$, whereas the paired competitor costs
$4\sin(\pi/6)=2$. The theorem includes comparison with every ambient
orthogonal transporter, so this paired cost is globally minimal in the
example.

Davis and Kahan's complex-phase example already shows how a phase change
can decrease full trace displacement \citep[Example~4.2]{davis1970rotation}.
Bolotnikov, Li, and Rodman formulate the ambient unitary minimal-distortion
problem and prove Q-norm and Frobenius optimality results
\citep[Problem~1.1 and Theorems~3.1, 4.2]{bolotnikov2001minimal}. The formula
above specifies the unrestricted real trace-norm minimum through singleton
and paired angle costs. The real field permits pairing between distinct
principal planes, which is why the answer is an optimization over matchings.


\section{Discussion}
\label{sec:discussion}

The formalization covers the original perturbation theory on real and
complex Hilbert spaces, including unbounded self-adjoint operators and
operator-ideal norms. The unbounded statements require explicit domain
conditions and bounded residuals. The operator-ideal statements require
finite norm values. These hypotheses determine which approximation problems
the estimates address.

The correction of the full-displacement claim
\citep[Proposition~4.4]{davis1970rotation} separates two extremal statements.
The valid comparison \citep[Proposition~4.3]{davis1970rotation}
gives squared-displacement minimality; its Q-norm consequence
is valid. Our real four-dimensional family exhibits full trace-norm failure
at every prescribed positive acute angle. In each example, the competitor has
two nonzero displacement singular values while the direct rotation has four.
The sums and squared sums of these values account for their different norm
comparisons. The earlier complex example supplied the construction; our
calculation makes its contradiction of the printed real claim explicit.

The matching theorem addresses the resulting optimization problem for any
equal-dimensional pair of finite-dimensional real subspaces. Its proof uses
principal-angle geometry, real orthogonal block decompositions, and
singular-value norms. These are also among the foundations required for the
Davis--Kahan formalization. The proof of the matching theorem uses additional
optimization arguments, including a passage to the trace-norm endpoint.

\section{Conclusion}
\label{sec:conclusion}

We formalized 28 of the 29 selected Davis--Kahan results and refuted the
remaining printed statement \citep[Proposition~4.4]{davis1970rotation}.
The sine-theta theorem includes
unbounded self-adjoint operators and symmetric operator-ideal norms with
their domain and finiteness assumptions. The direct-rotation calculation
identifies the failed singular-value comparison and gives an explicit real
counterexample and a valid Q-norm consequence.

We also formalized the operator theory used in these results: principal
angles and direct rotations, Borel spectral calculus, approximation numbers
and symmetric ideals, and Sylvester separation. Parts of this development
were subsequently adapted into Tau Ceti. Using these foundations, we proved
an exact weighted-matching formula for the real trace-norm transport
minimum.

\section*{AI assistance disclosure}
ChatGPT and Claude assisted with the formalization and manuscript preparation.
Formalizations were carried out by LLMs with human oversight.
The new matching result was derived via GPT 6.1 sol and reviewed by the human
  authors.

\section*{Code and data availability}
\ifDKAnonymous
Artifact information is omitted for anonymous review.
\else
The formalization, result-to-Lean mapping, provenance records, and manuscript
sources are available at
\url{https://github.com/AIQ-Kitware/aiq-dkps-formalization}.
\fi

\ifDKAnonymous\else
\section*{Acknowledgments}

This material is based upon work supported by the Defense Advanced Research
Project Agency (DARPA) under Contract No. HR001125CE017. Any opinions, findings
and conclusions or recommendations expressed in this material are those of the
author(s) and do not necessarily reflect the views of the Defense Advanced
Research Project Agency (DARPA).
\fi

\bibliographystyle{unsrtnat}
\bibliography{references}

\clearpage
\appendix

\section{Symbols and conventions}
\label{app:notation}

We write $\sym{A}$ for a reference operator and
$\sym{Acurrent}=\sym{A}+\sym{H}$ for the perturbed operator. The
sine-theta theorem uses coordinate maps for trial and exact subspaces.
The transport results name the subspaces $\sym{Usource},\sym{Vtarget}$
and their projections $\sym{Psource},\sym{Qtarget}$ directly.

In the colored version, blue identifies exact or target data, amber trial
or source data, cyan the exact complementary block, magenta angles, green
residuals, and violet spectral separation. The residual $\sym{R}$ and
direct rotation $\sym{Rrotation}$ have distinct roles and colors. The
notation and definitions also apply in monochrome.

\begingroup
\small
\setlength{\tabcolsep}{4pt}
\renewcommand{\arraystretch}{1.15}
\begin{longtable}{@{}>{\raggedright\arraybackslash}p{0.29\textwidth}>{\raggedright\arraybackslash}p{0.65\textwidth}@{}}
\toprule
\textbf{Symbol} & \textbf{Definition or interpretation}\\
\midrule
\endfirsthead
\multicolumn{2}{@{}l}{\textit{Symbols and conventions (continued)}}\\[2pt]
\toprule
\textbf{Symbol} & \textbf{Definition or interpretation}\\
\midrule
\endhead
\bottomrule
\endfoot
$\sym{A},\sym{H},\sym{Acurrent}$ & Reference operator, additive perturbation, and perturbed operator $\sym{Acurrent}=\sym{A}+\sym{H}$.\\
$\sym{Eambient};\ \sym{Ftrial},\sym{Hexact},\sym{Gunwanted}$ & Ambient, trial-coordinate, exact-coordinate, and complementary-coordinate Hilbert spaces.\\
$\sym{E0}:\sym{Ftrial}\to\sym{Eambient}$ & \dkSymbolDescription{E0}\\
$\sym{A0}$ & Self-adjoint trial operator. A Rayleigh--Ritz choice is $\sym{A0}=\like{E0}{E_0^*}\sym{Acurrent}\sym{E0}$.\\
$\sym{ReferenceBlock0},\sym{ReferenceBlock1}$ & Reference blocks on $\DKRange(\sym{Psource})$ and $\DKKernel(\sym{Psource})$ in Section~\ref{sec:foundations}.\\
$\sym{F0}:\sym{Hexact}\to\sym{Eambient}$ & \dkSymbolDescription{F0}\\
$\sym{F1}:\sym{Gunwanted}\to\sym{Eambient}$ & \dkSymbolDescription{F1}\\
$\sym{Lambda1}$ & Complementary self-adjoint block, with $\sym{Acurrent}\sym{F1}=\sym{F1}\sym{Lambda1}$ on its domain.\\
$\sym{SpectralMeasure},\sym{SpectralSet}$ & Projection-valued measure and Borel spectral set; $\sym{SpectralMeasure}(\sym{SpectralSet})$ is the corresponding spectral projection. These symbols differ from the ambient space $\sym{Eambient}$ and complementary block $\sym{S}$.\\
$\sym{Pexact}=\sym{F0}\like{F0}{F_0^*}$ & Projection onto the desired exact subspace in the sine-theta theorem.\\
$\sym{Theta0},\sym{sinTheta0}$ & Trial-coordinate angle operator and its positive sine, with $\sym{sinTheta0}=(I-\like{E0}{E_0^*}\sym{Pexact}\sym{E0})^{1/2}$.\\
$\sym{Theta}=\operatorname{diag}(\sym{Theta0},\sym{Theta1})$ & Full angle operator; $\sym{Theta1}$ compares the orthogonal complements. Its operator norm is the largest spectral angle.\\
$\sym{CrossedSourceDefect},\sym{CrossedTargetDefect}$ & Source directions perpendicular to the target, and target directions perpendicular to the source. Section~\ref{sec:foundations} uses the condition $\dim\sym{CrossedSourceDefect}=\dim\sym{CrossedTargetDefect}$.\\
$\sym{S}=(I-\sym{Pexact})\sym{E0}$ & Complementary block, with modulus $|\sym{S}|=\sym{sinTheta0}$.\\
$\sym{X}=\like{F1}{F_1^*}\sym{E0}$ & Coordinates of the trial component in the unwanted space; $\sym{S}=\sym{F1}\sym{X}$.\\
$\sym{SylvesterLeft},\sym{SylvesterRight},\sym{SylvesterRHS}$ & Generic left and right Sylvester coefficients and right side: $\sym{SylvesterLeft}\sym{X}-\sym{X}\sym{SylvesterRight}=\sym{SylvesterRHS}$. The historical equation writes the right coefficient as $\like{SylvesterRight}{A}$.\\
$\sym{R}=\sym{Acurrent}\sym{E0}-\sym{E0}\sym{A0}$ & Bounded residual agreeing with this difference on $\DKDom(\sym{A0})$.\\
$\sym{delta}>0$ & \dkSymbolDescription{delta}\\
$[\sym{betaEndpoint},\sym{alphaEndpoint}]$ & Interval in the interval/exterior gap hypothesis.\\
$\DKRange(T),\DKKernel(T),\DKDom(T)$ & Range, kernel, and domain of an operator.\\
$\sym{Usource},\sym{Vtarget}$ & Source and target subspaces in the displacement problem.\\
$\sym{Psource},\sym{Qtarget}$ & Orthogonal projections onto $\sym{Usource},\sym{Vtarget}$ in that problem.\\
$\sym{Rrotation},\sym{Wcompetitor}$ & Direct rotation and an arbitrary orthogonal or unitary transporter carrying $\sym{Usource}$ onto $\sym{Vtarget}$.\\
$\sym{IntertwinerT}$ & The intertwiner $\sym{Qtarget}\sym{Psource}+(I-\sym{Qtarget})(I-\sym{Psource})$, whose unitary polar factor is the acute direct rotation.\\
$\sym{theta},\like{theta}{\theta_j}$ & Individual principal angles, in $[0,\pi/2]$.\\
$\sym{Xframe},\sym{Yframe}$ & Orthonormal coordinate matrices for the source and target in the real counterexample.\\
$\sym{Jturn},\sym{Cturn},\sym{Hmix}$ & Planar quarter-turn, rotation through $2\sym{theta}$, and orthogonal mixing matrix used in that counterexample.\\
$\sym{CounterPrincipalOne},\sym{CounterPrincipalTwo};\ \sym{CounterMovingPlane},\sym{CounterFixedPlane}$ & Principal planes of the direct rotation and moving/fixed planes of the competitor in \cref{fig:prop44-displacement-planes}.\\
$\sym{CounterChordA},\sym{CounterChordB}$ & Displacement singular values $\sqrt{2-\sqrt2}$ and $\sqrt2$ in that counterexample.\\
$\sym{Kdisplacement}=I-\sym{Wcompetitor}$ & Full displacement operator in the failed Ky Fan comparison.\\
$\sym{OmegaOne},\sym{OmegaTwo}$ & Principal-plane projections used to test that comparison.\\
$|T|=(T^*T)^{1/2}$ & Positive modulus of an operator.\\
$s_j(T)$ & Ordered singular values; approximation numbers provide the corresponding comparisons for bounded operators in Hilbert space.\\
$\sym{NormN}(T)$ & \dkSymbolDescription{NormN}\\
$\|T\|_{\mathrm{op}}$ & Operator norm; in finite dimension, the largest singular value.\\
$\|T\|_F$ & Frobenius norm, $\bigl(\sum_j s_j(T)^2\bigr)^{1/2}$.\\
$\|T\|_1$ & Trace norm, $\sum_j s_j(T)$.\\
$\|T\|_{(j)}$ & Ky Fan norm, the sum of the $j$ largest singular values.\\
$\sym{NormQ}\text{-norm}$ & A norm satisfying $\sym{NormN}(T)^2=\sym{NormM}(T^*T)$ for a unitarily invariant norm $\sym{NormM}$.\\
$\sym{Matching},\like{SingletonCost}{s_i},\like{PairCost}{e_{ij}}$ & Matching of positive principal-angle labels, with singleton and paired costs in \cref{eq:nuclear-matching-costs}.\\
$\sym{MatchingEdgeWeight};\ \like{TraceCost}{C(M)}$ & Edge weight $s_i+s_j-e_{ij}$; total trace-norm displacement cost of the matching.\\
\end{longtable}
\endgroup


\section{Reading the Lean signatures}
\label{app:lean-guide}
\label{app:reading-lean}

A Lean theorem declaration has the form
\[
  \texttt{theorem name (parameters) (hypotheses) : conclusion := proof}.
\]
Parameters introduce the mathematical objects under discussion.  A hypothesis
such as \texttt{(hA : IsSelfAdjoint A)} supplies the corresponding fact;
bundled predicates such as \texttt{IsTrialResidual} and
\texttt{IsExactSpectralDecomposition} collect related assumptions.  Braces mark
arguments normally inferred by Lean, while square brackets request typeclass
instances supplying ambient algebraic, analytic, or topological structure.  The proposition after the final colon is the theorem's conclusion
\citep{demoura2021lean4}.

In \Cref{lst:dk-sintheta-headline}, \texttt{RCLike} is the Mathlib abstraction
covering $\mathbb{R}$ and $\mathbb{C}$.  A type
\texttt{ContinuousLinearMap K F E} is a continuous linear map.  A type
\texttt{LinearPMap K E E} is a linear partial map whose
domain is a submodule; this is the representation used for possibly unbounded
operators.  \texttt{N.Mem T} states that the ideal gauge associated with $\sym{NormN}$ is
finite on $T$, and \texttt{N.gaugeReal T} is the corresponding real value.  The
suffix \texttt{Real} refers to the codomain of the gauge and not to the scalar
field.

\section{Formalized Davis--Kahan results}
\label{app:dk-inventory}

Table~\ref{tab:dk-full-inventory} lists the \DKResultCount{} source targets:
the four headline theorems \citep[Section~2]{davis1970rotation} and every named theorem,
proposition, lemma, and corollary established in the paper. The table keeps
compact IDs for reference. The artifact maps each ID to the corresponding
Lean declarations.
\ifDKAnonymous\else
The \href{https://github.com/AIQ-Kitware/aiq-dkps-formalization/blob/main/dev/davis-kahan-1970-formalization-result-inventory.json}{result-to-Lean mapping}
contains that correspondence.
\fi

\begin{longtable}{>{\raggedright\arraybackslash}p{0.24\linewidth}>{\raggedright\arraybackslash}p{0.68\linewidth}}
\caption{The \DKResultCount{} Davis--Kahan results selected for formalization.}
\label{tab:dk-full-inventory}\\
\toprule
ID & Result \\
\midrule
\endfirsthead
\toprule
ID & Result \\
\midrule
\endhead
\path{S2-sin-theta} & Single-angle sine theorem \\
\path{S2-tan-theta} & Single-angle tangent theorem \\
\path{S2-sin-two-theta} & Double-angle sine theorem \\
\path{S2-tan-two-theta} & Double-angle tangent theorem \\
\path{DK-3.1-prop} & Acute direct rotation existence and uniqueness \\
\path{DK-3.2-prop} & Nonacute existence criterion \\
\path{DK-3.3-prop} & Principal square-root characterization \\
\path{DK-3.4-prop} & Square as a direct rotation \\
\path{DK-3.1-thm} & Classification of pairs of subspaces \\
\path{DK-3.1-cor} & Compact classification by angle eigenvalues \\
\path{DK-3.5-prop} & Angle commutation and eigenspace geometry \\
\path{DK-3.2-cor} & Reversal symmetry \\
\path{DK-4.1-prop} & Pointwise and singular-value extremality of the direct rotation \\
\path{DK-4.1-cor} & Unitarily invariant norm minimality of restricted displacement \\
\path{DK-4.2-prop} & Basis-angle square-sum extremality \\
\path{DK-4.3-prop} & Squared-displacement unitarily invariant norm minimality \\
\path{DK-4.4-prop} & Full-displacement optimality claim \\
\path{DK-5.1-thm} & Banach-space Sylvester lower bound \\
\path{DK-5.2-thm} & Semibounded self-adjoint Sylvester theorem \\
\path{DK-5.1-lem} & Strong-cutoff convergence of singular values \\
\path{DK-6.1-lem} & Direct-sum norm comparison and converse \\
\path{DK-6.2-lem} & Reflection-pinch contraction \\
\path{DK-6.1-prop} & Sine proof, ambient limitation, and symmetric sine theorem \\
\path{DK-6.1-thm} & Generalized sine theorem \\
\path{DK-6.2-thm} & Pairwise-gap square-norm sine theorem \\
\path{DK-6.3-thm} & Generalized tangent theorem and supporting machinery \\
\path{DK-6.3-lem} & Finite-rank near-maximizer off-diagonal block estimate \\
\path{DK-8.1-thm} & Branch selection and spectral repulsion \\
\path{DK-8.2-thm} & Smallness selects the acute branch \\
\bottomrule
\end{longtable}

A single source result may use several Lean declarations. The IDs identify
source statements independently of the organization of the formal library.
The printed claim \citep[Proposition~4.4]{davis1970rotation}
is refuted by the counterexample in
Section~\ref{sec:prop44}; its Q-norm consequence is stated separately there.

\input{appendix_result_summaries.tex}


\end{document}
, without editing.
% EVERY \sym, \like, and \dkrole obeys this switch, including the glossary.
% No color is required to identify a mathematical object: the glyphs, indices,
% definitions and adjacent prose also distinguish them.
%
% Families (base hue): exact/target=blue, trial/source=amber, exact unwanted=
% cyan, angles and their operators=magenta, residual/displacement=green,
% gaps=violet. Norm gauges=slate; provenance retains brown/steel/vermilion.
% P_ex is the exact-target projection in sine-theta; P is the source
% projection in transport. R denotes a residual and D a direct rotation.
% Semantic keys retain their roles even when the printed notation changes.
\newif\ifDKMathColors
\DKMathColorstrue
\ifdefined\DKMonochromeMath
  \DKMathColorsfalse
\fi

% Family anchors (print-compatible versions of the Manim light palette).
\definecolor{dkfg}{HTML}{1B1F24}
\definecolor{dkmuted}{HTML}{6B7280}
\definecolor{dkwanted}{HTML}{1F6FD1}
\definecolor{dktrial}{HTML}{A66705}
\definecolor{dkunwanted}{HTML}{007C91}
\definecolor{dksine}{HTML}{D12A7B}
\definecolor{dkresid}{HTML}{087D43}
\definecolor{dkgap}{HTML}{6E4BD8}
\definecolor{dkold}{HTML}{705519}
\definecolor{dkcurrent}{HTML}{054669}
\definecolor{dkperturb}{HTML}{BB491B}
\definecolor{dkrefuted}{HTML}{C62828}

% Related yet distinguishable shades. These are NOT new unrelated categories:
% objects inherit their parent family; a different shade signals a different
% mathematical object (e.g., a frame versus its subspace or its projection).
\definecolor{dkwantedframe}{HTML}{1254A9}
\definecolor{dkwantedprojection}{HTML}{08618D}
\definecolor{dkwantedrotation}{HTML}{4248B0}
\definecolor{dkwantedspectrum}{HTML}{345F8F}
\definecolor{dkwantedvector}{HTML}{1468AA}
\definecolor{dktrialframe}{HTML}{905600}
\definecolor{dktrialprojection}{HTML}{79511B}
\definecolor{dktrialcompetitor}{HTML}{A04D10}
\definecolor{dktrialspectrum}{HTML}{79591C}
\definecolor{dktrialvector}{HTML}{9B522A}
\definecolor{dkunwantedframe}{HTML}{006477}
\definecolor{dkunwantedspectrum}{HTML}{00546B}
\definecolor{dkunwantedprojection}{HTML}{145B6D}
\definecolor{dkunwantedvector}{HTML}{256879}
\definecolor{dkangleoperator}{HTML}{923478}
\definecolor{dksineoperator}{HTML}{AA275A}
\definecolor{dksineblock}{HTML}{B43F70}
\definecolor{dksylvester}{HTML}{804074}
\definecolor{dkanglemax}{HTML}{B32185}
\definecolor{dkanglecos}{HTML}{824276}
\definecolor{dkresidvector}{HTML}{086038}
\definecolor{dkdisplacement}{HTML}{226C50}
\definecolor{dkgaporiginal}{HTML}{4E4D89}
\definecolor{dkgapstat}{HTML}{7039A7}
\definecolor{dkgapinterval}{HTML}{574A82}
\definecolor{dkambient}{HTML}{36516A}
\definecolor{dknorm}{HTML}{475875}
\definecolor{dknormaux}{HTML}{66526E}
\definecolor{dkintertwiner}{HTML}{4A6280}
\definecolor{dkturn}{HTML}{47647D}
\definecolor{dkmix}{HTML}{505B65}

% Coloring primitives; all descendants (including indices and adjoints) can
% be colored as ONE semantic expression by writing \like{key}{expression}.
\newcommand{\dkrole}[2]{\ifDKMathColors\textcolor{dk#1}{#2}\else#2\fi}
% An optional plain-text reminder supplies definitions and types on hover,
% while the final argument stays short enough for the printed notation table.
\newcommand{\dkDeclareSymbol}[5][]{%
  \expandafter\def\csname dkgloss@#2\endcsname{#5}%
  \if\relax\detokenize{#1}\relax
    \expandafter\pdfstringdef\csname dkpdfgloss@#2\endcsname{#5}%
  \else
    \expandafter\pdfstringdef\csname dkpdfgloss@#2\endcsname{\detokenize{#1}}%
  \fi
  \expandafter\def\csname dkglyph@#2\endcsname{#3}%
  \expandafter\def\csname dkrole@#2\endcsname{#4}%
  \expandafter\def\csname dksym@#2\endcsname{\dkrole{#4}{#3}}%
}
\DeclareRobustCommand{\sym}[1]{%
  \ifcsname dksym@#1\endcsname
    \dkSymbolTooltip{#1}{\csname dksym@#1\endcsname}%
  \else
    \PackageError{dk-math-colors}{Unknown symbol key #1}{See math_symbol_colors.tex}%
  \fi
}
\DeclareRobustCommand{\like}[2]{%
  \ifcsname dkrole@#1\endcsname
    \dkSymbolTooltip{#1}{\dkrole{\csname dkrole@#1\endcsname}{#2}}%
  \else
    \PackageError{dk-math-colors}{Unknown role key #1}{See math_symbol_colors.tex}%
  \fi
}

% A short description belongs to the semantic key, not the bare letter.
% The same descriptions are available to the notation appendix.
\DeclareRobustCommand{\dkSymbolDescription}[1]{\csname dkgloss@#1\endcsname}
\newif\ifDKSymbolTooltips
\DKSymbolTooltipstrue
\ifdefined\DKNoSymbolTooltips\DKSymbolTooltipsfalse\fi
\newsavebox{\DKTooltipBox}
\newcount\DKTooltipCount
% Register only shipped widgets through hyperref's existing form-field
% machinery. Poppler/Okular need the catalog's field list, not just page
% annotations. The deferred aux write also discards measuring-only copies.
\makeatletter
\newcommand{\dkRegisterTooltipField}{%
  \expandafter\HyField@@AddToFields\expandafter{\number\pdffeedback lastannot}%
}
\makeatother
\AtBeginDocument{\ifDKSymbolTooltips\Form[NeedAppearances=false]\fi}
\AtEndDocument{\ifDKSymbolTooltips\endForm\fi}
% Native LuaLaTeX PDF widgets: no navigation action or JavaScript. Avoid
% pdfcomment's soulpos dependency, which previously broke local Vim builds.
% Measure in the CURRENT math style, then emit the original expression as-is:
% boxing the visible math would change spacing, scripts and line breaking.
\newcommand{\dkSymbolTooltip}[2]{%
  \ifDKSymbolTooltips
    \begingroup
      \DKSymbolTooltipsfalse % suppress nested annotations while measuring
      \ifmmode
        \edef\DKTooltipStyle{\ifcase\mathstyle
          \displaystyle\or\displaystyle\or\textstyle\or\textstyle\or
          \scriptstyle\or\scriptstyle\or\scriptscriptstyle\or\scriptscriptstyle\fi}%
        \sbox{\DKTooltipBox}{$\DKTooltipStyle#2$}%
      \else
        \leavevmode\sbox{\DKTooltipBox}{#2}%
      \fi
      \global\advance\DKTooltipCount by 1
      \pdfextension annot width\wd\DKTooltipBox
        height\ht\DKTooltipBox depth\dp\DKTooltipBox
        {/Subtype /Widget /FT /Btn /Ff 65536 /F 512
         /T (dk-symbol-\the\DKTooltipCount) /TU (\csname dkpdfgloss@#1\endcsname)
         /H /N /BS << /W 0 >> /Border [0 0 0]}%
      \dkRegisterTooltipField
      #2%
    \endgroup
  \else
    #2%
  \fi
}
\pdfstringdefDisableCommands{%
  \def\sym#1{\csname dkglyph@#1\endcsname}%
  \def\like#1#2{#2}%
}

% Declarations: [optional detailed hover reminder], key, glyph, role, short description.
% This registry includes local proof notation. The notation appendix selects
% the symbols useful for reference; registration does not add a printed row.
% Manim's original symbol keys, now split into semantic subshades.
% Ambient and perturbed operators (provenance is independent of geometry).
\dkDeclareSymbol[Reference self-adjoint operator A on the ambient Hilbert space; the perturbed operator is Ahat = A + H. In a single-operator eigenproblem, A itself is the operator being approximated. An unbounded A acts only on domain(A).]
  {A}{A}{old}{Reference self-adjoint operator before perturbation.}
\dkDeclareSymbol{A1}{A_1}{old}{Complementary block of the reference operator.}
\dkDeclareSymbol{a1}{a_1}{old}{Reference eigenvalue or complementary spectral value.}
\dkDeclareSymbol[Reference block A_0: the restriction of A to the range of the reducing source projection P, expressed in source coordinates. In Theorem 8.1 it lies at or below alpha; in Theorem 8.2 its spectrum lies in the enlarged central interval. The perturbed blocks are denoted Lambda_0 and Lambda_1.]
  {ReferenceBlock0}{A_0}{old}{Reference operator restricted to the source subspace in Section 8.}
\dkDeclareSymbol[Reference block A_1: the restriction of A to the kernel of the reducing source projection P, expressed in complementary coordinates. In Theorem 8.1 it lies at or above alpha + delta in quadratic-form order.]
  {ReferenceBlock1}{A_1}{old}{Reference operator restricted to the source complement in Section 8.}
\dkDeclareSymbol[Self-adjoint ambient operator Ahat being studied; in the perturbation setting Ahat = A + H. Its action on a trial vector is compared with the proposed eigenvalue times that vector.]
  {Ahat}{\widehat A}{current}{Operator under study; in the perturbation setting, the reference operator plus its perturbation.}
\dkDeclareSymbol[Self-adjoint ambient operator Ahat being studied. In the perturbation setting Ahat = A + H. The trial residual is R = Ahat E_0 - E_0 A_0, and the exact complementary block satisfies Ahat F_1 = F_1 Lambda_1 on its domain.]
  {Acurrent}{\widehat A}{current}{Operator under study; in the perturbation setting, the reference operator plus its perturbation.}
\dkDeclareSymbol[Additive self-adjoint perturbation H = Ahat - A. A bound on its norm can control the residual of an eigenspace of one operator relative to the other.]
  {H}{H}{perturb}{Additive self-adjoint perturbation of the reference operator.}
\dkDeclareSymbol{H0}{H_0}{perturb}{Wanted diagonal block of the perturbation.}
\dkDeclareSymbol{H1}{H_1}{perturb}{Unwanted diagonal block of the perturbation.}
\dkDeclareSymbol{B}{B}{perturb}{Off-diagonal perturbation or coupling block.}
\dkDeclareSymbol{b}{b}{perturb}{Scalar perturbation or coupling coefficient.}
% Trial/source family: subspace, frame, compressed operator, projector, vector.
\dkDeclareSymbol{U}{U}{wanted}{Desired exact subspace in the local geometric notation.}% Manim's U remains wanted in local slide notation.
\dkDeclareSymbol{V}{V}{trial}{Trial subspace in the local geometric notation.}% Manim's V remains trial in local slide notation.
\dkDeclareSymbol[Source subspace U to be carried onto the target V by an orthogonal or unitary map W. In the transport sections P projects onto U and Q onto V. The matching theorem also names the source subspace K.]
  {Usource}{U}{trial}{Source subspace to be transported onto the target subspace.}
\dkDeclareSymbol[Trial frame E_0: F -> E, a linear isometry from trial-coordinate space F into ambient space E. Its range is the trial subspace; E_0* returns trial coordinates and E_0* E_0 = I_F. In rank one, E_0(x) = x v.]
  {E0}{E_0}{trialframe}{Isometric embedding of the trial subspace; its adjoint returns trial coordinates.}
\dkDeclareSymbol[Self-adjoint trial operator A_0 on trial-coordinate space F. The residual compares Ahat E_0 with E_0 A_0. Rayleigh--Ritz chooses A_0 = E_0* Ahat E_0 when this compression is defined; the theorem allows other trial operators.]
  {A0}{A_0}{trialspectrum}{Self-adjoint operator on trial coordinates; its eigenvalues are the trial spectral values.}
\dkDeclareSymbol[Trial vector: a proposed approximate eigenvector. In the rank-one geometry v has unit length and E_0(x) = x v. Its residual for trial scalar rho is r = Ahat v - rho v.]
  {v}{v}{trialvector}{Trial vector or vector in the trial subspace.}
\dkDeclareSymbol[Scalar trial operator in rank one, A_0(x) = rho x. Figure 1 uses the Rayleigh value rho = <v, Ahat v> for unit v; then rho v is the projection of Ahat v onto the trial line and r is perpendicular to v.]
  {rho}{\rho}{trialspectrum}{Trial eigenvalue or scalar trial operator.}
\dkDeclareSymbol[Approximate eigenvalue assigned to a trial vector v. For the local operator A, its residual is A v - mu v. For unit v, the Rayleigh--Ritz choice is mu = <A v, v>.]
  {mu}{\mu}{trialspectrum}{Trial spectral value.}
\dkDeclareSymbol{a0}{a_0}{trialspectrum}{Wanted trial eigenvalue or trial spectral value.}
\dkDeclareSymbol[Ambient orthogonal or unitary transporter W with W(U) = V, equivalently W P W* = Q. Its full displacement is I - W on the whole ambient space. The real matching theorem minimizes its trace norm over every such orthogonal W.]
  {Wcompetitor}{W}{trialcompetitor}{Orthogonal or unitary competitor carrying the source subspace onto the target.}
\dkDeclareSymbol[Orthogonal projection P onto the source subspace U in the transport problem; P = P* = P squared. A feasible orthogonal or unitary W satisfies W P W* = Q. Restricted displacement (I - W)P tests only source vectors.]
  {Psource}{P}{trialprojection}{Orthogonal projection onto the source subspace in the transport problem.}
\dkDeclareSymbol{ProjSource}{P_U}{trialprojection}{Orthogonal projection onto the source subspace in the transport problem.}
\dkDeclareSymbol[Vector x in the source subspace U. For principal vectors, unit x_j in U and y_j in V are chosen with <x_i, y_j> = 0 for i different from j and <x_j, y_j> = cos(theta_j).]
  {SourceVector}{x}{trialvector}{Vector in the source subspace.}
\dkDeclareSymbol[Orthonormal coordinate frame X for the source subspace, so X* X = I. Together with a target frame Y, the singular values of X* Y are the principal cosines. In the four-dimensional counterexample X is the vertical stack of I_2 and 0.]
  {Xframe}{X}{trialframe}{Orthonormal coordinate matrix for the source subspace in the counterexample.}
\dkDeclareSymbol[Hilbert space F of trial coordinates. E_0 embeds F isometrically into ambient space E; A_0 and the positive sine operator sin(Theta_0) act on F, while the residual R maps F into E.]
  {Ftrial}{F}{trialframe}{Hilbert space of trial coordinates.}
\dkDeclareSymbol[Ambient Hilbert space E over the real or complex scalars. It contains the trial subspace range(E_0) and the orthogonal exact subspaces range(F_0) and range(F_1). The source theorem permits infinite-dimensional separable E.]
  {Eambient}{E}{ambient}{Ambient Hilbert space containing the trial and exact subspaces.}
% Exact/wanted family: the target subspace, its frame, projection and rotation.
\dkDeclareSymbol[Target subspace V in the transport problem: feasible maps satisfy W(U) = V, or equivalently W P W* = Q. The matching theorem also names this subspace L.]
  {Vtarget}{V}{wanted}{Target subspace in the transport problem.}
\dkDeclareSymbol[Exact wanted frame F_0: F_ex -> E, a linear isometry onto the desired exact spectral subspace. F_0* returns its coordinates. P_ex = F_0 F_0* is its ambient orthogonal projection; in Figure 1, F_0(h) = (h, 0).]
  {F0}{F_0}{wantedframe}{Isometric coordinates for the desired exact spectral subspace; star denotes the adjoint.}
\dkDeclareSymbol[Self-adjoint operator block on coordinates for the desired exact spectral subspace. Its spectrum consists of the wanted exact spectral values; in the corresponding decomposition Ahat F_0 = F_0 Lambda_0 on the operator domain.]
  {Lambda0}{\Lambda_0}{wantedspectrum}{Self-adjoint block on the desired exact spectral subspace.}
\dkDeclareSymbol{u}{u}{wantedvector}{Vector in the desired exact subspace.}
\dkDeclareSymbol{f0}{f_0}{wantedvector}{Component or coordinate vector in the desired exact subspace.}
\dkDeclareSymbol{lambda0}{\lambda_0}{wantedspectrum}{Wanted exact eigenvalue or wanted spectral value.}
\dkDeclareSymbol{lambdau}{\lambda_u}{wantedspectrum}{Eigenvalue associated with a desired exact vector.}
\dkDeclareSymbol[Exact wanted orthogonal projection P_ex: E -> E, with P_ex = F_0 F_0*. It retains the desired exact spectral component. Its complementary projection is P_ex_perp = I - P_ex = F_1 F_1*; the trial leakage is S = P_ex_perp E_0.]
  {Pexact}{P_{\mathrm{ex}}}{wantedprojection}{Orthogonal projection onto the desired exact subspace in the sine-theta theorem.}
\dkDeclareSymbol[Exact complementary orthogonal projection P_ex_perp = I - P_ex = F_1 F_1*: E -> E. It retains the unwanted exact component. In Figure 1, P_ex_perp v = S(1), whereas P_ex_perp r is the unwanted component of the residual.]
  {PexactPerp}{P_{\mathrm{ex}}^\perp}{unwantedprojection}{Orthogonal projection onto the exact complementary subspace.}
\dkDeclareSymbol{ReflectionSigma}{\Sigma}{wantedprojection}{Reflection associated with the desired exact spectral projection.}
\dkDeclareSymbol[Orthogonal projection Q onto the target subspace V in the transport problem. A feasible orthogonal or unitary W satisfies W P W* = Q, where P projects onto the source U.]
  {Qtarget}{Q}{wantedprojection}{Orthogonal projection onto the target subspace in the transport problem.}
\dkDeclareSymbol{ProjTarget}{P_V}{wantedprojection}{Orthogonal projection onto the target subspace in the transport problem.}
\dkDeclareSymbol[Canonical direct rotation D from the source U to the target V for an acute pair. It is the unitary polar factor of T = QP + (I - Q)(I - P), and acts by the principal angle on each principal plane. R is reserved for the trial residual.]
  {Rrotation}{D}{wantedrotation}{Canonical direct rotation carrying the source subspace onto the target.}
\dkDeclareSymbol[Direct rotation D between the ranges of orthogonal projections P and Q, including the nonacute case. Relative to range(P) and kernel(P), its diagonal blocks are positive and its off-diagonal blocks are opposite adjoints. It exists when the two crossed defect spaces have equal dimension, and is unique when both defects vanish.]
  {SourceDirectRotation}{D}{wantedrotation}{Direct rotation, including the nonacute case in the source-result table.}
\dkDeclareSymbol[Orthonormal coordinate frame Y for the target subspace, so Y* Y = I. In the counterexample Y is the vertical stack of cos(theta) I_2 and sin(theta) J_0; hence X* Y = cos(theta) I_2 and both principal angles equal theta.]
  {Yframe}{Y}{wantedframe}{Orthonormal coordinate matrix for the target subspace in the counterexample.}
\dkDeclareSymbol[Vector y in the target subspace V. For principal vectors, unit x_j in U and y_j in V are paired so that <x_j, y_j> = cos(theta_j); distinct pairs have zero cross inner product.]
  {TargetVector}{y}{wantedvector}{Vector in the target subspace.}
\dkDeclareSymbol[Hilbert space F_ex of exact wanted coordinates. The isometry F_0: F_ex -> E identifies it with the desired exact spectral subspace. This space differs from H, the additive perturbation.]
  {Hexact}{F_{\mathrm{ex}}}{wantedframe}{Hilbert space of coordinates for the desired exact spectral subspace.}% Coordinate Hilbert space (not perturbation).
% Exact/unwanted family; the trial error operators are MAGENTA, not cyan.
\dkDeclareSymbol{Uperp}{U^{\perp}}{unwanted}{Orthogonal complement of the desired exact subspace in the local notation.}
\dkDeclareSymbol[Exact complementary frame F_1: G -> E, a linear isometry onto the orthogonal complement of range(F_0). F_1* returns unwanted coordinates; F_1 F_1* = I - P_ex. In Figure 1, F_1(g) = (0, g).]
  {F1}{F_1}{unwantedframe}{Isometric coordinates for the exact complementary subspace; star denotes the adjoint.}
\dkDeclareSymbol[Self-adjoint unwanted block Lambda_1 on G, satisfying Ahat F_1 = F_1 Lambda_1 on its domain. Its spectrum is separated from the trial spectrum of A_0 by the specified positive delta. Figure 1 uses the scalar block 0.7.]
  {Lambda1}{\Lambda_1}{unwantedspectrum}{Self-adjoint operator on exact complementary coordinates: the unwanted spectral block.}
% Spectral measure and set (distinct from the ambient space E and sine block S).
\dkDeclareSymbol[Projection-valued spectral measure E associated with a self-adjoint operator A: A = integral lambda dE(lambda). For a Borel set S, E(S) is the orthogonal projection onto the corresponding spectral subspace.]
  {SpectralMeasure}{E}{wantedprojection}{Projection-valued spectral measure of a self-adjoint operator.}
\dkDeclareSymbol[Borel subset S of the real spectral axis. The spectral projection E(S) selects the part of a self-adjoint operator with spectral values in S; its range is a reducing subspace.]
  {SpectralSet}{S}{wantedspectrum}{Borel spectral set used to select a spectral subspace.}
% Generic Sylvester coefficients and right side, by source/target role.
\dkDeclareSymbol[Left coefficient B of the Sylvester map X -> B X - X C; B acts on the output space of X. For the sine bound, B = Lambda_1 on the unwanted-coordinate space G.]
  {SylvesterLeft}{B}{unwantedspectrum}{Left coefficient of the Sylvester equation; corresponds to the unwanted spectral block.}
\dkDeclareSymbol[Right coefficient C of the Sylvester map X -> B X - X C; C acts on the input space of X. For the sine bound, C = A_0 on trial space F. The historical equation also denotes this coefficient by A.]
  {SylvesterRight}{C}{trialspectrum}{Right coefficient of the Sylvester equation; corresponds to the trial operator.}
\dkDeclareSymbol[Right-hand side Y in B X - X C = Y, with the same input and output spaces as X. For the sine bound, Y = F_1* R: F -> G, the residual expressed in unwanted exact coordinates.]
  {SylvesterRHS}{Y}{resid}{Right-hand side of the Sylvester equation, obtained by projecting the residual.}
\dkDeclareSymbol{w}{w}{unwantedvector}{Vector in an unwanted or complementary spectral subspace.}
\dkDeclareSymbol{f1}{f_1}{unwantedvector}{Component or coordinate vector in the unwanted exact subspace.}
\dkDeclareSymbol{lambda1}{\lambda_1}{unwantedspectrum}{Unwanted exact eigenvalue or complementary spectral value.}
\dkDeclareSymbol{lambdaw}{\lambda_w}{unwantedspectrum}{Eigenvalue associated with an unwanted exact vector.}
\dkDeclareSymbol[Hilbert space G of exact complementary coordinates. F_1: G -> E embeds it onto the unwanted exact subspace; Lambda_1 acts on G and X = F_1* E_0 maps trial coordinates into G.]
  {Gunwanted}{G}{unwantedframe}{Hilbert space of exact complementary coordinates.}
\dkDeclareSymbol{E1}{E_1}{unwantedframe}{Isometric coordinates for the complementary trial subspace.}
\dkDeclareSymbol{Vperp}{V^{\perp}}{unwanted}{Orthogonal complement of the trial subspace in the local notation.}
% Angle and error geometry, including distinct rectangular and positive maps.
\dkDeclareSymbol[Individual principal angle, between 0 and pi/2. For orthonormal subspace frames X and Y, the singular values of X* Y are the principal cosines. In Figure 1, theta is the angle between the unit trial vector v and the exact wanted line.]
  {theta}{\theta}{sine}{Principal angle between the subspaces under comparison.}
\dkDeclareSymbol[Full positive angle operator Theta = diag(Theta_0, Theta_1) for the source/target decompositions. Theta_0 compares the source with the target; Theta_1 compares their orthogonal complements. Its operator norm is the largest spectral angle. For positive c, the conditions Theta <= c and Theta < c mean operator norm at most c and strictly below c, respectively.]
  {Theta}{\Theta}{angleoperator}{Full angle operator, the diagonal sum of the source and complementary angle blocks.}
\dkDeclareSymbol[Directed principal-angle operator Theta_0 on source or trial coordinates. For a trial frame E_0 and exact wanted projection P_ex, cos(Theta_0) = sqrt(E_0* P_ex E_0), and sin(Theta_0) = sqrt(I_F - E_0* P_ex E_0). In the two-projection geometry it comes from the positive cosine compression P Q P restricted to range(P).]
  {Theta0}{\Theta_0}{angleoperator}{Directed principal-angle operator acting on source or trial coordinates.}
\dkDeclareSymbol[Complementary directed angle operator Theta_1, obtained from the projections I-P and I-Q by the same cosine-compression construction as Theta_0. It acts on the source complement, or its isometric coordinates. The full angle operator is diag(Theta_0, Theta_1).]
  {Theta1}{\Theta_1}{angleoperator}{Directed angle operator comparing the orthogonal complements.}
\dkDeclareSymbol[Right-angle source defect D_0 = range(P) intersect kernel(Q). These are source vectors orthogonal to the target. The full-angle conclusion in the Section 8 smallness criterion assumes dim(D_0) = dim(D_1); in Lean their identification is CrossedDefectsEquivalent.]
  {CrossedSourceDefect}{D_0}{fg}{Source directions perpendicular to the target subspace.}
\dkDeclareSymbol[Right-angle target defect D_1 = kernel(P) intersect range(Q). These are target vectors orthogonal to the source. Equality of its Hilbert dimension with that of D_0 is assumed for the full-angle conclusion of Theorem 8.2.]
  {CrossedTargetDefect}{D_1}{fg}{Target directions perpendicular to the source subspace.}
\dkDeclareSymbol[Positive complementary cosine block C_1 = cos(Theta_1). The reference complement is identified with the perturbed complement by the direct rotation D, so A_1 and Lambda_1 act on common coordinates. Theorem 8.1 compares A_1 - alpha I with C_1 (Lambda_1 - alpha I) C_1 as quadratic forms.]
  {ComplementCosineBlock}{C_1}{anglecos}{Cosine of the complementary angle block in the spectral-repulsion inequality.}
\dkDeclareSymbol[Scalar sine of a principal angle. For a unit principal trial vector v and exact wanted projection P_ex, sin(theta) = ||(I - P_ex)v||. In Figure 1 it is the length of the pink leg S(1).]
  {sintheta}{\sin\theta}{sine}{Sine of an individual principal angle: the length of the complementary component of a unit principal trial vector.}
\dkDeclareSymbol[Positive sine operator sin(Theta_0) = sqrt(I_F - E_0* P_ex E_0): F -> F. It equals |S| for S = (I - P_ex)E_0: F -> E. Its eigenvalues in finite dimensions are the principal sines; the sine theorem bounds its ideal norm by N(R)/delta.]
  {sinTheta0}{\sin\Theta_0}{sineoperator}{Positive sine operator on trial coordinates; the modulus of the rectangular complementary error.}
\dkDeclareSymbol[Rectangular complementary error S = (I - P_ex) E_0: F -> E. It records how much each trial vector leaks outside the exact wanted subspace. Its modulus |S| = sqrt(S* S) is sin(Theta_0): F -> F. In rank one S(1) = P_ex_perp v.]
  {S}{S}{sineblock}{Rectangular complementary error: the trial frame projected onto the unwanted exact subspace.}
\dkDeclareSymbol[Complementary-coordinate map X = F_1* E_0: F -> G, with S = F_1 X. It satisfies Lambda_1 X - X A_0 = F_1* R on domain(A_0). Isometric multiplication gives N(X) = N(S); the separate modulus identity abs(S) = sin(Theta_0) gives N(S) = N(sin(Theta_0)), whenever these norms are finite.]
  {X}{X}{sylvester}{Solution of the Sylvester equation; for the sine bound, the trial frame in unwanted exact coordinates.}
\dkDeclareSymbol{alphaAngle}{\alpha}{anglemax}{Largest principal angle.}
\dkDeclareSymbol[Scalar c = cos(theta) in the counterexample. In the general two-projection decomposition, the same role denotes a positive cosine operator C; the cosine and sine blocks commute and satisfy C squared + S squared = I.]
  {CosAngle}{c}{anglecos}{Cosine of a principal angle or the corresponding positive cosine block.}
\dkDeclareSymbol[Scalar s = sin(theta) in the counterexample. In the general two-projection decomposition, the same role denotes a positive sine operator S; it commutes with the positive cosine block C and satisfies C squared + S squared = I.]
  {SinAngle}{s}{sineoperator}{Sine of a principal angle or the corresponding positive sine block.}
\dkDeclareSymbol[Positive small-angle parameter in the local subspace comparison. It controls the allowed angular deviation; it is an angle parameter rather than a perturbation-operator norm.]
  {EpsilonAngle}{\varepsilon}{anglemax}{Angle parameter controlling the maximum subspace deviation.}
\dkDeclareSymbol{PhiAngle}{\varphi}{angleoperator}{Auxiliary angle or angle operator in the two-subspace comparison.}
\dkDeclareSymbol{Bblock}{B}{sineblock}{Off-diagonal block in the two-subspace decomposition.}
\dkDeclareSymbol{Cblock}{C}{sylvester}{Complementary or Sylvester block in the two-subspace decomposition.}
\dkDeclareSymbol{Ablock}{A}{anglecos}{Positive cosine-related block in the two-subspace decomposition.}
% Residual/error-displacement family.
\dkDeclareSymbol[Residual vector r = Ahat v - rho v = R(1) for the rank-one frame E_0(x) = x v. Figure 1 draws it from rho v to Ahat v. For unit v and the Rayleigh value rho = <v, Ahat v>, <v, r> = 0; a general trial scalar need not give orthogonality.]
  {r}{r}{residvector}{Residual vector of an approximate eigenpair.}
\dkDeclareSymbol[Bounded residual R: F -> E, agreeing with Ahat E_0 - E_0 A_0 on domain(A_0). The domain condition requires E_0(domain(A_0)) inside domain(Ahat). R = 0 means that E_0 intertwines A_0 with Ahat; its norm measures the defect.]
  {R}{R}{resid}{Bounded residual of the trial eigenproblem, agreeing with the operator difference on the trial domain.}
\dkDeclareSymbol[Full displacement Delta_W = I - W on the ambient space. In the failed Ky Fan comparison W = W_(pi/4); Delta_W has singular values sqrt(2), sqrt(2), 0, 0. Multiplying by a principal-plane projection restricts its input to that plane.]
  {Kdisplacement}{\Delta_W}{displacement}{Full displacement of a transporter from the identity.}
\dkDeclareSymbol[Squared full displacement Dcal(W) = (I-W)* (I-W), a positive operator. Proposition 4.3 compares its unitarily invariant norm for a direct rotation D with that for every feasible transporter W.]
  {SquaredDisplacement}{\mathcal D}{displacement}{Squared full displacement as a function of the transporter.}
% Matching proof: the gradient and its spectral level are not a space or gap.
\dkDeclareSymbol[Skew-symmetric spectral gradient G of Phi_p. It commutes with W and represents first variations by d/dt Phi_p(W exp(tB)) at t = 0 equal to -tr(G B) for skew-symmetric B. Feasible stationarity gives J_src G J_src = -G.]
  {SpectralGradient}{G}{displacement}{Skew-symmetric spectral gradient of the displacement power sum in the matching proof.}
\dkDeclareSymbol[Positive spectral level g of the skew-symmetric gradient G: on this invariant subspace, -G squared = g squared I. Consequently G/g squares to -I and supplies the complex structure used in the real matching proof.]
  {GradientLevel}{g}{displacement}{Positive gradient level; the gradient squared is minus this level squared times the identity.}
\dkDeclareSymbol[Complex structure I_g = G/g on the positive gradient-level space V_g. It is a real isometry with square -I. Multiplication by a+ib means a times identity plus b times I_g; source reflection J_src reverses I_g and acts as conjugation.]
  {GradientComplexStructure}{\mathcal I_g}{displacement}{Complex multiplication constructed from a positive gradient level.}
\dkDeclareSymbol[Real skew-symmetric matrix B, with B transpose = -B, giving the orthogonal variation W exp(tB). If B commutes with the source projection P, this curve preserves the constraint W(U) = V.]
  {SkewDirection}{B}{ambient}{Skew-symmetric direction for a first variation of an orthogonal transporter.}
% Spectral separation family: gap and original gap are NOT identical.
\dkDeclareSymbol[Positive spectral separation in the stated configuration. In the sine theorem it separates the trial operator from the exact complementary block; Theorem 8.1 uses it for the ordered reference gap, and Theorem 8.2 for the perturbed central-band separation. Figure 1 uses delta = abs(rho - 0.7).]
  {delta}{\delta}{gap}{Positive spectral separation in the stated gap configuration.}
\dkDeclareSymbol{g}{g}{gaporiginal}{Original spectral gap of the reference operator; distinct from the certified sine-theta separation.}
\dkDeclareSymbol{Delta}{\Delta}{gapstat}{Population or statistical spectral gap.}
\dkDeclareSymbol[Spectral endpoint alpha. In the interval/exterior configuration, one spectrum lies between beta and alpha, the other at or below beta - delta or at or above alpha + delta. Theorem 8.2 places the selected perturbed block in the central interval. In Theorem 8.1, alpha bounds the lower reference block and alpha + delta bounds the upper block from below.]
  {alphaEndpoint}{\alpha}{gapinterval}{Upper endpoint of the central band, or lower-block cut in the ordered gap.}
\dkDeclareSymbol[Lower endpoint beta of the finite interval from beta to alpha in the gap hypothesis. One spectrum lies in that interval; the other lies at or below beta - delta or at or above alpha + delta.]
  {betaEndpoint}{\beta}{gapinterval}{Lower endpoint of the finite interval in the interval/exterior gap hypothesis.}
% Spectral/projector and auxiliary matrix objects, context-scoped.
\dkDeclareSymbol{P}{P}{fg}{Orthogonal projection; its source or exact-target role is specified locally.}% Neutral historical Manim key; prefer Pexact/Psource.
\dkDeclareSymbol{Q}{Q}{fg}{Orthogonal projection; its role is specified locally.}% Neutral historical Manim key; prefer Qtarget.
\dkDeclareSymbol[Two-projection intertwiner T = QP + (I - Q)(I - P), satisfying TP = QT and T* T = I - (P - Q) squared. For ||P - Q|| < 1 it is invertible; its unitary polar factor T |T| inverse is the canonical direct rotation.]
  {IntertwinerT}{T}{intertwiner}{Two-projection intertwiner; its unitary polar factor is the acute direct rotation.}
\dkDeclareSymbol[Planar quarter-turn J_0 = ((0, -1), (1, 0)), with J_0 squared = -I_2 and J_0 transpose = -J_0. It defines the rotation C = cos(2 theta) I_2 + sin(2 theta) J_0 used to construct the real competitor.]
  {Jturn}{J_0}{turn}{Planar quarter-turn, with square minus the identity.}
\dkDeclareSymbol[Partial isometry V_pol in the polar decomposition S = V_pol |S|, where |S| = sqrt(S* S). It is isometric on the orthogonal complement of kernel(S) and zero on kernel(S). For an invertible operator its polar factor is unitary.]
  {PolarU}{V_{\mathrm{pol}}}{intertwiner}{Partial isometry in the polar decomposition; unitary when the operator is invertible.}
\dkDeclareSymbol[Planar rotation C = cos(2 theta) I_2 + sin(2 theta) J_0 through twice the principal angle. The competitor W_theta is orthogonally conjugate to C direct-sum I_2, so it rotates one two-plane and fixes the other.]
  {Cturn}{C}{turn}{Planar rotation through twice the principal angle in the counterexample.}
\dkDeclareSymbol[Orthogonal mixing matrix H_mix = (1/sqrt(2)) times the block matrix ((I_2, I_2), (I_2, -I_2)). It constructs W_theta = H_mix (C direct-sum I_2) H_mix transpose.]
  {Hmix}{H_{\mathrm{mix}}}{mix}{Orthogonal mixing matrix in the counterexample.}
\dkDeclareSymbol[Source reflection J_src = 2P - I: it acts as +I on the source subspace and -I on its orthogonal complement, so its square is I. In the matching proof stationarity gives J_src G J_src = -G for the skew-symmetric spectral gradient G.]
  {Jreflect}{J_{\mathrm{src}}}{trialprojection}{Reflection equal to twice the source projection minus the identity.}
\dkDeclareSymbol[Skew angle generator J_ang for a direct rotation D = exp(J_ang Theta). On each principal plane it is the oriented quarter-turn; it vanishes on the common fixed directions. It commutes with the full angle operator Theta.]
  {AngleGenerator}{J_{\mathrm{ang}}}{turn}{Skew angle generator in the direct-rotation result summaries.}
\dkDeclareSymbol{Dblock}{D}{mix}{Auxiliary diagonal or block matrix in the counterexample.}
\dkDeclareSymbol[Orthogonal projection onto the first principal plane span(e_0, e_3) in the four-dimensional counterexample. For Delta_W = I - W_(pi/4), the Ky Fan 2-norm of Delta_W Omega_1 equals 2.]
  {OmegaOne}{\Omega_1}{turn}{First principal-plane projection in the failed Ky Fan comparison.}
\dkDeclareSymbol[Orthogonal projection onto the second principal plane span(e_1, e_2) in the four-dimensional counterexample. For Delta_W = I - W_(pi/4), the Ky Fan 2-norm of Delta_W Omega_2 equals 2, although the trace norm of Delta_W is only 2 sqrt(2).]
  {OmegaTwo}{\Omega_2}{intertwiner}{Second principal-plane projection in the failed Ky Fan comparison.}
\dkDeclareSymbol{Dplanar}{D}{wantedrotation}{Planar direct-rotation block.}
\dkDeclareSymbol{Hreflection}{H}{trialcompetitor}{Planar reflection carrying the source line to the target.}
% Unitarily invariant gauge family; N, M are NOT a norm versus its input.
\dkDeclareSymbol[Unitarily invariant norm N, or a symmetric operator-ideal norm on its domain of finiteness. It satisfies N(U T V) = N(T) for unitaries and N(B T) <= ||B||op N(T). In the sine theorem it is normalized by N(u v*) = ||u|| ||v|| and has Ky Fan dominance where both norms are finite. Both R and sin(Theta_0) must belong to its ideal.]
  {NormN}{N}{norm}{Unitarily invariant norm on its domain of finiteness; in the sine theorem, normalized by $N(uv^*)=\|u\|\|v\|$.}
\dkDeclareSymbol[Sine-theta separation G_delta(B,C): finite interval/exterior separation in either spectral order, or ordered separation B <= a I and C >= (a+delta) I, or the reverse. The ordered alternatives allow both spectra to be unbounded.]
  {SineGap}{\mathcal G_\delta}{gap}{Interval/exterior separation or either ordered separation, as in the sine-theta theorem.}
\dkDeclareSymbol[Auxiliary unitarily invariant norm M witnessing the Q-norm property N(T) squared = M(T* T). For the Frobenius norm N, one can take M to be the trace norm.]
  {NormM}{M}{normaux}{Auxiliary unitarily invariant norm used to define a Q-norm through the squared modulus.}
\dkDeclareSymbol[Q-norm property: a unitarily invariant norm N is a Q-norm when N(T) squared = M(T* T) for some unitarily invariant norm M. The Frobenius norm is an example, with M the trace norm; the repaired full-displacement theorem uses this norm class.]
  {NormQ}{Q}{normaux}{Q-norm: its square is a unitarily invariant norm of the squared modulus.}
\dkDeclareSymbol[Singleton trace cost s_i = min(4 sin(theta_i/2), 2) for angle label i. It transports one real principal plane independently, choosing a rotation of cost 4 sin(theta_i/2) or a reflection of cost 2.]
  {SingletonCost}{s}{norm}{Minimum trace displacement for one real principal plane: choose a rotation or a reflection.}
\dkDeclareSymbol[Paired trace cost e_ij = 4 sin((theta_i + theta_j)/2) for distinct angle labels i,j. It is attained by coupling their principal planes in a real four-dimensional block. The edge weight s_i + s_j - e_ij is the difference between the cost of two singletons and this paired cost.]
  {PairCost}{e}{norm}{Trace displacement cost of coupling two principal planes in a four-dimensional block.}
\dkDeclareSymbol[Scalar phase cost f_p(t) = (2 sin(t/2)) to the power p, for 0 <= t <= pi. A real planar rotation through t contributes two copies, so its displacement power cost is 2 f_p(t).]
  {MatchingPhaseCost}{f_p}{norm}{Displacement power cost for one scalar phase.}
\dkDeclareSymbol[Singleton power cost a_i(p) = min(2 f_p(theta_i), 2 to the power p). These are the rotation and reflection choices for one principal plane. At p=1 it equals the trace-cost weight s_i.]
  {SingletonPowerCost}{a}{norm}{Minimum displacement power cost for one angle.}
\dkDeclareSymbol[Paired power cost b_ij(p) = min over d <= tau <= h of 2(f_p(h+tau) + f_p(h-tau)), where h=(theta_i+theta_j)/2 and d=abs(theta_i-theta_j)/2. This is the positive-determinant pair minimum; unrestricted block optimization also compares two singletons. At p=1 it equals e_ij.]
  {PairPowerCost}{b}{norm}{Positive-determinant displacement power cost for a pair of angles.}
\dkDeclareSymbol[Matching M on the finite general graph of positive principal-angle labels. Each label belongs to at most one edge; unmatched labels use singleton transporters. Minimizing paired and singleton costs is equivalent to maximum-weight matching with edge weights g_ij = s_i + s_j - e_ij. The empty matching is allowed.]
  {Matching}{M}{normaux}{Matching of positive principal-angle labels: each label is used in at most one pair.}
\dkDeclareSymbol[Displacement power sum Phi_p(W) = sum_j s_j(I - W)^p. The matching proof first minimizes it for 1 < p < 2, then lets p decrease to 1, where it becomes the trace norm ||I - W||_1.]
  {PowerCost}{\Phi}{norm}{Sum of the p-th powers of displacement singular values.}
% Auxiliary matching-proof notation: neutral ink, with the same hover registry.
\dkDeclareSymbol[Off-diagonal gradient block G_01 mapping the source complement K perpendicular into K. Stationarity makes the diagonal blocks of G zero, so G = ((0,G_01),(-G_01 transpose,0)) in the source/complement decomposition.]
  {GradientOffDiagonal}{G_{01}}{fg}{Complement-to-source block of the skew-symmetric gradient.}
\dkDeclareSymbol[Orthogonal group O(n): real n by n matrices W satisfying W transpose times W = I. Its elements preserve Euclidean inner products and lengths. The notation O(2) refers to the same group in dimension two.]
  {OrthogonalGroup}{O(n)}{fg}{Group of real orthogonal matrices.}
\dkDeclareSymbol[Feasible transporters F = {W in O(n) : W(K) = L}. Equivalently W P_K W transpose = P_L. This nonempty compact set is the optimization domain throughout the matching proof.]
  {FeasibleTransport}{\mathcal F}{fg}{Orthogonal maps carrying the source subspace onto the target.}
\dkDeclareSymbol[Singular value s_j(T), the nonnegative square root of an eigenvalue of T transpose times T, counted with multiplicity. Summing s_j(I-W) gives the trace displacement; summing its p-th powers gives Phi_p(W). These singular values differ from the singleton weights s_i.]
  {TransportSingularValues}{s_j}{fg}{A displacement singular value in the matching proof.}
\dkDeclareSymbol[Scalar singular-value magnitude s between zero and two for the displacement I-W of an orthogonal matrix. Uniform convergence of s to the power p to s on this interval permits the trace-norm limit. This variable differs from the rotation parameter s used earlier in the proof.]
  {DisplacementSize}{s}{fg}{Scalar displacement singular-value magnitude in the uniform limit.}
\dkDeclareSymbol[All-singleton cost C_0 = sum_i s_i, before any two angle labels are paired. A matching has cost C_0 minus the sum of its selected edge weights.]
  {SingletonBaseline}{C_0}{fg}{Trace displacement when every angle uses its singleton transporter.}
\dkDeclareSymbol[Matching power cost C_p(M): sum of b_ij(p) over selected pairs plus a_i(p) over unmatched labels. Its minimum equals the minimum of Phi_p over feasible transporters. At p=1 it is C(M).]
  {MatchingPowerCost}{C_p}{fg}{Total paired and singleton power cost of a matching.}
\dkDeclareSymbol[Invariant orthogonal block V_a preserved by the source projection, target projection, and minimizing transporter. The displacement cost adds over these blocks, and their principal-angle lists combine into the original list.]
  {MatchingBlock}{V_a}{fg}{Simultaneous invariant block in the matching proof.}
\dkDeclareSymbol[Gradient-level space V_g: the eigenspace of -G squared for a positive eigenvalue g squared. It is preserved by G, W, and the source reflection J_src. Multiplication by i on it is G/g.]
  {GradientLevelSpace}{V_g}{fg}{Positive gradient-level eigenspace.}
\dkDeclareSymbol[Eigenphase varphi of W on V_g viewed as a complex space with multiplication by i = G/g. It lies between zero and pi and satisfies f_p prime(varphi) = g. There are at most two such phases.]
  {GradientEigenphase}{\varphi}{fg}{One of the at most two phases at a positive gradient level.}
\dkDeclareSymbol[Orthogonal projection E onto the first phase space of W on V_g. Here W = exp(i varphi_1) E + exp(i varphi_2) (I-E). In coordinates where the source reflection is conjugation, E = X + iY.]
  {GradientPhaseProjection}{E}{fg}{Projection onto the first phase space in a gradient-level space.}
\dkDeclareSymbol[Real part X of the phase projection E = X + iY. It is a real symmetric matrix with spectrum in the closed interval from zero to one, satisfying X squared - Y squared = X and XY + YX = Y. This X differs from the source-coordinate and Sylvester matrices elsewhere in the paper.]
  {ProjectionRealPart}{X}{fg}{Real symmetric part of the phase projection.}
\dkDeclareSymbol[Imaginary part Y of the phase projection E = X + iY. It is a real skew-symmetric matrix. It maps the lambda eigenspace of X to the 1-lambda eigenspace and constructs the small invariant planes.]
  {ProjectionImaginaryPart}{Y}{fg}{Real skew-symmetric imaginary part of the phase projection.}
\dkDeclareSymbol[Eigenvalue lambda in the closed interval from zero to one of the real part X of the phase projection. It is local to the block construction and differs from the original operator's spectral parameter.]
  {ProjectionEigenvalue}{\lambda}{fg}{Eigenvalue of the real part of the phase projection.}
\dkDeclareSymbol[Plane coupling eta = sqrt(lambda (1-lambda)) for 0 < lambda < 1. If u is a unit lambda eigenvector of X, then v = Yu/eta completes its invariant real plane.]
  {ProjectionCoupling}{\eta}{fg}{Normalization used to construct an invariant real plane.}
\dkDeclareSymbol[Pair mean angle h = (theta_i + theta_j)/2. Together with d = (theta_i - theta_j)/2 after ordering theta_i >= theta_j, it parametrizes the two-angle block cost.]
  {PairMeanAngle}{h}{fg}{Mean of the two ordered principal angles.}
\dkDeclareSymbol[Pair half-difference d = (theta_i - theta_j)/2 after ordering theta_i >= theta_j. It satisfies 0 <= d <= h and h+d <= pi/2.]
  {PairHalfDifference}{d}{fg}{Half the difference of the two ordered principal angles.}
\dkDeclareSymbol[Scalar pair parameter tau in the closed interval from d to h. The paired power cost is the minimum of 2(f_p(h+tau) + f_p(h-tau)) over this interval. At p=1 a minimum is attained at tau=h.]
  {PairPhaseParameter}{\tau}{fg}{Scalar parameter in the two-angle power-cost minimum.}
\dkDeclareSymbol[Attaining rotation angle gamma = arccos(cos(tau)/cos(d)). Choosing U_0 = U_1 = R_gamma realizes phase parameters x=h and y=tau in the four-dimensional block.]
  {PairAttainingAngle}{\gamma}{fg}{Rotation angle realizing a chosen scalar pair parameter.}
\dkDeclareSymbol[Four-dimensional direct rotation D = ((C,-S),(S,C)) in source and complementary principal coordinates, with C the diagonal principal cosines and S the diagonal principal sines. Every transporter on this block is D diag(U_0,U_1).]
  {PairDirectRotation}{D}{fg}{Direct rotation on the two principal planes.}
\dkDeclareSymbol[Diagonal cosine block C = diag(cos(theta_i), cos(theta_j)) in the four-dimensional direct rotation D. This matrix differs from the scalar displacement cost C(M).]
  {PairCosineBlock}{C}{fg}{Diagonal principal-cosine block.}
\dkDeclareSymbol[Diagonal sine block S = diag(sin(theta_i), sin(theta_j)) in the four-dimensional direct rotation D. This matrix is local to the matching proof.]
  {PairSineBlock}{S}{fg}{Diagonal principal-sine block.}
\dkDeclareSymbol[Free planar orthogonal matrices U_0 and U_1 acting on the two-dimensional source and its two-dimensional complement, respectively. Multiplying D by diag(U_0,U_1) parametrizes every transporter on this four-dimensional block.]
  {PairFreeOrthogonal}{U}{fg}{Free source or complementary planar orthogonal factor.}
\dkDeclareSymbol[Phase parameter x in the closed interval from h to pi-h for a positive-determinant four-dimensional transporter. Its eigenvalues are exp(plus or minus i(x+y)) and exp(plus or minus i(x-y)).]
  {PairPhaseX}{x}{fg}{First phase parameter of a four-dimensional transporter.}
\dkDeclareSymbol[Phase parameter y in the closed interval from d to pi-d for a positive-determinant four-dimensional transporter. Together with x it gives rotational phases x+y and x-y, folded into the closed interval from zero to pi for the displacement cost.]
  {PairPhaseY}{y}{fg}{Second phase parameter of a four-dimensional transporter.}
\dkDeclareSymbol[Scalar a = cos(x) in the four-dimensional phase calculation. Together with b = cos(y), it expresses the trace as 4ab and determines the reciprocal characteristic polynomial.]
  {PairPhaseCosineX}{a}{fg}{Cosine of the first phase parameter in the trace calculation.}
\dkDeclareSymbol[Scalar b = cos(y) in the four-dimensional phase calculation, paired with a = cos(x). These scalars are local to the characteristic-polynomial calculation.]
  {PairPhaseCosineY}{b}{fg}{Cosine of the second phase parameter in the trace calculation.}
\dkDeclareSymbol[Rotation parameter r: the two free orthogonal factors use planar rotation angles r+s and r-s. It appears in cos(y) = cos(d) cos(r) for two rotations.]
  {PairRotationR}{r}{fg}{Half the sum of the two free planar rotation angles.}
\dkDeclareSymbol[Rotation parameter s: the two free orthogonal factors use planar rotation angles r+s and r-s. It appears in cos(x) = cos(h) cos(s) for two rotations. This scalar differs from the singular values and singleton costs.]
  {PairRotationS}{s}{fg}{Half the difference of the two free planar rotation angles.}
\dkDeclareSymbol[Planar reflection J_2 = diag(1,-1). Every two-dimensional orthogonal matrix of determinant -1 can be written as a planar rotation times J_2.]
  {PairPlanarReflection}{J_2}{fg}{Fixed planar reflection used to parametrize determinant -1 factors.}
\dkDeclareSymbol[Real planar rotation R_t = ((cos(t),-sin(t)),(sin(t),cos(t))). Its two displacement singular values are both 2 times abs(sin(t/2)). In the matching proof a rotation through t in the closed interval from zero to pi contributes 2 f_p(t).]
  {MatchingPlanarRotation}{R_t}{fg}{Planar rotation through the indicated angle.}
% Statistical application / operator provenance.
\dkDeclareSymbol{Sigma}{\Sigma}{old}{Population covariance matrix.}
\dkDeclareSymbol{SigmaHat}{\widehat\Sigma}{current}{Estimated or perturbed covariance matrix.}
\dkDeclareSymbol{YWSV}{V}{wanted}{Orthonormal coordinates for the population eigenspace in the statistical application.}
\dkDeclareSymbol{YWSVhat}{\widehat V}{trialframe}{Estimated eigenspace frame in the statistical application.}

% Figure-specific objects retain their own meanings when bare letters overlap.
\dkDeclareSymbol[Signed rotation angle phi of the displayed map on its invariant real two-plane. At theta = pi/4 the direct rotation turns its two principal planes by opposite signed angles of magnitude pi/4; the competitor turns its moving plane by pi/2 and fixes its other plane.]
  {PlaneRotationAngle}{\varphi}{sine}{Signed rotation angle of the orthogonal map on the displayed invariant plane.}
\dkDeclareSymbol[Unit vector q in the displayed invariant plane. The map T is the direct rotation D or the competitor W in that panel; Tq is the image and (I - T)q = q - Tq is the displacement.]
  {FigureInput}{q}{ambient}{Unit input to the orthogonal map shown in this invariant-plane panel.}
\dkDeclareSymbol[Image Tq of the panel unit input q under T = D or W. Both q and Tq lie on the unit circle; the green displacement arrow joins Tq to q, so it represents q - Tq.]
  {FigureImage}{Tq}{wantedvector}{Image of the unit input under the orthogonal map shown in this panel.}
\dkDeclareSymbol[Displacement (I - T)q = q - Tq, drawn from the image Tq to the unit input q. On a plane rotated through phi its length is 2 |sin(phi/2)|; it vanishes on a fixed plane.]
  {DisplacementVector}{(I-T)q}{displacement}{Difference between a unit input and its image under the displayed orthogonal map.}
\dkDeclareSymbol[First principal plane Pi_1 = span(e_0, e_3) for the source-target pair in the real four-dimensional counterexample. At theta = pi/4 the direct rotation acts on this plane by a quarter of pi, with displacement singular value sqrt(2 - sqrt(2)).]
  {CounterPrincipalOne}{\Pi_1}{fg}{First principal plane in the real counterexample: the span of e0 and e3.}
\dkDeclareSymbol[Second principal plane Pi_2 = span(e_1, e_2) for the source-target pair in the real four-dimensional counterexample. The direct rotation acts by the same angle magnitude as on Pi_1, with the opposite sign in the displayed plane coordinates.]
  {CounterPrincipalTwo}{\Pi_2}{fg}{Second principal plane in the real counterexample: the span of e1 and e2.}
\dkDeclareSymbol[Moving plane M = range(W - I) = span(e_0 + e_2, e_1 + e_3) for the competitor. W rotates it by 2 theta; at theta = pi/4 the two displacement singular values contributed by this plane are both sqrt(2).]
  {CounterMovingPlane}{M}{ambient}{Moving plane of the competing transporter: the range of W minus the identity.}
\dkDeclareSymbol[Fixed plane Z = kernel(W - I) = span(e_0 - e_2, e_1 - e_3) for the competitor. W acts as the identity here, so this plane contributes two zero displacement singular values.]
  {CounterFixedPlane}{Z}{ambient}{Fixed plane of the competing transporter: the kernel of W minus the identity.}
\dkDeclareSymbol[Direct-rotation displacement singular value a = 2 sin(pi/8) = sqrt(2 - sqrt(2)) at theta = pi/4. Each of its two principal planes contributes two copies, giving trace displacement 4a.]
  {CounterChordA}{a}{displacement}{Direct-rotation displacement singular value in the counterexample: sqrt(2 - sqrt(2)).}
\dkDeclareSymbol[Competitor displacement singular value b = 2 sin(pi/4) = sqrt(2) at theta = pi/4. Its moving plane contributes two copies; the fixed plane contributes two zeros, giving trace displacement 2b.]
  {CounterChordB}{b}{displacement}{Competing-transporter displacement singular value in the counterexample: square root of 2.}
\dkDeclareSymbol[Singular value s_j(I - T) of the full ambient displacement, ordered from largest to smallest. Each rotating real two-plane contributes two values 2 |sin(phi/2)|; fixed directions contribute zeros. Their sum is the trace cost.]
  {DisplacementSingular}{s_j}{displacement}{An individual singular value of the full displacement from the identity.}
\dkDeclareSymbol[Edge weight g_ij = s_i + s_j - e_ij, the difference between the cost of two singletons and the paired cost. A matching has total cost C_0 minus the sum of its edge weights, with C_0 = sum_i s_i. Maximizing this sum gives the minimum transport cost; edges of negative weight may be omitted.]
  {MatchingEdgeWeight}{g_{ij}}{current}{Edge weight: the cost of two singletons minus their paired cost.}
\dkDeclareSymbol[Total trace-norm displacement cost. For a transporter W, C(W) = ||I - W||_1, the sum of its displacement singular values. For a matching M, C(M) is the sum of its paired and unmatched singleton costs.]
  {TraceCost}{C}{norm}{Trace-norm displacement cost of the indicated orthogonal transporter or matching.}

% Spectral values have their own key, distinct from an operator block.
\dkDeclareSymbol[Spectral parameter lambda of a self-adjoint operator. In a finite-dimensional eigenvalue list lambda_j(A) it denotes an individual eigenvalue; in A = integral lambda dE(lambda), it is the real integration variable.]{SpectralValue}{\lambda}{wantedspectrum}{Real spectral parameter or individual eigenvalue of a self-adjoint operator.}

\setlength{\emergencystretch}{3em}
\numberwithin{equation}{section}
\renewcommand{\thesubsection}{\thesection(\alph{subsection})}

% Short wrappers use the paper's symbol registry for glyphs and hover definitions.
\newcommand{\W}{\sym{Wcompetitor}}
\newcommand{\Psrc}{\sym{Psource}}
\newcommand{\Ptgt}{\sym{Qtarget}}
\newcommand{\Ggrad}{\sym{SpectralGradient}}
\newcommand{\Jsrc}{\sym{Jreflect}}
\newcommand{\Fset}{\sym{FeasibleTransport}}
\newcommand{\Phicost}{\like{PowerCost}{\Phi_p}}
\newcommand{\fphase}{\sym{MatchingPhaseCost}}
\newcommand{\Vlevel}{\sym{GradientLevelSpace}}
\newcommand{\glev}{\sym{GradientLevel}}
\newcommand{\Icomp}{\sym{GradientComplexStructure}}
\newcommand{\Ephase}{\sym{GradientPhaseProjection}}
\newcommand{\Xreal}{\sym{ProjectionRealPart}}
\newcommand{\Yimag}{\sym{ProjectionImaginaryPart}}
\newcommand{\lambdablock}{\sym{ProjectionEigenvalue}}
\newcommand{\etablock}{\sym{ProjectionCoupling}}
\newcommand{\hmean}{\sym{PairMeanAngle}}
\newcommand{\dhalf}{\sym{PairHalfDifference}}
\newcommand{\taupair}{\sym{PairPhaseParameter}}
\newcommand{\phaseX}{\sym{PairPhaseX}}
\newcommand{\phaseY}{\sym{PairPhaseY}}
\newcommand{\Dpair}{\sym{PairDirectRotation}}
\newcommand{\Cpair}{\sym{PairCosineBlock}}
\newcommand{\Spair}{\sym{PairSineBlock}}
\newcommand{\Ufree}[1]{\like{PairFreeOrthogonal}{U_{#1}}}
\newcommand{\Rplane}[1]{\like{MatchingPlanarRotation}{R_{#1}}}
\newcommand{\single}[1]{\like{SingletonPowerCost}{a_{#1}(p)}}
\newcommand{\paired}[2]{\like{PairPowerCost}{b_{#1#2}(p)}}
\newcommand{\why}{\par\smallskip\noindent\textbf{Justification.} }
\newcommand{\rehearse}[1]{\par\smallskip\noindent\textit{Explain aloud:} #1\par}

\title{Theorem 5: a worked proof of the real matching formula}
\author{}
\date{}
\begin{document}
\maketitle

This note follows the six steps and lettered inferences in the manuscript.
Each inference states what has been established; its justification gives
the calculation or theorem that permits it. The longer calculations are
included here so the argument can be studied offline.

The theorem concerns equal-dimensional subspaces
$\like{Usource}{K},\like{Vtarget}{L}\subseteq\mathbb R^n$.
An orthogonal \emph{transporter} is a matrix $\W$ with
$\W(\like{Usource}{K})=\like{Vtarget}{L}$.
The objective is its full ambient trace displacement $\|I-\W\|_1$.
For positive principal angles $\like{theta}{\theta_1},\ldots,
\like{theta}{\theta_m}$, put
\[
 \like{SingletonCost}{s_i}=\min\{4\sin(\like{theta}{\theta_i}/2),2\},
 \qquad
 \like{PairCost}{e_{ij}}
   =4\sin((\like{theta}{\theta_i}+\like{theta}{\theta_j})/2).
\]
The theorem states that the minimum displacement is the minimum, over
matchings $\sym{Matching}$ of the angle labels, of
\[
 \like{TraceCost}{C(M)}
 =\sum_{\{i,j\}\in\sym{Matching}}\like{PairCost}{e_{ij}}
   +\sum_{i\text{ unmatched}}\like{SingletonCost}{s_i}.
\]

There are two bounds to prove. Constructing a transporter from a matching
proves that the unrestricted minimum is \emph{at most} the matching minimum.
To prove the reverse bound, we must control a minimizer that initially may
mix any number of principal planes. Steps 2--3 supply that control for a
differentiable auxiliary cost. Step 6 passes its minimum value to the
trace-norm limit.

\section{Choose a differentiable cost and a global minimizer}

\subsection{The cost attains a global minimum}
\textbf{Inference.} For every $1<p<2$, the feasible set and cost
\[
 \Fset=\{\W\in\sym{OrthogonalGroup}:
              \W(\like{Usource}{K})=\like{Vtarget}{L}\},
 \qquad
 \Phicost(\W)=\sum_{j=1}^n
       \like{TransportSingularValues}{s_j}(I-\W)^p
\]
have a global minimizer. We fix one and call it $\W$.

\why Equal dimensions let us choose orthonormal bases of the two subspaces,
extend them to ambient orthonormal bases, and map the first full basis onto
the second. Thus $\Fset$ is nonempty. Write $\Psrc=P_K$ and $\Ptgt=P_L$
for the source and target projections. For orthogonal $\W$,
\[
 \W(\like{Usource}{K})=\like{Vtarget}{L}
 \quad\Longleftrightarrow\quad
 \W\Psrc\W^T=\Ptgt.
\]
Indeed, the matrix on the left of the projection equation is the orthogonal
projection onto $\W(\like{Usource}{K})$.
The projection equation is closed under limits. The orthogonal group is
closed because $\W^T\W=I$ is a closed equation, and bounded because every
column has length one. Hence $\Fset$ is compact by Heine--Borel.
Singular values vary continuously with a matrix, so $\Phicost$ is continuous.
The extreme value theorem now gives a minimizer.

The reason for choosing $p>1$ is differentiability, including where a
displacement singular value is zero. The trace norm at $p=1$ lacks this
property. Appendix~\ref{sec:study-differentiability} derives the matrix
derivative used here. We select a minimizer of $\Phicost$ for each $p$;
we do not assume the trace-norm minimizer also minimizes this other cost.
The upper bound $p<2$ selects the auxiliary range already formalized in
Lean. Passing to exponent one only requires such a range immediately
above one; the particular upper endpoint is not used in the limit.

\subsection{Invariant blocks split the cost and angle data}
\textbf{Inference.} If an orthogonal decomposition
$\mathbb R^n=\bigoplus_a\sym{MatchingBlock}$ is preserved by
$\Psrc,\Ptgt,\W$, its block costs add and its block principal-angle lists
combine into the original principal-angle list.

\why In an adapted basis these matrices are block diagonal. Thus
$I-\W$ is block diagonal, and so is $(I-\W)^T(I-\W)$.
The eigenvalue list of this last matrix is the union, with multiplicity,
of the block eigenvalue lists. Taking nonnegative square roots and summing
$p$th powers gives
\[
 \Phicost(\W)=\sum_a\Phicost(\W|_{\sym{MatchingBlock}}).
\]
Likewise, $\Psrc\Ptgt\Psrc$ is block diagonal. Its restriction to the source
has eigenvalues $\cos^2\theta_j$, so the block angle lists combine with
multiplicity. Restricting to the source is necessary: the ambient
compression also has zeros from source-complement directions.

For comparison, the principal planes of the two subspaces already exist
before we study $\W$. Choose principal unit vectors $u_i,v_i$ with
$\langle u_i,v_j\rangle=\delta_{ij}\cos\theta_i$ and, when $\theta_i>0$, set
\[
 q_i=\frac{v_i-\cos\theta_i\,u_i}{\sin\theta_i}.
\]
The principal-vector inner products give
$\langle u_i,q_j\rangle=0$ and $\langle q_i,q_j\rangle=\delta_{ij}$.
Hence $\Pi_i=\operatorname{span}(u_i,q_i)$ are mutually orthogonal planes.
These planes describe the pair of subspaces. They need not individually
be preserved by the minimizing transporter. Steps 2--3 find a simultaneous
decomposition that also preserves that transporter.
\rehearse{Why does a decomposition of the subspaces alone fail to prove
that arbitrary mixing cannot improve the minimum?}

\section{Differentiate along feasible rotations}

\subsection{Source-preserving rotations give feasible curves}
\textbf{Inference.} If $B^T=-B$ and $B\Psrc=\Psrc B$, then
$\W\exp(tB)\in\Fset$ for every real $t$.

\why Transposition respects the exponential power series, so
\[
 \exp(tB)^T=\exp(-tB)=\exp(tB)^{-1}.
\]
Thus the exponential is orthogonal. Every power of $B$ commutes with
$\Psrc$, so its exponential does too. For a source vector $v$,
\[
 \Psrc\exp(tB)v=\exp(tB)\Psrc v=\exp(tB)v.
\]
Therefore $\exp(tB)$ maps the source into itself. Its inverse $\exp(-tB)$
does the same, giving equality of the image and source. Composing with
$\W$ still maps the source onto the target.

Since the selected $\W$ minimizes the cost, the differentiable scalar
function $t\mapsto\Phicost(\W\exp(tB))$ has a minimum at zero. Its
derivative there is zero. This necessary condition is stationarity.

\subsection{The first variation is a trace pairing}
\textbf{Inference.} There is a real skew-symmetric matrix $\Ggrad$ such that,
for every skew-symmetric $B$, including those that mix the rotation planes,
\begin{equation}\label{eq:study-first-variation}
 \left.\frac{d}{dt}\Phicost(\W\exp(tB))\right|_{t=0}
     =-\operatorname{tr}(\Ggrad B).
\end{equation}

\why First define the scalar cost and planar quarter-turn
\[
 \fphase(\theta)=(2\sin(\theta/2))^p,
 \quad 0\le\theta\le\pi,
 \qquad J_0=\begin{pmatrix}0&-1\\1&0\end{pmatrix}.
\]
A rotation $\Rplane{\theta}=\cos\theta\,I+\sin\theta\,J_0$ satisfies
\[
 (I-\Rplane{\theta})^T(I-\Rplane{\theta})
       =4\sin^2(\theta/2)I.
\]
Its two displacement singular values agree, so its cost is $2\fphase(\theta)$.
Every real orthogonal matrix admits an orthogonal decomposition into
rotation planes with $0<\theta<\pi$ and scalar $+1,-1$ lines.
Define $\Ggrad=\fphase'(\theta)J_0$ on each such plane, and zero on the
scalar lines. These blocks are skew-symmetric and commute with the
corresponding blocks of $\W$; hence $\Ggrad^T=-\Ggrad$ and
$\Ggrad\W=\W\Ggrad$.

We must justify the formula when $B$ mixes these blocks. Let
\[
 \mathcal L(B)=
   \left.\frac{d}{dt}\Phicost(\W\exp(tB))\right|_{t=0}.
\]
The derivative of $\W\exp(tB)$ at zero is $\W B$.
Differentiability of the cost therefore makes $\mathcal L$ a linear
functional of $B$. Differentiability is the hypothesis that permits the
averaging argument below; continuity alone would not suffice.

Let $S_\varepsilon$ independently multiply each rotation plane or scalar
line by a sign $\varepsilon_a\in\{-1,+1\}$. It is orthogonal and commutes
with $\W$. Consequently,
\[
 \W\exp(tS_\varepsilon B S_\varepsilon)
     =S_\varepsilon\W\exp(tB)S_\varepsilon.
\]
Orthogonal conjugation preserves the displacement singular values, so
$\mathcal L(S_\varepsilon B S_\varepsilon)=\mathcal L(B)$.
If there are $k$ blocks, the average is
\[
 2^{-k}\sum_\varepsilon S_\varepsilon B S_\varepsilon
      =\operatorname{diag}(B_{11},\ldots,B_{kk}).
\]
An off-diagonal block $B_{ab}$ acquires sign $\varepsilon_a\varepsilon_b$;
half the choices give each sign. A diagonal block acquires sign one.
By linearity, the derivative at $B$ equals the derivative at this block
diagonal average. The averaged curve is used to calculate a derivative;
it need not preserve the source subspace.

On a rotation plane, a real $2\times2$ skew-symmetric matrix is $bJ_0$.
The matrices $\Rplane{\theta}$ and $\exp(tbJ_0)$ multiply to
$\Rplane{\theta+tb}$, so the derivative of the block cost is
$2b\fphase'(\theta)$. Meanwhile,
\[
 -\operatorname{tr}(\fphase'(\theta)J_0\,bJ_0)
   =-b\fphase'(\theta)\operatorname{tr}(-I_2)
   =2b\fphase'(\theta).
\]
On a scalar line the diagonal block of a skew-symmetric matrix is zero.
Adding the block derivatives proves \cref{eq:study-first-variation}.
Since $\Ggrad$ is block diagonal, $\operatorname{tr}(\Ggrad B)$ also
depends only on the diagonal blocks of $B$.
\rehearse{What permits averaging the direction? Why may the averaging use
directions that do not preserve the source? Where does the factor two arise?}

\subsection{Stationarity forces the gradient to be off-diagonal}
\textbf{Inference.} Relative to source and source complement,
\[
 \Ggrad=\begin{pmatrix}0&\sym{GradientOffDiagonal}\\
                   -\like{GradientOffDiagonal}{G_{01}^T}&0\end{pmatrix},
 \qquad \Jsrc=2\Psrc-I,
 \qquad \Jsrc\Ggrad\Jsrc=-\Ggrad.
\]

\why Choose $B=\Psrc\Ggrad\Psrc$. Transposition gives $B^T=-B$ and
$\Psrc^2=\Psrc$ gives $B\Psrc=\Psrc B=B$, so Step 2(a) makes the
variation feasible. Step 2(b) and stationarity give
$0=-\operatorname{tr}(\Ggrad\Psrc\Ggrad\Psrc)$.
Cyclicity of trace and $\Psrc^2=\Psrc$ identify this with
$-\operatorname{tr}(B^2)$: inserting the extra end projection in $B^2$
does not change the trace. Since $B^T=-B$,
\[
 -\operatorname{tr}(B^2)=\operatorname{tr}(B^TB)
     =\sum_{a,b}B_{ab}^2=\|B\|_F^2.
\]
Thus $\Psrc\Ggrad\Psrc=0$. Repeating with $I-\Psrc$ gives the other
diagonal block zero. In the same source-adapted coordinates,
$\Jsrc=\operatorname{diag}(I,-I)$; multiplication gives
$\Jsrc\Ggrad\Jsrc=-\Ggrad$.

There are two decompositions in use. The rotation planes in Step 2(b)
diagonalize $\W$ and define $\Ggrad$. The source/complement decomposition
in this step diagonalizes $\Psrc$ and $\Jsrc$. They need not coincide.
The gradient equations connect them.

\section{Obtain blocks with one or two source directions}

\subsection{Positive gradient levels carry complex coordinates}
\textbf{Inference.} Each eigenspace $\Vlevel$ of $-\Ggrad^2$ for an
eigenvalue $\glev^2>0$ is preserved by $\Ggrad,\W,\Jsrc$ and admits
complex multiplication $\Icomp=\Ggrad/\glev$.

\why Skew-symmetry gives
\[
 (-\Ggrad^2)^T=-\Ggrad^2,
 \qquad \langle-\Ggrad^2x,x\rangle=\|\Ggrad x\|^2\ge0.
\]
Thus the real spectral theorem applies. The commutation $\Ggrad\W=\W\Ggrad$
implies $\Ggrad^2\W=\W\Ggrad^2$, while anticommutation with $\Jsrc$
implies $\Jsrc\Ggrad^2\Jsrc=\Ggrad^2$. An operator commuting with
$-\Ggrad^2$ preserves each eigenspace: apply the two commuting operators
to an eigenvector to check its image has the same eigenvalue.

On $\Vlevel$, $\Icomp^2=-I$ and
$\Icomp^T\Icomp=-\Ggrad^2/\glev^2=I$. Define
$(a+ib)x=ax+b\Icomp x$. These identities make it complex scalar
multiplication compatible with the Euclidean metric. A real plane
$\operatorname{span}(x,\Icomp x)$ becomes one complex line.
Commutation makes $\W$ complex-linear and unitary. Anticommutation gives
\[
 \Jsrc((a+ib)x)=(a-ib)\Jsrc x.
\]
Since $\Jsrc^2=I$ and $\Jsrc$ preserves lengths, it is a conjugation.

To choose coordinates where this is ordinary coordinate conjugation,
let $F=\{x:\Jsrc x=x\}$. Every $x$ has the unique decomposition
\[
 x=\frac{x+\Jsrc x}{2}
     +\Icomp\left(-\Icomp\frac{x-\Jsrc x}{2}\right).
\]
Both parenthesized real components lie in $F$.
Thus an orthonormal real basis of $F$ supplies a complex orthonormal basis
of $\Vlevel$. In particular, $F$ has half its original real dimension.

\subsection{There are at most two phases on a positive level}
\textbf{Inference.} On the complex space $\Vlevel$, the eigenphases of
$\W$ satisfy $\fphase'(t)=\glev$, $0<t<\pi$, and have at most two values.

\why On an original rotation plane, $\Ggrad=\fphase'(t)J_0$ with
$\fphase'(t)>0$ for $0<t<\pi$. On its level, $\glev=\fphase'(t)$ and
$\Icomp=J_0$, so $\W$ acts as multiplication by $e^{it}$.
Direct differentiation gives
\[
 \fphase'(t)=p(2\sin(t/2))^{p-1}\cos(t/2),
\]
\[
 \fphase''(t)=p(2\sin(t/2))^{p-2}
   \bigl((p-1)\cos^2(t/2)-\sin^2(t/2)\bigr).
\]
The second derivative is positive before
$\tan^2(t/2)=p-1$ and negative after it. The first derivative is
strictly increasing and then strictly decreasing, so a positive
horizontal line meets its graph at most twice. The spectral theorem gives
\[
 \W|_{\Vlevel}=e^{i\varphi_1}\Ephase+e^{i\varphi_2}(I-\Ephase)
\]
for a complex orthogonal projection $\Ephase$. If there is only one phase,
take both phase values equal and choose $\Ephase=0$.
Two phase values can have arbitrarily large multiplicities, so a further
argument is needed to get small real blocks.

\subsection{The phase projection splits into small common blocks}
\textbf{Inference.} $\Vlevel$ splits orthogonally into real blocks of
dimension two or four preserved by $\W,\Psrc,\Ptgt$. They carry one or
two principal-angle entries, respectively.

\why In the conjugation coordinates from Step 3(a), write
$\Ephase=\Xreal+i\Yimag$ with real matrices. Its adjoint and idempotence
equations give
\[
 \Xreal^T=\Xreal,\quad \Yimag^T=-\Yimag,\quad
 \Xreal^2-\Yimag^2=\Xreal,\quad
 \Xreal\Yimag+\Yimag\Xreal=\Yimag.
\]
For a real vector $x$, the orthogonal-projection inequality gives
$0\le\langle\Xreal x,x\rangle
 =\operatorname{Re}\langle\Ephase x,x\rangle\le\|x\|^2$.
Thus the eigenvalues of the symmetric matrix $\Xreal$ lie in $[0,1]$.

Choose a unit eigenvector $u$ with eigenvalue $\lambdablock$. Then
\[
 \Xreal(\Yimag u)=(1-\lambdablock)\Yimag u,
 \qquad
 \Yimag^2u=-\lambdablock(1-\lambdablock)u.
\]
For $0<\lambdablock<1$, put
$\etablock=\sqrt{\lambdablock(1-\lambdablock)}$ and
$v=\Yimag u/\etablock$. Skew-symmetry gives
\[
 \langle u,\Yimag u\rangle=0,
 \qquad \|\Yimag u\|^2=-\langle u,\Yimag^2u\rangle=\etablock^2.
\]
Therefore $u,v$ are orthonormal and
\[
 \Xreal u=\lambdablock u,\quad
 \Xreal v=(1-\lambdablock)v,\quad
 \Yimag u=\etablock v,\quad \Yimag v=-\etablock u.
\]
Their real coordinate plane is preserved by both matrices. At
$\lambdablock=0$ or $1$, the norm calculation gives $\Yimag u=0$, so
the coordinate line spanned by $u$ suffices. The case
$\lambdablock=1/2$ still uses the two-vector construction: $u$ and $v$
are orthogonal even though they have the same eigenvalue of $\Xreal$.

The orthogonal complement of a chosen coordinate block is preserved too:
if $x$ is perpendicular to it, test $\langle\Xreal x,z\rangle$ using
symmetry and $\langle\Yimag x,z\rangle$ using skew-symmetry.
Finite-dimensional induction exhausts the real coordinate space.
Complexifying these coordinate lines or planes gives complex dimensions
one or two, hence original real dimensions two or four. These complex
spans are preserved by $\Ephase$ and coordinate conjugation $\Jsrc$,
hence by $\W$ and $\Psrc=(I+\Jsrc)/2$. Orthogonality and invariance
also give invariance under $\W^T$, so
$\Ptgt=\W\Psrc\W^T$ preserves them.
The fixed space of conjugation has half the real dimension. This is the
source space in the block, so its rank is one or two and its angle list
has one or two entries.
\rehearse{Why are two phase values insufficient by themselves? Why does
a real coordinate plane here become a four-dimensional block in the
original real space?}

\subsection{The zero-gradient space gives fixed directions and planes}
\textbf{Inference.} On $\ker\Ggrad$, a minimizing transporter decomposes
into fixed lines and two-dimensional one-angle blocks.

\why The scalar derivative is strictly positive on $(0,\pi)$ and zero
at the scalar phases $0,\pi$. Thus the restriction of $\W$ has only
eigenvalues $+1,-1$ and is a symmetric reflection $I-2N$, where $N$ is
the orthogonal projection onto its $-1$ eigenspace.
The two-projection decomposition for $\Psrc,N$ splits the space into
common invariant lines and planes; it follows from the same principal-vector
construction used in Step 1(b), allowing unmatched directions when ranks
differ. On the nontrivial planes the source rank is one. The same blocks
preserve $\W=I-2N$ and $\Ptgt=\W\Psrc\W^T$.

On an invariant line, $\Psrc$ is either zero or identity, and so is
$\Ptgt$ with the same choice. Replacing the scalar action $-1$ by $+1$
therefore preserves transport and lowers the power cost from $2^p$ to zero.
Global minimality excludes every such $-1$ line.

\section{Calculate the one- and two-angle costs}

\subsection{One angle chooses a rotation or a reflection}
\textbf{Inference.} The minimum on a real principal plane is
\[
 \single{i}=\min\{2\fphase(\theta_i),2^p\}.
\]

\why A planar rotation carrying an unoriented source line onto its target
rotates through $\theta_i$ modulo $\pi$. The smaller displacement cost is
$2\fphase(\theta_i)$ because $0\le\theta_i\le\pi/2$ and
$\fphase$ is increasing on $[0,\pi]$. The alternative folded angle is
$\pi-\theta_i\ge\theta_i$. Every planar reflection has eigenvalues
$1,-1$ and displacement singular values $0,2$, so its cost is $2^p$.
Reflections carrying the source line onto the target exist, for example
reflection in an angle-bisecting line. These two determinant classes
exhaust $O(2)$.

\subsection{A positive-determinant pair reduces to one scalar minimum}
\textbf{Inference.} Order the two angles so that $\theta_i\ge\theta_j$ and set
\[
 \hmean=(\theta_i+\theta_j)/2,\qquad
 \dhalf=(\theta_i-\theta_j)/2.
\]
The positive-determinant minimum on their four-dimensional sum is
\begin{equation}\label{eq:study-pair-cost}
 \paired{i}{j}=\min_{\taupair\in[\dhalf,\hmean]}
       2\bigl(\fphase(\hmean+\taupair)+\fphase(\hmean-\taupair)\bigr).
\end{equation}

\why We separate the matrix parametrization, its phase calculation,
the lower bound, and attainment of that bound.

\paragraph{Parametrize every transporter.}
Use principal coordinates with the two source directions first:
\[
 \Dpair=\begin{pmatrix}\Cpair&-\Spair\\\Spair&\Cpair\end{pmatrix},
 \quad \Cpair=\operatorname{diag}(\cos\theta_i,\cos\theta_j),
 \quad \Spair=\operatorname{diag}(\sin\theta_i,\sin\theta_j).
\]
This direct rotation maps the source onto the target. For any other
transporter $\W$, the orthogonal matrix $\Dpair^T\W$ preserves the
source and its complement, so
\[
 \W=\Dpair\operatorname{diag}(\Ufree{0},\Ufree{1}),
 \qquad \Ufree{0},\Ufree{1}\in\like{OrthogonalGroup}{O(2)}.
\]
This parametrization includes mixing between the principal planes.
If $\det\W=1$, the two free factors are both rotations or both reflections.

\paragraph{Calculate the phases.}
For rotations put $\Ufree{0}=R_{r+s}$ and $\Ufree{1}=R_{r-s}$.
Define $\phaseX,\phaseY\in[0,\pi]$ by
\[
 \cos\phaseX=\cos\hmean\cos s,
 \qquad \cos\phaseY=\cos\dhalf\cos r.
\]
For reflections write $\Ufree{0}=R_{r+s}J_2$,
$\Ufree{1}=R_{r-s}J_2$, with $J_2=\operatorname{diag}(1,-1)$, and define
\[
 \cos\phaseX=-\sin\dhalf\cos r,
 \qquad \cos\phaseY=\sin\hmean\cos s.
\]
In both cases the characteristic polynomial gives eigenvalues
$e^{\pm i(\phaseX+\phaseY)},e^{\pm i(\phaseX-\phaseY)}$.
Here is an explicit way to check that matrix calculation. Set
$u=\cos\phaseX$, $v=\cos\phaseY$. Block multiplication gives
\[
 \operatorname{tr}\W=4uv,\qquad
 \operatorname{tr}(\W^2)=16u^2v^2-8(u^2+v^2)+4.
\]
To reproduce the second identity, expand
\[
 \operatorname{tr}(\W^2)=
 \operatorname{tr}((\Cpair\Ufree{0})^2)
 +\operatorname{tr}((\Cpair\Ufree{1})^2)
 -2\operatorname{tr}(\Spair\Ufree{1}\Spair\Ufree{0}).
\]
For $A=\operatorname{diag}(a,b)$,
\begin{align*}
 \operatorname{tr}((AR_\alpha)^2)&=(a+b)^2\cos^2\alpha-2ab,\\
 \operatorname{tr}((AR_\alpha J_2)^2)&=(a-b)^2\cos^2\alpha+2ab,\\
 \operatorname{tr}(AR_\beta AR_\alpha)
   &=(a^2+b^2)\cos\alpha\cos\beta-2ab\sin\alpha\sin\beta,\\
 \operatorname{tr}(AR_\beta J_2 AR_\alpha J_2)
   &=(a^2+b^2)\cos\alpha\cos\beta+2ab\sin\alpha\sin\beta.
\end{align*}
Substitute $\alpha=r+s$, $\beta=r-s$, the principal cosines or sines,
and the definitions of $\hmean,\dhalf$. The coefficient of $z^2$ in the
characteristic polynomial is
$((\operatorname{tr}\W)^2-\operatorname{tr}(\W^2))/2$.
Orthogonality and determinant one make the polynomial reciprocal, yielding
\[
 \det(zI-\W)=z^4-4uvz^3+
       \bigl(4(u^2+v^2)-2\bigr)z^2-4uvz+1.
\]
Multiplying the two quadratic factors for phases
$\phaseX+\phaseY$ and $\phaseX-\phaseY$ gives precisely this polynomial.
These calculations justify the phase formulas without a choice of
individual eigenvectors.

In the rotation case the cosine bounds give
\[
 \hmean\le\phaseX\le\pi-\hmean,
 \qquad \dhalf\le\phaseY\le\pi-\dhalf.
\]
In the reflection case they first give
$\pi/2-\dhalf\le\phaseX\le\pi/2+\dhalf$ and
$\pi/2-\hmean\le\phaseY\le\pi/2+\hmean$.
These intervals lie in the same displayed rectangle because
$\hmean+\dhalf=\theta_i\le\pi/2$.

\paragraph{Reduce the rectangle to an edge.}
The displacement cost is
\[
 2\left[\fphase(\min\{\phaseX+\phaseY,2\pi-\phaseX-\phaseY\})
                    +\fphase(|\phaseX-\phaseY|)\right].
\]
Replacing $(\phaseX,\phaseY)$ by $(\pi-\phaseX,\pi-\phaseY)$
preserves the rectangle and cost, so assume $\phaseX+\phaseY\le\pi$.
Choose
\[
 \taupair=\max\{\dhalf,\hmean-|\phaseX-\phaseY|\}.
\]
We have $\dhalf\le\taupair\le\hmean$ and
$\hmean-\taupair\le|\phaseX-\phaseY|$ by this choice.
For the other phase, the rectangle gives
$\hmean+\dhalf\le\phaseX+\phaseY$ and
\[
 2\hmean\le2\max\{\phaseX,\phaseY\}
              =\phaseX+\phaseY+|\phaseX-\phaseY|.
\]
Both candidates for $\hmean+\taupair$ are therefore at most
$\phaseX+\phaseY$. Monotonicity of $\fphase$ shows that replacing the
phases by $\hmean+\taupair,\hmean-\taupair$ cannot increase cost.

\paragraph{Attain every point of the edge.}
For any $\taupair\in[\dhalf,\hmean]$, choose
\[
 \Ufree{0}=\Ufree{1}=R_\gamma,
 \qquad \gamma=\arccos(\cos\taupair/\cos\dhalf).
\]
The ratio belongs to $[0,1]$ because
$0\le\dhalf\le\taupair\le\pi/2$, and its denominator is positive since
$\dhalf\le\pi/4$. In the rotation parametrization, $r=\gamma$, $s=0$,
so $\phaseX=\hmean$, $\phaseY=\taupair$.
The scalar cost is continuous on a compact interval and attains its
minimum. The rectangle bound and this construction prove
\cref{eq:study-pair-cost} in both directions.

\subsection{The unrestricted pair also compares two singletons}
\textbf{Inference.} The minimum over both determinant classes is
\[
 \min\{\paired{i}{j},\single{i}+\single{j}\}.
\]

\why For negative determinant, one free factor is a rotation and the
other a reflection. With $c_i=\cos\theta_i\le c_j=\cos\theta_j$,
their contributions to $\operatorname{tr}\W$ have the forms
$(c_i+c_j)\cos\alpha$ and $(c_i-c_j)\cos\beta$, in either order. Thus
\[
 |\operatorname{tr}\W|\le c_i+c_j+|c_i-c_j|=2c_j.
\]
A real orthogonal $4\times4$ matrix of determinant $-1$ has eigenvalues
$1,-1,e^{it},e^{-it}$, allowing $t=0,\pi$ for repeated real eigenvalues.
Its trace is $2\cos t$, so $t\ge\theta_j$ and its cost is at least
$2^p+2\fphase(\theta_j)$. Reflecting on the larger-angle principal
plane and rotating on the smaller-angle plane attains this bound.

There are four choices when the two planes are transported independently:
two rotations, two reflections, or either mixed choice. The paired
positive-determinant minimum is at most the first two costs: evaluate
\cref{eq:study-pair-cost} at $\taupair=\dhalf$ and
$\taupair=\hmean$, respectively. For the second comparison,
$2\fphase(\theta_i+\theta_j)\le2\cdot2^p$.
The cheaper mixed choice is exactly the negative-determinant minimum.
The minimum of all four independent choices is
$\single{i}+\single{j}$. Comparing it with $\paired{i}{j}$ therefore
includes both determinant classes.
\rehearse{Where did we account for a four-dimensional block that is
cheaper as two singletons? Why does a paired block sometimes fix two
directions rather than rotate both principal planes independently?}

\section{Prove equality of the minimum costs}

For a matching $\sym{Matching}$, define
\[
 \sym{MatchingPowerCost}(\sym{Matching})
  =\sum_{\{i,j\}\in\sym{Matching}}\paired{i}{j}
     +\sum_{i\text{ unmatched}}\single{i}.
\]

\subsection{A matching constructs a transporter of the same cost}
\textbf{Inference.}
$\min_{\W\in\Fset}\Phicost(\W)
 \le\min_{\sym{Matching}}\sym{MatchingPowerCost}(\sym{Matching})$.

\why Use the mutually orthogonal principal planes from Step 1(b).
For an edge $\{i,j\}$, use the attaining paired transporter on
$\Pi_i\oplus\Pi_j$ from Step 4(b). On an unmatched plane, use the
attaining rotation or reflection from Step 4(a). Use identity on all
remaining common source/target and ambient directions. A matching uses
each label at most once, so these blocks are mutually orthogonal.
Their orthogonal sum is feasible and has exactly the sum of their costs.
The unrestricted minimum is no larger than the cost of any such
constructed transporter, hence no larger than the best matching cost.

\subsection{A minimizing transporter constructs a matching of the same cost}
\textbf{Inference.} The reverse inequality holds, and hence
\begin{equation}\label{eq:study-power-matching}
 \min_{\W\in\Fset}\Phicost(\W)
   =\min_{\sym{Matching}}\sym{MatchingPowerCost}(\sym{Matching}),
 \qquad 1<p<2.
\end{equation}

\why Steps 2--3 decompose the chosen global minimizer into fixed
directions and blocks carrying one or two angle entries. Each restriction
must minimize its local transport cost. Otherwise replace it by a cheaper
local transporter: the full map is still orthogonal and feasible, and
block additivity strictly lowers the global cost, a contradiction.
Step 4 assigns a singleton to a one-angle block. For a two-angle block,
choose the paired cost or two singleton costs according to which attains
its minimum. Step 1(b) identifies these entries with the original angle
list, including multiplicities. Each occurrence belongs to one block,
so the selected pairs form a matching.

Zero-angle entries do not lower the cost of a positive entry: their
singleton cost is zero, while pairing $\theta$ with zero forces
$\dhalf=\hmean=\theta/2$ and costs $2\fphase(\theta)$, at least the
singleton minimum. Two zero entries also cost zero. Remove such pairings
and keep zero-angle directions fixed. If this strictly lowered the cost,
Step 5(a) would construct a cheaper feasible transporter, contradicting
global minimality. The result is therefore a matching on exactly
the original positive labels with the minimizer's cost. Its cost is at
least the minimum matching cost, proving the reverse inequality.

The decomposition used to extract this matching need not be the same
basis used to realize it. The equality of the angle lists and the local
cost formulas are what permit the construction back in the original
principal coordinates.
\rehearse{State the two inequalities separately. Which one is proved by
a construction, and which one uses the structure of a global minimizer?}

\section{Pass to the trace norm and evaluate the weights}

\subsection{The minimum values converge}
\textbf{Inference.} Both sides of \cref{eq:study-power-matching} converge
to their corresponding trace-norm minimum values as $p$ approaches $1$
from above.

\why Every singular value of $I-\W$ lies in $[0,2]$, since
$\|I-\W\|_{\mathrm{op}}\le\|I\|_{\mathrm{op}}+\|\W\|_{\mathrm{op}}=2$.
The map $(p,s)\mapsto s^p$, with $0^p=0$, is continuous on
$[1,3/2]\times[0,2]$, hence uniformly continuous. Therefore
\[
 \epsilon_p=n\sup_{0\le s\le2}|s^p-s|\longrightarrow0,
 \qquad
 |\Phicost(\W)-\|I-\W\|_1|\le\epsilon_p
 \quad\text{for every }\W\in\Fset.
\]
If two functions on the same nonempty set differ by at most $\epsilon$,
their attained minima differ by at most $\epsilon$: evaluate each
function at a minimizer of the other. This proves convergence of the
left-hand minimum. It also explains why pointwise convergence alone is
insufficient for this inference.

The singleton costs converge because they are minima of two continuous
expressions. For each pair, the scalar costs in Step 4(b) converge uniformly
on the fixed compact interval $[\dhalf,\hmean]$, so the same argument
gives convergence of their minima. There are finitely many matchings.
If $C_p(M)\to C(M)$ for each of finitely many $M$, then
$\max_M|C_p(M)-C(M)|\to0$, and the minimum values converge as well.
We pass minimum \emph{values} to the limit; no convergence of the chosen
minimizers is required.

\subsection{The limiting weights are the stated trace costs}
\textbf{Inference.} At $p=1$,
\[
 a_i(1)=\min\{4\sin(\theta_i/2),2\}=s_i,
 \qquad b_{ij}(1)=4\sin((\theta_i+\theta_j)/2)=e_{ij}.
\]

\why The singleton formula follows by substitution. For a pair,
the sum-to-product identity gives
\begin{align*}
 2\bigl(f_1(\hmean+\taupair)+f_1(\hmean-\taupair)\bigr)
 &=4\bigl(\sin((\hmean+\taupair)/2)
                   +\sin((\hmean-\taupair)/2)\bigr)\\
 &=8\sin(\hmean/2)\cos(\taupair/2).
\end{align*}
This is nonincreasing for $\dhalf\le\taupair\le\hmean$, so its minimum
is attained at $\taupair=\hmean$ and equals $4\sin\hmean$.
The attaining construction has rotational phases $2\hmean$ and zero:
it rotates one real plane through $\theta_i+\theta_j$ and fixes the
other real plane. Its displacement has two singular values
$2\sin\hmean$ and two zeros.
Together with Step 6(a), this proves the stated trace-norm equality.
Compactness of $\Fset$ also gives attainment at $p=1$.

Finally, define the all-singleton cost and edge weights
\[
 \sym{SingletonBaseline}=\sum_i s_i,
 \qquad \sym{MatchingEdgeWeight}=s_i+s_j-e_{ij}.
\]
Replacing two singleton costs by one paired cost gives
\[
 C(M)=\sym{SingletonBaseline}
       -\sum_{\{i,j\}\in M}\sym{MatchingEdgeWeight}.
\]
Thus minimizing displacement is maximum-weight matching on a finite
general graph. Unmatched vertices remain allowed. A matching of maximum
cardinality is a different optimization objective.
\rehearse{Why can the exponent-one theorem follow without differentiating
the trace norm? What is the logical role of uniform convergence?}

\appendix
\section{Why the displacement power cost is differentiable}
\label{sec:study-differentiability}

The first-variation argument uses differentiability to obtain a linear
functional of the direction $B$. Here is the finite-dimensional matrix
calculus behind that assertion, including repeated and zero singular values.

\subsection*{A trace derivative for a real symmetric matrix}
Let $A$ be real symmetric and let $h:\mathbb R\to\mathbb R$ be continuously
differentiable on a neighborhood of its spectrum. Then
\[
 D\bigl[\operatorname{tr}h(A)\bigr](H)
       =\operatorname{tr}(h'(A)H)
 \quad\text{for symmetric }H.
\]
One way to justify this without differentiating individual eigenvectors
is to group eigenvalues by their distinct values at $A$.
Fix such a value $\lambda$, its eigenspace projection $P_\lambda$, and
multiplicity $m$. For a small symmetric perturbation $H$, Weyl's bound
places exactly $m$ nearby eigenvalues at $\lambda+O(\|H\|)$.
This bound follows from the min--max formula because every Rayleigh
quotient changes by at most $\|H\|$.

Choose orthonormal eigenvectors $v_1,\ldots,v_m$ for that cluster of
$A+H$. The gap between $\lambda$ and the other distinct eigenvalues of
$A$ implies $(I-P_\lambda)v_j=O(\|H\|)$: project the eigenvector
equation onto the complementary space, where $A-\lambda I$ has a
bounded inverse. Hence
\[
 \langle Av_j,v_j\rangle=\lambda+O(\|H\|^2).
\]
The cluster projection $P_\lambda(H)=\sum_jv_jv_j^T$ converges to
$P_\lambda$, as the following identity makes explicit:
\[
 \|P_\lambda(H)-P_\lambda\|_F^2
   =2\sum_{j=1}^m\|(I-P_\lambda)v_j\|^2=O(\|H\|^2).
\]
Both projections have rank $m$; expanding the squared Frobenius norm
gives the identity. Summing the eigenvector equations gives the cluster sum
\[
 \sum_{j=1}^m\lambda_j(A+H)
   =m\lambda+\operatorname{tr}(P_\lambda H)+o(\|H\|).
\]
For each eigenvalue, scalar differentiability gives
$h(\lambda+\delta)=h(\lambda)+h'(\lambda)\delta+o(|\delta|)$.
The Weyl bound controls all these remainders by $o(\|H\|)$.
Summing first over a cluster, then over the finitely many distinct
eigenvalues, proves the trace derivative. If $A$ has just one distinct
eigenvalue, its cluster is the whole space and no complementary inverse
is needed. Repeated eigenvalues are handled by the cluster sum.

\subsection*{Apply the trace derivative to a symmetric dilation}
For a real square matrix $D$, set
\[
 \mathcal A(D)=\begin{pmatrix}0&D\\D^T&0\end{pmatrix}.
\]
If $D=U\Sigma V^T$ is a singular value decomposition, conjugation by
$\operatorname{diag}(U,V)$ reduces this matrix to
$\bigl(\begin{smallmatrix}0&\Sigma\\\Sigma&0\end{smallmatrix}\bigr)$.
Its eigenvalues are the positive and negative singular values of $D$,
including zeros. Therefore, with $h(x)=|x|^p$,
\[
 F_p(D):=\sum_j s_j(D)^p
       =\tfrac12\operatorname{tr}h(\mathcal A(D)).
\]
For $p>1$, $h$ is continuously differentiable with
$h'(x)=p\operatorname{sign}(x)|x|^{p-1}$ and $h'(0)=0$.
The preceding trace derivative applies, even when $D$ is singular.
The same block diagonalization gives
\[
 h'(\mathcal A(D))=
 p\begin{pmatrix}0&U\Sigma^{p-1}V^T\\
                       V\Sigma^{p-1}U^T&0\end{pmatrix}.
\]
For a matrix direction $H$, the chain rule and block multiplication yield
\[
 DF_p(D)[H]
       =p\operatorname{tr}\bigl((U\Sigma^{p-1}V^T)^T H\bigr).
\]
No negative power of a zero singular value occurs.
Taking $D=I-\W$ and $H=-\W B$ proves that the first variation used in
Step 2(b) exists and is linear in $B$. The sign-averaging and planar
calculation there identify it with $-\operatorname{tr}(\Ggrad B)$.

\section{Reconstruct the proof without the text}

Try writing the following chain and supplying a justification for each arrow:
\[
 \text{global minimizer}
 \ \Longrightarrow\ \text{stationarity}
 \ \Longrightarrow\ \Jsrc\Ggrad\Jsrc=-\Ggrad,
 \qquad \Ggrad\W=\W\Ggrad;
\]
\[
 \begin{aligned}
 \text{levels of }-\Ggrad^2
 &\ \Longrightarrow\ \text{complex coordinates and at most two phases}\\
 &\ \Longrightarrow\ \text{real blocks of dimension two or four};
 \end{aligned}
\]
\[
 \begin{aligned}
 \text{block minima and additive costs}
 &\ \Longrightarrow\ \text{two matching inequalities}\\
 &\ \Longrightarrow\ \text{trace-norm limit}.
 \end{aligned}
\]
Be able to identify where each of the following enters: equal subspace
dimensions, finite ambient dimension, $p>1$, $p<2$, global rather than
merely stationary optimality, the real scalar field, and uniform convergence.

For a concrete reconstruction, use the three angles $30^\circ,45^\circ,75^\circ$.
There are four matchings. Their costs are approximately $4.5660$ for no
pair, $4.4350$ for $\{1,2\}$, $4.7041$ for $\{1,3\}$, and $4.4994$ for
$\{2,3\}$. Explain why the best construction uses a four-dimensional
paired block and a two-dimensional reflection, and which part of the proof
establishes that unrestricted six-dimensional mixing cannot do better.

\end{document}
